% 本文件由 LaTeX 排版 Agent 根据有效 translation.json 自动生成。
% author=Origen; work_id=0412; title=论原理

\chapter{论原理}

% structure: p0001 source=h1 target=omitted reason=page-title-already-rendered-by-parent

序言\newline{}卷一\newline{}卷二\newline{}卷三\newline{}卷四\par

\SourceRecord{Translated by Frederick Crombie.；Ante-Nicene Fathers；Vol. 4.；Edited by Alexander Roberts, James Donaldson, and A. Cleveland Coxe.；Buffalo, NY: Christian Literature Publishing Co.,；1885.；<http://www.newadvent.org/fathers/0412.htm>.}

% structure: p0001 source=h1 target=section reason=page-title

\section{\WorkTitle{论首要原理}（前言）}

% structure: p0002 source=h2 target=omitted reason=duplicate-page-title

1. 凡相信并确信恩宠与真理是由耶稣基督而来，并知道基督就是真理，正如祂亲口所说：\BibleQuote{我是真理，}的人，他们获取激发人度好而幸福生活的知识，不是从别处，而是从基督的言语和教训。我们所谓基督的言语，不仅指祂成为人、居住在肉身内时所说的话；因为早在那时以前，基督，\OriginalTerm{圣言}{Word of God}，已经在梅瑟和先知们内。因为若没有天主的圣言，他们怎能预言基督呢？我们若无意将本论文限于一切可能的简练范围，就不难从圣经中证明，梅瑟和先知们所说所做的一切，都是由于充满基督的圣神。因此，我认为引用保禄在\WorkTitle{致希伯来人书}中的这一见证就足够了，他说：\BibleQuote{因着信德，梅瑟长大以后，拒绝被称为法郎公主的儿子；宁愿同天主的子民一起受苦，也不愿暂时享受罪恶的快乐；把基督的凌辱看得比埃及的宝藏更为宝贵。}此外，基督升天后，祂在宗徒们内说话，保禄用这些话表明：\BibleQuote{或者你们寻求基督在我内说话的凭据？}\par

2. 然而，许多自称相信基督的人，不仅在琐碎小事上，而且在最重要的问题上，例如关于天主、主耶稣基督或圣神，甚至关于其他受造的存在物，即各种德能和圣善的力量上，彼此分歧；因此，似乎有必要首先就这些问题中的每一个确定明确的界限，制定毋庸置疑的规则，然后再转向对其他问题的研究。因为当我们不再从所有声称拥有真理的人（尽管许多希腊人和异族人声称要传扬真理）中，从他们认为错误的意见中寻求真理，而是相信基督是天主子，并确信必须从祂那里学习真理之后，既然有许多人自以为持有基督的意见，但其中一些人又与他们的前辈意见不同，然而，教会从宗徒们开始有序传承并保留至今的教导仍然存在，只有与教会传统毫无差异的，才应被接受为真理。\par

3. 应当知道，圣宗徒们在宣讲基督信仰时，对于某些他们认为对每个人都有必要的问题，甚至对那些在探究神圣知识方面似乎有点迟钝的人，也极其清晰地表达了自己；然而，他们将其陈述的根据留待那些配得上圣神卓越恩赐的人，特别是藉着圣神本身获得语言、智慧和知识恩赐的人去探究；至于其他问题，他们只陈述事实如此，而对存在的方式或起源保持沉默；显然是为了使后来更热忱的、爱好智慧的后继者有一个锻炼的题材，可以展示他们才能的果实——我指的是那些准备使自己成为智慧合格而相称的领受者的人。\par

4. 宗徒教导中清楚传达的特别要点如下：\par

\Emph{第一}，有一位天主，祂创造并安排了一切，在无物存在时，祂呼唤万物进入存在——这位天主从创世和世界奠基之初——是一切义人的天主：亚当、亚伯尔、舍特、厄诺士、哈诺客、诺厄、色勒、亚巴郎、依撒格、雅各伯、十二位圣祖、梅瑟和先知们；这位天主在末世，正如祂事先藉着先知预告的，派遣了我们的主耶稣基督，首先召唤以色列归向祂，其次，在以色列人民不忠之后，召唤外邦人。这位公义而良善的天主，我们主耶稣基督的父，亲自赐下了法律、先知和福音，祂也是宗徒们以及旧约和新约的天主。\par

\Emph{第二}，耶稣基督自己，祂来了（到世界上），是在万有之先由父所生；祂在创造万有时服侍了父——\BibleQuote{因为万有都是藉着祂而造}——祂在末世，脱下（祂的荣耀），成为人，道成肉身，虽是天主，却成为人，并保持祂所是的天主；祂取了与我们相似的身体，只在这一点上不同，即由童贞女和圣神所生；这位耶稣基督真正地诞生了，真正地受苦了，并不是仅仅表面上承受了（众人）共同的死亡，而是真正地死了；祂真正地从死者中复活了；复活后，祂与门徒交谈，并升了天。\par

然后，\Emph{第三}，宗徒们说到圣神与圣父圣子同享尊荣和威严。但关于祂，并没有清楚区分祂应被视为受生或固有的，还是也是天主子或不是：因为这些问题需要我们尽己所能从圣经中探究，并要求仔细研究。而且，这个圣神感动了每一位圣人，无论是先知还是宗徒；在旧约时代的人中，和在基督降临时受感的人中，并不是有不同的圣神，这一点在全教会中都教导得非常清楚。\par

5. 在这些之后，宗徒的教导也指出：灵魂自身具有实体和生命，在离开世界之后，将按其所应得而受赏，注定要获得永生和真福的继承——如果其行为为此获得它——或者被交付于永火和惩罚——如果其罪债使之沉沦于此。此外，还有从死者中复活的时候，那时这身体，如今\BibleQuote{播种在腐败中，将要在不朽中复活}，并且\BibleQuote{播种在羞辱中，将要在光荣中复活}。教会教导中也明确界定了这一点：每个理性的灵魂都拥有自由意志和自主选择；它必须与魔鬼及其天使以及敌对势力进行斗争，因为他们企图用罪过困住它；但若我们正直而明智地生活，就应努力挣脱那类重担。由此也推知，我们理解自己并不受必然性所辖制，以致被迫无论如何、甚至违背意愿去行善或作恶。因为若我们是自己的主人，有些影响力或许驱使我们犯罪，另一些则帮助我们获得救恩；然而我们并不被任何必然性强迫去行是或行非——那些认为星辰的运行和运动是人类行为的原因的人所持的正是这种观点，不仅认为超越自由意志影响的事件如此，连那些置于我们能力之内的事件也是如此。至于灵魂，它是藉由一种称为\OriginalTerm{传殖说}{traducianism}的过程来自种籽，以致其理性或实体可被视为置于身体本身的精粒之中，还是有任何其他起源；而这起源本身，是藉由生育与否，还是从外部赋予身体与否，在教会教导中并未足够清晰地加以区分。\par

6. 关于魔鬼及其天使，以及敌对势力，教会的教导业已指明这些存在确实存在；但他们是什么，或如何存在，并未足够清晰地解释。然而，大多数人持有这种意见：魔鬼曾是一位天使，他成为背教者后，带领尽可能多的天使与他一同堕落，这些天使至今被称为他的天使。\par

7. 这也是教会教导的一部分：世界是在某一时间受造并开始存在的，并且将因邪恶而被毁灭。但世界之前有什么，或世界之后将有什么，尚未为众人所确知，因为教会教导中没有明确的陈述。\par

8. 最后，圣经是由天主的圣神所写成的，并且具有一种意义，不仅是一见之下显浅的含义，还有另一种被多数人忽略的意义。因为所写的文字是某些奥迹的形态和神圣事物的图像。关于这一点，整个教会有一个共同意见：全部法律确实是属神的；但法律所传达的属神意义并非为所有人所知，仅为那些在智慧与知识的言语上领受了圣神恩宠的人所知。\par

术语\OriginalTerm{无形体}{\GreekText{ἀσώματον}}，即无形体的，不仅在其他许多著作中，而且在我们自己的圣经中也已废弃不用。若有人从题为\WorkTitle{伯多禄的教义}的小册子中向我们引用这句话，其中救主似乎对祂的门徒说：\Quote{我不是一个无形体的魔鬼}，我首先要回答：那著作不包括在教会书籍之中；因为我们能证明它既不是伯多禄所写，也不是由天主圣神所感动的人所写。但即使这一点被承认，那里的\OriginalTerm{无形体}{\GreekText{ἀσώματον}}一词并不传达希腊和外邦作者在哲学家讨论无形体本性时所指的含义。因为在那本小册子中，他使用短语\Quote{无形体的魔鬼}来表示魔鬼身体的形态或轮廓，无论它是什么，都不像我们这粗重可见的身体；但根据小册子作者的意思，必须理解为祂没有魔鬼所具有的身体，魔鬼的身体自然细微，薄如空气（因此许多人不考虑或称之为无形体），而是具有实在可触摸的身体。按照人类的习惯，凡不是那种性质的，都被简单人或无知者称为无形体；就好像有人说我们呼吸的空气是无形体，因为它不是一种可以抓握、持有或抵抗压力的物体。\par

9. 然而，我们将探究希腊哲学家称为\OriginalTerm{无形体}{\GreekText{ἀσώματον}}或\Quote{无形体}的事物在圣经中是否以另一个名称出现。因为还需要研究如何理解天主本身——是如同有形体的，具有某种形状，还是与形体不同的性质——我们的教导中对此没有清楚指明。同样的问题还须针对基督和圣神，以及每一个灵魂和每一个具有理性本质的存在进行研究。\par

10. 这也是教会教导的一部分：存在某些天使和某些善良的力量，他们在完成人的救恩中作祂的仆人。然而，他们何时受造，或属何种本性，或如何存在，并未清楚说明。关于太阳、月亮和星辰，它们是有生命的还是无生命的，没有明确的宣示。\par

因此，每个人都必须按照\Quote{以知识之光光照自己}的训示，使用这类元素和基础，如果他想要形成一个连贯的系列和真理体系，与所有这一切的理性相和谐，以便通过清晰而必要的陈述，查明每个个别主题的真理，并如我们所说，形成一套教义体系，通过例证和论证——无论是他在圣经中发现的，还是他通过严密追索后果并遵循正确方法所演绎出来的。\par

\SourceRecord{Translated by Frederick Crombie.；Ante-Nicene Fathers；Vol. 4.；Edited by Alexander Roberts, James Donaldson, and A. Cleveland Coxe.；Buffalo, NY: Christian Literature Publishing Co.,；1885.；<http://www.newadvent.org/fathers/04120.htm>.}

% structure: p0001 source=h1 target=section reason=page-title

\section{论原理（第一卷）}

% structure: p0002 source=h2 target=subsection reason=relative-heading-rank; base=h2

\subsection{第一章 论天主}

\Emph{我知道}有些人会试图说，即使按照我们自己的圣经的宣告，天主也是形体，因为在梅瑟的著作中他们发现写著：\BibleQuote{我们的天主是吞噬的火；}并且在\BibleBook{若望}{福音}中：\BibleQuote{天主是神，朝拜祂的人应当以心神和真理朝拜祂。}按照他们，火和神被视为无非是形体。现在，我想问这些人，对于那处宣告天主是光的经文，他们有什么可说的？正如若望在他的书信中写道：\BibleQuote{天主是光，在祂内毫无黑暗。}确实，祂是那光，照亮所有能接受真理之人的全部理智，正如在\BibleBook{}{圣咏}第三十六篇中所说：\BibleQuote{在祢的光明中我们才能见光。}因为还能提什么别的天主之光，\BibleQuote{在其中任何人见光，}除了天主的一种影响，人借此被光照，要么透彻地看见万有的真理，要么认识那被称为真理的天主本身？这就是这句话的意思：\BibleQuote{在祢的光明中我们才能见光；}即，在祢的圣言和智慧——即祢的圣子——内，我们将在祂内看见祢圣父。因为祂被称为光，难道祂被认为与太阳的光有什么相似吗？或者怎么能有丝毫理由想象，从那形体的光，任何人能获得知识的原因，并达到对真理的理解？\par

2. 那么，如果他们同意我们关于光的本性的论证——理性本身已经证明——并承认天主不能被理解为像光那样的形体，那么类似的推理也适用于\BibleQuote{吞噬的火}这一说法。因为天主作为火会吞噬什么？祂会被认为吞噬物质性的东西，如木、草、禾秸吗？如果祂是火，吞噬这类材料，这还有什么配得上天主的荣耀？但让我们思考：天主确实吞噬并彻底摧毁；祂吞噬罪恶的思想、邪恶的行为和罪恶的欲望，当它们潜入信徒的内心时；并且，祂与祂的圣子一同居住在那能接受祂的圣言和智慧的灵魂中，按照祂自己的宣告：\BibleQuote{我和父要来，并要在他那里作我们的住所？}祂在吞尽他们所有的恶习和情欲之后，使他们成为圣洁的殿宇，配得上祂自己。此外，对于那些因\BibleQuote{天主是神}这一说法而认为祂是形体的人，我认为应如下回答。圣经的习惯是，当它要指称任何与这粗糙的固体物体相对立的东西时，便称其为精神，如同在\BibleQuote{文字叫人死，精神却叫人活}这句话中，毫无疑问，\InnerQuote{文字}指的是属肉体的事物，而\InnerQuote{精神}指的是属理智的事物，我们也称之为\OriginalTerm{属神的}{spiritual}。宗徒又说：\BibleQuote{直到今日，每逢诵读梅瑟时，还有帕子盖在他们的心上；但无论何时，当他们转向主时，帕子就除去了；主的神在哪里，那里就有自由。}因为只要一个人没有转向属神的理解，帕子就盖在他的心上；这帕子，即粗俗的理解，圣经本身也被说成或被认为是被遮盖的：这就是帕子盖在梅瑟脸上对百姓说话的意思，即当律法被公开诵读时。但如果我们转向主——那里有天主的话，并且圣神启示属神的知识——帕子就除去了，我们就会以揭开的帕子在圣经中看见主的荣耀。\par

3. 既然许多圣人都分享圣神，因此不能理解祂为一个形体，被分割成有形的部分由每位圣人分享；但祂显然是一种圣化的德能，凡配得因祂的恩宠而被圣化的人，都说有分于祂。为使我们的论述更容易理解，我们举一个极不相同的事物为例。有许多人参与医学科学或技艺：我们是否因此认为，那些这样做的人取用了某种称为医学的形体的部分，这形体摆在他们面前，并以这种方式参与同一事物？或者我们不应理解为，所有以敏锐和训练有素的头脑理解这门技艺和学问本身的人，可说是分享了医术？但这些在将医学与圣神作比较时，并不能被视为完全平行的例子，因为它们只是用来证明：为许多个体所分享的事物不一定被视为形体。因为圣神与医学的方法或科学有很大不同，圣神是一个理智的存在体，并以独特的方式自立和存在，而医学根本不属于那种性质。\par

4. 但我们必须转向福音本身的言词，其中宣告说：\BibleQuote{天主是神}，我们必须表明如何按我们所说的一致来理解这一点。因为我们当探究救主是在什么场合说了这些话，是在谁面前说的，以及所探讨的主题是什么。我们毫无疑义地发现，祂是向撒玛黎雅妇人说了这些话，对她说——她按撒玛黎雅人的观点认为应在\Place{革黎斤山}朝拜天主——\BibleQuote{天主是神}。因为那撒玛黎雅妇人以为祂是犹太人，便问祂：天主应在耶路撒冷朝拜，还是在这座山上？她说：\BibleQuote{我们的祖先在这座山上朝拜，而你们却说，应该朝拜的地方是在耶路撒冷。}因此，对这撒玛黎雅妇人的意见——她以为天主应按不同地方的特权而被正确或适当地朝拜，或由犹太人在耶路撒冷，或由撒玛黎雅人在革黎斤山——救主回答说，凡愿跟随主的人必须放下对各地方的一切偏爱，并且如此表达：\BibleQuote{时候将到，那时真正的朝拜者既不在耶路撒冷，也不在这座山上朝拜圣父。天主是神，朝拜祂的人，应当以心神和真理朝拜祂。}并且请注意祂如何合乎逻辑地将心神与真理联结：祂称天主为神，为将祂与形体区分；祂称祂为真理，为将祂与阴影或肖像区分。因为那些在耶路撒冷朝拜的人，既非以真理也非以心神朝拜天主，而是屈从于天上事物的阴影或肖像；那些在革黎斤山上朝拜的人也是如此。\par

5. 那么，在尽可能驳斥了可能暗示我们要以任何程度将天主视为有形体的每一种观念之后，我们继续要说，按严格真理，天主是不可思议且无法度量的。因为无论我们能以知觉或思考获得关于天主的何种知识，我们都必须相信祂比我们所感知到的远为优越。因为，就好像我们看见一个人不能忍受一线光或一盏极小的灯焰，而我们想要让这样一位视力无法接受比我们所述更大光量的人认识太阳的辉煌和灿烂，难道不是有必要告诉他，太阳的光辉是无法言喻、无法估量地比他所见的一切光更美好、更荣耀吗？同样，我们的理解力，当被血肉的枷锁禁锢，并因其对这类物质实体的参与而变得更为迟钝和愚钝时，尽管与我们本性的身体相比，它被认为远为优越，但在其努力审视和观察无形事物时，它几乎只占一线光或一盏灯的位置。但在所有理性的、即无形的存有中，有谁如此超乎一切——如此不可言喻、不可估量地超乎一切——如同天主？祂的本性不能被任何人性的理解力（即使是最纯洁、最明亮的）所把握或看见。\par

6. 但是，如果我们用另一个比喻来使事情更清楚，这并不显得荒谬。我们的眼睛常常不能直视光本身的性质——即太阳的实体；但当我们看到它的光辉或它的光线可能透过窗户或一些接受光的小孔射入时，我们可以反思那光体的供应和源头是多么巨大。同样，上智安排的事工和这整个世界的计划，就是天主性质的某种射线，是与祂的实在和存有相比较而言的。因此，由于我们的理解力本身不能看见天主本身，它就从祂事工的美和祂造物的合宜中认识世界的圣父。因此，天主不应被思想为有形体的或存在于形体中的，而是一种单一的理智本性，在自身内不容许任何附加物；因此，祂不能被相信在祂内有较大和较小之分，而是祂在各方面都是\OriginalTerm{一}{\GreekText{Μονάς}}，并且可以说是\par

\SourceRecord{Translated by Frederick Crombie.；Ante-Nicene Fathers；Vol. 4.；Edited by Alexander Roberts, James Donaldson, and A. Cleveland Coxe.；Buffalo, NY: Christian Literature Publishing Co.,；1885.；<http://www.newadvent.org/fathers/04121.htm>.}

% structure: p0001 source=h1 target=section reason=page-title

\section{\WorkTitle{论原理（卷二）}}

% structure: p0002 source=h2 target=subsection reason=relative-heading-rank; base=h2

\subsection{第一章 论世界}

1. 虽然前卷中的所有讨论都涉及世界及其安排，但现今似乎理应专门重新讨论有关世界本身的几个要点，即其开端与终结，或在开端与终结之间发生的天主眷顾的那些安排，或那些被认为在世界受造之前已然发生、或将在终结之后发生的事件。\par

在此探讨中，首先明显的一点是：世界在其一切多样且变化的条件中，不仅由理性及更神圣的本性、以及各种身体组成，还包括无言的动物、野兽、飞禽和水中的一切活物；其次，由地方组成，即天或诸天、地或水，以及居于中间的、他们称之为\OriginalTerm{以太}{æther}的空气，以及从地而出或生于地的一切。既然世界中存在如此多的多样性，且理性存有本身之间也有如此大的差异（据信所有其他差异和多样性皆由此而生），那么，除了这一原因，还能为世界的存在指任何其他原因吗？尤其是在我们考虑到那终结——据前卷所示，万物都要借以恢复其原初状态——之时？如果这在逻辑上似乎成立，那么，如上所述，对于世界中如此巨大的多样性，我们还能想象出什么其他原因，除非是那些从原初的统一与和谐（他们最初由天主所造时处于其中）中堕落者运动与偏斜的多样性和变化？他们被逐出那善的状态，被不同的动机与欲望的困扰所牵制，被引向各种方向，从而根据其不同的倾向，将自身本性单一且不可分割的善转变成了各种类型的心智？\par

2. 然而，天主藉其智慧不可言喻的技艺，将一切事物（无论它们如何被造）转化为某种有用的目的，并为了万物的共同利益，将那些在心灵构造上彼此差异如此之大的受造物召回至同一种劳作与意图的和谐中；因此，尽管它们受不同动机的影响，却仍共同完成一个世界的圆满与完美，而且心智的这种多样性本身趋向于一个完美的终点。因为有一种能力抓住并维系着世界的一切多样性，并将不同的运动导向同一工作，以免如此浩大的世界工程因灵魂的分裂而瓦解。为此，我们认为，万有之父天主，为了确保祂所有受造物的救恩，藉其圣言与智慧那不可言喻的计划，如此安排了其中的每一事物：使每一个灵体（不论称作灵魂、理性存有或其他）都不被强迫——违反其意志的自由——走向任何别的途径，只能走向其自身心智的动机所引导的途径（以免如此便似乎夺去了运用自由意志的能力，而这必然改变存有自身的本性）；并且，这些存有不同的意志将被适当且有益地适应于一个世界的和谐——其中一些需要帮助，另一些能够提供帮助，还有一些则造成奋斗与竞争，以使那些在进步中的人，他们的勤勉更值得赞许，且经过奋斗的困难所确立的胜利后获得的等级地位更为稳固。\par

3. 尽管整个世界被安排成不同等级的职分，但其状况不应被视为带有内在矛盾与不和谐；正如我们同一个身体有许多肢体，并由同一个灵魂维系，同样，我认为整个世界也应被视为某种巨大无边的活物，藉天主的权能与理性（如同一个灵魂）得以维持。这一点，我想在圣经中也由先知的宣告所表明：\Quote{我岂不充满天地吗？——主说}；以及\Quote{天是我的宝座，地是我的脚凳}；还有救主的话，祂说我们不可\Quote{指着天起誓，因为天是天主的宝座；也不可指着地起誓，因为地是祂的脚凳}。同样，保禄在雅典人面前的演讲中也有同样的话，他说：\Quote{我们生活、行动、存留都在祂内}。因为我们如何生活、行动、存留于天主内，岂不是因祂以权能包容并维系整个世界吗？天如何是天主的宝座，地如何是祂的脚凳（如救主亲自所宣称），岂不是因祂的权能充满天地间的一切事物（如主自己的话）？而万有之父天主以其权能的圆满充满并维系世界（依据我们所引用的段落），我想没有人会难以承认。既然前述讨论已表明理性存有不同的运动及其各异的看法带来了世界中的多样性，我们就要审视：这个世界是否可能像其开端一样也有终结。因为毫无疑问，其终结必在诸多多样性与变异中寻求；而在世界的终结中所发现的这种多样性，又将为接续现今世界的另一个世界的多样性提供根据与契机。\par

4. 如果现在，在我们讨论的过程中，已经确定了这些事情是这样，那么似乎接下来我们应该考虑有形实体的本性，因为世界中的多样性不能没有物体而存在。从事物本身的本质来看，物体本性能够接受多样性和各种变化，因此它能够经历一切可能的转变，例如，木材变成火，火变成烟，烟变成空气，油变成火。食物本身，无论是人的还是动物的，不也表现出同样的变化根据吗？因为我们所吃的任何东西，都转化为我们身体的结构。但是，水如何变成土或空气，空气又如何变成火，或火变成空气，或空气变成水，虽然不难解释，但在此场合仅提一下就够了，因为我们的目的是讨论物体质料的本性。因此，我们所谓质料，是指置于物体之下的东西，即通过赋予和植入属性，使物体得以存在的东西；我们提到四种属性——热、冷、干、湿。这四种属性被植入到\OriginalTerm{质料}{\GreekText{ὕλη}}（或说质料，因为质料就其自身本性而言，是没有上述这些属性的）中，就产生了不同种类的物体。虽然这种质料，如上所述，按其自身固有本性是没有属性的，但它从未被发现没有属性而存在。我不明白为何如此多杰出的人物曾认为，这种质料——它是如此伟大，具有这样的性质，足以满足天主愿意存在的世界上所有物体，并且作为造物主的助手和仆人，接受祂希望赋予的一切形式和种类，并接受祂愿意赐予的任何属性——是非受造的，也就是说，不是由天主自己——万物的造物主——所造的，而是它的本性和能力是偶然的结果。我惊异于他们指责那些否认天主的创造能力或祂对世界的上智安排的人，并控告他们不虔敬，认为如此伟大的世界可以没有工匠或监督者而存在；而他们自己却同样犯有不虔敬的罪，说质料是非受造的，且与非受造的天主同永。按照这种观点，如果我们为了论证而假设质料不存在，正如这些人所主张的，他们说天主在没有任何东西存在时不能创造任何东西，那么无疑天主就会无所事事，因为没有质料可以运作，他们认为这质料不是由天主安排，而是偶然提供的；并且他们认为，这种偶然发现的东西，足以使天主用于如此宏大的事业，彰显祂大能的力量，并通过接纳祂全部智慧的计划，被区分并形成一个世界。现在，在我看来，这是非常荒谬的，是那些完全不知道非受造本性的能力和智慧的人的意见。但是，为了更清楚地了解事物的本质，让我们暂且假设质料不存在，而天主在以前没有任何东西存在时，使祂所愿意的事物存在了，那么我们为什么认为天主会创造比祂用自己的能力和智慧所产生的那质料更好或更大或不同种类的质料，以便使先前不存在的东西得以存在呢？祂会创造一个更糟、更劣等的质料，还是与他们所谓的非受造的质料相同的质料？现在我认为任何人都很容易看出，没有比那实际接受世界形式和种类的质料更好或更差的质料，如果不具备这些，它们就不能接受这些形式和种类。那么，称那如果被相信是由天主形成的话，无疑会被发现就是他们所称为非受造的那样的东西为非受造的，难道不是不虔敬吗？\par

5. 但是，为使我们相信圣经的权威确实如此，请听在\BibleBook{玛加伯}{下书}中，当七位殉道者的母亲鼓励她的儿子忍受酷刑时，如何确认了这一真理；因为她說：\BibleQuote{我儿，我恳求你，观看天和地，以及其中所有的一切，看到这些，就知道天主从无中创造了这一切。}\BibleRef{玛下 7:28} 在\WorkTitle{牧者书}中，也在第一条诫命中，这样说道：\Quote{首先，要相信只有一位天主，祂创造并安排了万物，并使万物从虚无中进入存在。} 或许\BibleBook{}{圣咏}中的这句话也与此相关：\BibleQuote{祂一发言，万有形成；祂一下令，万物受造。}\BibleRef{咏 33:9} 因为\BibleQuote{祂一发言，万有形成}这句话似乎是指现存事物的本质；而另一句\BibleQuote{祂一下令，万物受造}，似乎是指塑造本质本身的属性。\par

% structure: p0009 source=h2 target=subsection reason=relative-heading-rank; base=h2

\subsection{二、论物体本性的永恒性}

1. 关于这个问题，有些人习惯于探究：既然圣父生出非受造的圣子并发出圣神，不是好像先前不存在，而是因为圣父是圣子或圣神的根源与泉源，且在祂们内不能理解任何先与后；那么是否也能理解在理性本性与物质之间存在着一种类似的合一或关系。为了使这一点能得到更充分、更透彻的考察，讨论通常首先指向这个探究：这承载生命并包含精神与理性心智运动的物质本性本身，是否将与它们同样永恒，还是将完全消逝并毁灭。为了更精确地确定这个问题，我们首先需要探究，理性本性在达到圣德与幸福的顶峰之后（在我看来这是极其困难、几乎不可能达到的）是否可能完全脱离形体而存在，还是必须永远与身体结合。因此，如果有人能说明为什么它们有可能完全不需要身体，那么似乎就会得出这样的结论：如同从无中受造、经过时间间隔之后产生的物质本性在不存在时被造出一样，当它所服务的目的不再存在时，它也应停止存在。\par

2. 但是，如果完全不可能坚持这一点，即除圣父、圣子、圣神以外的任何其他本性能够不依赖身体而存在，那么逻辑推理的必然性迫使我们去理解：理性本性确实在起初被造，但物质实体只是在思想与理解中才与它们分离，并且似乎是为它们或以后才被形成的，而且它们从未、也并非不依赖它而生活；因为无形体的生活理当被视为唯独圣三的特权。如我们以上所述，这个世界的物质本性，具有接受一切转化的本性，当它被拖向较低等级的存在时，被塑造为较粗糙、较坚固的身体状态，以区分那些可见而多样的世界形式；但当它成为更完美、更幸福存在者的仆人时，它便在天体光辉中闪耀，并以属神的身体的衣饰装饰天主的众天使或复活之子；由此将充满那一个世界的多样而多变的状态。但如果有人愿意更充分地讨论这些事，就有必要以所有的敬畏和敬畏天主的心，更专注、更勤奋地查考\WorkTitle{圣经}，以察看在其中的隐秘和隐藏意义是否可能揭示关于这些事的任何东西；并且，在收集了许多关于这一点的见证之后，通过圣神向那些配得的人显示，或许能在其深奥神秘的语言中发现某些东西。\par

% structure: p0012 source=h2 target=subsection reason=relative-heading-rank; base=h2

\subsection{第三章 论世界的起始及其原因}

下一个探究的主题是：在现今存在的世界之前是否还有另一个世界；如果有，是否如现今这个一样，或是有些不同或更差；或者根本没有世界，而是像我们所理解的那样，在一切终结之后，当国度被交于天主父的时候，这却可能是另一个世界的终结——即这个世界在其后开始的终结；并且，各种理智本性的堕落是否促使天主产生了这个多样而多变的世界状况。我认为这一点也必须以类似的方式加以探究，即在这个世界之后，是否会有某种保存和修正的体系，确实是严厉的，并且给那些不愿服从天主的话的人带来许多痛苦，但这是一个过程，通过教导和理性训练，那些在今世致力于这些追求的人，在心灵被净化之后，得以进步而能够获得神圣智慧；在这之后，所有事物的终结将立即到来，并且为了那些需要它的人的纠正和进步，将再次出现另一个世界，要么与现今这个相似，要么比它更好，或大为逊色；并且那个世界，无论它是什么，在之后会持续多久；是否会有一个没有世界存在的时间，或者是否曾经有一个根本没有世界的时间；或者是否曾经有过或将来会有多个世界；或者是否会发生一个与另一个完全相似、无法区分的情况。\par

2. 为使此事更清楚起见，即形体性的物质是否能在时间间隔中存在，以及正如它在受造之前并不存在，它是否也能再次归于非存在，让我们首先看看，是否有人可以没有身体而生活。因为如果一个人可以没有身体而生活，那么一切事物也可以不要身体；我们在先前的论述中已经表明，万物都趋向同一终向。\newline{}现在，如果一切事物都可以没有身体而存在，那么无疑将没有任何形体性的实体，因为将不再需要它。但我们该如何理解宗徒在那些段落中的话呢？他在讨论死人复活时说：\BibleQuote{这可朽的必须穿上不朽，这必死的必须穿上不死。当这可朽的穿上了不朽，这必死的穿上了不死时，那所记载的话就应验了：\InnerQuote{死亡被胜利吞灭了！死亡啊，你的胜利在哪里？死亡啊，你的刺已被吞灭：死亡的刺就是罪，而罪的力量就是法律。}}\newline{}那么，宗徒似乎暗示了某种这样的意义。因为他所用的措辞，\Quote{这可朽的}和\Quote{这必死的}，仿佛带着触摸或指出什么的手势，除了指形体性的物质，还能指别的什么呢？这身体的物质，现在虽是可朽的，将来要穿上不朽，当一个完善的灵魂，一个带有不朽标记的灵魂开始居住在其中时。\newline{}如果我们说完善的灵魂是身体的衣服（因天主的圣言和祂的智慧，现在被称为不朽），你不要感到惊奇，因为耶稣基督自己，祂是灵魂的主和创造者，也被称为圣人们的衣服，正如宗徒所说：\BibleQuote{穿上主耶稣基督。}因此，正如基督是灵魂的衣服，同样，灵魂被称为身体的衣服，其理由也相当明白，因为灵魂是身体的装饰，遮盖并隐藏其必死的本性。\newline{}那么，\Quote{这可朽的必须穿上不朽}这句话，就好像宗徒说过：\Quote{身体这可朽的本性必须接受不朽的衣服——一个本身含有不朽性的灵魂，}因为它已穿上了基督，祂是天主的智慧和圣言。但当这身体（我们将来要在更光荣的状态下拥有它）成为生命的分享者时，它除了不死之外，也将成为不朽的。因为凡是必死的也必然是可朽的；但可朽的却不一定也是必死的。我们说一块石头或一块木头是可朽的，但我们不说它因此也是必死的。然而，由于身体分有生命，因为生命可以从它分离且确实分离，我们因此称它为必死的，并且按另一种意义也称它为可朽的。\newline{}因此，圣宗徒以非凡的洞察力，论及形体性物质的一般第一因（灵魂不断使用这种物质，不论它被赋予何种品质——现在确实是属血肉的，但将来更精微、更纯洁，被称为属神的），他说：\BibleQuote{这可朽的必须穿上不朽。}其次，论及身体的特殊原因，他说：\BibleQuote{这必死的必须穿上不死。}现在，不朽和不死除了是天主的智慧、圣言和正义——这些塑造、遮盖并装饰灵魂——之外，还会是什么呢？因此才有这话：\BibleQuote{这可朽的将穿上不朽，这必死的将穿上不死。}因为虽然我们现在可能大有进步，但我们仍是按部分认识，按部分说预言，如同在镜子里观看，模糊不清，那些我们似乎理解的东西，这可朽的尚未穿上不朽，这必死的也尚未穿上不死；而我们在身体中的这种训练无疑要延长到更长的时期，直到我们所披戴的这些身体本身，因天主的圣言、祂的智慧和完美的正义，而获得不朽和不死，因此才说：\BibleQuote{这可朽的必须穿上不朽，这必死的必须穿上不死。}\par

3. 然而，那些认为理性受造物可以在任何时候脱离身体而存在的人，可能会提出以下问题。如果这可朽坏的将穿上不朽，这必死的将穿上不死，而死亡在末后被吞灭，这就表明，唯有物质本性——死亡能在其上运作——将被毁灭，而身在身体中的人其心智的敏锐似乎被有形物质的本质所迟钝。但如果他们脱离身体，他们将完全摆脱这类扰乱带来的烦恼。但由于他们不能立即摆脱一切身体的包裹，他们应被视为居住在更精致、更纯洁的身体中，这些身体具有不再被死亡胜过、不再被它的刺所伤的特性；这样，最终随着物质本性的逐渐消失，死亡被吞灭，甚至在末了被消灭，其一切刺也因灵魂所接受的天主恩宠而完全迟钝，这恩宠使灵魂得以获取不朽和不死。于是所有人都有理由说：\BibleQuote{死亡，你的胜利在哪里？死亡，你的刺在哪里？死亡的刺就是罪过。} 如果这些结论似乎成立，那么我们必须相信我们未来的状态是无形的；而如果这一点被承认，且所有人都说屈服于基督，那么（无形状态）也必须给予所有屈服于基督的人；因为所有屈服于基督的人最终都将屈服于天主圣父，基督据说要将王权交给祂；因此看来那时身体的需要也将终止。如果终止，身体物质将归于无有，正如它先前并不存在一样。\par

现在让我们看看，对于作出这些断言的人，有什么可以说的。因为似乎必然的结果是，如果身体本性被消灭，它必须再次被恢复和创造；因为理性本性——自由意志的能力从未被夺走——似乎有可能在某些运动的影响下，通过主自身的特殊行动，再次进入某种运动，免得它们若是始终占据不变的状态，便不明白是藉着天主的恩宠而非自己的功劳而被安置在那最后的幸福境界中；这些运动无疑又将伴随着各种不同的身体，世界总是因这些身体而得以装饰；世界除了由各种不同事物构成之外，永远不会由其他任何东西构成——这种效果没有物质身体是无法产生的。\par

4. 现在我不明白，那些断言有时会出现与世界完全相似、毫无差异的另一个世界的人，究竟有什么证据可以维持他们的立场。因为若说有与世界完全相似的世界，那么亚当和厄娃将再次做他们先前所做的事：第二次将有同样的洪水，同样的梅瑟将再次带领近六十万人的民族离开埃及；犹达斯也将第二次出卖主；保禄将再次看守那些用石头砸死斯德望的人的衣服；此生所做的一切都将重复——我认为，如果灵魂受自由意志的推动，并根据其意志的力量保持前进或后退，这种情形无法用任何推理建立。因为灵魂并非被驱动着在一种循环中，经过许多世代后又回到同一轮次，从而行事或渴望这样或那样；无论他们自由意志的目标指向哪里，他们就会将行动引向那里。因为这些人所说的，与一个人断言如果将一\OriginalTerm{麦斗}{medimnus}谷粒倒在地面，第二次谷粒的落点与第一次完全相同，以致每一粒谷粒第二次都紧挨着先前抛下的位置，并且一\OriginalTerm{麦斗}{medimnus}以同样的次序和同样的标记散开——这当然是不可能的，即使连续倾倒无数年代也做不到。因此，在我看来，世界第二次以同样的次序、同样数量的出生、死亡和行为恢复是不可能的；但可能存在各种不同的世界，带有并非无关紧要的变化，以致另一个世界的状态可能由于某些明确的原因比这个世界更好，或更差，或再次处于中间状态。但关于这数目或尺度是什么，我承认自己无知，尽管若有人能告诉我，我乐意学习。\par

5. 然而，这个被称为时代的现今世界，据说是许多时代的终结。圣宗徒教导说，在现今之前的那个时代里，基督并未受苦，再往前的时代中也没有；我不知道自己能否列举出他未曾受苦的先前时代的数目。不过，我要根据保禄的某些陈述来说明我是如何领会这一点的。他说：\BibleQuote{可是如今，在时代的终结，他显现了一次，以自己为祭献，除去了罪过。}因为他说他一次被作为祭品，并在时代的终结显现以除去罪过。现在，关于这个现今的时代——据说它乃是为其他时代的终结而形成的——之后，还会有其他的时代接续而来，这从保禄本人那里我们已清楚得知，他说：\BibleQuote{为在未来的时代中，显示他在基督耶稣内对我们所施的恩宠的极丰富。}他没有说\Quote{在未来的时代}，也没有说\Quote{在未来的两个时代}，由此我推断，他的话语暗示了许多时代。如果还有什么比时代更伟大的东西，以至于在被造物中可理解为某些时代，但在那超越并高于可见受造物的其他存在物中（或许还有更伟大的时代），(这或许将在万物复原时发生，届时整个宇宙将达至完美的终点)，那么那个一切事物达至终结的时期，或许应被理解为某种比时代更伟大的东西。然而，此处圣经的权威打动了我，因为经上记载：\BibleQuote{一个时代甚至更多。}这个\Quote{更多}无疑意味着某种比时代更伟大的东西；试看救主的那句话：\BibleQuote{我愿我在哪里，这些人也同我在一起；就如我与你是一体，这些人也在我内合而为一。}这话岂不似乎传达了某种比一个时代和许多时代更伟大的东西，甚至可能比万世万代更伟大的东西——即当万物不再处于时代之中，而是天主在万有之内的时候。\par

6. 当我们尽己所能讨论了关于世界本质的这些要点之后，似乎不难探究\Quote{世界}一词的含义，它在圣经中常常显示出不同的意义。因为我们在拉丁文中称为\OriginalTerm{世界}{mundus}的，在希腊文中被称为\OriginalTerm{世界}{\GreekText{κόσμος}}，而\OriginalTerm{世界}{\GreekText{κόσμος}}不仅指世界，也指装饰。最后，在依撒意亚书中，当责备熙雍女子们的首领时，他说：\BibleQuote{你们将有秃头代替金头饰，因了你们的作为。}他用同一个词表示装饰和世界，即\OriginalTerm{世界}{\GreekText{κόσμος}}。因为大司祭的服装据说包含了世界的结构，如我们在\BibleBook{智慧}{篇}中读到：\BibleQuote{因为长袍上绣着整个宇宙。}我们这片土地及其居民也被称为世界，如圣经所说：\BibleQuote{整个世界都居于邪恶者手下。}然而，宗徒的门徒\Person{克莱孟}提到了那些被希腊人称为\OriginalTerm{对跖者}{\GreekText{᾿Αντίχθονες}}的人以及地球的其他部分，我们中无人能接近那里，那里的人也无人能来到我们这里，他也称这些为世界，说：\Quote{海洋对人类是不可逾越的；而在海洋的另一边，是那些世界，它们由统治天主的同样安排所治理。}那被天地所限的宇宙也称为世界，正如\Person{保禄}所宣告的：\BibleQuote{因为这世界的局面正在逝去。}我们的主和救主还指出了另一个不同于这个可见世界的世界，那确实难以描述和认识。他说：\BibleQuote{我不属于这世界。}因为，仿佛他属于另一个世界，他说：\BibleQuote{我不属于这世界。}关于这个世界，我们先前已经说过，解释是困难的；原因在于，免得给任何人提供机会，以为我们认为存在某些希腊人称为\Quote{理念}的图像；因为谈论一个仅仅存在于想象中或转瞬即逝的思想世界中的无形世界，确实不符合我们的著作；而他们如何声称救主来自那里，或者圣人将去那里，我不得而知。然而，毫无疑问，救主所指出的某些事物比现今这个世界更辉煌、更卓越，他激励并鼓励信徒以此为志向。但救主所指的那个世界是否因位置、本性或荣耀而与此世界截然分离，还是它在荣耀和品质上更优越，却仍限于这个世界的界限内（我认为后者更有可能），仍不确定，在我看来也不适合人类思考。但从\Person{克莱孟}似乎暗示的话来看，他说：\Quote{海洋对人类是不可逾越的，而在它后面的那些世界，}他用复数提及后面的世界，并暗示它们由至高天主的同一眷顾所管理和统治，他似乎为我们提供了某种观点的萌芽，即整个存在的事物——天上的、超天上的、地上的和地下的——通常被称为一个完美的世界，在这个世界之内或由其包含，如果有其他世界的话，也必须被设想为包含在里面。因此，他希望将太阳、月亮以及称为行星的其他天体的球体各自称为世界。不，即使是他们称为\OriginalTerm{不动天}{\GreekText{ἀπλανῆ}}的那个卓越球体，他们也希望恰当地称之为世界。最后，他们引用\BibleBook{巴路克}{先知书}来支持这一主张，因为书中更清楚地指出了七个世界或七个天。然而，在称为\OriginalTerm{不动天}{\GreekText{ἀπλανῆ}}的那个球体之上，他们将存在另一个球体，他们说，正如我们的天包含其下所有事物一样，它以其巨大的规模和难以描述的广度，将各球体的空间一同包含在其更为宏伟的圆周之内；因此所有事物都在它里面，如同我们这片土地在天下。这也被认为在圣经中被称为福地、活人之地，它有自己的更高之天，并且救主说圣人的名字被写在或已写在其上；由那更高的天所包围和封闭的这块土地，是救主在福音中应许给温良和怜悯的人的。因为他们认为我们这片从前被称为\Quote{干地}的土地，其名称来源于那块土地的名称，正如这片天也因那更高之天的名称而被称为穹苍一样。但我们已经在必须探究\BibleQuote{起初天主创造了天地}这一宣告的含义的地方，更详细地讨论了这些意见。因为除了第二天之后所造的那个\Quote{穹苍}或后来被称为\Quote{地}的那片\Quote{干地}之外，还显示了另一天和另一地。当然，有些人关于这个世界所说的话——它因受造而有朽坏性，但并未朽坏，因为造它并持守它以免朽坏掌控它的天主的旨意比朽坏更强而有力——可以更正确地应用于我们称之为\Quote{不动天}的球体，因为它由于天主的旨意而完全不受朽坏的影响，原因在于它没有接纳任何朽坏的原因，因为它是圣人和完全洁净者的世界，而不是像我们这个世界这样的恶人的世界。此外，我们必须看看，也许\Person{保禄}所说的正是这一点：\BibleQuote{我们并不注目那看得见的，而是那看不见的；因为看得见的是暂时的，看不见的才是永恒的。因为我们知道，如果我们地上帐棚式的住所被拆毁了，我们必由天主获得一所房舍，一所非人手所造的，永远在天上的住所。}当他在另一处说：\BibleQuote{因为我看见你手指所造的天，}并且天主藉先知的口论及一切可见的事物说：\BibleQuote{我的手创造了这一切，}他宣称那非人手所造的、永远在天上的住所是应许给圣人的，无疑是指出了可见与不可见之受造物的区别。因为\Quote{那些看不见的}和\Quote{那些不可见的}两种表达不应被理解为一回事。因为不可见的事物不仅不被看见，而且不具备可见的性质，即希腊人称为\OriginalTerm{无形体}{\GreekText{ἀσώματα}}的，也就是无形体的；而\Person{保禄}所说的\Quote{那些看不见的}确实具有被看见的性质，正如他所解释的，但由于尚未被应许给他们的人所看见，所以未被看见。\par

7. 既然我们已经尽己所能，概要地勾勒了关于万物的终结和最高福乐这三种见解，那么就让每一位读者自己细心勤奋地判断，其中哪一种可以被认可和采纳。因为有人说过，我们必须设想：要么，在万物都归于基督、并通过基督归于天主圣父之后，天主成为万物之中的万有时，一种无形体的存在是可能的；要么，尽管万物都已归于基督、并通过基督归于天主（就精神体是理性本性而言，他们与天主也成了一个精神），此时形体本身也与最纯净、最卓越的精神体结合，并根据接受者的品级或功绩改变为以太状态（如宗徒所说：\BibleQuote{我们也要改变}），而发出光辉；要么，至少当可见之物的形态消逝，一切朽坏都被抖落并清除，当这个世界所占据的整个空间——据说行星的球体就在其中——被留在下方和身后时，就到达了虔诚良善者固定不变的居所，它位于那被称为\Quote{不动的}\OriginalTerm{不动的}{\GreekText{ἀπλανής}}球体之上，如同在一片美地，在生命之地，这地将由温良的人承受；这地所属的那个天（以其更宏伟的广度环绕并包含这地本身）才真正并首要地被称为天，万物的终结与圆满可以稳妥而极其确实地置于其中——也就是说，这些事物经过被逮捕和因其罪过而受的净化惩戒，在履行并清偿一切义务之后，堪当在那地中获得居所；而那些顺从了天主的话、并从此因服从而显出配得上智慧的人，则被说成堪当那重天或诸天的国度；于是，预言得到了更相称的应验：\BibleQuote{温良的人是有福的，因为他们要承受土地}；\BibleQuote{神贫的人是有福的，因为天国是他们的}；以及圣咏中的宣告：\BibleQuote{他必高举你，你必承受土地}。因为下降到这地被称作\Quote{下降}，而上升到那高处则被称作\Quote{高举}。因此，通过圣人从土地到诸天的离去，似乎开辟了一条道路；以致他们与其说是安居在那地，不如说是带着意向居住在那里，即为了在也到达那完美境界时，继承天国。\par

% structure: p0021 source=h2 target=subsection reason=relative-heading-rank; base=h2

\subsection{第四章 法律与先知的天主，和我们的主耶稣基督的父，是同一的天主。}

1. 既然我们已尽力按序简短地整理了这些要点，接下来，按照我们起初的意图，要驳斥那些认为我们的主耶稣基督之父，与赐予梅瑟法律答复、或派遣先知、又是我们祖先亚巴郎、依撒格和雅各伯的天主，是不同的神的人。因为在这信德条文上，我们首先必须坚定立场。然后，我们必须思考福音中屡次出现、并且附于我们主和救主所有行动之后的那句话：\BibleQuote{这应验了某位先知所说的话}，显然，这些先知是那创造世界之天主的先知。因此我们得出结论：派遣先知的那位自己预言了关于基督所要预告的一切。毫无疑问，是圣父自己，而非另一位不同于祂的神，说出了这些预言。此外，救主或祂的宗徒经常引用旧约中的例证，表明他们承认古人的权威。救主在劝勉祂的门徒行善时也命令说：\BibleQuote{你们应当是成全的，如同你们的天父是成全的一样；因为祂使太阳升起，光照恶人，也光照善人；降雨给义人，也给不义的人。}即便是悟性薄弱之人也最明显地看出，祂在向门徒提出的可效法的对象，无非是那创造天并赐下雨水的天主。再者，祈祷者应当使用的这句话：\BibleQuote{我们的天父}，除了表示应当在世界的更美好部分、亦即祂的受造物中寻求天主之外，还能有何含义呢？此外，祂关于起誓所定下的卓越原则，说我们不应该\BibleQuote{指着天起誓，因为天是天主的宝座；也不可指着地起誓，因为地是祂的脚凳}，这与先知的话\BibleQuote{天是我的宝座，地是我的脚凳}最清晰地吻合吗？并且当祂将那些卖羊、牛、鸽子的从圣殿赶出去，并倒出兑换银钱者的桌子，且说：\BibleQuote{把这些东西拿出去！不要使我父的殿宇成为商场}，祂无疑称祂为祂的父，撒罗满曾为祂的名建造了一座宏伟的圣殿。此外，\BibleQuote{你们没有读过天主向梅瑟所说的话吗？我是亚巴郎的天主，依撒格的天主，雅各伯的天主；祂不是死人的天主，而是活人的天主}这句话，最清楚地教导我们：祂称圣祖（因为他们圣洁且活着）的天主为活人的天主，正是那曾在先知中说\BibleQuote{我是天主，除我以外没有别的神}的同一位。因为如果救主知道法律上所写的天主是亚巴郎的天主，并且就是那位说\BibleQuote{我是天主，除我以外没有别的神}的同一位，而承认那位（按异端者的想法）不知道在祂之上还有别的神的存在者，正是祂的父，那么祂就荒谬地宣称一位不认识更大之神的为祂的父。但如果祂说\BibleQuote{除我以外没有别的神}不是出于无知，而是出于欺骗，那么承认祂的父犯了虚言罪，就更加荒谬了。由此得出的结论是：祂知道除了万物的建立者和创造者天主之外，没有别的父。\par

2. 若要从福音中各处段落收集证明，以显示法律与福音的天主是同一的，未免过于繁琐。我们不妨略略提及《WorkTitle{宗徒大事录}》，其中\Person{斯德望}及其他宗徒向那创造天地、并藉祂圣先知的口发言的天主献上祈祷，称祂为\BibleQuote{亚巴郎的天主、依撒格的天主、雅各伯的天主}，是那\BibleQuote{带领祂的子民从埃及地出来的}天主。这些说法无疑清晰地引导我们的理智去相信造物主，并在那些虔敬而忠信地如此思想祂的人心中，培育对祂的爱慕；这正符合救主亲口所说的，当有人问祂法律中哪条诫命最大时，祂回答说：\BibleQuote{你应全心、全灵、全意，爱上主你的天主。第二条是：你应当爱近人如你自己。}又加上说：\BibleQuote{全部法律和先知，都系于这两条诫命。}那么，祂向那正在受祂教导、并即将进入门徒职务的人，为何特别推崇这条诫命，超越其他一切诫命？这条诫命无疑是要在他心中激发对那法律所宣告的天主的爱，因为法律本身就是以这些话语来表达的。但即便有所有这些最明显的证据，我们仍可以假定救主所说的\BibleQuote{你应全心、全灵、全意，爱上主你的天主}等语是指某个别的、不认识的天主。那么，如果法律和先知（如他们所说）来自造物主，即来自不同于祂所称为\Quote{善}的那一位的天主，那么祂随后加上的\BibleQuote{全部法律和先知，都系于这两条诫命}这句话又怎能合乎逻辑呢？因为那与天主疏远且无关的事物，又怎能依赖祂呢？当\Person{保禄}说：\BibleQuote{我感谢我所侍奉的天主，就是从我祖先起，以纯洁的良心在天主前所侍奉的}，他清楚表明他所皈依的并非某个新的神，而是基督。因为\Person{保禄}的其他祖先还能指谁呢？难道不是那些他所说\BibleQuote{他们是希伯来人吗？我也是。他们是以色列人吗？我也是}的祖先吗？不仅如此，他写给罗马人的书信的序言，对于懂得理解\Person{保禄}书信的人，不是清楚地显明了同样的道理，即他所宣讲的是哪一位天主吗？因为他说：\BibleQuote{保禄，耶稣基督的仆人，蒙召作宗徒，被选拔去传天主的福音，这福音是天主藉自己的先知在圣经上，所预先应许的，论及祂的儿子，按肉身说，是从达味后裔生的，按圣德的神力，藉着从死者中的复活，被立为具有大能的天主子，我们的主耶稣基督}等等。此外，还有下面这段：\BibleQuote{牛在打谷的时候，不可笼住它的嘴。天主岂是关心牛吗？还是全然为我们说这话？的确，这话是为我们写的，因为耕田的应当怀着希望去耕，打谷的也应当怀着有份的希望去打谷。}由此他清楚表明，那为我们（即为了宗徒）颁布法律的天主说：\BibleQuote{牛在打谷的时候，不可笼住它的嘴。}祂所关心的不是牛，而是宣讲基督福音的宗徒。在其他段落中，\Person{保禄}也接受法律的应许，说：\BibleQuote{孝敬你的父亲和你的母亲，这是一条附有应许的诫命，为使你得到幸福，并在天主你的主所赐给你的地上，得享长寿。}由此他无疑表明，法律、法律的天主以及祂的应许，都令他喜悦。\par

3. 但是，那些拥护此异端的人，有时惯于用某些欺骗性的诡辩误导简单之人的心灵，因此，我认为提出他们惯常作出的断言并驳斥其欺骗和谬误并非不当。以下是他们的宣称。经上记载：\BibleQuote{从来没有人见过天主。}但是，梅瑟所宣讲的那位天主，却被梅瑟本人及他之前的祖先看见；而救主所宣告的那位，却从未被任何人看见。因此，让我们问他们和我们自己：他们承认其为天主并声称与造物主不同的那位，是可见的还是不可见的？如果他们说是可见的，那么除了被证明违背圣经的宣告（圣经论到救主说：\BibleQuote{祂是不可见的天主的肖像，是一切受造物的首生者}）之外，他们也将陷入断言天主是有形体的荒谬之中。因为除非借助形状、大小和颜色（这些是形体的特殊属性），没有东西能被看见。如果天主被宣布为形体，那么祂也将被发现是由物质组成的，因为每一个形体都由物质构成。但如果祂由物质构成，而物质无疑是可朽坏的，那么按照他们的说法，天主就易于朽坏！我们将向他们提出第二个问题：物质是被造的，还是非受造的（即不是被造的）？如果他们回答说是非受造的（即非受造的），我们将问他们是否一部分物质是天主，而另一部分是世界？但如果他们说是被造的，那么毫无疑问，他们承认他们所宣称的天主是被造的！——这个结果肯定既不是他们的理性也不是我们的理性所能接受的。但他们会说：天主是不可见的。那么你将怎么办？如果你说祂在本性上是不可见的，那么祂也不该对救主是可见的。相反，基督的天主圣父却被说成是被看见的，因为祂说：\BibleQuote{谁看见圣子，就看见圣父。}这当然会给我们很大压力，除非我们更正确地理解为认识，而不是看见。因为谁理解了圣子，也就会理解圣父。这样，梅瑟也必须被认为看见了天主，不是用肉眼看见祂，而是用内心的视觉和理智的感知理解祂，而且只是在一定程度上。因为很明显，那位回答梅瑟的祂说：\BibleQuote{你不能看见我的面容，但能看见我的背面。}当然，这些话应以适合神圣话语的神秘意义来理解，要摒弃和蔑视那些无知之人所编造的有关天主的正面和背面的老妇之谈。确实，不要让任何人以为我们说了不敬的话，说甚至对救主而言圣父也是不可见的。请思考我们在对付异端者时所使用的区分。因为我们已解释过，看见和被看见是一回事，认识和被认识或理解和被理解是另一回事。因此，看见和被看见是形体的属性，它当然不适用于圣父、圣子或圣神之间的相互关系。因为天主圣三的本性超越视觉的尺度，而赋予在形体中者（即所有其他受造物）彼此看见的属性。但是，对于无形体且主要是理智的本性，除了认识或被认识之外，没有其他适当的属性，正如救主亲自宣告说：\BibleQuote{除了父，没有人认识子；除了子和子所愿意启示的人，也没有人认识父。}那么很明显，祂没有说：\BibleQuote{除了子，没有人看见父}，而是说：\BibleQuote{除了子，没有人认识父。}\par

4. 现在，如果因为旧约中出现的那些语句，如当天主被说成愤怒或后悔，或当任何其他人类情感或情欲被描述时，（我们的对手）认为他们有了反驳我们的根据（我们主张天主是完全不受苦难的，应被视为完全摆脱此类情感），那么我们必须向他们表明，类似的陈述甚至在福音的比喻中也能找到；例如，当说那栽种了一个葡萄园、租给佃农的人，佃农杀死了派去的仆人，最后甚至杀死了儿子，那人在愤怒中夺走了他们的葡萄园，并将邪恶的佃农交与毁灭，又将葡萄园转给了别人，那些人会在季节里按时交出果实。同样，关于那些市民，当家长出发去为自己取得一个国家时，他们派使者跟在他后面说：\BibleQuote{我们不愿意这人统治我们；}因为家长获得了国家后回来，愤怒地命令他们在自己面前被处死，并放火烧了他们的城。但是，无论我们在旧约或新约中读到天主的愤怒，我们都不照字面理解这些语句，而是在其中寻求属灵的意义，以便我们按天主应被思考的方式来思考祂。关于这些点，在解释第二篇圣咏的经文\BibleQuote{那时祂要在怒中对他们说话，在烈怒中惊扰他们}时，我们已经尽我们微薄的能力表明了这样的语句应如何理解。\par

% structure: p0026 source=h2 target=subsection reason=relative-heading-rank; base=h2

\subsection{第五章 论正义与良善}

1. 如今，既然这一看法对某些人有分量——即我们所谈论的那异端的首领们自认为已建立了一种区分，据此他们宣称正义是一回事，良善是另一回事，并将这一区分甚至应用于神圣事物，声称我们的主耶稣基督的圣父确实是一位良善的天主，但并非正义的天主，而法律和先知的天主是正义的，但并非良善的；——我认为有必要尽可能简短地回应这些说法。这些人，然后，认为良善是某种如此的情感，以至于将恩惠施予所有人，即使接受者不配或不值得任何恩待；但在我看来，他们在这一点上没有正确应用他们的定义，因为他们认为，对于遭受任何痛苦或灾祸的人，并没有施予恩惠。另一方面，他们将正义视为一种按其应得报偿每个人的品质。但这里，他们又没有正确解释自己定义的意义。因为他们认为，将恶事降于恶人、将恩惠降于善人是正义的；也就是说，按照他们的观点，正义的天主似乎并不愿善待恶人，反而对他们怀有一种仇恨。他们收集这类例子，每当在旧约圣经中找到一段历史，例如，关于洪水的惩罚，或其中所描述的那些人的命运，或索多玛和哈摩辣被火与硫磺雨毁灭，或全体百姓因他们的罪而倒毙在旷野，以致离开埃及的人中没有一个进入应许之地，除了若苏厄和加肋布。而从新约中，他们收集了怜悯和虔诚的言语，救主通过这些言语训练门徒，并似乎藉此宣告，除了天主圣父之外，没有一个是良善的；他们因而敢于称救主耶稣基督的圣父为良善的天主，而说\OriginalTerm{世界之神}{God of the world}是另一位，他们乐意称他为正义的，但并不也是良善的。\par

2. 首先，我认为他们必须表明——若是可能，且符合他们自己的定义——造物主在惩罚中乃是公义的，祂按应得惩罚：或是在洪水时代灭亡的人，或是索多玛的居民，或是离开埃及的人。因为我们有时看到有人犯下比上述被毁灭之人更邪恶、更可憎的罪行，却并未看到每个罪人都为自己的恶行付出代价。他们是否会说，那位一度公义的天主后来变成了善？或者他们更认为，祂如今仍是公义的，只是忍耐着人类的过犯，而那时祂甚至不是公义的，因为祂将无辜的婴孩与残暴不虔的巨人一同灭绝？这些就是他们的见解，因为他们不知道如何领会超越字面的事物；否则他们会指出，将罪罚临到子孙头上直到三、四代，以及后裔的子女身上，这按字面如何是公义的。然而，我们却不按字面理解这些事；而是如同厄则克耳讲述比喻时所教导的，我们探究比喻本身所包含的内在含义。此外，他们也应当解释这一点：那位惩罚属世之人和魔鬼的天主，如何是公义的，并按各人的功德报答各人？因为按照他们的说法，如果这些受罚者本性邪恶败坏，他们本不能行任何善事。既然他们称祂为审判者，祂似乎不是审判行为，而是审判本性；如果败坏的本性不能行善，善良的本性也不能作恶。其次，如果他们称为善的那一位对待所有人都是善的，那么祂无疑也善对待那些注定要灭亡的人。为何祂不拯救他们？如果祂不愿拯救，祂就不再是善的；如果祂愿意却无能为力，祂就不是全能的。为何他们不听从福音中我们的主耶稣基督的圣父为魔鬼及其天使预备了火？那惩罚性且令人悲哀的过程，按他们的观点，又怎会是善天主的作为呢？连救主本身，善天主的圣子，也在福音中抗议并宣告：\Quote{如果提洛和漆冬所行的奇迹曾在你们中间发生，他们早就披上苦衣，坐在灰中，悔改了。}而当祂接近那些城市并进入其境内时，请问，祂为何避开进入那些城市，并向他们显示大量的奇迹？如果确定他们见到奇迹后会披上苦衣、坐在灰中悔改，祂为何不做？但祂没有这样做，祂无疑舍弃了那些福音的语言表明并非邪恶或败坏本性的人去灭亡，因为福音宣称他们能够悔改。再者，在福音的某个比喻中，君王进来观看宴席上的宾客，看见一个人没有穿着婚宴礼服，就对他说：\Quote{朋友，你没有婚宴礼服，怎么到这里来了？}然后吩咐仆人：\Quote{捆起他的手和脚，把他丢在外面的黑暗里；那里要有哀号和切齿。}让他们告诉我们，那位进来观看宾客的君王是谁？当他发现一人穿着不洁的衣服时，便命令仆人捆住他，扔到外面的黑暗里。祂与他们所称为公义的那位是同一的吗？那么祂为何命令仆人好坏都请来，而不指示仆人查问他们的功德？这样的做法所指示的，并非如他们所言，是一位按人的功德施行赏罚的公义天主，而是一位对所有人显示无分别的善的天主。既然这事必须被理解为善天主——即基督或基督的圣父——那么他们还能提出什么反对天主审判的公义呢？还有什么比他们控告律法的天主更为不义的呢？即祂命令那位被祂仆人邀请（仆人奉差去呼召好人坏人）的人，因穿着不洁的衣服而被捆起手脚，扔到外面的黑暗里？\par

3. 如今，我们从圣经权威所汲取的，应当足以驳斥异端者的论点。然而，如果我们根据理性本身与他们稍作讨论，似乎也并非不合宜。那么，我们问他们，是否知道在人看来什么是德行与邪恶的基础，以及是否看来我们可以谈论天主中的德行，或者如他们所认为的，在这两位神中的德行。也请他们回答这个问题：他们是否认为良善是一种德行；既然他们无疑会承认如此，那么他们对不义会说什么？他们绝不可能，在我看来，愚蠢到否认正义是一种德行的地步。因此，如果德行是一种福，而正义是一种德行，那么毫无疑问，正义就是良善。但若他们说正义不是一种福，那它必是恶或是无关紧要之事。现在我认为回答那些说正义是恶的人乃是愚昧，因为我将显得要么在回答无意义的话，要么在回答心智失常的人。那能赏善人以福的事，如何能显为恶呢？正如他们自己也都承认的？但若他们说那是无关紧要之事，那么既然正义如此，节制、明智及其他一切德行也都是无关紧要之事。对保禄所说的，我们又当如何作答呢？他说：\BibleQuote{若有什么德行，若有什么可称赞的，你们就该思念这些事，就是你们在我身上所学习的、所接受的、所听见的、所看见的？} 因此，让他们借着查考圣经，认识每一种德行，不要自欺，说那按各人的功过赏报人的天主，是因憎恨邪恶而以恶报恶，而不是因为犯罪者需要更严厉的医治，且因祂对他们施以那些有望改善，但暂时似乎带来痛苦感受的措施。他们不读那有关在洪水中被毁灭之人的希望的记载；关于这希望，伯多禄自己在他的第一封书信中这样说：\BibleQuote{基督，固然，在肉身内被杀，但在圣神内却得以生活；祂藉着圣神，曾去给那些在监狱中被拘禁的、不信从的灵魂宣讲，他们在诺厄的时代，当方舟被建造时，曾等待天主的容忍，其中少数，即八个灵魂，藉着水而得救。这水所预表的圣洗，如今也拯救你们。} 至于索多玛和哈摩辣，让他们告诉我们，他们是否相信先知的话出自那位创造主天主——就是那位据说曾降下火雨与硫磺毁灭它们的天主。厄则克耳先知论到它们说了什么？\BibleQuote{索多玛，必恢复原先的状况。} 然而，为什么在惩罚那些应受罚的人时，祂不是为他们的益处而惩罚他们呢？——祂也对加色丁说：\BibleQuote{你有火炭，坐上去；它们将成为你的帮助。} 此外，关于那些在旷野中倒下的人，让他们听听第七十八篇圣咏的记载，其标题属阿撒夫；因为他说：\BibleQuote{当祂击杀他们时，他们才寻求祂。} 这一切证明，法律和福音的天主是同一的、唯一的天主，祂是既正义又良善的天主，祂公道地施恩，慈爱地惩罚；因为没有正义的良善，或没有良善的正义，都不能显示天主性（真正的）尊严。\newline{}\par

我们将添加以下评论，是因他们的诡辩而被迫作出的。如果正义与良善是不同的事，那么，既然恶是善的对立面，不义是正义的对立面，不义毫无疑问将是与恶不同的东西；并且，按你们的意见，既然义人不是善人，那么不义的人也不是恶人；同样，既然善人不是义人，那么恶人也不是不义的人。然而，谁看不出这荒谬呢：一位良善的天主，有一个邪恶者与之对立；而一位正义的天主（他们说比良善者低劣），却无人与之对立！因为没有一个可称为不义者，正如有一个被称为邪恶的撒殚。那么，我们该怎么办呢？让我们放弃我们所维护的立场吧，因为他们将无法坚持一个恶人不是不义的人，一个不义的人不是恶人。如果这些品质不可分割地存在于这些对立面中，即不义在邪恶中，或邪恶在不义中，那么毫无疑问，善人将与义人不可分离，义人与善人也不可分离；因此，正如我们在恶意与不义中说同一邪恶，我们也可以认为良善与正义的德行是同一回事。\newline{}\par

4. 然而，他们又引用圣经的话，提出他们那个著名的问题，肯定说经上记载：\BibleQuote{坏树不能结出好果子，因为树可由其果实认出。}那么，他们的立场是什么？法律是何种树，由其果实，即其诫命的言辞，得以显明。因为如果法律被发现有善，那么毫无疑问，颁布法律的那一位就被相信为善的天主。但如果它是公义的而非善的，那么天主也将被视为公义的立法者。宗徒保禄说话毫不迂回，他说：\BibleQuote{法律是善的；诫命是圣的、公义的和善的。}由此清楚可见，保禄并未学习那些将公义与善分离的人的语言，而是受那同时是圣、善、公义的天主教导，并由祂的圣神光照；他藉同一圣神发言，宣称法律的诫命是圣的、公义的、善的。为更清楚地表明善在诫命中比公义和圣德更为优越，他在重复自己话时，不用那三个形容词，而独用\Quote{善}一词，说：\BibleQuote{那么，善事成了我的死亡吗？绝对不可！}因为他知道善是德行的\OriginalTerm{类}{genus}，而公义和圣德是属于该类的\OriginalTerm{种}{species}；在前面的经节中他合称\OriginalTerm{类}{genus}和\OriginalTerm{种}{species}，而在重复时则只强调\OriginalTerm{类}{genus}。但在接下来的经文里他说：\BibleQuote{罪恶藉着善在我内造成了死亡}，在此他以类的方式总结了他先前具体解释的内容。同样也应如此理解这句宣言：\BibleQuote{善人从他心内的善库里取出善物来，恶人从他恶库里取出恶物来。}因为在此他也假设了善与恶中存在一个\OriginalTerm{类}{genus}，无疑指出在善人身上有公义、节制、明智、虔诚以及一切可被称为或理解为善的事物。同样他也说，一个恶人必定是不公义的、不洁的、不圣的，以及一切单独构成恶人的事物。因为没有人会认为一个人没有这些邪恶的特征就是恶的（也确实不能如此），所以同样可以肯定，没有这些德行，没有人会被视为善的。然而，他们仍保留了主在福音中的那句话，他们以为这是特别赐予他们的盾牌，即：\BibleQuote{除了唯一的天主圣父，没有一个是善的。}他们声称这句话是专指基督的父，而这父不同于创造万有的天主，他没有给那创造者以善的称谓。现在让我们看看在旧约中，先知们的天主，即言语的创造者和立法者，是否不被称为善。在圣咏中出现了哪些说法？\BibleQuote{天主对以色列，对心地正直的人，何其美善！}以及\BibleQuote{愿以色列现在都说他是美善的，他的仁慈永远常存。}耶肋米亚哀歌中的话：\BibleQuote{上主对等候他的人，对寻求他的灵魂，是美善的。}因此，正如天主在旧约中常被称为善，同样，我们的主耶稣基督的父在福音中被称作公义。最后，在若望福音中，我们主亲自在祈祷父时说：\BibleQuote{公义的父啊！世界没有认识你。}恐怕他们会说，这是由于祂取了人身，才称世界的创造者为\Quote{父}，并称祂为\Quote{公义的}；但紧接着的话就堵住了他们的借口：\BibleQuote{世界没有认识你。}然而，按照他们的观点，世界只认识那唯一善的天主。因为世界无疑认识它的创造者，主亲自说，世界爱属于自己的事物。那么明确的是，他们视为善的那一位，在福音中被称作公义。任何人都可以闲暇时收集更多证据，包括新约中我们的主耶稣基督的父被称为公义的段落，以及旧约中天地的创造者被称为善的段落；这样异端者被众多见证所定罪，或许终有一时会羞愧。\par

% structure: p0032 source=h2 target=subsection reason=relative-heading-rank; base=h2

\subsection{第六章 论基督的降生成人。}

1. 在对此诸点略加说明之后，现在时候到了，我们要重新探讨我们的主和救主降生成人的问题，即祂如何或为何成为人。因此，我们已尽微弱之所能，借默观祂自身的工程而非凭自己的感受，思考了祂的天主性本质，并仍以肉眼看见了祂有形的受造物，同时以信德看见无形的受造物，因为人的软弱既不能用肉眼看见一切事物，也不能用理智完全理解——我们人类比其他任何理性受造物更为软弱和脆弱（因为那些在天上或被认为存在于天上之上的受造物更为优越）——，所以我们还须寻找一个介于一切受造物和天主之间的存有，即一位中保，宗徒保禄称之为\BibleQuote{一切受造物中的首生者}。再者，看到圣经中关于祂威严的声明，即祂被称为\BibleQuote{不可见的天主的肖像，一切受造物中的首生者}，并且\BibleQuote{在祂内万有都是受造的，天上的、地上的，可见的、不可见的，或是宝座，或是宰制，或是率领者，或是掌权者，一切都是借着祂、并在祂内受造的；祂先于万有，万有都依赖祂而存在}，祂是万有之首，唯独以天主父为首，因为经上记载：\BibleQuote{基督的头是天主}；又明明看见经上记载：\BibleQuote{除了圣子，没有人认识父；除了圣父，也没有人认识圣子}（因为谁能认识智慧是什么，除非是那把它召叫出来的那位？或者，谁能清楚理解真理是什么，除非是真理之父？谁能确切探究祂的圣言以及天主自身的普遍本性——这种本性出自天主——除非是那与圣言同在的独一真神？），我们应当确信：这个圣言或理性（如果可这样称呼）、这个智慧、这个真理，除了父以外，无人知晓；关于祂，经上记载：\BibleQuote{我以为世界本身也容纳不下所能写的书}，即关于天主圣子的光荣和威严的书。因为那些属于救主光荣的细节，是不可能全部写下来的。在思考了关于天主圣子存有的如此重要的问题之后，我们深深惊叹：如此卓越超群的本性，竟会舍弃祂威严的境况而成为人，并居住在人中间，正如祂嘴唇上所倾注的恩宠所证明的，以及祂的天父所见证的，并由祂所行的各种征兆、奇事和异能所承认的；祂在借身体显现之前，曾派遣先知作祂的前驱和祂来临的使者；在升天之后，祂使祂的圣宗徒——这些愚昧无知的人，从税吏或渔夫中被拣选出来，却充满了祂天主性的德能——走遍世界，以便从各种族、各民族中聚集一群虔诚信仰祂的人。\par

2. 但在祂所行的一切奇妙大能事迹中，这一件完全超越了人类的赞叹，超出了凡人脆弱理智的理解或感受：即那神圣威严的至高权能、那圣父的圣言、那天主的智慧——万有，可见的与不可见的，都是藉着它受造的——竟能被相信存在于那位在犹大显现的人的限度之内；不仅如此，天主的智慧竟能进入妇女的子宫，作为婴儿诞生，发出像小孩啼哭般的哀号！而后还记载说，祂在临终时极其忧伤，正如祂自己所说的：\BibleQuote{我的心灵忧伤至死}；最后，祂遭受了人类视为最耻辱的死亡，尽管祂在第三天复活了。因此，我们在祂身上看到一些如此人性的事，似乎与普通人的软弱毫无区别；又看到一些如此神圣的事，只能恰当地归于那原始而不可言说的天主性本质。人类理解的狭窄无法找到出路，反而被极度惊讶所压倒，不知该往何处退避、抓住什么或转向何方。若想到一位神，却看见一个会死的人；若想到一个人，却看见祂从坟墓中回来，推翻了死亡的权势，满载着战利品。因此，必须怀着极大的敬畏和虔诚来思量这一景象，以便在同一个存有身上清楚显示两种本性的真理；这样，在那神圣而不可言说的实体中，就不会看到任何不配或不体面的事，也不会以为所发生的事情是虚幻表象的错觉。用人的耳朵听这些话，用言语来解释它们，远远超出了我们地位、理智和语言的能力。我认为甚至圣宗徒的权能也不足以做到；不，这个奥秘的解释或许超出了所有天上势力的理解范围。因此，关于祂，我们将以最简短的词句陈述我们信经的内容，而不是人类理性通常提出的断言；这并不是出于轻率，而是由我们论述的性质所要求，向你们提出的更像是\Quote{推测}而非明确的肯定。\par

3. 因此，天主的独生子——正如前文讨论所表明的，藉着祂，按照圣经的观点，一切可见的与不可见的受造物都是藉祂而造的——祂既造了万有，也爱祂所造的。因为祂自己既是无形天主的无形肖像，祂就无形地将自己的一部分分施给所有理性的受造物，使各人按照对祂爱慕的程度，获得恰好相称的一份。但由于按照自由意志的能力，各个灵魂具有差异和多样性，以致一个以更热烈的爱依恋其存在的根源，另一个则以更微弱、更冷淡的爱依恋，那\OriginalTerm{灵魂}{anima}——关于它，耶稣曾说：\BibleQuote{没有人能夺去我的\OriginalTerm{生命}{animam}}——从受造之初就不可分离、不可分解地与祂结合，因为祂是天主的智慧与圣言、真理与真光，那灵魂完全地接受了祂，并进入了祂的光明与荣耀中，与祂卓越地成为一体，正如宗徒对那些应效法它的人所应许的：\BibleQuote{那与主结合的，便是与祂成为一神。}因此，这灵魂的实体介乎天主与肉身之间——因为天主的本性没有中间的工具就无法与身体相混合——于是，正如我们所说的，那天主之人便诞生了，那实体作为中介，其本性并不难以接受身体。另一方面，那灵魂作为理性的存在，接受天主也并不违反其本性，它如上所述，已经完全进入了那圣言、智慧和真理之中。因此，它连同它所取的身体，理应被称为天主子、天主的德能、基督、天主的智慧，这或是因为它完全在天主子内，或是因为它完全将天主子接受到自己内。再者，藉着祂万有受造的天主子，被称为耶稣基督和人子。因为天主子也被说成死了——这是指那可以接受死亡的本性；祂被称为人子，被宣告要在天父的光荣中与圣天使一同来临。为此，在整个圣经中，不仅神性本性用人的语言来表达，人性也用神性尊贵的名称来装饰。实在，关于祂，比任何其他更可确证：\BibleQuote{他们二人要成为一体，不再是一对，而是一体。}因为天主的圣言与灵魂成为一体，比男人与妻子成为一体更是如此。但除了那灵魂，谁更适宜与天主成为一神呢？它因爱如此紧密地与天主结合，以致可以公正地说它与他成为一神。\par

4. 祂的爱德之完美和配得的爱慕之真诚，为它形成了这种与天主不可分离的结合，以至于那灵魂的蒙受并非偶然，也不是出于个人的偏爱，而是作为对其德行的赏赐。请听先知这样对它说：\BibleQuote{你喜爱正义，憎恨邪恶；因此天主，你的天主，用喜乐的油傅抹了你，胜过你的同伴。}作为它爱德的赏报，它被用喜乐的油傅抹；也就是说，基督的灵魂与天主的圣言一同成了基督。因为用喜乐的油傅抹，无非就是充满圣神。当说\BibleQuote{胜过你的同伴}时，意思是圣神的恩宠赐给它，不像赐给先知那样，而是天主的圣言本身的本质圆满地临在于它内，正如宗徒所说：\BibleQuote{在祂内住有整个天主性圆满地有形有体。}最后，为此祂不仅说了\BibleQuote{你喜爱正义}，还加上\BibleQuote{你憎恨邪恶}。因为憎恨邪恶，圣经论及祂说：\BibleQuote{祂没有犯过罪，口中也找不到诡诈，}以及\BibleQuote{祂在各方面都曾受过试探，与我们相似，却没有罪。}而且主自己也说：\BibleQuote{你们谁能指证我有罪？}祂又论及自己说：\BibleQuote{看，这世界的元首来了，他在我身上一无所获。}所有这些（段落）都表明，在祂内没有罪的意识；先知为更清楚地表明罪的意识从未进入祂内，说：\BibleQuote{孩子还不知道呼唤父亲或母亲之前，他就转离邪恶了。}\par

5. 现在，如果我们在上文已表明基督拥有理性的灵魂，这使任何人感到困难——因为我们在所有讨论中经常证明，灵魂的本性既能向善也能向恶——那么这困难将以如下方式解释。的确，祂的灵魂本性与其他所有灵魂相同，这不可怀疑，否则它就不能被称为灵魂，若非真是灵魂。但由于选择善恶的能力在所有人能力范围之内，基督所拥有的这灵魂选择了喜爱正义，以至于与其爱德的无限量相应，它坚定不移、不可分离地依附于正义，以致决心的坚定、爱德的无限量和爱火的不熄灭，消灭了一切改变和变化的\OriginalTerm{感觉}{sensum}；而那本来依赖意志的东西，经过长期习惯的力量变成了本性；因此我们必须相信，在基督内有一个人的理性的灵魂，却不认为它有丝毫犯罪的感觉或可能性。\par

为更充分地说明这一问题，使用一个比喻并非不合理，尽管在如此困难的主题上难以找到恰当的比喻。然而，如果我们可以无冒犯地说，金属铁既能受冷也能受热。那么，如果一块铁持续放在火中，通过所有毛孔和脉络吸收热量，并且火不间断，铁从未离开，它就完全转化成了火本身；我们还能说这个本质上是一块铁的东西，在火中并不断燃烧时，有可能接受寒冷吗？相反，由于更符合真理，我们难道不更应该说（正如我们在熔炉中常见的那样）它已完全成为火，因为其中除了火之外什么也看不见？如果有人试图触摸或抓住它，他所体验到的将不是铁的作用，而是火的作用。如此，那灵魂就像火中的铁，持续放在圣言中，持续放在智慧中，持续在天主内，它在其所做、所感、所理解的一切事上都是天主，因此既不能被称为可转变的，也不能被称为可变化的，因为它不断受热，由于与天主圣言的结合而拥有了不变性。最终，必须认为所有圣人都从天主圣言获得了一些温暖；而在这一灵魂中，神圣的火本身就应被认为是安息的，从它那里一些温暖可能传递给了他人。最后，这句话：\BibleQuote{天主，你的天主，以喜乐油傅你，胜过你的同伴，}表明那灵魂以一种方式以喜乐油被傅油，即以天主圣言和智慧；而他的同伴，即圣先知和宗徒，以另一种方式。因为据说他们\BibleQuote{在香膏的芬芳中奔跑；}而那灵魂就是盛装那香膏本身的器皿，所有配得先知和宗徒都分享了它的芬芳。因此，正如香膏的本质是一回事，它的芬芳是另一回事，同样，基督是一回事，他的同伴是另一回事。正如盛装香膏本身的器皿决不能容纳任何恶臭；而享受其芬芳的人，如果稍微离开它的香气，就可能接受降临到他们身上的任何恶臭：同样，基督作为器皿本身，其中盛装着香膏的本质，是不可能接受相反气味的，而他的\BibleQuote{同伴}将按他们与器皿的远近程度，分受并接收他的香气。\par

事实上，我认为耶肋米亚先知也理解在他内天主智慧的性质，这智慧也是祂为拯救世界而取的那同一智慧，他说：\BibleQuote{我们鼻孔的气息，就是主基督，我们曾说过，在他的荫庇下，我们将在万民中生活。}既然我们身体的影子与身体不可分离，并不可避免地执行和重复它的动作和姿态，我认为，他意在指出基督灵魂的工作及其不可分离的运动，这运动按照他的动作和意志完成一切，称此为主基督的影子，我们在这影子下要在万民中生活。因为在此取人性的奥迹中，万民得以生活，他们通过信德效法这奥迹而获得救恩。当达味说：\BibleQuote{上主，求你记念我的凌辱，他们因你的基督而加给我的凌辱，}在我看来也指示同一件事。保禄说\BibleQuote{你们的生命已与基督一同藏在天主内，}以及在另一处\BibleQuote{你们既然寻求基督在我内说话的凭据，}时，还有别的意思吗？他现在说基督藏在天主内。这一表达的含义，除非被表明如同我们上文指出的先知在\BibleQuote{主基督的影子}这话中所指的那样，或许超出人类心智的理解。但我们在圣经中也看到许多其他关于\BibleQuote{影子}一词意义的陈述，例如\Quote{路加福音}中著名的句子，加俾额尔对玛利亚说：\BibleQuote{圣神要临于你，至高者的能力要庇荫你。}宗徒谈到法律时说，那些在肉身受割损的，\BibleQuote{事奉的是天上事物的模型和影子。}别处又说：\BibleQuote{我们在世上的生命岂不是影子吗？}如果不仅地上的法律是影子，而且我们在世上所有的生命也是影子，我们在万民中生活在基督的影子下，那么我们必须看看所有这些影子的真理是否会在那个启示中被认识，那时不再通过镜子，模糊不清，而是面对面，所有圣人都配得看见天主的荣耀和万物的原因与真理。宗徒已经通过圣神接受了这真理的凭据，他说：\BibleQuote{虽然我们曾按血肉认识基督，但如今不再这样认识他了。}\par

以上是我们在处理如此困难的主题如基督的降生成人和天主性时想到的一些想法。如果有人确实能发现更好的见解，并能用圣经中更清楚的证明来确立他的论点，那么他的意见应优先于我的被接受。\par

% structure: p0041 source=h2 target=subsection reason=relative-heading-rank; base=h2

\subsection{第七章 论圣神}

正如在开始时按照情况所需进行的那些最初讨论之后，关于圣父、圣子和圣神，我们认为应该回过头来，表明同一位天主是世界的创造者和奠基者，以及我们主耶稣基督的父，即法律、先知和福音的天主是同一的；其次，关于基督，应当表明那先前被证明是天主圣言和智慧的那一位如何降生成人；现在，我们以尽可能简洁的方式回到圣神的话题上。\par

现在，是时候凭我们所能，对圣神讲几句话了。我们的主和救主在\BibleBook{若望}{福音}中称圣神为\OriginalTerm{护慰者}{Paraclete}。因为正如同一的天主自己、同一的基督，同样，在基督来临前相信天主的人中，或藉着基督投奔天主的人中，即在先知和宗徒内，也是同一的圣神。的确，我们听说过某些异端者竟敢说有两个天主和两个基督，但我们从未听说有人宣讲过两个圣神的道理。他们怎能从圣经中为此辩护？如果确能发现对圣神的任何界定或描述，他们又能提出什么区别来区分圣神与圣神呢？即便我们向\Person{马尔西昂}或\Person{瓦伦提诺}让步，认为可以在天主性问题加以区分，将善天主的本性描述为一种，义天主的描述为另一种，他又能编造或发现什么，以引入圣神的区分呢？我认为，他们无法发现任何能表明任何区分的东西。\par

2. 我们以为，每位理性的受造物毫无区别地领受圣神，如同领受天主的智慧和圣言一样。然而，我注意到，圣神主要的降临是在基督升天之后宣告于人，而非在祂进入世界之前。因为在此之前，圣神的恩赐只赐予先知和少数个人——若民间有配得的人；但在救主来临之后，经上记载，\Person{岳厄尔}先知的预言应验了：\BibleQuote{到末日，我要将我的神倾注在所有血肉的人身上，他们必要说预言。}\BibleRef{岳 2:28}这类似那句名言：\BibleQuote{万民都要事奉祂。}\BibleRef{咏 72:11}藉着圣神的恩宠，与其它众多成果一同，这最荣耀的后果清楚地显明：关于先知书或\Person{梅瑟}法律中所记载的事，当时只有少数人，即先知们本人，以及几乎没有任何一个全民族中的人，能够超越字面意义看到更伟大、即属灵的意义；但现在有无数信徒确信，无论是割礼、安息日的休息、动物的流血祭献，还是天主在这些事上对梅瑟的答复，都不应按字面理解。这种理解方式无疑是由圣神的能力启发给所有人的。\par

3. 正如有许多理解基督的方式——祂虽是智慧，但并不在所有的人中起智慧的作用或拥有智慧的能力，而只在那些专心于祂内智慧之研究的人中；祂虽被称为医生，但并非对所有人都如此，而只对那些认识到自己软弱患病、投奔祂的怜悯以得康复的人——同样，我认为圣神也是如此，在祂内包含一切的恩赐。因为藉着圣神，有人领受了智慧的言语，有人领受了知识的言语，有人领受了信德；如此，对于每个能够领受祂的人，圣神自己成为那种特质，或被人理解为那配得参与之人所需要的特质。那些在福音中听到祂被称为\OriginalTerm{护慰者}{Paraclete}的人，没有看出这些分别和差异，也没有适当考虑祂因何工作或作为而被命名为\OriginalTerm{护慰者}{Paraclete}，就将祂比作某些普通的神灵，并因此试图扰乱基督的教会，在弟兄间激起不小的纷争；然而福音显示祂具有如此的力量和威严，以致说宗徒们尚不能领受救主想要教导的事，直到圣神降临，祂倾注到他们灵魂中，光照他们认识天主圣三的本性和信仰。但这些人因理智无知，不仅自己无法逻辑地陈述真理，甚至不能留意我们所提出的；他们对祂的神性怀有不配的想法，自陷于错误和欺骗，被错误之神败坏，而非受圣神的教导所训诲，如宗徒所宣称的：\BibleQuote{随从魔鬼的道理，禁止嫁娶，毁灭许多人，并禁食某些食物，借严格的持守炫耀，引诱无辜者的灵魂。}\BibleRef{弟前 4:1-3}\par

4. 因此我们必须知道，护慰者就是圣神，祂教导那不能用言语表达、可以说是不可言传的真理，以及\BibleQuote{是人不能说出来的}\BibleRef{格后 12:4}，即不能用人类语言指明的真理。我们认为，宗徒用\Quote{不能}一词代替\Quote{不可能}；正如他在那段话中说：\BibleQuote{凡事我都可行，但不都有益处；凡事我都可行，但不都造就人。}\BibleRef{格前 10:23}因为那些在我们能力范围内、我们可以拥有的事，他说对我们来说是可行的。但是护慰者，即被称为圣神的那位，因祂安慰的工作而得名；希腊文 \OriginalTerm{安慰}{\GreekText{παράκλησις}} 在拉丁文中称为 consolatio。因为如果有人因认识祂不可言说的奥秘而堪当分享圣神，他无疑会获得内心的安慰和喜乐。因为当他借着圣神的教导，认识一切发生之事的原因——即它们如何发生、为何发生——他的灵魂便绝不会受困扰，也不会感到任何悲伤；他也不会因任何事而惊慌，因为他紧靠天主的圣言和祂的智慧，借着圣神称呼耶稣为主。既然我们提到了护慰者，并尽力解释了关于祂应有的理解；既然我们的救主在若望书信中也被称为护慰者，他说：\BibleQuote{如果我们有人犯罪，我们在父那里有一位护慰者，就是正义的耶稣基督，祂为我们赎罪。}\BibleRef{若一 2:1-2}那么让我们思考，护慰者这个称谓用在救主身上和用在圣神身上是否含义不同。现在，护慰者当指着救主时，似乎意味着中保。因为在希腊文中，护慰者具有双重含义——中保和安慰者。因此，鉴于随后的话：\BibleQuote{祂为我们赎罪}，护慰者这个名号在救主的情况下似乎应理解为中保，因为祂因我们的罪在父面前代求。至于圣神，护慰者必须理解为安慰者，因为祂将慰藉赐予那些祂公开启示属灵知识领悟的灵魂。\par

% structure: p0047 source=h2 target=subsection reason=relative-heading-rank; base=h2

\subsection{第八章 论\OriginalTerm{灵魂}{Anima}}

1. 现在我们论述的顺序要求我们在讨论了前面的主题之后，对灵魂进行一般的探究；并且从次要的点开始，上升到更重要的点。现在，我想无人怀疑所有生物——甚至水中生物——都有灵魂。因为所有人的普遍意见都维持这一点；而且圣经的权威为此提供了佐证，经上说：\BibleQuote{天主造了大鲸，以及水生物种各样有生命的动物。}\BibleRef{创 1:21}这也被那些用特定词语定义灵魂的人从理智的普遍直觉所确认。因为灵魂被定义如下：一种 \OriginalTerm{具有想象力和冲动的}{\GreekText{φανταστικὴ} \GreekText{καὶ} \GreekText{ὁρμητική}} 实体，在拉丁文中虽然不太恰当，可译为 \OriginalTerm{感觉的和能动的}{sensibilis et mobilis}。这当然可以恰当地用于所有生物，甚至那些水中生物；对于有翅膀的动物，也可证明同一 \OriginalTerm{灵魂}{anima} 的定义成立。圣经也为第二种意见增添了权威，经上说：\BibleQuote{你们不可吃血，因为一切肉身的生命就是它的血；你们不可与肉同吃生命。}\BibleRef{肋 17:14}这最清楚地表明，每种动物的血就是它的生命。如果有人问，对于蜜蜂、黄蜂、蚂蚁以及其他水中生物，如牡蛎、蛤蜊和所有其他无血却明显是活物的，怎么还能说\BibleQuote{一切肉身的生命就是血}，那么我们必须回答，在那类生物中，在其他动物中由红血力量所发挥的机能，是由它们体内的液体发挥的，尽管颜色不同；因为颜色无关紧要，只要物质有生命即可。驮畜或小型牲畜有灵魂，这普遍同意，毫无疑义。然而，圣经的观点是清楚的，当天主说：\BibleQuote{大地要生出各种动物，各按其类：四足兽、爬虫和地上的各种野兽。}\BibleRef{创 1:24}至于人，虽然无人怀疑或需要探究，但圣经宣告：\BibleQuote{天主将生命的气息吹在他的脸上，人就成为了一个有灵的生命。}\BibleRef{创 2:7}现在需要探究的是天使品级是否有灵魂，或者本身就是灵魂；以及其他的神圣和天上的能力，以及相反种类的。的确，我们在圣经中找不到任何权威说天使或任何其他作为天主仆役的神圣神灵有灵魂或被称为灵魂，然而许多人觉得它们有生命。但关于天主，我们发现经上写道：\BibleQuote{我要把我的灵魂放在那吃了血的人的灵魂上，我要将他从他的百姓中剪除。}\BibleRef{肋 17:10}还有另一处：\BibleQuote{你们的月朔、安息日和集会，我都不悦纳；你们的斋戒、节期和庆日，我的灵魂憎恨。}\BibleRef{依 1:13-14}在第二十二篇圣咏中，关于基督——正如福音所见证，这篇圣咏确是关于祂的——有以下的话：\BibleQuote{上主，求你不要远离我的助佑，求你注视我的防卫；天主，求你从刀剑下拯救我的灵魂，从狗爪下拯救我的独一者。}\BibleRef{咏 22:20-21}此外还有许多关于基督在肉身内居住时的灵魂的其他见证。\par

2. 然而，降生成人的本性将使得对基督灵魂的任何探究成为多余。因为正如祂真实地拥有肉身，祂也同样真实地拥有灵魂。确实，既感觉到又陈述圣经中所谓天主的灵魂应如何理解，是困难的；因为我们承认那本性是单纯的，没有任何掺杂或附加。无论以何种方式理解，它暂时似乎被称为天主的灵魂；而关于基督则毫无疑义。因此，在我看来，在理解或断言关于圣天使和其他天上威能者类似的事情时，似乎并无荒谬之处，因为灵魂的那一定义似乎也适用于他们。因为谁能理性地否认他们是\Quote{有感觉和能够活动的？}但如果那定义——灵魂被说成是一种理性地\Quote{有感觉和能够活动的}实体——是正确的，那么同样的定义似乎也适用于天使。因为在他们之内，除了理性的感觉和运动之外，还有什么呢？那些被同一概念所包括的存在者，无疑具有相同的实体。保禄确实暗示有一种属灵魂的人，他说这样的人不能领受圣神的事，反而认为圣神的道理是愚拙的，并且不能领悟那些属神的事。在另一处，他说：\BibleQuote{所种的是属灵魂的身体，复活起来的是属神的身体}，指出在义人的复活中，将没有任何属于灵魂本性的东西。因此我们探究：是否存在着某种实体，就其是\OriginalTerm{灵魂}{anima}而言，是不完善的。但它是不完善是因为它从完善中堕落，还是因为天主如此创造它，这将在每个主题开始依次讨论时加以探究。因为如果属灵魂的人不领受圣神的事，并且由于他是属灵魂的，不能领受对更善本性——即神圣本性——的理解，那么也许正是为此缘故，保禄为了更清楚地教导我们，什么是我们借以能领悟属于圣神的事——即属神的事——的，他将理智而非灵魂与圣神结合并联系起来。因为我认为，当他这样说时，他指出了这一点：\BibleQuote{我要以心神祈祷，也要以理智祈祷；我要以心神歌咏，也要以理智歌咏}。他并没有说\Quote{我要以灵魂祈祷}，而是以心神和理智。他也没有说\Quote{我要以灵魂歌咏}，而是以心神和理智。\par

3. 但也许有人会问：如果理智是藉着心灵祈祷和歌唱，并且就是它接受成全和救恩，那么伯多禄怎么说：\Quote{你们收到你们信德的结局，即你们灵魂的救恩}？如果灵魂既不祈祷也不以心灵歌唱，它怎能盼望救恩？或者，当它到达幸福之境时，它就不再被称为灵魂了吗？让我们看看是否可以这样回答：正如救主来拯救那失落的，那么从前被称为失落的，当它被得救时就不再是失落的；同样地，这个被拯救的东西也许也称为灵魂，当它被置于救恩的状态时，它将从圣言那里得到一个表示其更完美状态的名称。但有些人认为还可以补充一点：正如那失落的东西在失落之前无疑存在过，那时它还是别的什么东西，而不是毁灭；同样，当它不再处于毁灭状态时也是如此。同样地，那据说已灭亡的灵魂，在尚未灭亡时曾经是某种东西，因此被称为灵魂；如果它再次脱离毁灭，它就可能第二次成为它灭亡前的样子，并被称为灵魂。但是，从灵魂这个名称在希腊文中的含义来看，一些好奇的研究者认为可能提示了一种重要的含义。因为在神圣语言中，天主被称为火，如\WorkTitle{圣经}说：\Quote{我们的天主是吞噬之火}。关于天使的实体，\WorkTitle{圣经}也说：\Quote{祂使自己的天使成为神风，使自己的仆役成为燃烧的火焰}；在另一处：\Quote{上主的天使在荆棘的火焰中显现。} 此外，我们受命要\Quote{在心灵上热忱}；这句话无疑表明天主的圣言是火热而燃烧的。\Person{耶肋米亚}先知也听到那位给他答复的说：\Quote{看，我已将我的话放在你口中成为火。} 既然天主是火，天使是火焰，众圣徒在心灵上热忱，那么相反地，那些从对天主的爱中堕落的人，无疑被认为在祂的爱中变冷，成了冷淡的。因为主也说：\Quote{由于罪恶增多，许多人的爱将冷淡。} 不，一切事物，无论是什么，在\WorkTitle{圣经}中与敌对势力相比时，魔鬼据说总是找到冷的；有什么比它更冷呢？在海中，海怪也被说成是王。因为先知暗示那蛇和海怪，当然是指其中一个恶灵，也在海中。但另一处先知说：\Quote{我要拔出我的圣剑攻击那飞行的蛇，攻击那蜿蜒的蛇，并杀死它。} 他又说：\Quote{即使他们躲避我的眼睛，下到海底，我必在那里命令蛇咬他们。} 在\WorkTitle{约伯传}中，它也被说成是水中一切的王。先知预言祸患将由北风向所有居住在地上的人燃起。现在，北风在\WorkTitle{圣经}中被描述为冷的，如\WorkTitle{智慧篇}所说：\Quote{那寒冷的北风}；\BibleRef{德 43:20} 这件事无疑也必须理解为指魔鬼。因此，如果那些神圣的东西被称为火、光、热忱，而那些相反的东西被称为冷；如果许多人的爱被说成冷淡；我们就必须探究，也许灵魂这个名称，希腊文称为\OriginalTerm{灵魂}{\GreekText{ψυχή}}，是否得名于从一种更美好的神圣状态变冷，并由此而来，因为它似乎从那种天然的神圣之热冷却了，因此被置于目前的状态，并得了现在的名称。最后，看看你是否容易在\WorkTitle{圣经}中找到灵魂受到称赞的地方：相反，它经常伴随着谴责的词语出现，如经文：\Quote{邪恶的灵魂毁灭拥有它的人}；\BibleRef{德 6:4} 以及：\Quote{那犯罪的灵魂，它必死亡。} 因为经文说了：\Quote{所有的灵魂都属于我；父亲的灵魂，儿子的灵魂，都属于我}，接下来似乎应该说：\Quote{行义的灵魂必蒙拯救}，和\Quote{犯罪的灵魂必死亡。} 但现在我们看见，祂将可责备的与灵魂联系在一起，而沉默了可称颂的事。因此，我们要看看，是否正如我们所说，由名称本身表明，它被称为\OriginalTerm{灵魂}{\GreekText{ψυχή}}，即\Emph{anima}，是因为它从公义的热忱和对神圣之火的参与中冷却了，但尚未失去恢复它起初热忱的力量。因此先知似乎也用这句话指出某种状态：\Quote{我的灵魂，回到你的安息中吧。} 从所有这些可以得出：理智离开了它的地位和尊严，就被造成或被称为灵魂；如果它得到修复和纠正，就回到理智的状态。\par

4. 如果情况如此，在我看来，理智的这种朽坏与堕落并非在所有受造物中相同，而是这种向灵魂的转化在不同个体中程度不一，有些理智甚至保留了部分先前的能力，而其他理智则保留甚少或毫无保留。因此，有些人在生命之初就展现出更活跃的理智，另一些人则心智较为迟钝，还有些人生来就全然愚钝、完全无法教导。然而，我们关于理智转化为灵魂的说法，或其他任何具有类似含义的陈述，读者必须仔细斟酌并自行决定，因为这些观点不应被视为我们以教义方式提出的主张，而仅仅是意见，以探究和讨论的风格处理。读者也应当考虑这一点：就救主的灵魂而言，福音中所记载之事，有些以\Quote{灵魂}之名归于祂，有些则以\Quote{灵}之名。因为当福音欲指出任何影响祂的苦难或扰动时，便以\Quote{灵魂}之名表明，如经上说：\Quote{现在我灵魂烦乱；}以及\Quote{我的灵魂忧闷至死；}以及\Quote{没有人夺取我的灵魂，而是我自己放下。} 祂交托的，不是灵魂，而是灵，在圣父手中；当祂说肉身软弱时，祂未说灵魂愿意，而是说灵愿意：由此看来，灵魂是介于软弱的肉身与愿意的灵之间的某种存在。\par

5. 但或许有人会用我们在陈述中已提醒过你们的那些反对意见之一来反驳我们，说：\Quote{那么，为何也提到天主的灵魂？} 对此，我们作如下回答：正如关于天主的一切有形之物的说法，如手指、手、手臂、眼睛、脚、口等，我们不应理解为人的肢体，而是藉这些身体肢体的名称来指称祂的某些能力；同样，我们也应认为，\Quote{天主的灵魂}这个称号所指的乃是别的什么。如果我们可以斗胆在这个主题上再多说一点，那么天主的灵魂或许可被理解为指天主的独生子。因为正如灵魂植入身体后，推动其中一切事物，并在它所作用的一切事物上施展力量；同样，天主的独生子，即祂的圣言与智慧，伸展并扩展到天主的每一种能力中，植入其中；或许为表明这一奥秘，天主在圣经中被称为或描述为一个身体。确实，我们必须考虑，是否正因如此，天主的灵魂可被理解为祂的独生子——因为祂亲自来到这苦难的世界，下到这涕泣之谷和这受辱之地，正如祂在圣咏中所说：\Quote{因为祢在苦难之地羞辱了我们。} 最后，我知道某些评注家在解释福音中救主的话\Quote{我的灵魂忧闷至死}时，认为这话指宗徒，祂称他们为自己的灵魂，因为他们比祂身体的其余部分更好。因为正如信徒的群众被称为祂的身体，他们说，宗徒作为比身体其余部分更好的，应被理解为祂的灵魂。\par

关于理性灵魂，我们已尽可能提出了这些观点，作为读者讨论的议题，而非确定无疑的教义命题。至于动物和其他哑巴受造物的灵魂，我们在上文已作了概括性的说明，就此足够。\par

% structure: p0054 source=h2 target=subsection reason=relative-heading-rank; base=h2

\subsection{第九章 论世界及理性受造物的运动，善恶及其原因。}

1. 但让我们现在回到本次讨论的既定顺序，就理解所能看见的创造之开端，即天主创造的起始。在那个起始中，我们应当设想，天主创造了如此众多的理性或理智受造物（无论称其为什么），我们先前称之为\Quote{理智}，其数量正是祂预见为足够的。确实，祂按祂自己预定的某个确定数目创造了它们：因为不可像某些人所想的那样，认为受造物没有界限，因为若没有界限，则既不能被领会，也不能被限定。倘若如此，受造物便不能被天主约束和管理。因为本性上，无限者也是不可领会的。此外，正如圣经所说：\BibleQuote{天主按数目和度量安排了一切；}因此，数目将正确适用于理性受造物或理智，使它们数目众多，足以被天主安排、管理和控制。而度量将恰当地适用于物质身体；我们应当相信，这度量是受造如此，如天主预知为装饰世界而足够的那样。这些就是我们所当相信的在天主创造之初（即在万有之前）所创造的东西。我们认为，这一点甚至体现在摩西以某种模糊用语引入的那个开端中，当时他说：\BibleQuote{起初天主创造了天地。}因为显然，这里所说的不是穹苍，也不是旱地，而是那原本的天和地，后来我们现今所见的天和地才由此得名。\par

2. 然而，既然那些我们前面说过在起初被造的理性本性，是在先前不存在时被造的，那么正因其不存在和存在的开端这一事实，它们必然是可变的和易变的；因为无论它们的实质中有什么能力，都不是本性地具有的，而是源于造物主的恩善。因此，它们之所是，既非它们自身所有，也非永恒持续，而是由天主赐予的。因为它并非始终存在；而凡是恩赐之物，也可能被收回而消失。被收回的理由在于灵魂的运动未按正义和适宜而行。因为造物主向祂所造的理智施予了自由自主行动的能力，借此，存在于它们内的善可以因它们自己意志的努力而成为它们自己的；但是，懈怠、对保存善的努力的厌烦、对更好事物的厌弃和疏忽，提供了背离善的开端。而背离善无非就是变为恶。因为可以肯定，缺乏善就是邪恶。由此，一个人越远离善，就越陷入邪恶。在这种状态中，每个理智按其作为，或多或少地忽略善，被拖入善的对立面，那无疑就是恶。由此可见，万物的造物主容许了多样性和差异的某些种子和原因，以便祂可以根据理智（即理性受造物）的多样性创造多样性和差异，而这种多样性必须被认为是由我们上述原因所导致的。我们现在想要解释的，就是所谓的多样性和差异。\par

3. 我们称之为世界的，是超乎诸天、在诸天中、在地上、或称为下界的地方，或任何存在的地方，连同其居民。这整个便称为世界。在这个世界中，有些存在被称为超诸天者，即居于更幸福的居所，披戴着光明而辉煌的身体；在他们之间显示出许多区别，例如宗徒说：\Quote{太阳的光荣是一种，月亮的光荣是另一种，星辰的光荣又是一种；因为星辰与星辰在光荣上有所区别。}有些存在被称为地上的，在他们之间，即在人之间，也有不小的差异；因为有些是野蛮人，有些是希腊人；而在野蛮人中，有些是野蛮凶悍的，有些性情较温和。他们中有些生活在完全认可的法律之下；另一些则在较普遍或严厉的法律之下；而有些则拥有非人道的野蛮习俗而非法律。他们中有些从出生时起就沦于卑微和服从，被当作奴隶抚养，处于主人、君王或暴君的统治之下。另一些则以更符合自由和理性的方式成长：有些身体健全，有些自幼患病；有些视力有缺陷，有些听力和言语有缺陷；有些生来如此，有些在出生后不久或至少成年时失去感官功能。何须我一一列举人类悲惨的全部情景，有些人得以幸免，而另一些人则陷入其中？每个人都可以自己权衡思量。还有一些不可见的权能，被委以管理地上事物；在他们之间也必须相信存在不小的差异，正如在人之间所见的那样。事实上，宗徒\Person{保禄}暗示存在某些较低的权能，并且在他们之间同样必定要寻求多样性的原因。至于哑巴动物、鸟类和水中的生物，似乎不必追问，因为可以肯定它们不应被视为首要等级，而是次要等级。\par

4. 既如此，既然一切受造之物都被说是藉着基督并在基督内造成的，正如宗徒保禄最清楚地指出的，当他说：\BibleQuote{因为一切受造之物，无论是天上的、地上的、有形的、无形的，或是宝座、掌权者、执政者、宰制者，都是在祂内并藉着祂受造的；一切都是藉着祂并为了祂受造的；}又如若望在他的福音中指明同一件事，说：\BibleQuote{在起初已有圣言，圣言与天主同在，圣言就是天主；这圣言在起初就与天主同在；万物是藉着祂造的；凡受造的，没有一样不是藉着祂造的；}并且圣咏中也写着：\BibleQuote{祢以智慧创造了万物；}——既如此，基督就是圣言和智慧，也是正义，那么毫无疑问，那些在圣言和智慧中受造的事物也被说是在基督的正义中受造的；这样，在受造之物中就不会有任何不义或偶然的事，而一切都被证明符合公平和正义的律法。那么，如此巨大的多样性和如此巨大的差异，怎会被理解为完全公正和正义的呢？我确信，人的能力或语言都无法解释，除非我们俯伏祈祷，恳求圣言、智慧和正义本身，祂是天主的独生子，祂以祂的恩宠倾注在我们的感官中，屈尊照亮那黑暗的，揭开那隐藏的，启示那秘密的；如果我们真配得上寻求、叩问、敲门，以致当我们在寻求时配得找到，叩门时配得敞开。那么，不是依靠我们自己的能力，而是依靠那创造万物的智慧，以及我们相信存在于祂所有受造物中的正义的帮助，尽管我们暂时无法说明，但信赖祂的仁慈，我们将努力考察和探究，世界上的这种巨大多样性和差异如何能与完全的正义和理性一致。当然，我指的只是普遍的理性；因为寻找每一个个别案例的特殊原因，要么是无知，要么是愚蠢的标志。\par

5. 现在，当我们说这个世界是在我们上面所解释的由天主创造的多样性中建立起来的，并且当我们说这位天主是良善、正义且最公义的，有许多人，尤其是那些来自马尔西翁、瓦伦提努斯和巴西里德斯学派的人，他们听说过灵魂有不同的本性，他们向我们提出异议：天主在创造世界时，将某些受造物安置在天上的住处，不仅赐予他们更好的居所，还授予他们更高贵、更光荣的地位；将执政者赐予另一些；赋予一些掌权者，一些宰制者；将天庭中最尊贵的席位赐予一些；使一些闪耀更辉煌的光荣，如星辰般灿烂；赐予一些太阳的光荣，另一些月亮的光荣，另一些星辰的光荣；使一颗星与另一颗星的光荣不同——这与天主的正义不符。并且，一言以蔽之，如果创造者天主既不愿着手、也没有能力完成一项良善完美的工作，那么在创造理性本性的存在时，即那些祂本身就是其原因的存在，祂为何使一些具有更高的等级，而另一些具有第二、第三或许多更低、更次等程度呢？其次，他们又就地上的存在向我们提出异议：有些人生来比其他人更幸运；例如，一个人由亚巴郎所生，出自恩许；另一个人则由依撒格和黎贝加所生，还在母胎中就取代了他的兄弟，并据说在出生前就被天主所爱。不仅如此，这个情况本身——尤其是有人生在希伯来人中，他在那里找到神圣法律的教导；另一个人生在希腊人中，他们自己也是智慧且学问不少的人；还有人生在埃塞俄比亚人中，他们惯于吃人肉；或生在斯奇提亚人中，弑亲是法律认可的行为；或生在陶鲁斯人中，那里陌生人被献为祭——这正是强烈反对的理由。因此他们的论点是：如果存在如此巨大的环境多样性，以及如此多样变化的出生条件，其中自由意志的能力没有发挥余地（因为没有人为自己选择出生在哪里、与谁一起或处于什么条件）；那么，如果这不是由灵魂本性的差异造成的，即一个恶本性的灵魂注定进入一个邪恶的民族，一个好灵魂进入一个正义的民族，那么剩下的结论只能是：这些事情必须被认为是受偶然和机遇支配的。如果这一点被承认，那么就不再相信世界是由天主创造的，或者由祂的眷顾所管理；其结果是，天主对每个人行为的审判看起来是无需期待的。在这件事上，确实，事物的真相只有祂能知晓，祂洞悉一切，甚至天主的深奥。\par

6. 然而，我们虽不过是人，却不愿因沉默助长异端者的傲慢，将在能力所及之内，就他们的异议给予我们想到的答复。我们已多次由圣经中所能引用的宣言表明，创造万有的天主是良善、公义、全能的。当祂起初创造祂所愿意创造的受造物时，即理性本性，祂创造他们的理由无非是出于自身，即祂自身的良善。因此，既然祂自身是那些将要受造之物存在的原因，在祂内既无任何变化或改变，亦无权能之缺失，祂创造祂所造的一切均为平等、相似，因为在祂自身内并无产生差异与多样性的理由。但既然这些理性受造物本身（正如我们多次指出、并将在适当处进一步表明的）被赋予自由意志的能力，这自由意志便激励各人或藉模仿天主而进步，或由于疏忽而致失败。而这，正如我们已说过的，便是理性受造物间差异的原因，其根源并非来自创造者的意志或判断，而是来自个体意志的自由。如今，天主认为按受造物的功绩来安排他们是公义的，便将这些不同的理智带入一个世界的和谐中，为要装饰仿佛一个居所，在其中不仅应有金器银器，也应有木器瓦器（有些作贵重的，有些作卑贱的），连同这些不同的器皿、或灵魂、或理智。依我之见，这便是那世界呈现多样性面貌的原因，而天主眷顾则继续按各人运动或情感及意向的多样性来治理每个个体。因此，创造者既不会在分配（因前述原因）各人按其功绩时显得不公；也不会认为每人出生的幸福或不幸福，或落在他身上的任何境遇是偶然的；也不会相信有不同的创造者或不同本性的灵魂存在。\par

7. 但在我看来，即使圣经本身对此奥秘的性质也并非全然缄默，正如宗徒保禄在讨论雅各伯与厄撒乌的事例时说：\Quote{孩子们还未出生，尚未行善或作恶，为使天主按拣选而定旨意得以坚立，不是出于行为，而是出于那召叫者，就有话对她说：\InnerQuote{年长的要服事年幼的}，正如经上所载：\InnerQuote{我爱了雅各伯，却恨了厄撒乌}。}随后他自问自答说：\Quote{那么我们可说什么呢？难道天主有不义吗？}为了使我们有探究这些事的机会，并查明这些事并非无缘无故发生，他自答说：\Quote{断乎不是！}因在我看来，关于雅各伯与厄撒乌所提出的同一个问题，也可针对天上、地上甚至地下的一切受造物提出。同样地，我认为正如他那里说\Quote{孩子们还未出生，尚未行善或作恶}，也可论及其他一切事物：\Quote{它们尚未} 受造，\Quote{也未行善或作恶，为使天主按拣选而定旨意得以坚立}，以致（如某些人所想的）一方面有些受造物被造为天上的，有些为地上的，另有些为地下的，\Quote{不是出于行为}（如他们所想的），\Quote{而是出于那召叫者}，那么，若事情如此，我们可说什么呢？\Quote{难道天主有不义吗？断乎不是！}因此，当仔细考察圣经中关于雅各伯与厄撒乌的记载时，并不认为天主有失公义，以致在它们出生或在此生做任何事之前就说\Quote{年长的要服事年幼的}；并不认为甚至在母胎中雅各伯就取代了他兄弟是不义的，只要我们觉得他按先前生活的功绩堪受天主的宠爱，以致配得在他兄弟之上；同样地，关于天上的受造物，如果我们注意到，多样性并非受造物的原始状态，而是由于先前存在的原因，创造者为各人按其功绩程度预备了不同的职务，其根据正是：每位作为由天主受造的理智或理性精神，按其心智的运动和灵魂的情感，为自己获得了或多或少的功绩，并成为天主所爱或所厌恶的对象；然而，其中一些拥有更大功绩者，却为了装饰世界的状态并履行对较低级受造物的职责，被命定与其他人一同受苦，以便藉此他们自己能参与创造者的忍耐，按宗徒的话说：\Quote{因为受造物被屈服于虚空之下，不是出于自愿，而是出于那使它屈服的那一位，在希望中。}因此，持守宗徒在论及厄撒乌与雅各伯出生时所说的话：\Quote{难道天主有不义吗？断乎不是！}我认为这句话应谨慎地应用于所有其他受造物，因为正如我们之前所说，创造者的公义应在一切事上彰显。依我之见，这一点最终会更清楚，如果各人——无论是天上的、地上的或阴间的受造物——都被说成在自身内、在其肉身出生之前就有其多样性的原因。因为万物都是由天主的圣言和祂的智慧所创造，并由祂的公义所安排。祂以慈悲的恩宠照顾众人，鼓励众人使用任何可助其痊愈的疗法，激励他们走向救恩。\par

8. 既然在审判之日，善人必与恶人分离，义人与不义之人分离，并且所有人将按天主的判决，依照他们各自应得的报应分派到他们所配得的各处地方，这是毫无疑义的；同样，我认为从前也有类似的情况，若天主愿意，我们将在随后说明。因为必须相信天主时时处处都是按祂的审判来行事和安排万物的。因为宗徒所说的话：\Quote{在大户人家，不但有金器银器，也有木器瓦器；有作贵重的，有作卑贱的。}以及他接着说的：\Quote{人若自洁，便成为贵重的器皿，成圣，合乎主用，准备行各种善事。}无疑指出了这一点：凡在今世自洁的人，必在来世准备好行各种善事；而不自洁的人，则按他不洁的程度成为卑贱的器皿，即不配。因此可以理解，从前也曾有理性器皿，无论洁净与否，即他们要么自洁了，要么没有；因此每个器皿按其洁净或不洁的程度，在今世得到了一处地方、或区域、或出生条件、或要履行的职务。这一切，直到最卑微的，天主以祂智慧的能力提供和辨别，并以祂管治的审判安排万物，依照最公正的报应，各人按其应得得到帮助或照顾。其中当然显明了公平的一切原则，而环境的不平等却维护了按功过报应的正义。但每个个别情况下功过的根据，只有天主自己，连同祂的独生圣言、祂的智慧及圣神，才能真正而清楚地认识。\par

% structure: p0063 source=h2 target=subsection reason=relative-heading-rank; base=h2

\subsection{第十章　论复活、审判、地狱之火及惩罚}

1. 但既然这篇论述提醒我们关于未来的审判和报应，以及罪人的惩罚，根据圣经的警告和教会的教导——即在审判之时，有永火、外面的黑暗、监狱、火炉及类似的其他惩罚为罪人预备了——让我们看看我们在这几点上应有怎样的看法。但为了按适当次序讨论这些主题，我认为我们应当先考虑复活的本质，以便知道那要来承受惩罚、安息或幸福的（身体）是什么；关于这个问题，在我们已撰写的其他有关复活的论文中，我们已更详细地讨论过，并表明了我们对此的看法。但现在，为了我们论文的逻辑顺序，从这些著作中重申几点并没什么不合理，尤其因为有些人因教会的信经而恼怒，好像我们对复活的信仰是愚蠢的，完全没有意义；这些人主要是异端者，我认为应当以如下方式回答他们。如果他们承认有死人的复活，就让他们回答我们：那死去的是什么？不是身体吗？那么，复活的就是身体。然后让他们告诉我们，他们认为我们是否要使用身体。我认为当宗徒保禄说：\BibleQuote{播种的是属生理性的身体，复活起来的是属神性的身体。}时，他们不能否认复活的是身体，或在复活中我们将使用身体。那么如何？如果我们确定使用身体，如果已腐朽的身体被宣告要复活（因为只有先前已腐朽的才能恰当地说复活），那么无人能怀疑它们会复活，以便我们在复活时再次穿上它们。这两件事是紧密相连的。因为如果身体复活，它们无疑复活来作我们的遮盖；如果我们需要被穿上身体，正如确实需要，我们应当只穿上我们自己的身体。但如果这些身体确实复活，而且它们是作为\Quote{属神性的}身体复活，那么毫无疑问，它们在摆脱腐败、除去必死性后，被称为从死者中复活；否则，某人从死者中复活以便再次死去就显得徒劳和多余。最后，如果人仔细考虑属生理性身体的特性，当它被种在地里时，又恢复属神性身体的特性，就能更清楚地理解这一点。因为正是从属生理性身体中，复活的能力和恩宠引出属神性身体，当它将其从卑贱的状况转化为荣耀的状况时。\par

2. 然而，异端者自以为博学智慧，我们要问他们：是否每一个身体都有某种形式，即按某种形状被塑造？若他们说，身体是按无形状塑造的，他们便显出自己是人类中最无知、最愚昧的。因为除非是完全不学无术的人，没有人会否认这点。但若他们理所当然地说，每一个身体当然按某种确定的形状被塑造，我们要问他们能否向我们指出并描述属神身体的形状？这是他们绝对无法做到的。此外，我们还要问他们关于复活者的差异。他们将如何证明这句话是真实的：\BibleQuote{鸟的肉是一种，鱼的肉是另一种；有属天的身体，也有属地的身体；属天的光荣是一种，属地的光荣是另一种；太阳的光荣是一种，月亮的光荣是另一种，星辰的光荣是另一种；这颗星与那颗星在光荣上有所不同；死者的复活也是如此。}那么，根据存在于天体之间的那种等级，让他们向我们展示复活者光荣的差异；若他们曾设法想出某种与天体差异相符的原则，我们便要他们借地上身体的比照来指定复活的差异。我们对这段话的理解是：宗徒欲描述在光荣中复活者（即圣人）之间的巨大差异，便借天体作比喻，说：\BibleQuote{太阳的光荣是一种，月亮的光荣是另一种，星辰的光荣是另一种。}他又欲教导我们在今世未洁净自身而将来复活者（即罪人）之间的差异，便借地上的事物作比喻，说：\BibleQuote{鸟的肉是一种，鱼的肉是另一种。}因为天上的事物堪与圣人相比，地上的事物堪与罪人相比。这些话是针对那些否认死者复活（即身体复活）的人而说的。\par

3. 现在我们将注意力转向我们中的某些（信徒），他们或因智力薄弱，或因缺乏合宜教导，对身体的复活持有非常低下卑贱的看法。我们要问这些人：他们如何理解血肉的身体将因复活的恩宠改变而成为属神的身体？如何理解那种在软弱中的将以能力复活？那种在羞辱中的将在光荣中复活？那种种在朽坏中的将转变为不朽的状态？因为，若他们相信宗徒所说的，那在光荣、能力和不朽中复活的身体已经成了属神的，那么再声称它可被血肉的情欲所缠累，就显得荒谬且与宗徒的本意相悖，因为宗徒明确宣称：\BibleQuote{血肉不能继承天主的国，朽坏也不能继承不朽。}然而，他们如何理解宗徒的话：\BibleQuote{我们都要改变？}这改变固然应按我们上面所教导的次序期待；在其中，无疑地，我们应当企盼一些配得上天主恩宠的事。我们相信，这事将按宗徒所描述的次序发生：将\BibleQuote{赤裸的谷粒，或任何其他果实}播种在地里，天主就\BibleQuote{随意给那谷粒一个身体}，一旦那谷粒死了。同样，我们的身体也应被设想像谷粒一样落入地里；（那包含身体实体、植入其中的胚芽）虽然身体死了，朽坏了，分散了，但靠着天主的话，那常在身体实体中安全的同一胚芽，将它们从地里提起，修复并复原它们，正如麦粒中的能力在其朽坏死亡后，将谷粒修复复原为具有茎和穗的身体。同样，对那些堪得承受天国产业的人，我们先前所提的身体恢复的胚芽，按天主的命令，从属地和血肉的身体中恢复出一个属神的身体，能居住在天上；而对那些功绩较低、境况更卑微、甚至在等级中最低并被完全弃绝的人，也按其生命和灵魂的尊贵程度，给予身体的光荣和尊严——尽管如此，即便是那些注定要受永火或严厉惩罚的人所复活的身体，也因复活的变化而如此不朽，以至于连严厉的惩罚也不能使其朽坏或瓦解。既然那将从死者中复活的身体有如此性质，现在让我们来看永火威胁的含义。\par

4. 我们在先知\Person{依撒意亚}书中发现，每个人所受惩罚的火被描述为他自己的；因为他说：\BibleQuote{走在你自己火的光中，并在你所点燃的火焰中}\BibleRef{依50:11}。这些话似乎表明，每个罪人都为自己点燃了己火的火焰，而不是被投入别人早已点燃或先于自身存在的一些火中。这火的燃料和食物就是我们的罪，宗徒\Person{保禄}称之为\BibleQuote{木、草、禾秸}\BibleRef{格前3:12}。并且我认为，正如丰盛的食物以及相反种类和数量的给养在身体中引发热病，且热病也有不同种类和持续时间，根据所收集的毒物为疾病提供材料和燃料的比例（这些不同毒物所收集的材料质量，决定了更急性或更慢性的疾病的原因）；同样，当灵魂为自己聚集了大量恶行和众多罪过时，在适当的时候，那一切邪恶的集合便会沸腾而受惩罚，被点燃而受责罚；那时，心灵本身或良心，藉着天主的德能，将罪过时刻印在其上的所有标记和形式记入记忆，便会看到一部历史，好像它所做的一切污秽、可耻和不洁的行为都展现在眼前：于是良心本身被困扰，被自己的刺扎透，成为控告自己和见证自己的证人。我认为这也是宗徒\Person{保禄}本人的意见，当他说：\BibleQuote{他们的思想互相控告或辩护，在天主藉耶稣基督审判人隐秘事的日子，照我的福音}\BibleRef{罗2:15-16}。由此可知，围绕灵魂的本体，由罪本身的有害情感产生了某些折磨。\par

5. 为了使对这事理的理解不致显得过于困难，我们可以从那些常降临到某些灵魂上的情欲的有害效果中汲取一些考量，例如当一个灵魂被爱火吞噬，或因热忱或嫉妒而消磨，或愤怒之情被点燃，或因巨大的疯狂或悲伤而耗尽；在某些情况下，有些人因无法忍受这些恶行的过度，便认为屈服于死亡比永远忍受这种折磨更可容忍。你或许会问：对于那些陷入上述恶习所产生邪恶中的人，他们在现世生命中未能为自己获得任何改善，并以此状态离开世界的人，在惩罚方面，让他们被这些有害情感——即愤怒、暴怒、疯狂或悲伤——的残余所折磨，是否就足够了？这些情感在生前未被任何良药减轻其致命毒性；还是说，这些情感改变后，他们将要受普遍惩罚的痛苦？现在我以为可以理解存在另一种惩罚：正如我们感受到，当身体的肢体从彼此支撑中松脱并撕裂时，会产生最剧烈的疼痛；同样，当灵魂被发现超脱了天主为善与有益的行为及观察而创造它的秩序、联系与和谐，并在其理性动态的联系中不能与自身和谐时，就必须认为它承受着自己分裂的责罚和折磨，并感受自己无序状态的惩罚。而当灵魂的这种分解与撕裂受到火的试验时，无疑会固化成一个更坚固的结构，并得以修复。\par

6. 还有许多其他事情我们未能察觉，只有我们灵魂的良医才知道。因为如果我们为了身体健康，由于饮食带来的不良影响，而认为有必要服用某种令人不快且痛苦的药物，有时甚至需要根据疾病的性质使用切割刀子的严厉手术；如果病毒的毒性甚至超出了这些疗法，最后必须用火焚烧；那么，我们就更应当明白，天主——我们的良医，渴望除去我们灵魂的缺陷，这些缺陷是它们从各种罪恶和过犯中染上的，应当采用这种惩罚性的措施，甚至对那些失去心智健全的人施加火的惩罚！这种做法的图画也出现在圣经中。在\BibleBook{申命}{纪}中，天主的话用发热、寒冷、黄疸、视力衰弱的痛苦、心智错乱和瘫痪、失明、肾脏虚弱等惩罚来威胁罪人。因此，如果有人闲暇时从整部圣经中搜集所有疾病的列举，这些疾病在针对罪人的威胁中被称为身体疾病的名称，他会发现要么是灵魂的恶习，要么是它们的惩罚，被形象地表示出来。现在要理解，正如医生对病人施用药物，以便通过细心的治疗使他们恢复健康，天主对待那些堕落犯罪的人也是如此，这由以下事实证明：天主愤怒的杯被命令通过耶肋米亚先知递给所有民族，让他们喝，陷入疯狂，并吐出来。这样做时，他威胁他们说，如果有人拒绝喝，他将不会被洁净。由此可以理解，天主报仇的愤怒对于灵魂的净化是有益的。另外，通过火施加的惩罚也被理解为具有治疗的目的，依撒意亚这样论及以色列说：\BibleQuote{主将洗去熙雍子女的污秽，并以审判之灵和焚烧之灵从他们中间清除血迹。}关于加色丁人，他这样说：\BibleQuote{你有火炭；坐在它们上面：它们将成为你的帮助。}在其他地方他说：\BibleQuote{主将在焚烧之火中圣化。}在玛拉基亚的预言中他说：\BibleQuote{主坐着吹气，并净化，并将倾泻出洁净的犹大子孙。}\par

7. 但福音中提到的那个命运，即不忠心的管家，据说要被分割，其中一部分与不信者放在一起，仿佛那不属于他们自己的部分被送到别处，无疑表明了对那些其灵魂（如我所见）与灵魂分离的人的某种惩罚。因为如果这圣神是神圣本性的，即被理解为圣神，我们将理解这是指圣神的恩赐：当藉着圣洗或圣神的恩宠，智慧之言、知识之言或任何其他恩赐被赐予一个人，却没有被恰当运用，即要么埋在地里，要么包裹在手帕里，那么圣神的恩赐当然会从他的灵魂中被收回，而剩余的部分，即灵魂的本质，将被分派与不信者同列，与那圣神——通过与主结合，它本应成为一灵——分离。现在，如果这不是指圣神，而是指灵魂自身的本质，那被称为其较好部分的是按天主的肖像和模样造的；而另一部分，即后来因自由意志的堕落而违反其原始纯洁状态的本性而接受的——这部分，作为物质的友人和爱侣，与不信者一同受罚。还有第三种理解方式：即每个信徒，即使教会中最卑微的，据说都有一位天使伴随，救主宣称这天使总是看到圣父的面，而这天使当然与他的守护对象是一体的；因此，如果后者因其不顺从而变得不配，天主的使者就被说成从他那里被取走，然后他的一部分——即属于他人性的部分——与神圣部分分离，被分派与不信者同列，因为它没有忠实地遵守天主所分配给他的天使的告诫。\par

8. 但外界的黑暗，按我的判断，与其说是指某种没有光线的黑暗大气，不如说是指那些深陷于无知黑暗之中、被置于理智之光范围之外的人。我们还必须看到，恐怕这也可能是这种表达方式的意思：正如圣人们在复活后将得到他们在今生神圣纯洁地居住过的那些身体，明亮而荣耀，同样，那些在今世喜爱错误黑暗和无知之夜的恶人，在复活后也可能被披上黑暗黑色的身体，以便那在今世占据他们内心思想的无知之雾，在未来成为身体的外部遮盖。关于监狱也应持有类似的观点。这些评论已尽可能简短，为的是暂时保持我们论述的顺序，但愿目前足够。\par

% structure: p0072 source=h2 target=subsection reason=relative-heading-rank; base=h2

\subsection{第十一章 论相反的应许。}

1. 现在让我们简短地看看，关于许诺，我们应持有什么观点。确定无误的是，没有任何生物能够完全静止不动，而喜欢各种运动、持续的活动和意志；我认为，这种本性在一切生物中显而易见。那么，理性动物，即人的本性，必然处于持续的运动和活动之中。的确，如果他忘记自己，不知道什么适合他，他所有的努力都指向服务于身体的用途，在他的所有行动中，他都沉迷于自己的享乐和肉欲；但如果他是一位致力于关心或提供公共利益的人，那么，或通过为国家的利益咨询，或通过服从长官，他努力去促进那似乎确实有利于公共福祉的事物。现在，如果任何人具有这样的本性，即明白还有比那些似乎有形的事物更好的东西，从而将他的劳动奉献给智慧和科学，那么他无疑会将全部注意力转向这类追求，以便通过探究真理，确定事物的原因和理由。因此，正如在这一生中，一个人认为最高善是享受身体快乐，另一个人认为是为了公共利益而咨询，第三个人认为专注于研究和学习；那么，让我们探究，在那真正的生命（即与基督一同隐藏在天主内的生命，也就是永生）中，我们是否也会有某种类似的秩序和存在状态。\par

2. 某些人，拒绝思考的劳苦，对律法的字面采取肤浅的看法，并且在某种程度上屈服于自己欲望和情欲的放纵，仅是字面的门徒，认为未来许诺的实现要在身体快乐和奢华中去寻找；因此，他们特别希望在复活之后，再次拥有那些永不缺乏吃喝以及履行血肉之躯所有功能的身体结构，而不遵循宗徒保禄关于复活属神身体的意见。因此，他们说，复活之后会有婚姻和生儿育女，他们想象着地上的耶路撒冷城将被重建，其根基用宝石铺设，城墙用碧玉建造，城垛用水晶；城墙由许多宝石构成，如碧玉、蓝宝石、玉髓、祖母绿、红玛瑙、条纹玛瑙、黄金石、绿玉、风信子石、紫水晶。此外，他们认为其他国家的人民将作为他们享乐的仆役赐给他们，他们要使用这些人作为耕田者或筑墙者，通过他们，他们毁坏和倒塌的城市将再次建立；他们认为他们将接受列国的财富来生活，并将控制他们的财富；甚至米甸和基达的骆驼也会来，带给他们金子、乳香和宝石。他们认为这些观点可以通过先知书中关于耶路撒冷的许诺来确立权威，也通过那些段落，如说事奉主的人将吃喝，但罪人将饥渴；义人将喜乐，但忧愁将占据恶人。并且从新约中，他们引用救主的言论，其中祂向门徒许诺关于酒的喜乐，说：\BibleQuote{从今以后，我不再喝这杯，直到我在我圣父的国里与你们同喝新酒。}他们还加上救主称那些现在饥饿口渴的人为有福，许诺他们将得到饱饫的声明；他们还引用许多其他圣经例证，但他们并未察觉到其意义应当以比喻方式理解。然后，又按照今生事物的形式，并按照这个世界等级或位阶的级次，或他们权力的大小，他们认为他们将成为国王和王子，像那些现在存在的世俗君主；主要似乎是因为福音中的那句话：\BibleQuote{有权管辖五座城。}简而言之，按照今生一切类似事物的方式，他们期望许诺中所期望的一切事情都得以实现，即现在存在的一切将再次存在。这就是那些虽相信基督，却以一种犹太人的意思理解圣经的人的观点，他们从中得出毫无神圣许诺价值的东西。\par

3. 那些按照宗徒的理解接受圣经描绘的人，怀有希望：圣人们确将进食，但那将是生命之粮，能以真理和智慧之食滋养灵魂，并光照理智，使之从神圣智慧的杯中畅饮，正如圣经所宣告的：\Quote{智慧预备了宴席，宰杀了牲畜，在杯中调了酒，并高声呼喊：到我这里来，吃我为你预备的饼，喝我调的酒。}藉着这智慧之食，理智被滋养，达到圆满完美的状态，如同人起初被造时的样式，恢复为天主的肖像和模样；因此，即使有人离开此生时并未完全受教，却行了蒙悦纳的善工，他将在那座圣人之城——耶路撒冷——领受教导，即他会被培育、塑造，成为活石，成为特选而宝贵的石头，因为他坚定恒忍地经受了一生的奋战和虔敬的考验；并将在那里获得对这里已预言之事的更真实、更清晰的认识，即\Quote{人生活不只靠饼，而也靠天主口中所发的一切言语。}那些治理较低者、教导他们、培育他们、训练他们归向神圣事物的，也要被理解为是君王和统治者。\par

4. 但如果这些看法未能给怀有如此希望之人心中注入相称的渴望，那么让我们稍作退后，暂且不顾心灵对这事实本身自然天生的渴求，而加以探究，以便最终能描述那生命之粮的样式、那酒的品质、以及权能的特质——一切皆与属灵的观点相符。在那些通常靠手工完成的技艺中，事物为何如此、为何具有特殊品质或特殊目的，其缘由成了心灵探究的对象，而实际作品则藉手工展现；同样，在天主所造的工程中，我们所见之事的原因和理据仍是隐藏的。正如我们的眼睛看到工匠的作品时，一察觉任何异常的技艺卓越，心灵便炽燃渴望知道其性质、如何形成或为何做成；同样，程度远胜、无可比拟地，心灵以难以言喻的渴望炽燃，想知道我们所见天主所行之事的缘由。我们相信，这渴望和热望无疑是由天主植入我们内的；正如眼睛天然寻求光明和视觉，我们的身体天然渴求饮食，我们的心灵也怀着一种相称而自然的渴望，要认识天主的真理和万物的起因。我们从天主领受了这渴望，并非为了永远不得满足或\Emph{能够}满足；否则，若真理之爱永无机会得偿所愿，则它似乎被天主徒然植入我们心中。因此，即使在此生，那些以极大辛劳致力于虔敬和宗教追求的人，虽然只得着神圣知识浩瀚宝藏中的些许碎片，但由于他们的心思意念专注于这些追求，且因渴望之切而超越自身，他们便获得极大益处；又因他们的心灵被导向寻求真理的研究和热爱，他们便更适于接受将来的教导；这就如同有人要绘画时，先用轻铅笔勾勒出未来图像的轮廓，为之后加入的特征预备标记，这初步的轮廓素描便为铺上真正的色彩铺路；同样，在一定程度上，藉着我主耶稣基督的铅笔，我们心版上可描画轮廓和草图。因此，或许经上所说：\Quote{凡有的，还要给他，且使他富裕。}由此证实，凡在此生拥有真理和知识轮廓的，将来必得着一幅完美图像的美。\par

5. 我想，这样的渴望曾由那说\BibleQuote{我进退两难，我渴望离世与基督同在，那是好得无比的}人表明出来；他知道当他回到基督那里时，他就更清楚地知道世上一切事的原因，不论关于人、人的灵魂、心智，或关于任何其他主题，例如：什么是运行的灵，什么是生命之灵，或是赐给信徒的圣神的恩宠是什么。然后他也将明白以色列是什么，或各民族差异的含义；以色列十二支派是什么意思，每个支派的个别意义是什么。然后，他也将明白司祭和肋未人的理由，以及不同司祭等级的理由，其预象曾在梅瑟身上，还有与天主有关的喜年和年之周的真实含义。他也将看见节庆日、圣日以及一切祭献和洁净的理由。他也会感知治愈癞病的理由，以及各种癞病的区别，以及治愈遗精者的理由。此外，他将知道善的影响力及其伟大和特质，以及相反类的影响力，前者对人的情感，以及后者对人引起的争竞。他也要观看灵魂的本性，以及动物的多样性（无论是水中的、鸟类的、还是野兽的），以及为什么每个属类又分为如此多的种类；造物主的意图或祂智慧的目的是隐藏在每一件事物中。他也会熟悉某些性质与某些根茎或草药相关联的理由，以及另一方面，恶效应被其他草药和根茎所避免的理由。而且，他将知道背逆天使的本性，以及他们为何有能力奉承那些没有以全部信德的力量鄙视他们的人，以及他们为何存在以欺骗和引人走入歧途。他也会学到天主眷顾对每一件事的判断；并且人的遭遇中没有一件是偶然或碰巧发生的，而是按照一个如此精心策划、如此宏大的计划，以至于它连头上的头发数目也不忽略——不仅是圣人们的，也许还是所有人类的——这眷顾的统治计划甚至延伸到关心两个麻雀卖一个德纳的事，无论那里的麻雀是比喻义还是字面义理解。由此我们可以推定，在离世之后，要等到只把地上事物的缘由指示给那些堪当领受的人，可能会经过不少时间，使他们借知晓这一切，并借圆满知识的恩宠，享有说不出的喜乐。此外，如果天地之间的气域并非没有居住者，而且是理性种类的，正如宗徒所说：\BibleQuote{你们从前在其中行事为人，随从这世界的潮流，随从空中权势的首领，即现今在悖逆之子身上运行的灵。}他又说：\BibleQuote{我们要被提到云里，在空中与主相遇，这样我们就常与主同在。}\par

6. 因此，我们应设想，圣人们将留在那里，直到他们辨认出在空中所行之事中的\Quote{双重模式}。当我说\Quote{双重模式}时，我的意思是：当我们在地上时，我们看见动物或树木，看见它们之间的差异，也看见人之间极大的不同；但尽管我们看见了这些事，却不明白其中的理由；而我们从可见的多样性中只获得这样的提示：我们应当考察并探究这些事物是按什么原则被创造或不同地安排的。而这样一种求知的热忱或欲望在我们地上时被激发，其圆满的理解和领悟将在死后赐予，如果结果确实如我们所期待的话。因此，当我们完全理解了它的本性时，我们将以双重的方式理解我们在世上所看见的事物。因此，对于这空中的居所，我们必须持有类似的观点。因此，我认为，所有离开此世的圣人们，将留在某个位于地上的地方，圣经称之为\OriginalTerm{乐园}{paradise}，作为一个教导之所，可以这么说，是灵魂的课堂或学校，他们在其中将受教导关于他们在世上所见的一切，并接受一些关于未来之事的资讯，正如他们此生时已获得某种程度的未来事件的预兆，虽然\BibleQuote{藉着镜子观看，模糊不清}，\BibleRef{格前 13:12}这一切都在适当的时间和地点更清楚、更明白地启示给圣人们。如果有人确实心地纯洁、思想神圣、知觉更为熟练，他将因更快的进步而迅速升到空中的一处地方，并到达天国，通过那些——可以这样说——在不同地方的\OriginalTerm{住所}{mansions}，希腊人称之为\OriginalTerm{球体}{spheres, i.e., globes}，但圣经称之为\OriginalTerm{诸天}{heavens}；在其中的每一个中，他将首先清楚地看见那里所行之事，其次将发现为何如此行事的原因：这样他将依次穿过所有等级，跟随那已进入诸天的，即天主子\Person{耶稣}，他曾说：\BibleQuote{我愿意，我在那里，这些人也同我在一起} \BibleRef{若 17:24}。祂谈到这地方的多样性时曾说：\BibleQuote{在我父的家里，有许多住处} \BibleRef{若 14:2}。祂自己是无所不在的，迅速地穿越万物；我们不再将祂理解为存在于那些祂曾为我们而被限制的狭窄界限中，即不是在于祂在地上与人同居时所占据的那个有局限的身体，根据这一点，祂可能被视为被围困在某个单一的地方。\par

7. 因此，当圣人们到达天上的住所时，他们将清楚地看见星辰的本性，逐一了解它们是否赋予生命，或其状态如何。他们也将领会天主事工的其他原因，那些原因祂自己将启示给他们。因为祂将如同对待孩子一样，向他们展示万物的原因和祂创造的能力，并解释为什么那颗星被放置在天空的特定区域，为什么它与另一颗星相隔如此遥远的距离；例如，如果它更近或更远，会有什么后果；或者如果那颗星比这颗更大，万物的总体将不会保持不变，而一切都会转变为另一种存在状态。因此，当他们完成了与星辰和天体运行相关的一切事务后，他们将来到那些看不见的事物面前，或那些我们仅知其名的事物，以及那些不可见的事物，宗徒\Person{保禄}曾告诉我们这些为数众多，尽管它们是什么，或它们之间有什么差异，我们甚至无法凭借软弱的智力去猜测。因此，\OriginalTerm{理性本性}{rational nature}通过每一步而成长，不是如在此生中在肉身、身体和灵魂中成长，而是在理解和感知能力上扩展，被提升为一个已经圆满的心智，达致\OriginalTerm{圆满知识}{perfect knowledge}，不再被那些肉体的感官所阻碍，而是在理智的增长中增进；并永远纯洁地、可以说是面对面地注视着万物的原因，它达到圆满：首先是它由此上升到真理，其次是它由此安住于真理之中，以问题和事物的理解以及事件的成因为其可享用的食粮。因为正如在此生中，我们的身体靠着早期充足的食物供应而增长到其应有的尺寸；但在达到适当的身高后，我们不再为增长而进食，而是为活着并藉之维持生命；同样，我认为，心智在达到圆满时，也进食并享用适当合宜的食粮，其程度既无不足也无多余。而在一切事上，这食粮应被理解为\OriginalTerm{默观和理解天主}{contemplation and understanding of God}，这食粮对于被造受造的本性来说，有适当合宜的度量；这度量应当被每一个开始看见天主（即藉着心灵的纯洁理解祂）的人所遵守。\par

\SourceRecord{Translated by Frederick Crombie.；Ante-Nicene Fathers；Vol. 4.；Edited by Alexander Roberts, James Donaldson, and A. Cleveland Coxe.；Buffalo, NY: Christian Literature Publishing Co.,；1885.；<http://www.newadvent.org/fathers/04122.htm>.}

% structure: p0001 source=h1 target=section reason=page-title

\section{\WorkTitle{论首要原理（第三卷）}}

% structure: p0002 source=h2 target=subsection reason=relative-heading-rank; base=h2

\subsection{鲁菲努斯序言}

读者，请你在祈祷中纪念我，使我们也能堪当成为圣神的效仿者。前两卷\WorkTitle{论首要原理}的翻译，我不但应你的请求，甚至是在你的催促下于四旬期内完成的；但当时你，我热心的弟兄玛卡里乌斯，不仅住在我附近，而且有比现在更多的闲暇可支配，所以我也更加努力；而我在讲解后两卷时花费了更长时间，因为你从城市遥远的边缘来得不那么频繁来催促我的工作。现在，如果你还记得我前一篇序言中警告过你的——某些人会愤慨，如果他们没听到我们说奥力振的坏话——我想你立刻就会经历到，这事已经发生了。但如果那些煽动人舌毁谤的魔鬼因那部作品而如此狂怒——在那部作品中他尚未完全揭露它们隐秘的行径——那么，你认为在这部作品中，他将揭露那些黑暗隐秘的途径，魔鬼通过这些途径潜入人心，欺骗软弱不稳定的灵魂，情况又会如何呢？你将立刻看到一切陷入混乱，骚动被激起，全城一片喧嚷，而那个试图藉福音之灯的光驱散无知之魔鬼黑暗的人被传来接受判决。然而，让那些渴望在神圣学问中受训、同时完整持守天主教信仰准则的人，对这些事不以为意。不过，我认为有必要提醒你，前几卷所遵循的原则在这几卷中也得到了遵守，即不翻译那些似乎与奥力振其他意见及我们自己的信仰相违背的段落，而是跳过这些段落，视其为他人篡改和伪造的。但是，如果他为了讨论和增进知识，在理性受造物（我们信仰的本质不依赖于这个主题）方面似乎表达了任何新意，也许当我们有必要按此顺序回答某些异端意见时，无论是在本卷还是在前几卷中，我都没有省略这些内容，除非他想在后续各卷中重复前几卷已经陈述的内容，那时为了简洁起见，我认为适当删减了一些重复。然而，如果有人出于增进知识的愿望而阅读这些段落，而不是为了吹毛求疵，那么最好由有学识的人来阐释它们。因为让语法学家解释诗歌的虚构和喜剧的荒谬表演，而认为没有导师和解释者就能学到那些关于天主、天上德能以及整个宇宙的论述——其中驳斥了外教哲学家或异端者的一些可悲错误——是荒谬的；其结果就是，人们宁愿鲁莽无知地谴责困难隐晦的事物，也不愿通过勤奋和研习来弄清它们的含义。\par

% structure: p0004 source=h2 target=subsection reason=relative-heading-rank; base=h2

\subsection{译自鲁菲努斯的拉丁文译本}

% structure: p0005 source=h3 target=subsubsection reason=relative-heading-rank; base=h2

\subsubsection{第一章 论意志的自由}

1. 当我们引领理智去默观那永恒无限的世界，并凝视其不可言喻的喜乐与福乐时，我们相信对天主应许也应持有某些类似的观点。但是，由于教会的宣讲包含对未来公正审判的信德，这种信德激励并劝勉人们度善生与德行生活，并尽一切可能避免罪恶；并且由此无疑表明，我们自己有能力投身于值得赞美或值得责备的生活，因此我认为有必要就意志的自由说几句话，因为许多作者已经以不平凡的风格论述了这个主题。为了更容易弄清什么是意志的自由，让我们探究意志和欲望的本质。\par

2. 在一切运动的事物中，有些运动的起因在自身之内，有些则从外部获得；凡是仅从外部运动的事物都是无生命的，如石头、木头，以及任何仅凭其物质的构成或身体的实体而持存的事物。的确，那种将物体因腐朽而分解视为运动的观点必须被排除，因为它与我们当前的目的无关。另一些事物运动的起因在自身之内，如动物、树木，以及一切由自然生命或灵魂所维系的事物；其中有些人认为金属的矿脉也应归入此类。火也被认为是自身运动的原因，泉水或许也是如此。在那些自身拥有运动起因的事物中，有些被称为从自身之外运动，有些则是被自身所运动。他们如此区分，因为那些从自身之外运动的事物虽然有生命，却没有灵魂；而那些有灵魂的事物则被自身所运动，当一种幻象，即欲望或刺激呈现在它们面前，激发它们朝向某物运动时。最后，在某些具有灵魂的事物中，有这样一种幻象，即一种意志或感觉，如同某种自然本能召唤它们，并唤醒它们作有秩序规律的运动；正如我们所见蜘蛛的情形，它们被一种幻象，即某种编织的意愿和欲望，以一种非常有秩序的方式激发起来进行织网的工作，无疑某种自然运动引发了这种工作的努力。甚至发现这种昆虫除了编织的自然欲望之外，没有其他感觉；同样，蜜蜂也具有建造蜂房和收集所谓空中蜂蜜的欲望。\par

3. 但是，既然理性的动物不仅自身拥有这些自然的运动，而且比其他动物更大程度地拥有理性的能力，借此它能判断并决定关于自然运动的事，并否决和拒绝一些运动，同时赞同和采纳另一些运动，那么，藉此理性的判断，人的运动可以被引导和导向一种可称赞的生活。由此得出，既然在人里面的这种理性的本性自身拥有区分善恶的能力，并且在区分时拥有选择它所赞同的事物的能力，那么它选择善时理应被视为值得称赞，而追随卑劣或邪恶之事时则应受谴责。这一点绝不容我们忽视：在某些不会说话的动物身上，其运动比其他动物更有规律，例如猎犬或战马，以至于有些人可能认为它们是被某种理性感觉所驱动。但我们必须相信，这与其说是理性的结果，不如说是某种自然本能的结果，这种本能在很大程度上是为了这类目的而赋予的。现在，正如我们开始指出的那样，鉴于理性的动物具有这样的本性，有些事可能从外部发生在我们人类身上；这些事接触我们的视觉、听觉或任何其他感官，可能激发并唤醒我们趋向善的运动或相反的运动；既然它们来自外部来源，阻止它们的到来并不在我们自己的能力范围之内。但是，决定并赞同我们应当如何使用那些如此发生的事情，这恰恰是我们内在的理性即我们自己的判断的职责；藉此理性的决定，我们使用从外部来到我们的激发，用于理性所赞同的目的，我们的自然运动由其权威决定，导向善行或相反。\par

4. 若有人现在要说，那些从外部原因发生在我们身上并唤起我们运动的事情，其性质是我们无法抵抗它们，无论它们激发我们行善或作恶，那么，持此观点的人不妨稍加注意自己的内心，仔细审视自己心灵的动态，除非他已发现，当对任何欲望的诱惑升起时，直到灵魂的赞同被获得，心灵的权威已对邪恶的暗示予以放任，否则一事无成；这样，似乎有两个方面基于某些可能的理由向居于我们内心法庭的法官提出请求，以便在陈述理由之后，执行的判决从理性的判断发出。举例说明：如果一个人决心过节制和贞洁的生活，使自己远离与妇女的一切玷污，而一个女人碰巧出现在他面前，煽动和引诱他违反自己的决心，那么这个女人并不是他犯过错的完全和绝对的原因或必然性，因为他有能力通过记住自己的决心来约束情欲的冲动，并通过德性的严厉告诫来抑制诱惑他的快乐；这样，所有放纵的感觉被驱除后，他的决心可以保持坚定和持久。最后，如果对某些受过学问训练、藉神圣锻炼而坚强的人，这类诱惑出现，他们立即记起自己是谁，并回想长期默想和教导的主题，并以更神圣教义的支持来强化自己，他们拒绝并击退一切快乐的诱惑，并通过置于他们内在的理性的干预来驱散敌对的情欲。\par

5. 既然这些立场通过某种自然的证据如此确立，那么把我们的行为的原因归咎于那些从外部发生在我们身上的事情，从而把责任从我们自己身上转移开（责任完全在于我们自己），这岂不是多余的吗？因为这就等于说我们就像木头或石头，本身没有运动，而是从外部接受运动的原因。这种断言既不真实也不恰当，仅仅是为否认自由意志而捏造出来的；除非我们真的认为自由意志在于：没有任何从外部发生在我们身上的事情能够激发我们行善或作恶。若有人将我们过失的原因归于身体自然的失调，这种理论证明与一切教导的理性相悖。因为，正如我们在许多人身上所见，他们在过放纵不贞的生活、成为奢侈和情欲的俘虏之后，如果碰巧被教导和训诫之言唤醒，踏上更好的生活道路，就会发生如此巨大的改变，以至于从奢侈邪恶的人转变为清醒、最贞洁和温和的人；同样，我们又看到，在安静诚实的人身上，与不安分和无耻的人交往后，他们的良好品行被恶言恶语败坏，变得像那些邪恶至极的人一样。这种情况有时也发生在成年人身上，以致他们在年轻时比年岁增长、能够更自由地生活时活得更贞洁。因此，我们推理的结果是表明：那些从外部发生在我们身上的事情并不在我们自己的能力范围内；但是，藉着我们内在的理性（它区分并决定应当如何使用这些事情），善用或恶用那些如此发生的事情，则是<Emph{是}>我们能力范围之内的。\par

6. 现在，为了以圣经的权威来证实理性推理——即我们生活得正确与否是我们自己的事，我们既不是被迫于外部原因，也不是如某些人所想，被迫于命运的安排——我们引用\Person{米该亚}先知的话作为见证：\BibleQuote{人哪，已通知你什么是善，或主要求你什么，无非是履行正义，爱好仁慈，并准备与上主你的天主同行。}\BibleRef{米 6:8}\newline{}\Person{梅瑟}也这样说：\BibleQuote{我已将生命与死亡的道路摆在你们面前：你们要选择善，并行走其中。}\BibleRef{申 30:15}\newline{}此外，\Person{依撒意亚}也宣告说：\BibleQuote{如果你们愿意，并听从我，你们将享用土地的出产。但如果你们不愿意，也不听从我，刀剑将吞噬你们；因为上主亲口说了这话。}\BibleRef{依 1:19-20}\newline{}在圣咏中也写道：\BibleQuote{如果我的百姓听了我的话，如果以色列行走在我的道路上，我必会制服他们的敌人。}\BibleRef{咏 81:14-15}\newline{}由此表明，百姓有能力去听，并行走在天主的道路上。救主也说：\BibleQuote{我对你们说：不要抵抗恶人；}\BibleRef{玛 5:39}\newline{}以及：\BibleQuote{凡向自己的弟兄发怒的，就要受审判；}\BibleRef{玛 5:22}\newline{}以及：\BibleQuote{凡注视妇女而贪恋她的，就已经在心里与她犯了奸淫；}\BibleRef{玛 5:28}\newline{}在颁布其他一些命令时——无非是要表明，遵守所命令的事出于我们自己的能力。因此，我们若违犯那些我们能够遵守的诫命，就理应受谴责。为此，祂自己又宣称：\BibleQuote{凡听了我这些话而实行的，我要告诉他像什么人：他像一个聪明的人，把房子建在磐石上；}\BibleRef{玛 7:24}\newline{}等等。同样，这段话：\BibleQuote{凡听了这些话而不实行的，就像一个愚笨的人，把房子建在沙土上；}\BibleRef{玛 7:26}\newline{}等等。甚至对右边的人说的话：\BibleQuote{来，我父所祝福的，你们来吧！\newline{}因为我饿了，你们给了我吃的；我渴了，你们给了我喝的；}\BibleRef{玛 25:34-36}\newline{}清楚地表明，这取决于他们自己：要么这些人因遵行所命令的而配得称赞，并领受所应许的；要么那些人因听了或接受相反的事而应受谴责，对他们说：\BibleQuote{可咒骂的，离开我，到那永火里去吧。}\BibleRef{玛 25:41}\newline{}我们也应注意，宗徒\Person{保禄}对我们说，我们对自己的意志拥有权能，我们自身就拥有获救或灭亡的原因：\BibleQuote{难道你轻看祂丰厚的慈爱、宽容和忍耐，不知天主的慈爱是引领你悔改吗？但你因固执和不悔改的心，为自己积蓄愤怒，直到审判和显示天主正义审判的日子；祂要按每个人的行为报应各人：对于那些恒心行善，寻求光荣、不朽和永生的人，他们将得到永生；但对于那些固执、不信真理而信从邪恶的人，将有愤怒、怨恨、患难和困苦，加给每一个作恶的人，先是犹太人，后是希腊人；但光荣、尊贵与平安，加给每一个行善的人，先是犹太人，后是希腊人。}\BibleRef{罗 2:4-11}\newline{}你还会在圣经中找到无数其他段落，清楚地表明我们拥有自由意志。否则，赐给我们诫命，遵守它们就能得救，违犯它们就要受谴责，就是自相矛盾的——如果我们没有被赋予遵守这些诫命的能力的话。\par

7. 但是，既然在圣经本身中发现一些表达出现在这样的上下文中，以至于可能从它们中理解出与此相反的意思，让我们把它们摆在我们面前，按照虔敬的规则加以讨论，并给出它们的解释，以便通过我们现在阐明的少数段落，其他类似且似乎排除对意志的任何控制的段落也能得到澄清。因此，那些天主在谈到\Person{法郎}时使用的表达，例如他频繁地说：\BibleQuote{我要使法郎的心硬}。因为如果他是由天主硬心的，并因此犯罪，那么他的罪的起因就不在于他自己。如果是这样，那么看起来\Person{法郎}并不拥有自由意志；而且，因此可以主张，依照这个例子，其他丧亡的人也不应将其毁灭归因于他们自己的自由意志。同样，在\BibleBook{厄则克耳}{先知书}中的那句话：\BibleQuote{我要从他们身上除去石心，赐给他们肉心，使他们遵行我的诫命，遵守我的道路}，也可能影响一些人，因为遵行他的道路或遵守他的诫命似乎是天主的恩赐，如果他除去那石心——那遵守他诫命的障碍——并赐予并植入一颗更好、更柔软的心，现在称为肉心。还要考虑我们的主和救主在福音中对那些问他为什么用比喻对群众说话的人的回答的性质。他的话是：\BibleQuote{使他们看见却看不见，听见却听不见，也不明白；免得他们悔改，罪得赦免。}此外，宗徒\Person{保禄}的话：\BibleQuote{不在乎那愿意的，也不在乎那奔跑的，只在乎那显示仁慈的天主}；在另一处：\BibleQuote{立志和行事都是由天主而来}；又在别处：\BibleQuote{因此，他愿意怜悯谁就怜悯谁，愿意使谁硬心就使谁硬心。这样，你将对我说：他为什么还指责人呢？谁能抗拒他的旨意呢？人哪，你是谁，竟敢向天主抗辩？受造之物岂能对造它的说：你为什么这样造我呢？难道陶匠没有权柄用同一团泥，造一个器皿为尊贵，另一个为卑贱吗？}——这些以及类似的宣告似乎对阻止许多人相信每个人都拥有对自己意志的自由，并使人觉得人得救或丧亡是天主旨意的结果，具有不小的影响。\par

8. 那么，让我们从对\Person{法郎}所说的话开始，据说他被天主硬心，以致他不让百姓离去；并且，与他的案例一起，宗徒\Person{保禄}的话也将被考虑，他说：\BibleQuote{因此，他愿意怜悯谁就怜悯谁，愿意使谁硬心就使谁硬心。}因为异端者主要依靠这些段落，断言救恩不在我们自己的能力之内，而是灵魂具有那样的本性，以致必须不可避免地要么丧亡要么得救；并且，邪恶本性的灵魂绝不可能变好，善良本性的灵魂也不可能变坏。因此他们坚持认为，\Person{法郎}也是由于败坏的本性而被天主硬心，天主使那些属地的本性的人硬心，却怜悯那些属神的本性的人。那么，让我们看看他们的断言是什么意思；首先，我们要求他们告诉我们，他们是否坚持认为\Person{法郎}的灵魂是属地的本性，如他们所说的丧亡。他们无疑会回答说是属地的本性。如果是这样，那么当他的本性反对他这样做时，相信天主或服从天主是不可能的。如果这是他本性上的状况，那么为什么还需要他的心动硬，并且不止一次，而是多次呢？除非确实因为他有可能被说服而屈服？没有人能被另一个人说成是硬了心，除非他自己原本不是刚硬的。如果他不是自己刚硬的，那么他也不是属地的本性，而是那种在神迹奇事面前可能屈服的人。但他对天主的计划是必要的，为了拯救群众，天主可以借着他抵抗众多奇迹并反抗天主旨意来彰显在他身上的能力，并且他的心动据此被说是硬了心。这就是我们对这些人的初步回答；通过这些，他们的断言可以被推翻，即他们以为\Person{法郎}是因邪恶的本性而毁灭。关于宗徒\Person{保禄}的话，我们也必须以类似的方式回答他们。因为，按照你们的观点，天主使谁硬心呢？就是那些你们称为败坏本性的人，并且我认为，如果他们不被硬心，他们就会做别的事。确实，如果他们因被硬心而走向毁灭，他们就不再是自然地丧亡，而是由于所遭遇的事。其次，天主对谁显示怜悯呢？就是对那些将要得救的人。那些因本性一次得救，从而自然地得到福乐的人，还需要第二次怜悯呢？除非从他们的情况可以看出，因为他们有可能丧亡，所以他们获得怜悯，以免丧亡，而得到救恩，并拥有善的国度。让我们以此回答那些杜撰并编造善或恶的本性（即属地的或属神的灵魂）的寓言的人，按照他们所说的，每个人据此要么得救要么丧亡。\par

9. 现在我们也必须回答那些认为法律的天主仅是公义而非良善的人；让我们问这样的人，他们认为法郎的心是怎样被天主硬化的——通过什么行为或什么预先安排。因为我们必须把握一个概念：在我们看来，天主既是公义又是良善，但按照他们的说法，仅只是公义。让他们向我们展示，一个连他们也承认是公义的天主，如何能公义地使一个人的心变硬，以便因那硬化而犯罪并丧亡。如果天主自己就是那些因不信（由于被硬化）而最终被祂以法官的权威定罪之人的毁灭原因，那么天主的公义又如何辩护？因为祂为什么指责法郎说：\BibleQuote{但既然你不让我的百姓走，看哪，我要击杀埃及中所有头生的，连你的长子}，以及天主通过梅瑟对法郎所说的一切其他话？因为每一个持守圣经记载的真实性、并希望表明法律和先知的天主是公义的人，都必须为这一切提供理由，并表明其中毫无贬损天主公义之处，因为尽管他们否认祂的良善，但承认祂是公义的审判者，是世界的创造者。然而，我们对那些断言这世界的创造者是恶毒的存在（即魔鬼）的人的回答方法则是不同的。\par

10. 但由于我们承认通过梅瑟说话的天主不仅是公义的，也是良善的，让我们仔细探究，一位公义和良善的神性如何能硬化法郎的心。让我们看看，是否能效法宗徒保禄的榜样，借助一些平行事例来解决这个难题：例如，我们能否表明，通过同一个行为，天主怜悯一个人，却硬化另一个人；祂并非有意或希望那被硬化的被硬化，而是因为在祂显示良善和忍耐的同时，那些以傲慢无礼对待祂的仁慈和宽容之人的心，因他们罪行的惩罚被延迟而硬化；而那些将祂的良善和忍耐作为悔改和革新的机会之人，则得到怜悯。然而，为了更清楚地说明我们的意思，让我们采用宗徒保禄在\BibleBook{}{希伯来书}中所用的比喻，他说：\BibleQuote{因为土地吸收了常下在它上面的雨水，并生出对耕种它的人有益的蔬菜，就会从天主那里得到祝福；但那长出荆棘和蒺藜的，就被废弃，临近诅咒，结局就是焚烧。}现在，从我们引用的保禄的话中，清楚地表明，通过天主同一个行为——即祂降雨在地上的行为——一部分土地，若得到精心耕种，就结出好果实；而另一部分，被忽略和无人照料，则长出荆棘和野蓟。如果有人仿佛以雨的口吻说：\Quote{是我，雨，造成了好的果实，也是我造成了荆棘和蒺藜的生长}，尽管这话听起来可能很刺耳，但却是真实的；因为除非雨水降下，否则不会长出果实、荆棘或蒺藜，而随着雨水的降临，大地就同时产生了这两者。现在，虽然由于雨水有益的作用，大地生出了两种植物，但植物的多样性并不能正确地归因于雨水；而应公正地将坏种子的责任归于那些人，他们虽然可以通过频繁的翻耕松土，通过多次耙地打碎土块，并除掉所有无用和有害的杂草，以及通过耕作所要求的一切劳动和辛劳来清理和预备田地以迎接即将到来的阵雨，却忽视了这样做，因此他们将收获荆棘和蒺藜，这是他们懒惰最合适的果子。因此，结果是，虽然雨水仁慈且无偏私地平均降在整个大地上，但通过雨水的同一个作用，被耕种的土地以祝福的方式为勤劳细心的耕种者产出有用的果实，而因农夫疏忽变得硬化的土地只长出荆棘和蒺藜。因此，让我们把天主所行的那些迹象和奇迹视为祂从上降下的阵雨；而人的意愿和渴望则视为耕种和未耕种的土地，这些土地就本质而言的确是相同的，正如每块土地都是相同的一样，但并非处于相同的耕种状态。由此可以得出，每一个人的意志，如果未受训练、野蛮、未开化，就会因天主的奇迹和奇事而变硬，变得越来越野蛮和多刺；或者，如果它被清除恶习并服从训练，就会变得更加柔顺，全心全意地顺服。\par

11. 但是，为了更清楚地确立这一点，不妨再举一个比喻：例如，有人说太阳既能使物体变硬，又能使之熔化，尽管硬化和熔化是性质相反的事。说太阳凭借同一热力既能融化蜡，又能烘干和硬化泥土，这并不错误；因为它的作用方式并非对泥土是一样，对蜡是另一样，而是因为泥土与蜡的性质不同，尽管它们按本性都是出于土。同样，天主的同一作为，藉着\Person{梅瑟}所行的神迹奇事施行出来，既显露出\Person{法郎}在极度邪恶中所怀的硬心，也显露出那些与以色列人混杂、并且记载与希伯来人一同离开埃及的其他埃及人的服从。至于经上说\Person{法郎}的心逐渐被制服，有一次他说：\BibleQuote{你们不可走得太远；你们要走三天的路程，但要把妻子、儿女和牲畜留下}，以及其他一些陈述，表明他逐渐向神迹奇事让步，这除了显示神迹奇事的能力对他产生了一定影响，但未达到应有程度之外，还说明了什么？因为如果硬心像许多人所理解的那样，他根本不会在少数事情上让步。但我认为，按照通常用法来解释\Quote{硬心}一词的比喻性或转义性质，并无不合理之处。因为那些以仁慈待奴隶而著称的主人，常常因自己极大的忍耐和纵容而使得奴隶变得傲慢无礼、一文不值时，对他们说：\Quote{是我把你弄成这个样子的；是我宠坏了你；是我的忍耐使你变得一无是处；你的乖僻和恶习都怪罪于我，因为我没有立即按你每次过失的报应惩罚你。}因为我们必须首先注意语言的比喻性或转义意义，然后才能理解表达的力量，而不是在我们未弄清其内在含义时就指责这个词。最后，宗徒\Person{保禄}显然谈论这类事，对那些仍留于罪恶中的人说：\BibleQuote{你藐视祂丰厚的恩慈、宽容和忍耐，不晓得天主的恩慈是领你悔改吗？你竟任着你那硬心和不悔改的心，为自己积蓄忿怒，以致在忿怒和显明天主公义审判的日子，忿怒临到你身上。}宗徒对那些在罪中的人说了这样的话。让我们把这些话应用在\Person{法郎}身上，看看它们是否也适切地描述了他：因为他按自己的硬心和不悔改的心，在忿怒的日子为自己积蓄和储存忿怒，因为若非有如此众多和伟大的神迹奇事，他的硬心就永远不会被宣告和显明出来。\par

12. 但是，如果我们所提出的证据似乎不够充分，宗徒的比喻似乎缺乏适用性，那么让我们加上先知权威的声音，看看先知们关于那些起初确实度着正义生活、配得领受天主仁慈的许多证据，但后来作为人而误入歧途的人说了些什么；先知自己也与他们认同，说：\BibleQuote{为什么，上主，祢使我们离开祢的道路？使我们的心顽固，不敬畏祢的名？求祢因祢仆人的缘故，为祢产业的支派而转回，使我们也能从祢的圣山得一点产业。}耶肋米亚也用了类似的言语：\BibleQuote{上主，祢欺骗了我们，我们被欺骗了；祢抓住了（我们），祢得胜了。}那么，那些祈求怜悯的人所说的\BibleQuote{为什么，上主，祢使我们的心顽固，不敬畏祢的名？}这句话，应当按照比喻的和道德的意义来理解，就好像有人说：\Quote{为什么祢这么久宽恕我们，当我们犯罪时不报应我们，反而遗弃我们，以致我们的邪恶增加，我们的犯罪自由因惩罚停止而扩展？}同样，一匹马若不持续感受骑手的刺，不被马嚼子擦伤嘴，它就会变硬。一个男孩若不常受责罚管教，也会成长为傲慢的青年，随时准备陷入恶习。因此，天主就遗弃和疏忽那些祂判定不配受惩罚的人：\BibleQuote{因为主所爱的，祂必管教；凡祂收纳的儿子，祂都鞭打。}由此我们可以推断，那些配得被主鞭打和管教的人，将被收纳为儿子的地位和爱情，以便他们也能通过忍受考验和磨难，能够说：\BibleQuote{谁能使我们与在基督耶稣内的天主的爱相隔绝？是困苦吗？是窘迫吗？是饥荒吗？是赤贫吗？是危险吗？是刀剑吗？}因为所有这些，每个人的决心得以显明和展示，他坚忍的坚定得以显明，与其说是对那在事情发生前就预知一切的天主，不如说是对理性属天的诸德能，他们在获得人类救恩的工作中分享一份，作为某种助手和天主的仆役。另一方面，那些尚未以如此坚定和爱情将自己奉献给天主，尚未准备好事奉祂，尚未预备灵魂接受考验的人，被说成是被天主遗弃，即不被教导，因为他们没有准备好接受教导，他们的训练或照顾无疑被推迟到以后。这些人确实不知道他们将从天主那里得到什么，除非他们首先产生受益的愿望；而最终将如此，如果一个人首先认识自己，感到自己的缺陷，并明白他应当或能够向谁寻求弥补自己的不足。因为一个人若事先不知道自己的软弱或疾病，就不能寻求医生；或者至少，在恢复健康后，那个人不会对他的医生心存感激，如果他没有首先认识到自己病情的危险性质。同样，除非一个人首先认清自己生活的缺陷和罪的邪恶本性，并用他自己的口通过告解将其宣告出来，否则他不能得到洁净或赦免，以免他不知道自己所有的是因恩赐而获得的，而将出自天主慷慨的认为是自己的财产，这种想法无疑产生心灵的傲慢和骄傲，并最终成为个人灭亡的原因。我们必须相信，魔鬼的情况就是如此，他将自己那无罪时持有的首位视为自己的，而非天主所赐；于是，在他身上应验了那句宣告：\BibleQuote{凡高举自己的，必被贬抑。}由此我看来，天主的奥秘向明智和通达的人隐藏，正如圣经所言：\BibleQuote{使任何血肉在天主面前不可自夸。}却启示给小孩子——即给那些在变为婴儿和小孩（即回到小孩子的谦卑和单纯）之后，然后进步的人；当他们达到成全时，记得他们得到幸福状态不是由于自己的功劳，而是由于天主的恩宠和怜悯。\par

13. 因此，一个应被遗弃的人，是因天主的判决而被遗弃；同时天主对一些罪人施以宽容，但这并非没有确定的行动原则。实际上，祂的忍耐本身对这些人有益，因为祂所眷顾的灵魂是不朽的；作为不朽和永恒的存在，即使未立即受到眷顾，也不被排除在救恩之外——救恩被推迟到更合适的时候。因为对于那些深受邪恶毒害的人，或许在较晚的时期获得救恩更为有益。正如医生有时虽能迅速掩盖伤口的疤痕，却暂时拖延治疗，期待更好更完全的康复——他们知道在处理伤口引起的肿胀时，延缓治疗并让恶性的体液暂时流出，比仓促地表面愈合更有益于健康；因为若关闭了病态体液的惯常出口，其毒素必将潜入肢体内里，渗透到内脏的深处，不仅引起身体的疾病，更会危及生命——同样，天主也深知人心隐秘、预知未来，以极大的忍耐容许某些事情发生，这些外在之事促使内在隐藏的私欲和恶习显露出来，以便借此洁净和医治那些因极度疏忽和漫不经心而在自己内接纳了罪恶之根和种子的人；当这些恶被驱逐到表面并暴露出来时，就能在一定程度上被抛弃和驱散。于是，一个人即使看似受严重痛苦折磨，四肢痉挛，也终有一天会得到缓解和停止受苦；在饱受痛苦之后，经历诸多苦难，他可能恢复其（原有的）状态。因为天主对待灵魂，不仅着眼于我们当下生命这短短六十年或更久，而是着眼于永恒无尽的时期，对不朽的灵魂施行祂的眷顾，如同祂自己是永恒不朽的一样。因为祂使祂按自己肖像和样式所造的理性本性成为不朽的；因此，不朽的灵魂并不因当下生命的短暂而被排除在天主的医治和治愈之外。\par

14. 但我们也要从福音中取一些关于我们刚才所提及之事的比喻，其中描述了一种岩石，其上只有薄薄一层表土，当种子落在上面时，据说很快就会发芽；但发芽后，随着太阳升到天上，它就枯萎了，并死去，因为它没有将根深深扎入土中。\newline{} 这岩石无疑代表了人的灵魂，因自己的疏忽而刚硬，因自己的邪恶而变成石头。因为天主并没有通过创造的行为赐给任何人一颗石心；而是每个人的心因自己的邪恶和悖逆而被称为石心。\newline{} 因此，如果有人要责怪农夫没有更快地将种子撒在石地上，因为撒在其他石地上的种子被看到迅速发芽，农夫肯定会回答说：\Quote{我撒这地更慢，是为了使它能保留所接受的种子；因为这地适合慢慢地播种，免得庄稼因过快发芽，仅从浅土的表层长出，而无法抵挡太阳的光线。}\newline{} 先前挑剔的人岂不会因农夫的理由和更高的知识而默许，并赞同那先前在他看来毫无理由的事，如今却视为基于理性而做？\newline{} 同样，天主，祂所有造物中技艺娴熟的农夫，无疑会隐藏并推迟那些我们认为本应更早获得健康的事物，到另一个时间，为的是使事物内在而非外在得以治愈。\newline{} 但如果现在有人向我们提出异议，说某些种子确实撒在了石地上，即落在刚硬如石的心上，我们应当回答，这也不是没有天主眷顾的安排；因为若非如此，就不会知道轻率的听和漠视的查究招致了什么谴责，也不会知道有序接受训练带来了什么益处。\newline{} 因此，灵魂得以认识自己的缺陷，并将责备归于自身，随之而保留并服从训练，即为了让它看到必须先除去自己的缺点，然后才能接受智慧的教导。\newline{} 因此，正如灵魂无数，它们的性情、目的、活动、欲望和冲动也各不相同，其多样性远非人心所能把握；因此，这种安排的技巧、知识和能力必须唯独留给天主，因为只有祂能知道每个灵魂的药方，并量定其痊愈的时间。\newline{} 所以，唯独祂，正如我们所说，认识各人的道路，并决定祂应如何引领法郎，好使祂的名通过他传遍全地，此前先用许多打击责罚他，最后将他淹死在海中。\newline{} 然而，这淹死并不意味天主对法郎的眷顾就此终结；因为我们不可想象，因为他被淹死，他就立刻完全消灭了：\BibleQuote{在天主的手中，有我们和我们的话语；一切智慧，以及技艺的学识，也都在祂手中}，正如圣经所宣告的。\newline{} 但这些论点我们已尽己所能加以讨论，处理了圣经中关于天主使法郎的心刚硬的那一章，并符合这句话：\BibleQuote{祂愿意恩待谁，就恩待谁；祂愿意使谁刚硬，就使谁刚硬。}\par

15. 现在让我们来看看\Person{厄则克耳}的那些段落，其中他说：\BibleQuote{我要从他们身上除掉石心，赐给他们肉心，使他们遵行我的律例，谨守我的典章。}因为如果天主在祂愿意的时候，除掉石心，赐予肉心，使祂的典章得以遵守，祂的诫命得以服从，那么看来除去邪恶就不在我们的能力之内。因为除掉石心似乎无非就是从任何人身上除去那使人硬心的邪恶，无论天主愿意从谁身上除去它。赐予肉心，为遵守天主的规诫，服从祂的命令，也无非是人变得顺服，不再抗拒真理，而是施行德行。那么，如果天主应许要这样做，而且在祂除掉石心之前，我们自己无法除掉它，结论就是除去邪恶不在我们的能力之内，而唯独在天主的能力之内。此外，如果在我们内形成肉心不是我们的行为，而是唯独天主的工作，那么度有德行的生活就不在我们的能力之内，而将完全显明是神圣恩宠的工作。这就是那些希望从圣经的权威证明没有任何事物在于我们自己的能力之人的主张。现在我们对这些人的回答是：这些段落不应如此理解，而应按以下方式理解。假设有一个人愚昧未受教导，他感到自己无知的耻辱，或因某人的劝勉所驱使，或因渴望效法其他智者所激励，便将自己交托给一个他确信会精心训练、胜任教导的人。那么，如果这个先前在无知中硬心的人，如我们所说，以全然的决心顺服一位师傅，并承诺在所有事情上服从他，师傅明见他决心的坚定性质，便会适当地应许除去一切无知，并在他的心智中植入知识；但这并非在门徒拒绝或抵抗他的努力时承担此事，而仅在他提供并约束自己在一切事上服从的情况下。同样，天主的圣言应许那些亲近祂的人，祂将除掉他们的石心，但不是从那些不听从祂圣言的人身上，而是从那些接受祂教导的规诫之人身上；正如我们在福音中看到病人走近救主，请求获得痊愈，并最终因而得医治。为使瞎子得以治愈、重见光明，他们的部分在于向救主恳求，并相信祂能治好他们；而祂的部分则在于恢复他们的视力。天主的圣言也正是以这种方式应许通过除掉石心（即除去邪恶）来赐予教导，使人能够遵行天主的规诫，并遵守法律的命令。\par

16. 接下来要面对的是救主在福音中说的那句话：\BibleQuote{他们看了，却看不见；听了，却听不明白，免得他们回头，罪得赦免。}对此，我们的对手会评论说：\Quote{如果那些更清楚地听的人必然被纠正和归正，且如此归正以致配得赦罪，而清楚听道并不在他们自己的能力之内，反倒取决于教导者更公开、更清楚地教导，而祂声明祂不向他们清楚地宣讲圣言，免得他们也许听了、明白了，就归正得救，那么当然，他们的得救并不取决于他们自己。如果是这样，那么无论得救或灭亡，我们都没有自由意志。}若不是后面还有\BibleQuote{免得他们也许回头，罪得赦免}这句话，我们或许更倾向于回答：救主不愿意那些祂预知不会成为善人的人明白天国的奥秘，因此才用比喻对他们说话；但既然有那句附加的话\BibleQuote{免得他们也许回头，罪得赦免}，解释就变得更加困难了。首先，我们需要注意，这段经文为反驳那些异端者提供了怎样的辩护。他们惯于从旧约中挑出任何在他们看来似乎断言造物主天主严酷残忍的表述，例如当祂被描述为有复仇或惩罚的情感，或任何此类情绪时，他们便否定造物主中存在良善；因为他们不以同样的心思和情感来判断福音书，也不留意福音书中是否有他们谴责和指责旧约的那些说法。因为显然，在所引的段落中，救主被表明（正如他们自己所承认的）没有清楚说话，正是为了使人不回头，回头后罪不得赦免。倘若仅仅按字面理解这些话，那么它们所含的肯定不亚于旧约中他们所挑剔的那些段落。如果他们认定新约中此类上下文中的任何表述需要解释，那么同样，旧约中那些受指责的表述也必然可以通过类似的解释而摆脱污名，从而表明新旧两约中的经文都出自同一位天主。现在让我们尽力回到所提出的问题上来。\par

17. 我们在先前讨论\Person{法郎}的案例时曾说过，有时过快治愈一个人反而没有益处，尤其当疾病被封在身体内部而更加猛烈发作的时候。因此，洞悉隐秘之事、在一切发生之前就已预知一切的天主，因祂的大善而延迟对这样的人的治疗，并将其康复推迟到更远的时期——可以这样说，祂通过不医治来医治他们，以免过于健康的状态使他们变得不可救药。因此，对于那些被视作\Quote{外面}的人，即我们的主和救主向他们说话的人，祂通过鉴察人心和肺腑，看出他们尚不能接受更清晰的教训，便以言语的覆盖来隐藏最深奥奥秘的含义，免得他们迅速悔改并得到医治（即快速获得罪的赦免）之后，又轻易滑回同样的疾病，因为他们发现这疾病可以毫无困难地被治愈。因为如果情况如此，无人能怀疑惩罚会加倍，邪恶的严重性也会增加：不仅那些看似被赦免的罪过重犯，而且当善德的殿堂被诡诈污秽之生灵践踏、内心充满隐藏的邪恶时，也就被亵渎了。对于那些尝过罪恶之不洁污秽食物、又尝过善德的甘美并将其甜美送入自己口中、却再次投身于罪恶的致命有毒供应的人，还能有什么补救呢？谁又会怀疑，拖延与暂时的放弃是更好的做法，以便将来有一天，他们若恰巧对邪恶感到厌倦，对其现在所喜悦的污秽变得厌恶，那时天主的圣言才最终恰当地向他们显明，神圣之物不致扔给狗，珍珠不致抛给猪——猪会践踏它们，并且反过来撕咬和攻击那些向他们宣讲天主圣言的人呢？因此，这些人就是被称作\Quote{外面}的人——毋庸置疑，这是相对于那些被称为\Quote{里面}、更清晰地听到天主圣言的人而言的。然而，那些\Quote{外面}的人确实听到了圣言，只是被比喻遮盖、被箴言隐藏。除了这些外面的人，还有另一些人，被称为\Place{提洛}人，他们完全听不到——救主知道，如果行在他们中间的神迹行在别人中间，他们早就披麻蒙灰悔改了；然而这些人连那些\Quote{外面}的人所听到的也听不到。我认为，这是因为这类人在邪恶中的等级远低于且差于那些被称为\Quote{外面}的人（即那些离\Quote{里面}之人不远、且配得听到圣言——尽管是借比喻——之人）；并且，也许他们的医治被延迟到那个日子，即审判之日，那时对他们而言，比那些神迹被记载下来并施行在他们面前的人更为容易忍受，以便他们最终在解除罪的负担后，能更轻松、更有耐力地踏上安全之路。这一点我希望读者铭记：对于如此艰难和晦涩的主题，我们的最大努力不在于清楚地找出问题的解答（每个人都会按圣神所赐的话语来做），而是以最明确的方式持守信德准则，努力表明天主眷顾公平管理一切，也以最正义的原则治理不朽的灵魂，按照每个人的功过和动机给予赏报；当前事物的安排并不局限于今世的生命，而是前存的功过状态始终为后续的状态提供基础，如此，基于永恒不变的正义律法和天主眷顾的统治，不朽的灵魂得以达到完美的顶峰。然而，如果有人反对我们的说法，认为某些邪恶败类的人故意回避宣讲的圣言，并问为什么圣言要传给那些人，以致被鄙视的\Place{提洛}人反而被优先对待（这种做法无疑增加了这些人的邪恶，使他们的定罪更加严厉——即他们听到圣言却不信），那么必须这样回答：天主是所有人思想之创造者，祂预见到人们对祂眷顾的抱怨，特别是那些说\Quote{我们怎能相信？因为我们既没有看见别人所见的，也没有听见别人所听的宣讲？这样看来，责任不在我们，因为那些听到圣言、见到神迹的人毫不延迟地成为信徒，是被神迹本身的力量所征服}的人——祂为了消除这类抱怨的借口，并表明破坏的根源并非天主眷顾的隐藏，而是人心意志的决定，便将祂恩宠的恩赐也赐予不配和不信的人，为要堵住一切人的口，使人心知道一切缺失在于自身，而不在于天主；同时，也能领悟和认识到：轻视已赐予的神圣恩惠的人，比那些不配获得或听到恩惠的人，所受的定罪更为沉重；并且，天主眷顾管理的一个特点及其最公正管理的标记是：有时祂对某些人隐藏看见或听到神圣能力奥秘的机会，免得他们在看见神迹的力量、认识并听到其智慧奥秘后，却加以藐视和冷漠，从而因他们的不敬而受更严厉的惩罚。\par

18. 现在让我们来看这句话：\BibleQuote{不在乎那愿意的，也不在乎那奔跑的，只在乎那显示仁慈的天主。}因为我们的对手断言，如果一个人的得救既不在于那愿意的，也不在于那奔跑的，而在于那显示仁慈的天主，那么我们的救恩就不在我们的权能之内。因为我们的本性是这样的：它容许我们或者得救，或者不得救；否则，我们的救恩就完全取决于祂的意愿，祂若愿意，就显示仁慈，并赐予救恩。现在，让我们首先询问这些人：渴望福乐是善行还是恶行？以善为终极目标而奔赴是否值得赞美？如果他们回答这样的行为应当受责备，他们显然是疯了；因为所有圣善之人都渴望福乐并奔赴福乐，当然无可指责。那么，一个未得救的人，如果他本性是恶，怎么会渴望福乐并奔赴福乐，却找不到它们呢？因为他们说，坏树结不出好果子，然而渴望福乐却是一颗好果子。而坏树的果子怎能是好的呢？如果他们断言渴望福乐并奔赴福乐是一种无关善恶的行为，即既非善也非恶，那么我们要回答说：如果渴望福乐并奔赴福乐是无关善恶的，那么它的对立面——渴望邪恶并奔赴邪恶——也应该是无关善恶的；然而事实是，渴望邪恶并奔赴邪恶绝不是无关善恶的，而是明显邪恶的行为。因此可以确定，渴望并追随福乐不是无关善恶的，而是一种有德的行为。\par

既然我们已用所给的答复驳斥了这些异议，让我们赶紧转入主题本身的讨论，其中说道：\BibleQuote{不是由于人愿意，也不是由于人奔跑，而是由于显示仁慈的天主。}在\WorkTitle{圣咏集}——在\WorkTitle{登阶之歌}中（这些歌归属于\Person{撒罗满}）——出现了以下陈述：\BibleQuote{除非主建造房屋，建造的人就徒然劳力；除非主守护城池，守警者就徒然警醒。}这些话确实并非指示我们停止建造或守护我们内在城池的安全；而是指出：凡不靠天主建造的，凡不靠他守护的，都是徒然建造，徒然守护。因为凡妥善建造和妥善保护的一切事，主都被视为建造或保护的原因。例如，倘若我们看到一座宏伟的建筑物，一座由美丽建筑技艺建造的华丽建筑群，我们岂不应当公正而恰当地说，这样的建筑不是靠人的能力，而是靠天主的帮助和力量建造的？然而，这样的陈述并不意味着人的努力和勤勉是懒惰的，毫无成就。或者，如果我们看到一座城池被敌人严密封锁，攻城的战车逼近城墙，城池受到围栅、武器、火器以及一切导致毁灭的战争器械的压迫，当敌人被击退和驱散时，我们岂不应当正确而恰当地说，天主打发了拯救给这座被解放的城池？然而，我们这样说，并不意味着哨兵的警觉、年轻人的机敏、卫兵的守护是缺失的。宗徒也应以类似的方式理解，因为单凭人的意志不足以获得救恩；任何凡人的奔跑也不能赢得天国的奖赏，获得在基督耶稣内天主召叫的奖品，除非我们这善意的意志、准备好的决心，以及我们内在的任何勤勉，都得到天主的帮助和供应。因此，宗徒极合逻辑地说：\BibleQuote{不是由于人愿意，也不是由于人奔跑，而是由于显示仁慈的天主；}正如我们关于农业所说的：\BibleQuote{我栽种了，\Person{阿颇罗}浇灌了，然而使之生长的是天主。所以，栽种的不算什么，浇灌的也不算什么，唯独使之生长的天主。}因此，正如一块田地带来丰盛良好的收成至完全成熟时，没有人会虔诚而合理地断言农夫造出了那些果实，而会承认那是天主产生的；同样，我们自己的成全也是如此，不是靠我们保持不动和懒惰，（而是靠我们的一些行动）：然而其成就不会归于我们，而归于天主，他是工程的首要原因和主要原因。因此，当一艘船克服了海上的危险，虽然结果是靠水手们极大的努力、航海技术的帮助、领航员的热情和仔细、顺风的推力以及星象的仔细观察而达成的，但任何有理智的人都不会把船只的安全——当它在波浪中颠簸、被巨浪折磨后最终安全抵达港口——归因于任何其他事物，而只归因于天主的仁慈。就连水手或领航员也不敢说：\Quote{我拯救了这船，}而是将一切归于天主的仁慈；并非他们认为自己没有为拯救船只贡献技能或劳力，而是因为他们知道，虽然自己贡献了劳力，但船只的安全是由天主确保的。同样，在我们人生的赛跑中，我们自己也必须付出劳力，并发挥勤勉和热情；但救恩作为我们劳力的果实，应希求自天主。否则，如果天主不要求我们的任何劳力，他的诫命就显得多余了。同样，\Person{保禄}责备某些人背离真理，称赞其他人坚守信德，也是徒劳；他向各教会传授某些规诫和制度，也是徒劳；我们自己渴望或追求善也是徒劳。但可以肯定，这些事并非徒然；可以肯定，宗徒们并非徒然给予训诲，主也并非毫无理由地制定法律。因此，我们反而要声明，异端者诋毁这些美好的声明是徒劳的。\par

19. 此后，接着这一点，即\BibleQuote{愿意和行事都是出于天主。}我们的对手主张，如果愿意出于天主，行事出于他，或者不论我们行善或行恶、愿善或愿恶，都是出于天主，那么我们就无\OriginalTerm{自由意志}{free-will}可言。对此我们必须回答：宗徒的话并没有说愿意行恶出于天主，或愿意行善出于他；也没有说行善或行恶出于天主；而是说了一个普遍陈述：愿意和行事都是出于天主。因为正如我们从天主获得这一品质——我们是人，我们呼吸，我们行动；同样，我们从天主获得（能力）借以愿意，就好比我们说我们的运动能力来自天主，或各个肢体的履行这些职责及其运动来自天主。由此，我当然不理解为：因为手移动，例如，不义地惩罚或偷窃，这行为就出于天主；而只是说运动的能力来自天主；而我们的责任是将那些我们执行的能力来自天主的运动，转向善或恶的目的。所以宗徒说的是：我们确实接收了意志的能力，但我们误用意志于善或恶的欲望。同样，我们也必须判断结果。\par

20. 但是关于宗徒的声明：\BibleQuote{因此，祂愿意怜悯谁，就怜悯谁；愿意使谁顽固，就使谁顽固。你必要对我说：那么，祂为什么还要指责人呢？谁能抗拒祂的旨意呢？然而，人啊！你是谁，竟敢向天主抗辩？被造的岂能对造它的说：你为什么这样造了我？难道陶匠没有权柄从同一泥团中，制造一个器皿作尊贵的器皿，另一个作卑贱的器皿吗？}\BibleRef{罗 9:18-21} 有人也许会争辩说，就如陶匠从同一\OriginalTerm{泥团}{lump}中制造一些器皿作\OriginalTerm{尊贵的器皿}{vessel unto honour}，另一些作\OriginalTerm{卑贱的器皿}{vessel unto dishonour}，同样，天主创造一些人是为了丧亡，另一些人是为了得救；因此，得救或丧亡并不在于我们自己的能力；按照这种推理，我们似乎没有\OriginalTerm{自由意志}{free-will}。我们必须这样回答持有这种意见的人：宗徒是否可能自相矛盾？如果对宗徒不能这样设想，那么，他们如何认为宗徒指责在格林多行淫乱的人、或犯了罪而不悔改自己所作的丑行、淫乱和污秽的人是对的？同样，祂又大大称赞行事正义的人，如奥乃息佛禄一家，说：\BibleQuote{愿主怜悯奥乃息佛禄的家庭，因为他屡次使我精神快慰，不以我的锁链为耻；而且他一到了罗马，便急切地访寻我，终于找到了。愿主在这一天使他获得主的仁慈。}\BibleRef{弟后 1:16-18} 然而，宗徒的庄重不允许指责应受指责的人（即犯了罪的人），并大大称赞因善工而应受称赞的人；然后，又好像没有人能行善或作恶一样，说这是造物主的行为，使人人都行善或作恶，因为祂制造一个器皿作尊贵的器皿，另一个作卑贱的器皿。祂又如何能加上那声明：\BibleQuote{我们众人都必须出现在基督的审判台前，为使每个人因自己的身体所行的，无论是善是恶，获得报应。}\BibleRef{格后 5:10} 因为如果一个人被造物主造来就是为了不能作恶，那么给他善的奖赏有何意义？或者，如果一个人因造物主的创造行为而不能行善，那么给他应得的惩罚又有什么公正？再者，这与另一处的声明：\BibleQuote{在大户人家中，不只有金器和银器，也有木器和瓦器，有些作尊贵的器皿，有些作卑贱的器皿。所以，人若自洁，脱离卑贱的事，必成为尊贵的器皿，圣洁的，为主所合用，预备行各样的善事。}\BibleRef{弟后 2:20-21} 岂不是相互矛盾吗？因此，那自洁的人被造成尊贵的器皿，而那不屑于洗净自己污秽的人，则被造成卑贱的器皿。按我看来，从这样的声明中，我们行为的起因绝不可归诸造物主。因为天主造物主建造一个器皿作尊贵的器皿，另一些器皿作卑贱的器皿；但那自洁于一切污秽的器皿，祂使之成为尊贵的器皿；而那以罪恶的污秽弄脏自己的器皿，祂使之成为卑贱的器皿。由此得出的结论是，每个人行为的起因是预先存在的；然后，每个人按自己的\OriginalTerm{应得之分}{deserts}，被天主造成尊贵的或卑贱的器皿。因此，每个单独的器皿自己向它的造物主提供了祂将它造为尊贵或卑贱的器皿的原因和机会。如果我们所声称的看来正确（事实上也是如此，并且与一切虔诚相协调），即每个器皿由天主预备为尊贵或卑贱的，是由于先前的原因，那么，当我们以同样的顺序和方法讨论更遥远的因由时，对灵魂的本性得出同样的结论，即认为这就是雅各在出生入世以前就被爱，而厄撒乌仍在母腹中就被恨的原因，这似乎并不荒谬。\par

21. 不，那从同一\OriginalTerm{泥团}{lump}制造一个器皿得尊荣，另一个得羞辱的声明并不会困扰我们；因为我们声称一切理性灵魂的本性都是相同的，就像同一泥团被描述为在\OriginalTerm{陶匠}{potter}手中一样。既然理性受造物的本性是同一的，天主就按照先前的功绩，从它里面创造并形成一些人得尊荣，另一些人得羞辱，正如陶匠从同一泥团中制造一些器皿得尊荣，另一些得羞辱。现在，关于宗徒所说的仿佛带责备语气的话：\BibleQuote{然而，人啊！你是谁，竟敢向天主抗辩？}\BibleRef{罗 9:20} 我认为，他的意思是，这样的责备并不指任何一个生活正直、公义、对天主有信心的人，即如梅瑟这样的人，经上记载说：\BibleQuote{梅瑟讲话，天主用声音回答他；} 正如天主回答了梅瑟，每一位\OriginalTerm{圣人}{saint}也向天主回答。但是，那\OriginalTerm{不信者}{unbeliever}，因自己生活和行为的不堪而在天主面前失去回答的信心，并且对这些事不是寻求学习进步，而是反对和抗拒，说得更明白些，就是能说出宗徒所指示的那些话的人，当他问：\BibleQuote{那么，祂为什么还要指责人呢？谁能抗拒祂的旨意呢？}\BibleRef{罗 9:19}——对这样的人，宗徒的责备才正当地指向：\BibleQuote{然而，人啊！你是谁，竟敢向天主抗辩？} 因此，这样的责备不适用于信徒和圣人，而是适用于不信者和恶人。\par

现在，对于那些引进不同本性灵魂的人，并且利用宗徒的这一宣告来支持他们自己的意见，我们必须如下答复：即使他们同意宗徒所说的，即从同一团泥中造出了那些成为\OriginalTerm{贵重的器皿}{vessels of honour}的和那些成为\OriginalTerm{卑贱的器皿}{vessels of dishonour}的，他们称之为要得救和要毁灭的本性者，那么将不再有不同本性的灵魂，而是所有人都有一本性。而如果他们承认同一个陶工无疑可以表示同一位造物主，那么得救者和丧亡者也就不会有不同的创造者了。现在，确实，让他们选择：他们愿意认为所指的是那创造恶人和败坏之人的慈善的造物主，还是那创造善人和备好得尊荣之人的不慈善的造物主？因为回答的必要性将迫使他们从这两者中选择其一。但按照我们的声明，我们说是由于先前的缘由，天主才造了贵重的器皿或卑贱的器皿，天主公义的认可毫不受限制。因为有可能，这个由于先前缘由在今世被造为贵重器皿的器皿，如果行为轻忽，会在另一个世界按其行为的功过转变为卑贱的器皿；同样，如果某人由于先前缘由在今生被他的造物主造为卑贱的器皿，却改过自新，洁净自己脱离一切污秽和恶习，他可能在新的世界成为贵重的器皿，圣化且合用，预备行各样的善事。最后，那些由天主在今世造为以色列人，却度了与其民族高贵不相称的生活，并从其出身的伟大堕落了的人，将在来世，因他们的不信，在某种程度上由贵重的器皿转变为卑贱的器皿；而另一方面，许多在今世被算作\OriginalTerm{埃及人}{Egyptian}或\OriginalTerm{厄东人}{Idumean}器皿的人，采纳了以色列人的信仰和行为，当他们行了以色列人的事，并进入主的教会，将在天主众子的启示中出现为贵重的器皿。由此，更符合虔诚规则的是相信，每个理性造物，按照其志向和生活方式，有时从恶转善，有时从善堕落；有些持守于善，另一些进步到更好的状况，总是上升到更高的事物，直到达到至高的等级；而另一些则留在邪恶中，或者，如果他们内在的邪恶开始进一步蔓延，他们就降到更坏的状况，沉入邪恶的最深处。因此我们也必须设想，可能有一些人起初确实以小的过犯开始，但倾注邪恶到如此程度，并如此精通邪恶，以至于在邪恶的量度上他们甚至与敌对的力量相等；再者，如果藉着多次严厉惩罚的施行，他们能够在未来某个时候恢复理智，逐渐试图为他们的伤口找到医治，他们可能在停止邪恶后被恢复到善的状态。因此我们认为，既然灵魂正如我们经常所说是不死不灭、永恒的，那么有可能在无数无终的时期中，在无量不同的世界里，它可以从至善降至至恶，或者从至恶恢复至至善。\par

但是，由于宗徒的话——他在论及贵重的器皿或卑贱的器皿时所说的\BibleQuote{所以，人若自潔，脫離這些事，就必作貴重的器皿，為聖所使用，準備作各樣的善事。}——似乎不将任何事归诸天主的权能，而全然归诸我们自己；而他在宣告\BibleQuote{陶工有權從同一團泥中，造一個貴重的器皿，另一個卑賤的器皿。}的那些话中，似乎又将一切归于天主——我们不应以为这些陈述是矛盾的，而是要将两种含义调和一致，并从两者中得出一个意义，即：我们既不应认为那些在我们自己能力内的事无需天主的助佑就能完成，也不应认为那些在天主手中的事无需我们的行为、意愿和意向的介入就能成就；因为我们没有能力如此地愿意或行任何事，以致不知道这个使我们能愿意或能做的能力本身是由天主赐予我们的，按照我们上述的区分。或者，当天主造器皿，一些为贵重，另一些为卑贱时，我们应当认为祂并非将我们的意愿、志向或功过视为贵贱的原因，仿佛它们是一种质料，祂从其中为我们每人造或贵重或卑贱的器皿；而其实是灵魂本身的运动，或理智的意向，可能自行向那知晓人心和思想意念的天主暗示，这器皿应当被造为贵重还是卑贱。但关于\OriginalTerm{自由意志}{freedom of the will}的问题，我们已经尽力讨论，这些就足够了。\par

% structure: p0030 source=h2 target=subsection reason=relative-heading-rank; base=h2

\subsection{希腊文译本}

% structure: p0031 source=h3 target=subsubsection reason=relative-heading-rank; base=h2

\subsubsection{第一章 论自由意志，兼对看似取消自由意志的圣经语句的解释和诠释。}

1. 由于教会的宣讲中包含关于天主公义审判的教义，这教义若被相信为真实，会激励听者度有德行的生活，并尽力避免罪恶，因为他们明确承认值得称赞和责备的事在我们自己的能力内，因此，让我们来单独讨论一些关于自由意志的要点——这是一个所有问题中最必要的。为了理解何谓自由意志，有必要阐明其概念，使这概念得到精确的界定后，课题就呈现在我们面前。\par

2. 在运动的事物中，有些自身内在拥有其运动的原因；另一些则仅从外部被推动。只有可移动的事物是从外部被推动的，例如木块、石头以及所有仅凭其构造而结合在一起的物质。且让我们把那将物体的流动称为运动的观点搁置一旁，因为它对于我们当前的目的并无必要。但动物和植物自身内在拥有其运动的原因，总之，一切凭借本性和灵魂而结合在一起的事物（据说金属也属于此类）也是如此。除此之外，火也是自动的，或许还有泉水。如今，在那些自身内在拥有其运动原因的事物中，据说有些是从自身之外被推动，有些则从自身之内被推动：无生命之物是从自身之外，有生命之物则是从自身之内。因为有生命之物是从自身之内被推动的，其中生起一种激起努力的\OriginalTerm{表象}{phantasy}。再者，在某些动物中，形成激起努力的表象，表象的本性有条理地激发努力，例如在蜘蛛中形成织网的表象；随后便尝试织网，表象的本性有条理地只激发这昆虫做此事。除了其表象本性之外，人们相信这昆虫别无所有。在蜜蜂中，则形成生产蜂蜡的表象。\par

3. 然而，理性的动物除了其\OriginalTerm{表象本性}{phantasial nature}之外，还拥有\OriginalTerm{理性}{reason}，它判断表象，拒绝一些而接受另一些，以便动物能按照它们被引导。因此，既然理性的本性中具有有助于沉思美德与恶习的手段，我们通过遵循它们，在观看善与恶之后，选择其一而避免另一，当我们致力于实践美德时，我们值得称赞，而行相反之事时则当受责备。然而，我们不可不知，万物所赋予的本性的大部分，在动物中具有程度不等的差异，或大或小；因此猎犬和战马的本能，可谓以某种方式接近于理性的能力。至于落入某一种在我们内里激起这样或那样表象的外在原因，显然并非取决于我们自身的事情；但决定我们以这种或那种方式使用所发生之事，则无非是我们内在理性的特权，理性随机会激发我们趋向于德性行为或转向其反面。\par

4. 然而，倘若有人主张，正是这种外在原因具有如此性质，以致它如此出现时无法抗拒，那么让他省察自己的情感和运动，看看是否有一个主导原则因某些貌似有理的论据而对某个对象表示赞同、同意和倾向。举例来说，一个决心守贞并避免肉体交合的男子，若有一个女子出现在他面前，激发他做出与其决心相悖的行为，这并非完全取消其决心的原因。因为此人完全沉溺于快感的奢侈和诱惑，不愿抵抗，也不愿持守其决心，便犯下了放纵之罪。另一人（同样情形发生在他身上，但他受过更多教导并操练了自己）确实遭遇了诱惑和勾引；但他的理性，因被提升至更高点并受精心训练，且在其对于德性行为的看法上得以坚固，或近乎坚固，便击退那激发，熄灭了欲望。\par

5. 情况既然如此，说我们从外部被推动，并通过宣称我们如同木块和石头一般被外在原因所拖曳，从而将责任推卸于自身之外，既不符合真理，也不合乎理性，而是那想要摧毁\OriginalTerm{自由意志}{free-will}概念之人的说法。因为如果我们向这样的人询问什么是自由意志，他会说它在于：当决意要做某事时，没有外在原因激发相反的事。然而，仅仅责备身体的本性则是荒谬的；因为\OriginalTerm{训导理性}{disciplinary reason}若抓住了那些最放纵和最野蛮的人，只要他们听从其劝诫，便会产生转变，以致向善的改变和变化极其广泛——最放荡的人往往变得比先前看似天性并非如此的人更好，最野蛮的人转变成如此温和的状态，以致那些从未如此野蛮的人相比之下显得野蛮，在他们内里产生了如此大的温良。我们又看见另一些最稳定可敬的人，因与恶习交往而从其可敬稳定的状态中被驱离，以致陷入放荡的习惯，往往在中年开始其邪恶，并在青春时期过后陷入混乱，而青春就其本性而言是不稳定的。因此，理性表明外在事件并不取决于我们，但以这种方式或相反方式使用它们是我们自身的事务，因为我们已接受理性作为法官和审查者，考察我们应当如何面对那些来自外部的事件。\par

6. 如今，我们理应度有德行的生活，而天主也对此有所要求，这并非依赖祂或任何其他事物，也不像有些人所想的那样依赖命运，而是我们自己的作为。先知\Person{米该亚}将证明这一点，他说：\BibleQuote{人哪，已通知了你，什么是善，或主要求你什么，无非是履行正义，爱好仁慈？}\Person{梅瑟}也：\BibleQuote{我已将生命之路和死亡之路摆在你面前：选择善事，并走在其中。}\Person{依撒意亚}也：\BibleQuote{如果你们愿意，并听从我，你们将享用地上的美物；但若你们不愿意，且不听从我，刀剑将吞灭你们：因为主的口说了。}在\BibleBook{}{圣咏}中：\BibleQuote{如果我的子民听从了我，以色列行走在我的道路上，我必将他们的敌人贬抑为无，并伸手攻击那些压迫他们的人；}这表明祂的子民有能力听从并行走在天主的道路上。救主也命令说：\BibleQuote{但我告诉你们：不要抵抗恶人；}以及：\BibleQuote{凡向弟兄发怒的，应受审判；}以及：\BibleQuote{凡注视妇女而有意贪恋她的，已在心中与她犯了奸淫；}以及祂所赐的任何其他诫命，都表明遵守所命令的事在于我们自己，而且我们若违犯，将理所当然地应受谴责。因此祂进一步说：\BibleQuote{凡听我的话而实行的，好比一个精明的人，把自己的房屋建在磐石上，}等等，等等；\BibleQuote{凡听我的话而不实行的，好比一个愚昧的人，把自己的房屋建在沙土上，}等等。当祂对右边的人说：\BibleQuote{我父所祝福的，来吧！}等等；\BibleQuote{因为我饿了，你们给了我吃的；我渴了，你们给了我喝的，}这极为明显地表明，祂向这些值得称赞的人赐予许诺。但与此相反，对另一些人，作为应受谴责的，祂说：\BibleQuote{可咒骂的，离开我，到永火里去！}我们也看看\Person{保禄}如何对我们谈论自由意志，以及我们自己是如何成为毁灭或得救的原因，他说：\BibleQuote{难道你轻视天主丰厚的仁慈、容忍和忍耐，不知道天主的仁慈是引领你悔改吗？但你可凭着你的硬心和不悔改的心，为自己积蓄忿怒，直到忿怒的日子，以及天主公义审判显现的日子；祂要照各人的行为报应各人：凡恒心行善，寻求荣耀、尊贵和不朽的，就有永生；但那好争论、不信从真理、却信从不义的人，就有忿怒、恼恨、患难、困苦，加给一切作恶的人，先是犹太人，后是希腊人；但荣耀、尊贵与平安，加给一切行善的人，先是犹太人，后是希腊人。}确实，圣经中有无数章节极其清楚地确立了自由意志的存在。\par

7. 然而，由于旧约和新约中的某些说法导致相反的结论——即遵守诫命而得救，或违反诫命而丧亡，并不在于我们自己——让我们逐一列举这些说法，并查看对它们的解释，以便从我们所列举的之中，任何人以类似方式选出所有看似否定自由意志的段落，并考虑对它们所作的解释。现在，关于\Person{法郎}的陈述困扰了许多人，天主曾多次论及他：\BibleQuote{我要使法郎的心硬。}因为如果他是由天主所硬化，并因硬化而犯罪，那么他就不是自己犯罪的原因；如果是这样，那么法郎也没有自由意志。有人会说，同样地，那些丧亡的人也没有自由意志，他们不会自己丧亡。\Person{厄则克耳}中的宣告也：\BibleQuote{我要除去他们的石心，给他们换上肉心，使他们遵行我的法令，遵守我的诫命，}这可能让人以为是天主赐予力量遵行祂的诫命，并遵守祂的法令，因为祂除去了障碍——石心，并植入更好的——肉心。也让我们看看福音中的段落——救主对那些问祂为什么用比喻对群众讲话的人的回答。祂的话是：\BibleQuote{他们看是看，却看不见；听是听，却不明白；免得他们悔改，并得到罪的赦免。}还有\Person{保禄}中的段落：\BibleQuote{这既不在于那愿意的，也不在于那奔跑的，而在于那显示仁慈的天主。}其他地方的宣告也：\BibleQuote{愿意和行事都是出于天主；} \BibleQuote{祂愿意怜悯谁，就怜悯谁；愿意使谁心硬，就使谁心硬。你也许会说：那么，祂为什么还要指责呢？谁能抗拒祂的意愿呢？} \BibleQuote{说服来自那召叫的，而不来自我们。} \BibleQuote{不，人哪，你是谁，竟敢与天主争辩？被造的能对那造它的说：你为什么这样造我？难道陶匠没有权柄用同一团泥，造一个贵重的器皿，又造一个卑贱的器皿吗？}现在，这些段落本身就足以困扰大众，好像人没有自由意志，而是天主随祂的意愿拯救和毁灭。\par

8. 那么，让我们从关于法老所说的话开始——他由天主心硬，以致他不释放百姓；同时也要审视宗徒的这句话：\BibleQuote{因此，祂愿意怜悯谁，就怜悯谁；愿意使谁心硬，就使谁心硬。}一些持不同意见者误用这些段落，他们自己也几乎摧毁了\OriginalTerm{自由意志}{free-will}，因为他们引入败坏的、不能得救的性体，以及另一些得救的、不可能失落的性体。他们说，法老由于是败坏的性体，因此被天主心硬，天主怜悯属神的，却使属土的心硬。现在让我们看看他们的意思。我们要问他们，法老是否属于属土的性体；当他们回答时，我们要说，属土的性体的人全然不服从天主；但若不服从，又何需他的心被心硬，且不止一次，而是多次呢？除非，既然他可能服从（若是如此，当被征兆和奇事紧迫时，他必定会服从，因为他不是属土的），天主需要他更大程度地不服从，以便为众人的得救彰显祂的大能作为，因此使他的心变硬。这首先将是我们的回答，为推翻他们关于法老属于败坏的性体的假设。关于宗徒的那句话，也必须给他们同样的答复。因为天主使谁心硬？是那些丧亡的人，好像他们若非被心硬就会服从？还是显然那些要得救的人，因为他们不属于败坏的性体？祂又怜悯谁？是那些要得救的人吗？那些一次就被预备好得救、且因自己的性体必将成为有福的人，何需第二次的怜悯呢？除非，既然他们可能招致毁灭，若未得到怜悯，他们将获得怜悯，以免招致他们可能遭受的毁灭，而处于得救之人的状态。这就是我们对这些人的回答。\par

9. 但对于那些自以为理解\Quote{心硬}一词的人，我们必须提出询问：他们说天主借着祂的作为使人心硬，是什么意思？祂这样做的目的是什么？因为让他们观察一位实际上既公义又良善的天主的概念；但如果他们不承认这一点，暂且向他们让步，承认祂是公义的；并让他们表明，那良善而公义的天主，或仅仅公义的天主，如何显得公义，在使那因心硬而丧亡的人心硬时；以及当人因他们的心硬和不服从而被祂管教时，公义的天主如何成为毁灭和不服从的原因。祂又为何责备他，说：\BibleQuote{你不让我的百姓去}；\BibleQuote{看哪，我要击杀埃及中所有头生的，连你的头生也在内}；以及凡借着\Person{梅瑟}干预而记载为天主对法老所说的一切话呢？因为凡相信圣经是真实的，并相信天主是公义的人，若他诚实，就必须努力表明天主在使用这些表达时，如何能被清楚地理解为公义的。但如果有人站着，露着头宣称世界的造物主倾向于邪恶，我们就需要用别的话来回答他们。\par

10. 但由于他们说他们视祂为公义的天主，而我们视祂为既是良善又是公义的天主，就让我们思考，这位良善而公义的天主如何能使法老的心变硬。那么，看看我们能否借着宗徒在\BibleBook{}{希伯来书}中所用的比喻来证明：借着同一种作为，天主怜悯一个人，同时使另一个人心硬，尽管祂并非有意使之心硬；而是（尽管）怀着良善的目的，心硬却作为此类人内在的邪恶原则的结果而随之发生，因此祂被称为使那心硬者心硬。他说：\BibleQuote{一块田地，时常吸收降在其上的雨水，并为那些耕种它的人长出合用的菜蔬，就从天主获得祝福；但那长出荆棘和蒺藜的，就被废弃，近于诅咒，它的结局就是焚烧。}因此，就雨水而言，只有一种作为；既然雨水只有一种作为，被耕种的土地就结出果实，而被忽略且贫瘠的土地则长出荆棘。现在，对于降雨的那位来说，说\Quote{我产生了果实和地里的荆棘}，似乎是亵渎；然而，尽管是亵渎，却是真实的。因为，若没有降雨，就不会有果实或荆棘；但若在适当的时候适度降下，两者就都产生了。如今，那吸收了常常降在其上的雨水，却仍长出荆棘和蒺藜的土地，就被废弃，近于诅咒。因此，雨水的祝福甚至降在了劣等土地上；但后者因被忽略和未耕种，就生出荆棘和蒺藜。因此，同样地，天主所做的奇事，就好像雨水；而不同的心意，就好像被耕种和被忽略的土地，二者都像土地一样，属于同一个性体。\par

11. 正如太阳发出声音说：\Quote{我融化又烘干}，融化和烘干是相反的事，但就这一点而言，他并未说谎；因为蜡因同一热量而熔化，泥巴被烘干。同样，借着梅瑟所行的同一作为，一方面证明了法郎的顽固——这是他邪恶的结果，另一方面也证明了那些跟随希伯来人离开的混杂埃及人群的屈服。简短的说法是，当法郎说：\BibleQuote{但你们不可走远：你们要走三天的路程，留下你们的妻子}，以及其他他在征兆面前逐渐让步所说的话，证明了奇迹也对他产生了一些影响，但并未完成一切（它们所能完成的）。然而，如果那多数人所认为的——法郎心硬——是由天主亲自造成的，那么连这些也不会发生。按照通常用法来软化这类表达并不荒谬：因为好主人常常对他们的奴隶说，当这些奴隶被主人的仁慈和宽容宠坏时：\Quote{是我使你变坏，如此重大的过失应归咎于我。} 因为我们必须注意这句话的性质和力量，而不应诡辩地争论，忽略表达的意义。因此，保禄清楚地研究了这些要点后，对罪人说：\BibleQuote{或者你藐视祂丰厚的慈爱、宽容和忍耐，不晓得天主的恩慈是领你悔改吗？但你竟任着你刚硬不悔改的心，为自己积蓄忿怒，直到忿怒的日子，并天主公义审判启示的日子。} 现在，让宗徒对罪人说的话也向法郎说，那么对他所作的宣告就被理解为特别合适，他是按照自己的刚硬和不悔改的心，为自己积蓄忿怒；因为除非有奇迹发生，而且是这样重大和重要的奇迹，他的刚硬就不会被证明或显明出来。\par

12. 但由于这样的叙述难以获得赞同，且被认为是牵强的，让我们也从先知的宣告中看看，那些虽然经历了天主的极大仁慈却未过道德生活、后来犯罪的人说了什么。\BibleQuote{上主，祢为何使我们远离祢的道路？祢为何使我们的心刚硬，以致不敬畏祢的名？求祢为祢仆人、为祢产业支派的缘故回转，使我们得以继承祢圣山的一小块地方。} 以及在耶肋米亚中：\BibleQuote{上主，祢欺骗了我，我被欺骗了；祢坚强，祢得胜了。} 因为那些恳求怜悯的人所说的\BibleQuote{祢为何使我们的心刚硬，以致不敬畏祢的名？}，其性质如下：\Quote{祢为何长久宽恕我们，不因我们的罪而惩罚我们，反而离弃我们，直到我们的过犯达到极点？} 如今祂让大多数人不受惩罚，一方面是为了让每个人的习性得以考验，只要这取决于我们自己；另一方面是为了让有德之人因所受的考验而显明；而其他人，虽然不被天主忽略——因为祂在万物存在之前就知晓一切——但从理性的受造物和他们自身而言，他们以后可以获得医治的方法，因为他们若不谴责自己，就不会知道这益处。对每个人有益的是，他能觉察自己的本性和天主的恩宠。因为一个人若不察觉自己的软弱和天主的恩惠，虽然他得到了益处，但未经验自己，也没有谴责自己，就会以为那由天上恩宠所赐的恩惠是他自己的成就。这种想象也会导致虚荣，成为跌倒的原因；我们设想，魔鬼就是这种情况，他把自己在无罪状态时所拥有的优先地位归功于自己。\BibleQuote{因为凡高举自己的，必被贬抑；} 而\BibleQuote{凡谦卑自己的，必被高举。} 请注意，正是为此，天主的事向明智和聪明的人隐藏了，正如宗徒所说：\BibleQuote{使一切有血肉的，在天主面前一个也不能自夸。} 而向婴孩启示了，即那些在童年之后走向更好事物的人，他们记得自己达到最大可能的幸福，与其说是出于自己的努力，不如说是出于（天主的）不可言说的美善。\par

13. 因此，那被遗弃的人并非无理由地被遗弃于天主的审判，而天主对某些罪人长期忍耐，是因为这对他们有益，考虑到灵魂的不死和永无尽期的世界，使他们不至于过快地被带入救恩状态，而是更缓慢地被引导，在经历许多邪恶之后。因为正如那些能够快速治愈一个人的医师，当他们怀疑身体中有隐藏的毒物时，反而做与治愈相反的事，通过他们治愈的愿望本身使这更加确定，认为最好在相当长的时间内让病人处于发炎和疾病之中，以便他更可靠地恢复健康，而不是显得快速治愈，而后导致复发，并且（因此）仓促的治愈只能持续一时；同样地，天主也知道人心中的隐秘，预见未来的事，在祂的长期忍耐中，允许（某些事发生），并通过外在发生的事拔出隐秘的邪恶，以便洁净那因疏忽而接受了罪之种子的人，使他在罪恶浮出表面时将其吐出，尽管他可能已深陷于邪恶，之后他可以在邪恶之后获得医治并被更新。因为天主治理灵魂，我可以说，不是就现世生命的五十年而言，而是就无限的时代而言：因为祂使思维原理本性不死，并与其自身相近；而理性灵魂并不像在此生一样被排除在医治之外。\par

14. 现在来吧，让我们使用福音中的以下比喻。有一块岩石，上面有少许表层土壤，种子落在上面，就迅速发芽；但发芽后，因为没有根，太阳升起时就晒焦枯干了。这块岩石就是人的灵魂，因疏忽而硬化，因邪恶而变成石头；因为没有人从天主那里领受一颗石心，而是因邪恶变成这样的。如果有人责怪农夫没有更早地将种子撒在岩石地上，当他看到其他岩石地接受了种子而生长茂盛时，农夫会回答说：\Quote{我要更缓慢地播种这地，撒下能够抓牢的种子，这种较慢的方法对这地更好，并且比那更快地接受种子、更表层的方法更可靠。} （那责怪的人）会赞同农夫，认为他说话合理，行事有技巧：同样，万物的伟大农夫也推迟那可能被视为过早的恩惠，以免浮于表面。但很可能有人会对此提出反对：\Quote{为什么有些种子落在仅有表层土壤的地上，而这灵魂就好比一块岩石？} 对此我们必须回答说：对这个灵魂来说，它渴望更好事物，但操之过急，且不是通过正确的途径，那么实现它的渴望更好，以便它因此自责，在很长一段时间后，才能忍受合乎自然的耕作。因为灵魂，可以这样说，是无数的；它们的情操、运动、目的、攻击和努力也是无数的，唯一的出色管理者知道季节、合适的帮助、途径和方法，就是万物的天主和父，祂知道祂如何借着如此重大的事件，以及借着在海中淹死，引导法老；祂对法老的眷顾并未因此终止。因为当他被淹死时，他并没有被消灭：\Quote{因为我们的所有和我们的言语，都在天主手中；一切智慧和对工艺的认识也都在祂手中。} \BibleRef{智 7:16} 这就是对以下陈述的一个适度的辩护：\Quote{法老的心变硬了}，以及\Quote{天主怜悯祂所愿意怜悯的，并使祂所愿意的硬心。} \BibleRef{出 33:19；罗 9:18}\par

15. 让我们也看看厄则克耳中的宣告，其中说：\BibleQuote{我要除去他们石心，给他们换上肉心，使他们遵行我的律例，谨守我的诫命。}因为如果天主愿意时，就除去石心，植入肉心，使人遵行祂的诫命和遵守祂的命令，那么除去邪恶就不在我们能力之内。因为除去石心无非就是除去那使人硬心的邪恶，从天主愿意除去它的人身上除去；而植入肉心，使人遵行天主的诫命，遵守祂的命令，这岂不就是要变得对真理稍微柔顺、不抗拒，并且能够实践德行吗？如果天主应许这样做，而祂除去石心之前，我们不自己放下石心，那么显然，除去邪恶不取决于我们自己；如果我们对在我们内产生肉心不做任何事，而是天主的作为，那么合乎德行地生活就不是我们自己的行为，而完全是\OriginalTerm{天主的恩宠}{divine grace}的结果。那些只从（圣经的）字句就否定自由意志的人会这样说。但我们要回答，说：我们应该这样理解这些段落：即，一个人，例如，他碰巧是愚昧无知的，在认识到自己的缺陷之后，或因老师的劝诫，或因其他方式，自愿把自己交付给他认为能引导他接受教育和德行的人；当他这样交付自己时，老师应许除去他的无知，植入教导，这并不是说他对自己的训练和避免无知毫无贡献，而是因为他把自己带来接受医治；同样，天主的圣言应许除去邪恶——它称之为石心——从那些来到它面前的人，不是若他们不愿意，而是（只有）若他们将自己交托给病人的医生，正如在福音中，病人来寻求救主，请求得到医治，就这样被治愈了。而且，我要说，盲人恢复视力，就他们的请求而言，是那些相信自己能被医治的人的行为；但就视力的恢复而言，是我们救主的工作。因此，天主的圣言应许，向那些来到它面前的人植入知识，通过除去那石心与硬心——即邪恶——好使人遵行天主的诫命，并遵守天主的命令。\par

16. 此后又有福音中的段落，救主说，祂用比喻对那外面的人说话，是为叫\BibleQuote{他们看是看见，却不了解；听是听见，却不明白；免得他们悔改，罪得赦免。}现在，我们的反对者会说：\Quote{如果某些人确凿无疑地因听到更清晰的言语而悔改，以致配得赦罪，而听到这些更清晰的言语并不取决于他们自己，而是取决于教导的人，而祂为此缘故没有更明确地向他们宣告，免得他们看见并明白，那么这样的人得救就不在他们能力之内；如果是这样，我们就不能拥有关于得救与灭亡的自由意志。}对此等论证，本可做出有效的反驳，若不是加上了一句：\BibleQuote{免得他们悔改，罪得赦免}——就是说，救主不愿那些将来不成为善良和有德行的人明白祂教导中更奥妙的部分，为此缘故用比喻对他们说话；但现在，因为\BibleQuote{免得他们悔改，罪得赦免}这句话，辩护就更困难了。首先，我们必须注意这段经文对异端者的意义，他们从旧约中搜寻那些段落，其中——如他们自己大胆断言的那样——显出世界创造者在报复和惩罚恶人方面的残忍，或无论他们愿意用什么名称来指称这种品质，他们这样说只是为了说良善不存在于创造者中；而他们对待新约的方式却不是这样，也不是本着坦诚的精神，而是忽略了与新约中他们认为该受谴责的旧约段落相似的地方。因为显而易见，按照福音，救主如他们所断言的那样，通过祂先前的话，显然不为此缘故清楚说话，好使人不悔改，且悔改后不至于成为配得赦罪的人：这个陈述本身并不逊于旧约中被反对的那些段落。如果他们设法为福音辩护，我们必须问问他们，他们对相同的问题采取不同处理方式，难道不该受责备吗？他们不因新约绊倒，反而设法为其辩护，却对旧约中类似的地方提出控告，而他们本应以同样方式为新约中的段落辩护。因此，我们将迫使他们，因为相似之处，将所有著作视为同一位天主的著作。来，让我们尽我们所能，对提交给我们的问题给出一个答案。\par

我们在研究\OriginalTerm{法郎}{Pharaoh}的题目时也曾肯定，有时候快速的治愈对受治者并非有利，因为如果他们被麻烦的疾病缠身后又轻易摆脱了纠缠他们的病患，就会轻视这种容易治愈的恶，不再警惕再次陷入其中，因而再次被卷入。因此，对于这样的人，永恒的天主，那洞察隐秘、在一切存在之前就知晓一切的主，依照他的美善，延迟赐予他们更迅速的帮助，并且，可以这样说，在帮助中不帮助，后者对他们更为有利。那么，很可能，我们正在谈论的所谓\Quote{外面的人}，按照我们的推测，已被\OriginalTerm{救主}{Saviour}预先见到不会在\OriginalTerm{悔改}{repentance}中保持稳定，因此如果他们更清楚地听到所说的话，就会被救主如此对待：不能清楚地听见他教导中更深奥的事，以免在迅速悔改并通过获得罪赦而痊愈后，轻视自己邪恶的伤口为轻微易治，再次迅速滑入其中。也许他们还要为他们先前违背美德所犯的过犯受罚，当他们离弃美德时，尚未填满（全部）时间；以便在神圣眷顾离弃他们，并被自己所播种的邪恶充满更多之后，他们后来能被召唤到更稳固的悔改；以免再次迅速陷入他们先前在傲慢对待美德要求并投身于更坏事物时所卷入的那些邪恶。因此，那些被称为\Quote{外面的人}（显然是与\Quote{里面的人}比较而言），他们离\Quote{里面的人}并不很远；当\Quote{里面的人}清楚听见时，他们自己则模糊地听见，因为对他们说话的是比喻；但他们确实听见了。然而，在那些被称为\Quote{外面的人}中，另有一些被称为\OriginalTerm{提洛人}{Tyrians}，虽然救主被预知如果他们临近他的边界，他们早就坐在麻衣和灰烬中悔改了，但他们甚至听不到那些被\Quote{外面的人}所听见的话（很可能在功德上远逊于那些\Quote{外面的人}），以便在另一个季节，当他们比那些没有接受圣言的人（其中他也提到了提洛人）更可容忍时，在更适当的时候听到圣言后，能获得更持久的悔改。但请你观察，除了我们寻求探究（真理）的愿望之外，我们是否不更努力地在每件事上对天主和他的基督保持虔诚的态度，因为我们千方百计证明，在如此伟大和如此特殊、涉及天主各种\OriginalTerm{眷顾}{providence}的事上，他关怀不朽的灵魂。确实，如果有人要追问那些被反对的事：为什么那些看见奇迹、听见圣言的人没有受益，而提洛人却会悔改，如果这些事行在他们中间、说给他们听；并应问说：为什么救主向这些人宣讲这些事，反而有害于他们，使他们的罪被算为更重？我们必须回答这样的人：那理解所有指责他眷顾之人的倾向——（声称）正是由于眷顾他们才不信，因为眷顾没有允许他们看见它使别人看见的事，也没有安排他们听见那些使别人听了受益的圣言——愿意证明他们的辩护没有理性根据，他赐予了那些指责他治理的人所要求的恩惠；以便在他们获得之后，仍被证明有最大的不虔敬，因为他们即使在那时也没有让自己受益，并停止这样的狂妄；并在这个具体点上获得自由后，得以知道天主在赐予某些人恩惠时，有时会延迟和拖延，不赐予看见和听见那些事的恩宠，而这些事在被看见和听见后，会使那些在如此伟大特殊的事之后仍不信的人的罪变得更重、更严重。\par

18. 现在我们来看这段经文：\BibleQuote{这样看来，不是由于愿意的，也不是由于奔跑的，而是由于显示慈悲的天主。}因为那些挑剔的人说：如果\BibleQuote{不是由于愿意的，也不是由于奔跑的，而是由于显示慈悲的天主，}那么救恩就不取决于我们自己，而取决于那造我们如此的那位的安排，或取决于那按自己旨意显示慈悲者的意愿。现在我们必须向这些人提出以下问题：渴望善事是美德还是恶行？而渴望奔跑以在追求善事中达到目标是值得称赞还是应受责备？如果他们说是应受责备的，他们就会给出荒谬的回答；因为圣人们渴望并奔跑，显然在此行动中并未做任何可责备之事。但如果他们说渴望善事、追求善事是美德，我们就要问他们：一个趋于败坏的本性如何渴望更好的事物？因为这就像一棵坏树结出好果子，因为渴望更好的事物是一种美德行为。他们也许会给出第三种回答：渴望并追求善事是无关紧要的事情之一，既不美好也不邪恶。对此我们必须说：如果渴望并追求善事是无关紧要的事，那么其反面也是无关紧要的事，即渴望并追求邪恶。但渴望并追求邪恶并不是无关紧要的事。因此，渴望并追求善事也不是无关紧要的事。因此，这就是我认为可以为\BibleQuote{不是由于愿意的，也不是由于奔跑的，而是由于显示慈悲的天主}这一说法作出的辩护。\Person{撒罗满}在\BibleBook{}{圣咏集}（因为\WorkTitle{上行之歌}是他的作品，我们将从中引用）中说：\BibleQuote{除非主建造房屋，建造的人徒然劳作；除非主守护城市，守夜的人徒然警醒：}这并不是劝阻我们建造，也不是教导我们不要守护我们灵魂中的城市，而是表明：凡不依靠天主而建造、且未从他那里得到守护的，是徒然建造、徒然警醒，因为天主理当被称为建造之主，万有之治理者，城市守护之统御者。因此，如果我们说这样的建筑不是建造者的工作，而是天主的；又说这样的城市之所以未受敌人侵害，不是由于守夜者的成功努力，而是由于统御一切的天主，那么我们并不为错，因为理解到人类的手段也起了一些作用，但应将益处感恩地归于使其完成的天主。同样，既然（单纯的）人类渴望不足以达到目的，那些如同运动员般奔跑的人并不能赢得天主在基督耶稣内召叫的奖赏——因为这些事是靠天主的助佑完成的——所以，说\BibleQuote{不是由于愿意的，也不是由于奔跑的，而是由于显示慈悲的天主}是恰当的。正如关于农事所记载的：\BibleQuote{我栽种了，\Person{阿颇罗}浇灌了，但使之生长的是天主。所以，栽种的无足轻重，浇灌的也无足轻重，唯有使之生长的天主才是一切。}我们不能敬虔地断言丰收是农夫或浇灌者的工作，而是天主的工作。同样，我们自身的成全也不是好像我们自己什么都不做；因为它不是由我们完成的，而是天主成就了其中大部分。为使这一论断更清楚地被相信，我们将从航海术中取一个比喻。因为与风的作用、空气的温和、星辰的光亮——这一切共同合作保证船员的安全——相比，航海术在将船引入港口中所占的比重能有多大呢？既然连水手们自己出于虔诚，也往往不敢断言他们救了船，而是将一切归于天主；并非好像他们什么都没做，而是因为天主眷顾所成就的远大于他们技艺所成就的。在我们得救的事上，天主所做的大于我们自己所做的无限倍；因此，我认为，才说\BibleQuote{不是由于愿意的，也不是由于奔跑的，而是由于显示慈悲的天主}。因为如果以他们所想象的方式解释\BibleQuote{不是由于愿意的，也不是由于奔跑的，而是由于显示慈悲的天主}这一陈述，那么诫命就是多余的；\Person{保禄}自己责备一些人堕落、称赞另一些人保持正直并为众教会制定律法也是徒然的；我们投身于渴望更好的事物也是徒然的，努力奔跑也是徒然的。但是，\Person{保禄}给出这样的劝告，责备一些人，称赞另一些人，并不是徒然的；我们投身于渴望更好的事物，追求卓越的事物，也不是徒然的。因此，他们没有很好地解释这段经文的含义。\par

19. 此外，还有一段经文说：\BibleQuote{愿意的和实行的，都是出于天主。}有些人断言：如果愿意和实行都出于天主，那么无论我们愿意恶事或实行恶事，这些（行动）都来自天主；若是这样，我们就没有自由意志。但另一方面，当我们愿意更好的事、实行更卓越的事时，既然愿意和实行都来自天主，那么就不是我们做了更卓越的事，而只是我们显似（做了），其实是天主赐予的；因此即使在这一点上，我们也没有自由意志。对此我们必须回答：宗徒的话并不断言愿意恶事是出于天主，或愿意善事是出于祂（同样地，实行较好或较差之事也是如此）；而是说一般意义上的愿意和一般意义上的奔跑（是出于祂）。因为正如我们从天主获得（属性）成为生物和人类，同样我们也获得一般意义上的愿意，以及可说是一般意义上的运动。并且，正如我们拥有生命和运动以及活动的能力（例如这些肢体，手或脚），我们却不能正确地说从天主获得了这种运动的形式，即我们用来打击、破坏或夺取他人财物的运动；而只是从祂接受了单纯的运动的能力，我们用它来达到较好或较坏的目的；同样，我们从天主获得了行动的能力（就我们是生物而言），并从造物主获得了愿意的能力，而我们运用愿意的能力以及行动的能力来追求最高贵或相反的目标。\par

20. 然而，宗徒的宣告似乎会将我们引向没有自由意志的结论；他在反驳自己时说道：\BibleQuote{因此，祂愿意恩待谁，就恩待谁；愿意使谁心硬，就使谁心硬。你会对我说：那么祂为什么还要指责人呢？有谁能抗拒祂的旨意呢？但是人哪，你是谁，竟敢向天主抗辩？受造之物岂能对造它的说：你为什么这样造了我？难道陶匠没有权柄从同一团泥中，造一个器皿作贵重的，另一个作卑贱的吗？}因为有人会说：如果陶匠用同一团泥造一些器皿作贵重，另一些作卑贱，而天主也这样造一些人得救，另一些人丧亡，那么得救或丧亡就不取决于我们自己，我们也没有自由意志。现在我们必须向这样解释这些经文的人提问：是否能想象宗徒自相矛盾？我想，没有人敢这么说。那么，如果宗徒没有自相矛盾，按照这种理解，他怎能合理指责格林多那位犯淫乱的人，或那些堕落、没有为所行的放荡和污秽悔改的人呢？他又怎能祝福那些他所称赞行为良好的人，就如他这样称赞敖尼斯佛洛的家：\BibleQuote{愿主赐仁慈给敖尼斯佛洛的家，因为他屡次使我精神愉快，不以我的锁链为耻；而且他一到罗马，就殷勤地寻找我，并找到了我。愿主在那日使他获得主的仁慈。}同一宗徒既责备罪人当受指摘，又称赞行善者当受赞许，但另一方面又说似乎不取决于我们自己，而是造物主造成了贵重的器皿和卑贱的器皿，这是不一致的。那么这句话又怎会正确：\BibleQuote{因为我们众人必须出现在基督的审判台前，为叫各人按他在肉身内所行的，或善或恶，领取相当的报应。}既然作恶的人因为被造为卑贱的器皿而达到这种邪恶程度，而行善的人因为从起初就被造为此目的而成为贵重的器皿，所以行善？另外，另一处陈述岂不也与这些人从我们引用的经文所得出的观点相冲突（即创造主应负责一个器皿贵重、另一个卑贱）？它说：\BibleQuote{在大户家庭中，不仅有金器银器，也有木器和瓦器；有作贵重的，也有作卑贱的。所以人若自洁，脱离卑贱的，成为贵重的，圣洁的，合主人使用的，预备行各样的善事。}因为如果自洁的人成为贵重的器皿，而任由自己不洁的人成为卑贱的器皿，那么就这些话而言，创造主完全没有责任。因为创造主造贵重的器皿和卑贱的器皿，并非从起初依据祂的预知，因为祂不凭预知事先定罪或称义；而是那些自洁的人成为贵重的器皿，那些任由自己不洁的人成为卑贱的器皿：这样，在形成贵重和卑贱器皿时，就有更古老的原因（在起作用），使得一个被造为前者，另一个被造为后者。但如果我们一旦承认在形成贵重器皿和卑贱器皿时存在着某些更古老的原因，那么回到灵魂的主题上，并（设想）雅各伯被爱、厄撒乌被恨有一个更古老的原因，这原因在雅各伯披上身体之前就已存在，在厄撒乌于黎贝加腹中成胎之前也已存在，这又有什么荒谬之处呢？\par

21. 同时，清楚显示的是，就底层本性而言，正如在陶匠手中的同一块泥，从中制成器皿，有的为尊贵，有的为卑贱；同样，每一个灵魂的本性都在天主手中，可以说，只有一个理性的团块，某些更古老的原因导致一些被造成尊贵的器皿，另一些被造成卑贱的器皿。但若宗徒的话含有责备，当他说：\Quote{然而，人啊，你是谁，竟敢向天主抗辩？}这教导我们，那在天主面前有自信、忠信、且生活有德行的人，不会听到\Quote{你是谁，竟敢向天主抗辩？}这样的话。例如梅瑟这样的人，\Quote{因为梅瑟说话，天主就发声回答他；}正如天主回答梅瑟，圣人也这样回答天主。但那没有这种自信的人，显然，要么因为他失去了它，要么因为他探究这些问题不是出于求知之爱，而是出于找错之心，因此说：\Quote{为什么祂还要责备？有谁能抗拒祂的旨意呢？}他就该受责备的话：\Quote{然而，人啊，你是谁，竟敢向天主抗辩？}\par

现在，对那些引入不同本性，并利用宗徒的声明来支持他们观点的人，我们必须作如下回答。如果他们主张那些丧亡者和那些得救者是由同一块泥形成的，并且得救者的创造者也是丧亡者的创造者，而祂是良善的，祂创造的不只是属灵的本性，还有属土的本性（因为这从他们的观点得出），那么，仍然有可能的是：那因先前某些义行而成为尊贵器皿的人，若以后没有类似行为，也没有做与尊贵器皿相称的事，在另一个世界就变成了卑贱的器皿；另一方面，那因比此生更古老的原因在此世成为卑贱器皿的人，在改良后，在新创造中可能成为\Quote{尊贵的器皿，圣洁的，合主人使用的，预备行各样的善事。}也许那些现在是以色列人，却生活得不配其出身，他们将失去其地位，如同从尊贵的器皿变为卑贱的器皿；许多现在的埃及人和依杜默亚人，他们接近以色列人，当他们结出更多果实的时候，将进入主的教会，不再被算作埃及人和依杜默亚人，而成为以色列人。因此，根据这个观点，是由于他们（不同的）意向，有些人从较坏的状态进步到较好的状态，另一些人从较好堕落到较坏；还有其他一些人，保持在一个有德行的进程中，或从善进步到更善；而另一些人则相反，停留在邪恶的进程中，或从坏变得更坏，随着他们的恶行流淌。\par

22. 但是，由于宗徒在一处并不把成为尊贵或卑贱的器皿归于天主，而是将全部归回给我们自己，说：\Quote{人若自洁，就必成为尊贵的器皿，圣洁的，合主人使用的，预备行各样的善事；}在别处又不把它归于我们自己，而似乎将全部归给天主，说：\Quote{陶匠有权从同一块泥中，造一个器皿为尊贵，另一个为卑贱；}既然他的陈述并不矛盾，我们必须调和它们，并从两者中提取一个完整的陈述。既不是我们自己的能力，离开天主的认识，迫使我们进步；也不是天主的认识（这样做），除非我们自己也为好的结果贡献一些东西；也不是我们自己的能力，离开天主的认识和使用与我们相称的能力的运用，使人成为尊贵或卑贱的器皿；也不是天主的意志单独将人塑造成尊贵或卑贱，除非祂将我们的意志视为一种容许变化的质料，倾向于更好或更坏的行为进程。这些观察就我们对自由意志的论述而言已经足够了。\par

% structure: p0055 source=h3 target=subsubsection reason=relative-heading-rank; base=h2

\subsubsection{第二章：论敌对势力}

1. 现在我们须依照圣经的记载，注意敌对势力或魔鬼本人如何与人类争斗，煽动并诱使人犯罪。首先，在 \BibleBook{}{创世纪}中，蛇被描述为诱惑了厄娃；关于此事，在名为 \WorkTitle{梅瑟升天录} 的作品（宗徒犹达在其书信中提及的一篇短论）中，总领天使弥额尔与魔鬼争论梅瑟遗体时，说蛇受魔鬼鼓动，是亚当和厄娃犯罪的原因。此事也引起一些人的探究，即：那从天上向亚巴郎说话的天使是谁，他说：\Quote{现在我知道你敬畏天主，且为我的缘故没有顾惜你的爱子，即你所爱的。}因为显然他被描述为一位天使，说他那时知道亚巴郎敬畏天主，并如圣经所载，没有顾惜他的爱子；虽然他没有说亚巴郎这样做是为了天主，而是为了他，即说话者的缘故。我们也必须查明，在 \BibleBook{}{出谷纪}中记载的那个想要杀死梅瑟（因为他正启程往埃及去）的是谁；后来，那被称为毁灭天使的又是谁；以及在 \BibleBook{}{肋未纪}中被称为\Quote{阿颇蓬沛乌斯}，即驱逐者的那一位，圣经论及他说：\BibleQuote{一阄归于上主，一阄归于阿颇蓬沛乌斯，即驱逐者。}在 \BibleBook{撒慕尔}{纪上}，也说有恶神窒息撒乌耳；在 \BibleBook{列王}{纪上}，先知米加雅说：\BibleQuote{我看见以色列的上主坐在祂的宝座上，天上的万军侍立在祂左右。上主说：谁去诱骗以色列王阿哈布，使他上去倒在辣摩特基肋阿得呢？这个这样说，那个那样说。随后有一个神出来，站在上主面前说：我去诱骗他。上主问他说：用什么方法？他说：我去，在他所有先知的口中成为说谎的神。上主说：你必能诱骗他，且会成功；你去，快这样做。因此，上主已将说谎的神放在你所有这些先知的口中；上主已说了关于你的凶言。}从最后这段引文清楚表明，有一个神出于自己的（自由）意志和选择，选定去诱骗（阿哈布）并施行谎言，以便上主误导君王走向死亡，因为他理应受苦。在 \BibleBook{编年}{纪上}中也说：\BibleQuote{撒殚，即魔鬼，起来反对以色列，怂恿达味统计人民。}在 \BibleBook{}{圣咏}中，也提到恶天使骚扰某些人。在 \BibleBook{}{训道篇}中，撒罗满说：\BibleQuote{如果长官的神气起来反对你，不要离开你的岗位，因为坚毅能制止许多过犯。}在 \BibleBook{}{匝加利亚}中，我们读到魔鬼站在耶叔亚右边，反抗他。依撒意亚说，上主的剑兴起攻击那龙，那弯曲的蛇。至于厄则克耳，在他的第二个神视中，他毫不含糊地预言了提洛王子的敌对势力，并说那龙住在埃及的河中。不仅如此，整部 \BibleBook{}{约伯传}所写的内容，除了魔鬼的活动（他要求获得权柄管辖约伯所有的一切、他的儿女、甚至他本人）之外，还有什么呢？然而魔鬼因约伯的忍耐而败退。在那卷书中，上主通过祂的回答，告诉我们许多关于那反对我们的龙的能力。以上是 \WorkTitle{旧约}中，就我们目前所能回忆的，关于敌对势力在圣经中被提及或被说成反对人类、后来受罚的记载。\par

现在我们也来看 \WorkTitle{新约}，其中撒殚接近救主，诱惑祂；那里也记载，恶神和不洁的魔鬼曾占据许多人，被救主从病人身上驱逐出去，这些人也被说成是因祂而得自由。甚至犹达斯，当魔鬼已把出卖基督的念头放在他心里之后，后来撒殚完全进入他里面；因为经上记载，在吃过饼后，\BibleQuote{撒殚进入了他里面}。宗徒保禄教导我们，不应给魔鬼留地步；却说：\BibleQuote{穿上天主的全副武装，使你们能抵抗魔鬼的诡计}；指出圣徒必须\BibleQuote{\Emph{战斗}，不是对抗血和肉，而是对抗率领者，对抗掌权者，对抗这黑暗世界的霸主，对抗天界中的邪恶势力}。他还说，救主甚至是被这世界的首领钉在十字架上，他们将要归于乌有；他说，他们的智慧他并不讲论。因此，全部这些事，圣经教导我们，有某些无形的仇敌与我们争战，圣经命令我们武装起来抵抗他们。由此，在信主基督的人中，较为简单的人认为，人所犯的一切罪，都是由于这些敌对势力持续努力影响罪人的心灵所致，因为在那无形的斗争中，这些势力被发现比人强。因为，举例来说，如果没有魔鬼，就没有人会迷失。\par

2. 我们既然更清楚地看明了这事的原因，便不持这一意见，并且考虑到那些显然因我们身体结构之必然结果而发生的罪。难道我们真要以为魔鬼是我们感到饥饿或口渴的原因吗？我想没有人敢这样主张。既然如此，若魔鬼不是我们感到饥饿和口渴的原因，那么，当每人到了青春期，这一时期唤起了天然热力的刺激时，两者又有什么区别呢？无疑将得出结论：正如魔鬼不是我们感到饥饿和口渴的原因，同样，他也不是那种在成熟时期自然产生的情欲——即性交欲望——的原因。诚然，这一原因并非总是由魔鬼推动，以致我们不得不认为，若无魔鬼存在，身体就不会有那种性交的欲望。其次，让我们思考：正如我们已经指出的，人类渴望食物并非出于魔鬼的煽动，而是出于某种自然的本能；那么，若无魔鬼，人类经验能否在取用食物时表现出如此节制，以致永不超越正当限度？即：没有人会取用不当，或超过理性所允许的限度；因此，人们在饮食上遵守适度与节制，便永不会犯错。我确实认为，即使没有魔鬼的煽动，人类也难以做到如此极大的节制，以致无人在取用食物时会超出应有的限度与克制，除非他因长期的习惯和经验学会了这样做。那么，情况究竟如何？在饮食之事上，即使没有魔鬼的挑动，我们也可能犯错，倘若我们碰巧不够节制或不够谨慎；那么，我们在性交的欲望中，或在对自然欲望的克制中，情况是否就不相似呢？我确实认为，同样的推理也适用于其他自然冲动，例如贪婪、愤怒、忧伤，以及一般因不节制之恶而超出自然适度限度的一切冲动。因此，有显而易见的理由支持这一意见：正如在善事上，人的意志本身软弱无力以成就任何善事（因为事事皆靠天主助佑才得以成全），同样，在相反性质的事上，我们从那些按自然使用的事物中接受了某些初始元素，可以说是罪的种子；但当我们过度放纵它们，未抵抗不节制的初次冲动时，敌对势力便抓住这初次过犯的机会，从各方面煽动和压迫我们，力图将我们的罪扩展到更广阔的领域，为我们人类提供罪的机会和起点，这些敌对势力将这些罪扩散到广大范围，若有可能，甚至超越一切限度。因此，当人们起初略略贪求金钱时，贪婪随情欲的增长而开始滋长，最终堕入贪婪。此后，当心智盲目接替情欲，敌对势力借其暗示催逼心智时，金钱不再只是被贪求，而是被偷窃、被强夺，甚至以流人血的方式获得。最后，一个确凿的证据表明，如此巨大的恶习出自魔鬼，这从以下事实容易看出：那些被过度情爱、或不可遏制的愤怒、或过分忧伤所压迫的人，所受的痛苦并不亚于那些被魔鬼身体折磨的人。因为某些历史记载说，有些人因情爱而陷入疯狂，有些人因愤怒，不少人因忧伤，甚至因过度喜乐；我认为，这源于此：那些敌对势力，即那些魔鬼，在不节制已为他们敞开心灵入口之后，便进入其中，完全占据了他们的感性本性，尤其是当对德行之荣耀的感觉没有唤醒他们去抵抗时。\par

3. 然而，有些罪并非来自敌对势力，而是源于身体自然的冲动，宗徒保禄在下面这段话中明确指出了这一点：\BibleQuote{肉身与圣神相争，圣神与肉身相争；二者彼此敌对，以致你们不能做你们所愿意的事。}因此，如果肉身与圣神相争，圣神与肉身相争，我们有时就必须与血肉搏斗——也就是说，作为人，我们按肉身生活，不能承受超乎人的试探；因为经上论到我们说：\BibleQuote{你们所受的试探，无非是人所常受的；但天主是忠信的，他决不许你们受试探超过你们的能力。}正如公共比赛的负责人不随意或偶然地让参赛者入场，而是经过仔细审查，以最公正的考虑——或按身材或按年龄——将这人配给那人，例如男孩配给男孩，男人配给年龄或体力相近的男人；同样，我们也当理解天主眷顾的安排，他以最公正的原则安排所有进入此世人生争斗的人，按照每个人能力的本性——这能力只有洞察人心的那一位知道——因此，一人与某一种肉身的试探争战，另一人与另一种；一人暴露在试探之下如此长的时间，另一人只如此长的时间；一人被肉身诱惑于这种或那种放纵，另一人被诱惑于另一种；一人必须抵抗这种或那种敌对势力，另一人必须同时与两个或三个争战；有时对抗这种敌对势力，有时对抗另一种；在完成某些行为之后，遭遇某一组敌人，在另一行为之后遭遇另一组。请看宗徒的话是否也暗示了这样的情况：\BibleQuote{天主是忠信的，他决不许你们受试探超过你们的能力，}也就是说，每人受试探的程度与他的力量或抵抗能力成正比。\newline{}现在，尽管我们说，按照天主公义的审判，每人受试探的程度与他的力量相当，但我们不应因此认为受试探的人必定在争斗中得胜；如同在赛场上竞技的人，虽然按公正原则配给对手，却不一定能获胜。除非竞技者的力量相等，否则胜利者的奖赏就不算公正；失败者也就不应受到公正的责备，因为他确实允许我们受试探，但不过\BibleQuote{超过我们的能力}：因为我们受试探是按我们的力量；经文并没有说，在试探中，他也必开一条出路，叫我们承当得住，而是叫我们能承当得住。但取决于我们自己是否用他所赐的这力量——或奋力，或软弱。因为毫无疑问，在任何试探中，我们都有忍耐的力量，只要我们善用所赐的力量。然而，拥有得胜的\Emph{能力}与实际得胜是不同的，正如宗徒本人以极其谨慎的措辞所表明的，他说：\BibleQuote{天主必开一条出路，叫你们\Emph{能}承当得住，}而不是说你们\Emph{必}承当得住。因为许多人没有坚持住试探，反而被它战胜。天主使我们能坚持住（试探）——否则似乎就没有争斗了——但使我们拥有坚持住的\Emph{能力}。而这赐给我们使我们能得胜的能力，可以按照我们自由意志的机能，或勤勉地使用——然后我们得胜，或懒惰地使用——然后我们失败。因为如果所赐的能力是必须得胜、永不失败，那么对于不能失败的人，争斗还有什么理由？若成功抵抗的能力被剥夺，胜利还有什么功德？但如果得胜的可能性平等地赐给我们所有人，并且如何使用这可能性——即勤勉或懒惰——在于我们自己，那么失败者将受到公正的责备，得胜者将得到应有的赞扬。现在，根据我们尽己所能讨论的这几点，我认为清楚显明了：有些过犯我们绝非在恶势力的压力下犯下；而另一些过犯，则是在他们的唆使下，我们被诱入过度而无节制的放纵。由此，我们应当探究那些敌对势力如何在我们的内心产生这些唆使。\par

4. 关于从我们心中发出的思念、对已做之事的回忆、或对任何事物及原因的默想，我们发现这些有时出自我们自己，有时是由敌对势力引发的，而且并不少见的是，它们是由天主或圣天使所提示的。这样的一种说法，除非有圣经的见证加以证实，否则或许会显得难以置信。那么，思念发自我们自己内心，达味在\BibleBook{圣咏}{集}中作证说：\Quote{人的思念将向祢认罪，其余的思想要向祢守节。}然而，这同样也为敌对势力所促成，撒罗满在\BibleBook{训道}{篇}中以下述方式指出了这一点：\Quote{如果当权者对你发怒，你不可离弃岗位，因为温和能平息大错。}宗徒保禄也将以这些话为此作证：\Quote{攻破人的计谋，以及各种为反对基督的智识所垒起的堡垒。}尽管如此，这乃是天主的作为，达味在\BibleBook{圣咏}{集}中宣告：\Quote{那以你为援助，在你心内设置朝圣之路的人，真有福气！}宗徒又说：\Quote{天主将这事放在提托心里。}某些思念是由善天使或恶天使提示给人的内心，这由陪伴托彼特的天使以及先知的话所表明，他说：\Quote{我遂问那与我谈话的使者说…}\WorkTitle{牧人书}也宣告同样的事，说每个人由两位天使陪伴；每当善念在我们心中升起，它们是由善天使提示的；而相反性质的思念，则是恶天使的煽动。巴尔纳伯在\WorkTitle{巴尔纳伯书信}中也宣告同样的事，他在其中说有两条路，一条光明之路，一条黑暗之路，并断言有某些天使被安置在这两条路上——天主的天使掌管光明之路，撒殚的天使掌管黑暗之路。然而，我们不可想象，从我们心中所提示的，无论是善是恶，会产生任何其他结果，除了一种内心的骚动和一种促动我们向善或向恶的刺激之外。因为，当恶势力开始煽动我们从恶时，我们完全有能力弃绝邪恶的暗示，抵抗卑劣的诱惑，不做任何可责之事。反之，当天主的能力召唤我们趋向更好的事时，我们也有可能不听从召唤；在这两种情况下，我们的自由意志都得到了保存。诚然，我们在前文说过，某些对善恶行为的回忆，是由天主眷顾的行动或敌对势力提示给我们的，正如在\BibleBook{艾斯德尔}{传}中所显示的：当阿塔薛西斯王不记得那位义人摩尔德开的功绩时，因夜间失眠而疲倦，天主使他心中想起要人给他诵读自己伟大事迹的编年史；于是，他记起了从摩尔德开那里所受的恩惠，便下令将他的敌人哈曼悬挂起来，而将显赫的荣耀赐予摩尔德开，并使整个圣洁的民族免于迫在眉睫的危险。另一方面，我们必须设想正是通过魔鬼的敌对影响，那提示才被引入大司祭和经师们的心中，他们来到比拉多面前说：\Quote{大人，我们记得那骗子活着的时候曾说过：三天以后我要复活。}犹达斯关于出卖我们的主和救世主的计谋，也不单是出于他内心的恶念。因为圣经作证：\Quote{魔鬼已使依斯加略人西满的儿子犹达斯决意出卖耶稣。}因此，撒罗满正确地命令说：\Quote{在一切之上，你要谨守你的心。}宗徒保禄也警告我们：\Quote{为此，我们必须更注重所听的道理，免得被漂去。}而当他说：\Quote{也不可给魔鬼留有余地，}他便以此命令指出，由于某些行为或一种内心的懈怠，便为魔鬼留出了空间，以致一旦他进入我们的心，若非完全占有我们，至少也要玷污灵魂，若他没有完全掌控它，就将他的火箭射向我们；有时我们被深深刺伤，有时只是被点燃。这些火箭很少并只在少数情况下被熄灭，以致找不到可以刺伤的地方，即当一个人被信德的坚固大盾所遮盖时。确实，\BibleBook{厄弗所}{书}中的宣告：\Quote{因为我们的战斗不是对抗血和肉，而是对抗率领者，对抗掌权者，对抗这黑暗世界的霸主，对抗天界里邪恶的鬼神。}必须如此理解，好像\Quote{我们}指的是：\Quote{我保禄，和你们厄弗所人，以及所有不是对抗血和肉而战斗的人：}因为这样的人必须与率领者、掌权者、这黑暗世界的霸主战斗，不像格林多人那样，他们的战斗还是对抗血和肉，并且所遇见的试探没有不是人所常受的。\newline{}ewline{}\par

5. 然而，我们不可认为每个人都必须与所有这些（仇敌）争战。因为任何人，即使是圣人，也绝不可能同时与它们全体进行较量。倘若那事竟以某种方式发生——既然它当然不可能如此——人性就不可能承受而不遭毁灭。但譬如，若五十名兵士说他们要与其他五十名交战，人们不会理解为其中一人必须对抗全部五十人，而是每人都会正确地说\Quote{我们的战事是对抗五十人}，全体对全体；同样，这也应理解为宗徒的意思：基督的所有运动员和兵士都必须与所列举的所有仇敌摔跤并搏斗——这较量确实必须针对所有仇敌，但由单个个体与单个力量较量，或至少按天主——这较量的正义裁判者——所定的方式进行。因为我认为人性的力量有限，即使有\Person{保禄}——关于他说过\BibleQuote{他是我所拣选的器皿}；或有\Person{伯多禄}——地狱的门不能战胜他；或有\Person{梅瑟}——天主的朋友：但他们中无一人能承受这些敌对势力同时发起的全面进攻而不自毁，除非那一位的威能单独在他内运作，他曾说\BibleQuote{你们放心，我已战胜了世界}。因此\Person{保禄}充满信心地高呼\BibleQuote{我依靠那加强我力量的基督，能应付一切}；又说\BibleQuote{我比他们众人更劳苦，其实不是我，而是天主的恩宠与我同在}。因此，既然这力量——绝非源于人——在他内运行并发言，\Person{保禄}才能说\BibleQuote{我深信无论死亡、生命、天使、率领者、掌权者、现时的事、未来的事、高深、权能、任何受造物，都不能使我们与天主的爱隔绝，这爱是在我们的主基督耶稣内的}。因为我不认为人性单凭自身就能与天使、高处的掌权者、深渊的掌权者以及任何其他受造物较量；但当它感受到主居住其内时，对天主助佑的信心将引领它说\BibleQuote{上主是我的光明，我的救援，我还畏惧谁呢？上主是我生命的保障，我尚恐惧何人？当仇敌前来攻击我，要吃我的肉时，我的敌人和仇人，反而绊倒仆地。虽有大军向我进攻，我的心毫不战栗；虽然战争向我勃发，我仍然满怀信赖}。由此我推断，一个人或许永不能凭自己战胜对抗的权势，除非他获得天主助佑的恩惠。因此，天使也被说成与\Person{雅各伯}摔跤。然而，我理解作者的意思是说，天使\Emph{与}\Person{雅各伯}摔跤和\Emph{对抗}他摔跤不是同一回事；而与\Person{雅各伯}摔跤的天使，是那与他同在以确保他安全的那位，他在知晓他的道德进步后，还赐给他以色列的名号，即他在斗争中\Emph{与他同在}，在较量中协助他；因为无疑还有另一位天使，\Person{雅各伯}是与他争战，并且必须与他进行较量。最后，\Person{保禄}没有说我们\Emph{与}率领者或\Emph{与}掌权者摔跤，而是\Emph{对抗}率领者和掌权者。因此，尽管\Person{雅各伯}摔跤，无疑是在\Emph{对抗}那些权势中的某一个，\Person{保禄}宣告这些权势反抗并攻击人类，尤其攻击圣徒。因此最后圣经论到他说\BibleQuote{他与天使摔跤，与天主角力获胜}，这样，较量由天使的帮助支持，但成功的奖赏将得胜者引向天主。\par

6. 我们确实不可认为这类斗争是靠体力与摔跤技艺进行的；而是灵与灵相争，正如\Person{保禄}所言：\BibleQuote{我们的斗争是对抗率领者、掌权者及这黑暗世界的统治者。}更要理解这类斗争的性质：例如，当损失与危险临到我们，或有人诽谤、诬告我们时，敌对势力的目的不仅在于让我们承受这些（考验），更在于藉此驱使我们陷入过分愤怒或悲伤，或陷入彻底的绝望；或至少——这是更重的罪——当我们在烦恼中疲乏不堪、被征服时，被迫抱怨天主，仿佛祂没有公正地治理人类生活；结果便是我们的信德可能被削弱，或望德落空，或被迫放弃我们信念的真理，或对天主怀有不敬的想法。关于\Person{约伯}，在魔鬼请求天主将权力赐予它处理\Person{约伯}的财物之后，便记载了类似的事。由此我们也学到，每当遭遇财产损失，并非出于偶然的攻击，也不是某个人被俘虏，或我们所爱之人的住所倒塌压死人是出于意外；因为对于所有这些事件，每位信徒都该说：\BibleQuote{除非从上赐给你，你对我毫无权柄。}请看，\Person{约伯}的房子是在魔鬼首先获得对付他儿子的权力之后才倒塌的；否则，那三队骑兵也不会突然冲来抢走骆驼、牛和其他牲畜，除非是受那灵煽动，他们已成为那灵之意志的仆人。那看似火或被认为雷击的东西，也不会落在祖宗的羊群上，除非魔鬼曾对天主说：\BibleQuote{你不是四面圈上篱笆，围护他的家和他所有的一切吗？但如今伸手触摸他所有的一切，看他是否当面否认你。}\par

7. 以上所有论述的结果是要表明：世上一切被视为中间性质的事件，无论是悲伤的还是其他的，确非由天主造成，却并非没有天主；祂不仅不阻止那些邪恶的敌对势力（它们渴求促成这些事）实现其目的，反而容许它们这样做，尽管只是在特定时机针对某些人，正如关于\Person{约伯}所说的，他在某个时期内被置于他人权下，房屋被不义之人洗劫。因此，圣经教导我们，视所有发生之事为天主所允许，知道没有天主便不会有任何事件发生。我们怎能怀疑这一点，即没有天主（的旨意）就没有事临到人？因为我们的主和救主耶稣基督宣称：\BibleQuote{两只麻雀不是卖一个铜钱吗？但你们天上的父若不许可，一只也不会掉在地上。}但论述的必要性使我们以冗长的离题话，讨论了敌对势力对人的斗争以及人生中更悲哀的事件，即它的诱惑——正如\Person{约伯}所言：\BibleQuote{人在地上的整个生命岂不是一个考验吗？}——以便清楚地显示这些事件发生的方式以及我们应如何看待它们。接下来我们要注意，人如何陷入虚假知识的罪中，或者敌对势力通常怀着什么目的就这些事挑起与我们的冲突。\par

% structure: p0064 source=h3 target=subsubsection reason=relative-heading-rank; base=h2

\subsubsection{第三章 论三重智慧}

1. 圣宗徒切愿将关乎学识与智慧的深奥真理训示我们，在\BibleBook{格林多}{前书}中说：\BibleQuote{我们在成全的人中也讲智慧，但不是今世的智慧，也不是今世有权势者的智慧，这些权势是要消灭的；我们讲的乃是天主奥秘的智慧，这智慧是隐藏的，天主在万世之前，就预定了，使我们获得光荣；今世有权势的人中，没有一个认识这智慧的，因为他们如果认识了，决不至于把光荣的主钉在十字架上。}\BibleRef{格前 2:6-8}在这段经文中，他描述不同种类的智慧，指出有一种今世的智慧，有一种今世有权势者的智慧，另有一种天主的智慧。但当他用\Quote{今世有权势者的智慧}这个说法时，我不认为他是指一切今世有权势者共有的智慧，而是指他们中某些人所独有的智慧。再者，当他说\BibleQuote{我们讲的天主奥秘的智慧，这智慧是隐藏的，天主在万世之前，就预定了，使我们获得光荣}\BibleRef{格前 2:7}时，我们必须探究，他的意思是否说：这就是那先前世代未曾显明、未使人知道的天主智慧，如今却启示给祂的圣宗徒与先知了；而且这也是救主来临前的那智慧，\Person{撒罗满}藉此获得了智慧；救主亲自声明祂所教导的比\Person{撒罗满}更大，说：\BibleQuote{看，这里有一位大于撒罗满的！}\BibleRef{玛 12:42}——这话表明，受教于救主的人，所学的超越\Person{撒罗满}的知识。因为若有人断言救主自己固然有更大的知识，但并未传授比\Person{撒罗满}更多的知识，那如何与下文一致：\BibleQuote{南方的女王在审判时，将起来定这世代人的罪，因为她从地极来，听\Person{撒罗满}的智慧；看，这里有一位大于撒罗满的！}\BibleRef{玛 12:42}因此，有一种今世的智慧，也可能有每位今世有权势者各自所附的智慧。但唯独论及天主的智慧，我们看出这一点：它往昔在古老先前的时代运作较小，后来却藉基督更圆满地启示出来。关于天主的智慧，我们将在适当处所探究。\par

2. 但现在，既然我们正在讨论敌对权势激动争斗的方式，藉此将虚假的知识引入人心，诱惑人灵，使人自以为发现了智慧，我认为有必要点名并区分今世的智慧与今世有权势者的智慧，藉此发现这些智慧之父。因此，照我上面所陈述，我认为在今世有权势者的智慧之外，另有一种今世的智慧；藉此智慧，属于今世的事物似乎被领悟和理解。然而，这种智慧本身全然不适于对神圣之事、或世界治理的计划、或任何其他重要主题、或良好与幸福生活的训练形成意见；而是完全涉及诗歌、文法、修辞、几何、音乐等技艺，或许医学也应归入这一类。凡此种种，我们都应视作今世智慧所包含的。另一方面，我们理解今世有权势者的智慧，是如埃及人所谓的隐秘玄奥哲学，以及加色丁人与印度人的占星术，他们自称知晓高超之事，还有希腊人中间关于神圣之事的多般分歧意见。因此，在圣经中我们发现各国都有有权势者；如在\BibleBook{达尼尔}{书}中我们读到波斯国的有权势者，希腊国的另一位有权势者，这些显然由经文性质表明不是人，而是某种权势。在\BibleBook{厄则克耳}{先知书}中，提洛的有权势者也明确被揭示为一种属灵的权势。因此，当这些及其他同类者，各自拥有自己的智慧，建立自己的意见和情感，看见我们的主与救主宣称并声明祂为此目的来到世界，要使一切伪称的学识意见被消灭，却不认识祂内在所隐藏的，就立即为祂设下网罗：因为\BibleQuote{世上的君王起来，首领聚在一起，反对上主和祂的基督。}\BibleRef{咏 2:2}但他们的网罗被发现，他们试图实行的计划在主的光荣被钉十字架时显明出来，因此宗徒说：\BibleQuote{我们在成全的人中也讲智慧，但不是今世的智慧，也不是今世有权势者的智慧，这些权势是要消灭的；今世有权势的人中，没有一个认识这智慧的，因为他们如果认识了，决不至于把光荣的主钉在十字架上。}\BibleRef{格前 2:6-8}\par

3. 我们确实必须努力查明，\OriginalTerm{今世的首领}{princes of this world}们试图灌输给人们的这种智慧，究竟是由\OriginalTerm{敌对势力}{opposing powers}引入他们心中，目的是为陷害伤害他们，还是仅为欺骗他们——即并非意图对人造成什么伤害，而是因为今世的首领们认为这些意见是真实的，他们渴望将自己相信为真理的传授给他人。这正是我倾向于采纳的看法。因为，举例来说，某些希腊作者或某异端教派的领袖，在领受了错误教义代替真理后，在自己心中得出结论说这就是真理，于是接下来就努力说服别人相信他们意见的正确；同样，我们设想今世的首领们的做法也是这样，在世上，某些属灵势力被指派统治某些民族，并因此被称为今世的首领。此外，除了这些首领，还有某些今世的\OriginalTerm{特殊作用}{special energies}，即属灵势力，它们带来某些效果，这些效果是它们凭借自由意志选择产生的；属于这些势力的，就有那些操练今世智慧的首领。例如，有一种特殊的作用和力量，是诗歌的灵感；另一种是几何学的；如此，每一种这样的技艺和职业都有一位单独的力量提醒我们。最后，许多希腊作家认为，没有疯狂就没有诗歌艺术；因此他们的历史中多次记载，那些他们称之为诗人的人突然充满了一种疯狂的精神。我们又如何说那些他们称为\OriginalTerm{占卜者}{diviners}的人，他们通过那些控制他们的魔鬼的作用，以精心构造的诗句给出回答？那些称为\OriginalTerm{术士}{Magi or Malevolent}或恶意之人，常常通过召唤魔鬼到幼童身上，使他们背诵出令人赞叹和惊讶的诗篇。现在，我们设想这些效果是这样产生的：正如圣洁无瑕的灵魂，以全部情感和纯洁将自己献给天主，保持自己不受恶灵的一切污染，经过长期斋戒的净化，并受到神圣宗教训练的熏陶，借此获得一份神性，赢得预言之恩和其他神圣恩赐；同样，我们设想那些将自己置于敌对势力之路的人，即故意仰慕并采纳其生活方式和习惯的人，也从它们那里获得灵感，成为它们的智慧和教义的分享者。其结果是，他们充满了那些他们已将自己臣服于其下的神灵的作用。\par

4. 至于那些关于基督教导与圣经规则所允许的不同的人，查明是否由于\OriginalTerm{敌对势力}{opposing powers}出于背信弃义的目的，为了阻止相信基督而编造出某些神话式的和不敬神的教义，或者是否他们听到基督的圣言后，既不能将其从良心的隐秘中抛出，也不能保持其纯洁神圣，于是通过\OriginalTerm{合用的器皿}{vessels convenient to their use}，即通过他们的先知，引入了各种与\OriginalTerm{基督真理的规则}{rule of Christian truth}相反的错误，这并非徒劳的任务。我们现在更倾向于认为，那些\OriginalTerm{背教和逃亡的势力}{apostate and refugee powers}，因自身心灵和意志的极恶而离弃了天主，或因嫉妒那些（在他们认识真理后）已预备好升到他们自己已坠落之同等地位的人，为了阻止任何这种进步，就发明了这些错误和虚假教义的欺骗。那么，有许多证据清楚表明，当人的灵魂存在于这身体内时，可以接纳来自不同善恶神灵的不同作用（即活动）。至于恶灵，其活动有两种模式：即，要么完全彻底占领心智，以致俘虏失去理解和感觉的能力，例如那些通常称为附魔的人，我们看见他们失去理智、疯狂（如福音中所记载被救主治愈的那些人）；要么通过邪恶的暗示，以各种思想败坏一个有知觉和理智的灵魂，劝说它去行恶，\Person{犹达斯依斯加略}就是一个例证，他在魔鬼的暗示下犯了叛卖的罪行，正如圣经所宣称的：\BibleQuote{魔鬼已经将出卖他的意念放在犹达斯依斯加略的心里}。\par

但是，一个人接受\OriginalTerm{善灵的作用}{energy of a good spirit}，即受激发和鼓励向善，并被感动向往天上的或神圣事物；正如圣天使和天主自己在先知们身上运作，用他们神圣的劝言激励和劝勉他们走向更好的生活，然而确实如此，它仍留在个人的意志和判断之内，即愿意或不愿意跟随对神圣和天上事物的召唤。从这明显的区别，可以看出灵魂是如何因着更好神灵的临在而被感动，即，如果它没有因即将到来的灵感而遭遇任何\OriginalTerm{心灵的扰乱或疏离}{perturbation or alienation of mind}，也没有失去对自己\OriginalTerm{意志的自由控制}{free control of its will}；例如，所有先知或宗徒的情况，他们都毫无心理扰乱地服务于神圣的回应。现在，通过善灵的建议，人的记忆被唤起，记起更好的事物，我们以前面的例子已表明过，当我们提到\Person{摩尔德开}和\Person{阿塔薛西斯}的时候。\par

5.　其次，我以为应当探究这个问题：为什么人的灵魂有时受到善灵的影响，有时受到恶灵的影响？我认为其原因比个人的肉身出生更早，正如（洗者）\Person{若翰}在他母亲的胎中跳跃和欢欣时所显示的，那时\Person{玛利亚}问候的声音传到了他母亲\Person{依撒伯尔}的耳中；又如先知\Person{耶肋米亚}所宣告的，他在母腹形成以前已被天主所认识，在出生之前已被祂圣化，并且尚在年幼时就领受了预言的恩宠。另一方面，也毫无疑问地显明，有些人从生命的一开始就被敌灵所附：亦即，有些人出生时就带有恶灵；另一些人，根据可信的记载，自幼就从事占卜。还有一些人从存在之初就受到被称为\OriginalTerm{派松}{Python}的魔鬼（即腹语灵）的影响。对于所有这些实例，那些主张世界上一切事物都在天主眷顾掌管之下的人（这也是我们自己的信仰），在我看来，除了坚持认为存在某些前生的原因，因此灵魂在肉身出生之前，在其感官性质或运动上已经犯下一定的罪债，从而被天主眷顾判定置于这种状况之外，无法给出其他回答以表明神圣治理之上毫无不义的阴影。因为灵魂无论是否在肉身内，总是拥有自由意志；而自由意志总是朝向善或恶。任何有理性的、有感觉的存在，即心智或灵魂，都不可能没有某种或善或恶的运动。这些运动很可能在灵魂在此世做任何事之前就提供了功过的依据；因此，基于这些功过或依据，灵魂在出生之时甚至出生之前，可以说就被天主眷顾分配去承受善或恶。\par

因此，对于那些似乎在人出生之后立即或甚至在他们见到光明之前就降临到人身上的事件，我们应持这样的看法。至于不同的神灵向灵魂（即人的思想能力）提出的建议，促使人行善或行恶，即使在这种情形下，我们也必须认为有时存在着某些先于肉身出生的原因。因为心智有时警醒，摒弃恶念，便呼求善灵的帮助；反之，若心智疏忽怠惰，则会因不够谨慎而为那些像强盗一样暗中埋伏、伺机冲入人思想的灵体留出余地，正如宗徒\Person{伯多禄}所说：\BibleQuote{我们的仇敌魔鬼，如同咆哮的狮子巡游，寻找可吞食的人}。因此，我们必须日夜以一切谨慎守护我们的心，不给魔鬼留地步；而要竭尽全力，使天主的仆役——即那些被派遣来服事那承受救恩之人的神灵——得以在我们内找到居所，乐于进入我们灵魂的客舍，居住在我们内，以他们的劝导引导我们；当然，若他们发现我们的心宅因德行与圣善的操练而装饰妥当。关于那些敌视人类的权势，我们尽其所能地说了这些，就该足够了。\par

% structure: p0072 source=h3 target=subsubsection reason=relative-heading-rank; base=h2

\subsubsection{第四章 论人的诱惑。}

1.　现在，我认为不应缄默不言的是人的诱惑这一主题；这些诱惑有时源自血肉和血统，或源自血肉的智慧，据说这智慧是与天主为敌的。至于某些人所声称的，即每个人仿佛有两个灵魂，这个说法是否正确，我们将在解释了那些据称比任何源于人的诱惑更强大的诱惑（即我们从率领者和掌权者、这世界的黑暗统治者、以及天上邪恶的属灵势力所承受的诱惑，或我们从恶灵和不洁魔鬼所遭受的诱惑）的本质之后，再作判定。在探究这一主题时，我认为应按逻辑方法追问：在我们这些由灵魂、身体和活力（\OriginalTerm{生命力}{spiritus vitalis}）组成的人里面，是否还有另一种要素，它自有其冲动，能唤起朝向邪恶的运动？因为这类问题常被一些人这样讨论：即，是否像我们内据说共存着两个灵魂，一个更神圣、属天，另一个较低劣；或者，是否仅仅因为我们依附于按其自身本性是死的、完全缺乏生命的身体结构（因为物质身体的生命来自我们，即来自我们的灵魂，而物质身体与精神相反并敌对），我们就被牵引、被诱惑去实践那些合于身体的恶行；或者，第三，是否（如某些希腊哲学家的意见）我们的灵魂虽在实体上是单一的，却由几个部分构成，一部分称为理性的，另一部分称为非理性的，而所谓的非理性部分又分为两种情感——贪欲和情欲。我们发现，上述关于灵魂的这三种意见，有人持有；但其中我们提到为某些希腊哲学家所采纳的那一种，即灵魂是\OriginalTerm{三分}{tripartita}的，我未看到圣经的权威对此有大力支持；而对于其余两种，在圣经中却发现了相当多的段落，似乎可以应用于它们。\par

2. 现在，我们先讨论某些人所持的见解，即在我们内有一个良善的、属天的灵魂，另一个是属地的、低等的；并且那更好的灵魂是从天上植入我们内的，就像\Person{雅各伯}还在母胎时，那灵魂使他赢得胜过他兄弟\Person{厄撒乌}的胜利，在\Person{耶肋米亚}身上，从出生就得以圣化，在\Person{若翰}身上，从母胎就被\OriginalTerm{圣神}{Holy Spirit}所充满。他们声称，那被称为低等灵魂的东西是与身体一起由身体的种子产生的，因此它不能离开身体而活着或存在，所以他们常说它被称为肉身。因为\BibleQuote{肉身相反圣神起贪欲}这句话，他们认为不适用于肉身，而是适用于这个灵魂，这灵魂正是\OriginalTerm{肉身的灵魂}{soul of the flesh}。此外，他们还想从这些话中为\BibleBook{肋未}{纪}的声明辩护：\BibleQuote{一切肉身的生命，都在于血}。因为血弥漫于整个肉身，使肉身产生生命，所以他们断言，这个被称为一切肉身的生命的灵魂，是包含在血内的。此外，肉身与圣神之争，圣神与肉身之争，以及进一步的声明\BibleQuote{一切肉身的生命，都在于血}，按这些作者的说法，不过是给\OriginalTerm{肉身的智慧}{wisdom of the flesh}另起一个名称，因为它是一种物质的精神，不受天主的法律约束，也不能受约束，因为它有属地的愿望和肉体的欲望。他们认为宗徒就是针对这一点说了这些话：\BibleQuote{我看见在我的肢体中有另一个法律，反对我理性的法律，并把我掳去，使我隶属于那在我肢体中存在的罪的法律。}如果有人反驳他们说，这些话是关于身体本性的，身体按其本性的特殊性而言是死的，但据说它有感觉或智慧，与天主为敌，或与圣神争斗；或者有人说，在某种程度上，肉身本身拥有一种声音，它喊叫反对忍受饥饿、干渴、寒冷，或由丰裕或贫穷带来的任何不适，——他们会试图削弱和损害这些论点，指出还有许多其他心灵的骚动，它们完全不是源自肉身，然而圣神却与它们争斗，如野心、贪财、争竞、嫉妒、骄傲等；看到人的心智或精神与这些进行某种斗争，他们将这一切邪恶的原因归诸于上面所提到的这个肉身的灵魂，它是通过\OriginalTerm{灵魂传递论}{traducianism}的过程从种子产生的。他们也惯于引用宗徒的声明来支持他们的断言：\BibleQuote{肉身的作为是明显的，它们是：淫乱、不洁、放荡、崇拜偶像、巫术、仇恨、纷争、嫉妒、忿怒、争吵、分裂、异端、派别、妒忌、醉酒、狂欢宴乐，以及类似的事；}他们断言，所有这些都不是源自肉身的习惯或快乐，因此所有这些活动都应被视为内在于那没有灵魂的实体，即肉身中。此外，声明\BibleQuote{弟兄们，你们看看你们的蒙召，在你们中间按肉身有智慧的并不很多，}似乎需要理解为有一种智慧，属肉的和物质的，另一种是属神的，前者确实不能被称为智慧，除非有一个肉身的灵魂，它对于称为肉身的智慧是智慧的。除了这些段落之外，他们还引用了以下的话：\BibleQuote{因为肉身相反圣神起贪欲，圣神相反肉身，以致我们不能做我们所愿意的。}那么，关于他所说的\BibleQuote{我们不能做我们所愿意的？}这些事是什么呢？他们回答说，肯定不能指圣神；因为圣神的意志不会受阻。也不能指肉身，因为如果它没有自己的灵魂，那么它肯定也不会有意志。剩下的就是，这个灵魂的意志是所指的，它能够有自己的意志，而且它肯定与圣神的意志相对立。如果是这样，那么就可以确定，灵魂的意志是介于肉身和圣神之间的某种东西，它无疑服从并服务于它所选择服从的那一个。如果它屈服于肉身的快乐，它使人成为属肉的人；但当它结合于圣神时，就产生属神的人，因此他们被称为属神的人。宗徒的话似乎就是这个意思：\BibleQuote{但是你们不在肉身内，而在圣神内。}\par

因此，我们必须弄清楚这\OriginalTerm{介于中间的意志}{intermediate will}是什么，除了那据说属于肉身或圣神的意志之外。因为人们确认，凡被称为\OriginalTerm{圣神之工}{works of the spirit}的，都是圣神意志（的产物），而凡被称为\OriginalTerm{肉身之工}{works of the flesh}的，都（来自）肉身的意志。那么，除了这些之外，那\OriginalTerm{灵魂的意志}{will of the soul}是什么？它另有一个名称，宗徒反对我们执行这意志时说：\BibleQuote{你们不能做你们所愿意的？}由此似乎要表明，它不应依附于这两者中的任何一个，即既不依附于肉身，也不依附于圣神。但有人会说，灵魂执行自己的意志比执行肉身的意志更好；同样，执行圣神的意志比执行自己的意志更好。那么，宗徒怎么说\BibleQuote{你们不能做你们所愿意的？}因为在肉身与圣神之间的争斗中，圣神绝不是必胜的，显然在许多个人身上，肉身占了上风。\par

3. 然而，既然我们所进入的讨论主题极为深奥，需要从各方面加以考虑，让我们看看是否可能确定这样一点：正如当圣神战胜了肉身时，灵魂跟随圣神是更好的，同样，如果前者（灵魂）跟随那与圣神抗争的肉身，而圣神正要将灵魂召回其影响之下，这看起来是一条更糟的途径，但灵魂受肉身管辖，仍可能显得比留在自己意志的力量下更为有利。因为据说它既不冷也不热，而是保持在某种温吞的状态，那么它的归化将是缓慢且有些困难的事。如果它确实紧贴肉身，那么，在饱尝并充满了那些它因肉身恶习而遭受的恶事，仿佛被奢侈和情欲的重担所累之后，有时它可能更容易、更迅速地脱离物质的污秽，转向对天上之事的渴慕和对属神恩宠的品味。而宗徒所说的这句话\BibleQuote{圣神与肉身相争，肉身与圣神相争，致使我们不能做我们所愿意的事}——无疑是指那些被指明超出圣神意愿和肉身意愿的事——其意义（若以另一种方式表达）是：一个人要么处于美德状态，要么处于邪恶状态，总比两者都不是更好；但是，灵魂在归向圣神并与圣神结合之前，当其依附于身体并默想属肉身的事物时，似乎既非处于善的状态，也非明显处于恶的状态，而是类似于一只动物。然而，对它来说，如果可能，通过依附圣神而成为属神的，这是更好的；但若不能做到，那么，即使跟随肉身的邪恶，也比置于自己意志的影响下而保留非理性动物的地位更为有利。\par

我们现在已经讨论了这些要点，出于考虑每个个别观点的愿望，比原计划更长一些，以免那些通常由询问者提出的观点被认为逃过了我们的注意——他们询问我们之内是否有任何别的灵魂，除了这一个属天的、理性的灵魂之外，而这个别的灵魂在本性上反抗后者（指这个属天的理性的灵魂），被称为肉身、肉身的智慧或肉身的灵魂。\par

4. 现在让我们看看那些主张我们内只有同一种灵魂的一个运动、一个生命，且灵魂的救恩与毁灭皆归因于其自身行为的人，通常如何回答这些陈述。首先，让我们注意灵魂的激动是何种性质：当我们感到自己内里被向不同方向牵引时，心中升起一种思想争战，某些可能性被提示给我们，使我们时而倾向这边，时而倾向那边，有时因此被定罪，有时又认可自己的行为。然而，说恶灵有一种变化不定、自相矛盾的判断，这并不奇怪，因为这在所有人身上都是明显的：每当对不确定之事进行商议时，人们考虑并咨询何者为更好、更有益的选择。因此，如果两种可能性相遇并提出相反观点，它们将心灵向相反方向拖曳，这并不令人惊讶。例如，若一个人因反思而信服并敬畏天主，就不能说是肉性与圣神相争；而是因真实与有益之事的不确定性，心灵被向相反方向牵引。同样，当认为肉性挑动情欲，而更好的劝诫抵制此类诱惑时，我们不应设想是某一个生命在抵抗另一个，而是身体本性的倾向——它急于排空并净化充满精液之处；同理，不应认为是某种敌对力量或另一个灵魂的生命在我们内激起口渴的欲望并驱使我们饮水，或导致我们感到饥饿并驱使我们满足它。但正如通过身体的本能运动，食物和水被渴望或拒绝，同样，随着时间的推移积聚在各种器官中的天然种子也急于被排出丢弃，甚至有时在没有任何刺激原因的情况下自行溢出。因此，当经上说\BibleQuote{肉性与圣神相争}时，这些人理解为：肉性的习惯、必要或快乐唤醒人，使他离开神圣和属灵的事物。因为因身体的需要被拉扯，我们不得闲暇从事于永远有益的属神之事。同样，当灵魂投身于神圣和属灵的事，并与圣神结合时，它被称为与肉性作战，不容许肉性因放纵而松弛，或因那些它自然愉悦的快乐而变得不稳。他们也以此方式理解\BibleQuote{肉性的智慧与天主为敌}这句话，并非肉体真有灵魂或自己的智慧。而是正如我们惯常不当使用语言说大地口渴、想喝水——这里\Quote{想}的词用得不恰当，而是借用说法（\OriginalTerm{借用词语}{catachrestic}）——同样可以说房屋想重建，以及许多类似表达；对肉性的智慧或\BibleQuote{肉性与圣神相争}也应如此理解。他们通常也将\BibleQuote{你兄弟的血由地上向我喊冤}与此相连。因为向主呼喊的并非所流之血本身；而是不当地说血在呼喊，要求惩罚流血之人。宗徒所说的\BibleQuote{我在我的肢体上发现了另一条法律，与我理智的法律交战}，他们理解为：那愿意献身于天主圣言的人，因身体的需要和习惯——这些像一种法律深深嵌入身体——而被分散、分裂、阻碍，不能奋力钻研智慧，以默观神圣的奥秘。\par

5. 至于以下被列为肉性之工的事：异端、嫉妒、争吵或其他恶行，他们如此理解经文：心灵因屈服于肉性的情欲而变得感觉迟钝，被罪恶的重量压迫，没有精细或属灵的感觉，被说成成了肉性，并从其中获得名称——它在那里表现出更强烈的意志力量。他们还进一步追问：\Quote{这个被称为肉性情欲的邪恶感觉，其创造者将是谁？}因为他们主张灵魂和肉体的创造者除了天主别无其他。若我们断言善良的天主在自己的受造物中创造了任何与祂自己为敌的东西，那显然是荒谬的。那么，既然经上写着\BibleQuote{肉性的智慧是与天主为敌}，且这被宣告为受造的结果，天主自己就似乎创造了一种与祂自己为敌的本性，这本性既不能服从祂也不能服从祂的法律，仿佛它是一种被这样形容的动物。若接受这种观点，它与那些主张所造灵魂本性不同、按其本性注定灭亡或得救之人的观点有何区别？然而，这仅是异端者的意见，他们不能以敬虔维持天主的正义，便虚构这等不敬的发明。现在，我们已尽己所能将每方关于不同观点可能提出的论证摆出，让读者自行选择他认为应优先采纳的。\par

% structure: p0080 source=h3 target=subsubsection reason=relative-heading-rank; base=h2

\subsubsection{第五章 世界在时间中有其开始。}

1. 现在，既然在教会的信条中有一条主要是基于我们对神圣历史真实性的信念，即这个世界是在某个时间被创造并开始存在的，并且按照为万物所规定的时期循环，它将要因腐朽而被毁灭，那么重新讨论与此相关的几点似乎并非荒谬。而且，就圣经的可信度而言，关于此事的声明似乎易于证明。甚至异端者，尽管在许多其他事上广泛对立，但在这一点上似乎意见一致，都服从圣经的权威。\par

关于世界的创造，有哪部分圣经能比梅瑟关于其起源所传述的记载给我们更多信息呢？虽然它包含比单纯历史叙述所显示更深刻的意义，并且包含许多要按属灵理解的事物，且在论述深奥神秘的主题时以文字作为某种遮蔽，但叙述者的语言表明所有有形之物都是在某个时间被创造的。至于世界的终结，\Person{雅各伯}是第一个提供信息的，他在对子女说话时说道：\BibleQuote{雅各伯的儿子们，你们到我这里来集合，我要告诉你们在末日，}或\BibleQuote{在末日以后将有什么事。}那么，如果有\BibleQuote{末日}，或一个\BibleQuote{紧随末日的时期}，那么有开始的日子必然要结束。\Person{达味}也宣告说：\BibleQuote{诸天要毁灭，而祢将永存；是的，它们全部都要像衣服一样陈旧；祢要像更换外衣一样更换它们，它们就被更换了；但祢永不改变，祢的岁月没有穷尽。}我们的主救主确实在\BibleQuote{起初创造者使他们成为男人和女人}这句话中亲自见证了世界是被创造的；再者，当祂说\BibleQuote{天地要过去，但我的话决不会过去}时，祂指出它们是可朽坏的，并且必然要终结。此外，宗徒在宣告\BibleQuote{受造之物被屈服于虚妄之下，并不是出于自愿，而是由于那使它屈服的那一位，怀着希望，因为受造之物本身也要从败坏奴役中被释放，进入天主儿女光荣的自由}时，明确宣告了世界的终结；他在再次说\BibleQuote{这世界的模样正在逝去}时也是如此。现在，通过他所说的\BibleQuote{受造之物被屈服于虚妄之下}这一表述，他表明了这个世界有一个开始：因为如果受造之物因某种希望而被屈服于虚妄，那么它当然是由于一个原因而被屈服的；既然是出于一个原因，它必然有一个开始：因为如果没有某种开始，受造之物就不可能被屈服于虚妄，也不可能希望从败坏的奴役中得自由，如果它从未开始服役的话。任何愿意闲暇时搜索的人，都会在圣经中找到许多其他经文，其中世界既被说有开始，也被寄望于终结。\par

2. 现在，如果有人在此反对我们圣经的权威或可信度，我们要问他：他是否断言天主能够或不不能够理解一切？断言祂不能，显然是不敬之举。那么如果他回答，正如他必须回答，天主理解一切，那么正因它们能被理解这一事实，就被认为有开始和终结，因为完全没有开始的东西根本不能被理解。因为无论理解延伸多远，当认为没有开始时，理解的能力就被无限地撤回和移除。\par

3. 但这是他们通常提出的反对意见：他们说：\Quote{如果世界在时间中开始，那么在世界开始之前天主在做什么？因为说天主的本性是不活动、不动的，或假设良善在某一时刻没有行善，全能者在某一时刻没有运用其大能，既是亵渎又是荒谬的。}这就是他们惯常对我们这个说法的反对，即这个世界在某个时间开始，并且按照我们对圣经的信仰，我们可以计算它已经存在的岁月。对于这些命题，我认为任何异端者都难以给出与其观点本性相符的回答。但我们可以根据宗教的标准给出一个合理的回答，即我们说，天主并非在创造这个有形世界时才首次开始工作；而是正如在它毁灭之后将有另一个世界一样，我们也相信在这个世界存在之前已经存在着其他世界。这两个观点都将由圣经的权威来证实。因为在此之后将有另一个世界，\Person{依撒意亚}教导说：\BibleQuote{将有新天新地，我要使它们存留在我眼前，上主说。}；而在此世界之前也有其他世界存在，\BibleBook{}{训道篇}指出：\BibleQuote{是什么已经存在？就是将要存在的。是什么已经被创造？就是将要被创造的。太阳之下毫无新事。谁能说话并宣告，看，这是新的？它已经在我們之前的世代中存在过。}通过这些证言，既确立了在我们之前已有世代，也确立了在我们之后将有其他世代。然而，不应认为多个世界同时存在，而是这当前世界终结之后，其他世界将开始存在；关于这一点，无需重复每个具体陈述，因为我们已经在前面几页中这样做了。\par

4. 这一点确实不可轻易放过：圣经以一个新的、独特的名称称呼世界的创造，称之为\OriginalTerm{奠定}{\GreekText{καταβολή}}，这个词极不恰当地被拉丁文译为\Quote{constitutio}；因为在希腊文中，\OriginalTerm{奠定}{\GreekText{καταβολή}}更意指\Quote{dejicere}，即抛下——这个词，正如我们已经指出的，被拉丁文短语\Quote{constitutio mundi}不恰当地翻译了，如在\BibleBook{若望}{福音}中，救主说：\BibleQuote{那时必有大灾难，从世界的开始以来没有这样的}；在这段经文中，\OriginalTerm{奠定}{\GreekText{καταβολή}}被译为\Quote{开始}（constitutio），应按上述解释理解。宗徒在\BibleBook{厄弗所}{书}中也用了同样的说法，说：\BibleQuote{天主在创造世界以前，在基督内拣选了我们}；他将这个\Quote{创造}称为\OriginalTerm{奠定}{\GreekText{καταβολή}}，应按同样的意义理解。因此，探究这个新术语的含义似乎值得；我确实认为，既然圣人们的终结和圆满将在那看不见的、永恒的（时代）中，我们必须（如前文多次指出的）通过默观那终结本身，断定理性受造物也有一个类似的开始。如果它们有一个如他们所盼望的终结那样的开始，那么他们无疑从最初就存在于那看不见的、永恒的东西之中。若果真如此，那么就存在从较高境界到较低状况的下降，这不仅发生在那些因自身运动多样而应得变化的灵魂身上，也发生在那些为了服务整个世界而被迫从较高、不可见的领域降到较低、可见的领域的灵魂身上，尽管是违背他们意愿的——\BibleQuote{因为受造物屈服于虚幻，并非出于自愿，而是出于那使它屈服的，并怀着希望}；这样，太阳、月亮、星辰和天使才能对世界以及那些因心智缺陷过甚而需要更粗重、更坚固的肉体的灵魂履行其职责；并且为了那些需要这种安排的人，这个可见的世界也被召唤存在。由此推论，通过这个词的使用，似乎指向一种全体共同分享的从较高到较低状况的下降。的确，所有受造之物都怀有获得自由的希望——从奴役的败坏中解脱出来——当天主的子女们，那些堕落或分散的，将聚集为一，或者当他们在世上完成了其他只有安排万物的天主才知道的职责时。我们确实应该设想，世界被创造得具有这样的性质和容量，不仅要包含所有那些注定要在世上受训练的灵魂，还要包含所有那些准备伺候、服务并协助他们的权能。因为许多声明都确定所有理性受造物具有同一性：唯有在这个基础上，天主在其一切对待他们的事上的正义才能得到辩护，因为每个人自身都有理由知道自己为何被置于这种或那种生活等级中。\par

5. 因此，天主后来安排的这个事物秩序（因为祂从世界的起源就清楚地洞察到影响那些人的原因和理由：或因心智缺陷而应进入肉体，或由于对可见之物的贪恋而被冲走，还有那些，无论出于自愿或不自愿，被那（使其屈服于希望者）强迫去服侍那些陷入这种状况的人），没有被一些人理解，他们未能看出这种多样化的安排是由天主设立的，其根源在于自由意志的先前原因，因此得出结论：世上一切事物要么受偶然运动支配，要么受必然命运支配，没有任何东西在我们自身意志的能力范围内。因此，他们也无法证明天主的眷顾是无可指责的。\par

6. 但正如我们所说，所有生活在此世的灵魂都需要许多服务者、管理者或助手；同样，在末世，当世界的终结已迫在眉睫，全人类濒临最后的毁灭，不仅那些受他人管理的人变得软弱，就连那些受托管理职责的人也是如此，那时不再需要这样的帮助或这样的保卫者，而需要作者和创造者本人的帮助，以恢复一方被败坏和亵渎的服从纪律，以及另一方被败坏的管理纪律。因此，天主的独生子，祂是圣父的圣言和智慧，当祂与圣父同享那在世界以前就有的荣耀时，祂舍弃了它，取了奴仆的形状，服从至死，为要教导那些只能通过服从才能获得救恩的人学习服从。祂也通过将一切仇敌置于祂脚下，恢复了被败坏的管理和统治的法则，借此（因为祂必须为王，直到把一切仇敌放在祂脚下，并消灭最后的仇敌——死亡）祂可以教导统治者本人节制其统治。既然祂来是为了恢复纪律，不仅包括统治的纪律，也包括服从的纪律，正如我们所说，祂首先在自己身上完成了祂愿意由他人完成的事，祂服从圣父，不仅至十字架的死，而且在世界终结时，将一切经由祂而服从圣父并获救的人包含在自己之内，祂自己也与他们一起，在他们之内，被称为服从圣父；万有在祂内存在，祂自己是万有之首，在祂内是获救者的救恩和圆满。因此，宗徒论到祂说：\BibleQuote{当一切事物都屈服于祂以后，圣子自己也要屈服于那使一切事物屈服于祂的，好叫天主成为万物之中的万有。}\par

7. 我实在不知道，异端者不理解宗徒在这些话中的意思，如何将\Quote{\OriginalTerm{屈服}{subjection}}一词视为对圣子的贬抑；因为如果对这个称号的恰当性提出质疑，可以通过相反的假设很容易确定。因为如果处于屈服不好，那么相反的好，即不处于屈服。现在，按照他们的看法，宗徒的言语似乎通过这些话说：\BibleQuote{当一切事物都被征服于祂之后，圣子自己也要屈服于那使一切事物屈服于祂的，}这表明现在不屈服于圣父的那位，将在圣父先征服一切事物于祂之后，屈服于祂。但我惊讶的是，怎么可能认为：当一切事物尚未被征服于祂时，祂自己并不屈服；而在一切事物都被征服于祂、祂成为万民之王并掌管万有时，倒被认为变得屈服了，尽管祂先前并不屈服？因为这样的人不明白，基督对圣父的屈服表明我们的幸福已达完美，祂所承担的工作已胜利完成，因为祂不仅净化了至高统治权对全部受造物的能力，而且将人类\OriginalTerm{服从}{obedience}和\OriginalTerm{屈服}{subjection}的原理以纠正和改善的状态呈献给圣父。那么，如果说圣子屈服于圣父的那种屈服被认为是美好而有益的，那么就有极其合理和逻辑的推论：敌人对天主子所表现的屈服也应当被理解为同样有益有用；正如当圣子被说成屈服于圣父时，表示全部受造物的完美恢复，同样，当敌人被说成屈服于天主子时，被征服者的得救和失落者的恢复就在其中。\par

这种屈服的实现，将按照特定的方式、经过特定的训练、并在特定的时期完成；不可想象屈服是由必然的压力所带来的（免得整个世界似乎被强制屈服于天主），而是藉由言语、理智和教导；藉由对更好道路的呼唤、最佳的训练体系，以及运用适当合宜的威吓，这些威吓公正地临到那些忽视自己救恩和益处的任何关怀或注意的人。简而言之，我们人类在训练奴隶或子女时，也因他们年幼无法运用理智而用威吓和恐惧来约束他们；但一旦他们开始明白何为善、有益、可敬时，对鞭打的恐惧消除后，他们便因言语和理智的劝服而顺从一切善事。然而，为了在所有理性受造物中保存自由意志，每个人应如何被治理——即，哪些人由天主的圣言找到并训练，好像他们已经预备好并能接受；哪些人被推迟到更晚的时间；哪些人完全被隐藏，其处境远离听闻圣言；哪些人，当圣言被传扬和宣讲时却轻视它，而被某种矫正和惩罚驱向救恩，其悔改在某种程度上是被要求和强夺的；哪些人被赐予某些救恩机会，有时仅凭一次回答就证明其信德，从而无疑获得了救恩——这些结果出于何种原因或场合，或天主的智慧在他们内看见什么，或他们意志的什么运动使天主如此安排这一切，只有祂和祂的独生子（万物藉祂而创造并恢复）以及圣神（万物藉祂而圣化，祂由圣父所发）知道。愿光荣归于祂，直到永远。阿们。\par

% structure: p0090 source=h3 target=subsubsection reason=relative-heading-rank; base=h2

\subsubsection{第六章 论世界的终结。}

现在，关于世界的终结和万物的完成，我们在前面几页中，已依照圣经的权威，尽其所能地陈述了我们认为足以教导的内容；这里只补充几点劝勉性的意见，因为探讨的顺序使我们回到这个主题。至善，就是整个理性本性所追求的目标，也被称为一切福乐的终结，许多哲学家定义如下：他们说，至善就是尽可能相似于天主。但我认为这个定义与其说是他们的发现，不如说是源自圣经的观点。因为这一点，在众哲学家之前，梅瑟已经指出，他在描述人的最初创造时这样说：\BibleQuote{天主说：\InnerQuote{让我们照我们的肖像，按我们的模样造人。}}然后他接着写道：\BibleQuote{天主于是照自己的肖像造了人，就是照天主的肖像造了人：造了他们男和女，并祝福了他们。}这里\BibleQuote{照天主的肖像造了他}这句话，没有提到\Quote{模样}一词，其含义无非是：人在最初受造时就获得了天主肖像的尊荣；但模样的完美则保留给完成之时——即人藉自身勤勉模仿天主而为自己获得它，起初通过天主肖像的尊荣而获致成全的可能性，最终通过完成必要的工作而达到天主模样的完美实现。事实正是如此，宗徒若望更清楚明确地指出了这一点，他宣告说：\BibleQuote{小孩子们，我们还不清楚我们将是什么；但若从救主那里给我们启示，你们必会毫不怀疑地说：我们将相似祂。}这句话极其肯定地指出，不仅万物的结局——他说对他来说依然未知——是当仰望的，而且相似天主也是当仰望的，将按照功行的圆满而赐予。主自己在福音中不仅宣告这些结果将要实现，而且藉祂的代祷——祂亲自为门徒向圣父求得——使这一切成就，祂说：\BibleQuote{父啊！我愿我在哪里，这些人也与我在哪里；并且如同你与我原是一体，他们也在我们内合而为一。}其中，如果我们可以这样表达的话，天主肖像本身似乎已经进步，从仅仅相似变为相同，因为在终结或终点，天主是\BibleQuote{万有之中的万有}。关于这一点，有些人质疑：物质身体的本质，即使被洁净、纯化并完全成为属神的，是否似乎会阻碍获得（天主）肖像的尊荣或合一的性质？因为一个属物质的本性不能与确实无形的神性有任何相似，也不能真实恰当地被称为与它合一，尤其当我们由信仰的真理得知：那唯一的一，即圣子与圣父，必须归于（天主）本性的特殊性。\par

2. 既然应许在终结时天主将成为万有中的万有，我们就不宜设想动物，无论是绵羊或其它牲畜，能达到那个终结，免得暗示天主甚至居住在动物，无论是绵羊或其它牲畜内；同样，也不宜设想木石达到那终结，免得说天主也在它们内。同样，任何邪恶之事也不应被认为能达到那终结，免得当说天主在万有中时，也被说成在邪恶的器皿内。因为如果我们现在断言天主无所不在且充满万有，基于无物可离天主而空虚，我们却并不说祂现在在祂所居者中是\Quote{万有}。因此，我们必须更仔细地考察那标志着圆满福乐和万物终结的是什么——它不但被称为天主在万有中，也被称为\Quote{万有中的万有}。那么就让我们探究天主在万有中将成的\Quote{万有}究竟是什么。\par

3. 我认为，天主被称为\Quote{万有中的万有}这一表述，意思是祂在每个个体中是\Quote{万有}。祂将在每个个体中成为\Quote{万有}的方式如下：当任何理性悟性从各种恶习的渣滓中得以净化，所有邪恶的乌云被彻底扫清后，所能感受、理解或思想的一切，将完全是天主；当它不再看到或持有任何非天主的东西，而是天主成为它一切活动的尺度与标准时；如此，天主就是\Quote{万有}，因为不再有善恶之分——邪恶无处存在；因为天主是万有，祂近处无恶；那常拥有善、天主为他就是万有者，也不再渴望吃知善恶树上的果子。因此，当终结复归于起始，万物的结局与其开端相比，那理性受造物的状态将被重建——它当初无需吃知善恶树上的果子；如此，当所有邪恶感被移除，个体得以净化清洁，那唯一至善的天主就成为他的\Quote{万有}——这不限于少数个体或相当数量，而是祂自身是\Quote{万有中的万有}。当死亡不再存在于任何地方，死亡的刺和一切邪恶都不复存在时，那时天主确实是\Quote{万有中的万有}。但有些人认为，理性受造物或本性的那种圆满与福乐，只能在我们上面所提及的同一状态中存留，即：万有都拥有天主，天主为他们成为万有——如果他们丝毫不因与身体本性的结合而受阻。否则他们认为，最高福乐的荣耀会被任何物质实体的混合所阻碍。但这个问题我们在前面几页已经更详细地讨论过了。\par

4. 现在，既然我们发现宗徒提到了一个属神身体，就让我们尽所能探究我们应如何构想这样的事物。就我们的理解所能把握而言，我们认为属神身体的本性应不仅为所有圣洁完美的灵魂所居住，也为所有将从腐败的奴役中获得解放的受造物所居住。关于身体，宗徒也曾说：\BibleQuote{我们有一所非人手所造，在天上永存的房屋}，\BibleRef{格后5:1}即指在诸福乐的居所。从这一陈述我们可以推测，那身体的性质是何等纯洁、精微、荣耀——如果我们将其与那些虽属天体、光芒璀璨、却是人手所造、肉眼可见的身体相比。但关于那身体，它被说成是非人手所造、在天上永存的房屋。因此，既然\BibleQuote{所见的是暂时的，所不见的是永远的}，\BibleRef{格后4:18}所有我们在地上或天上看见的、能被看见的、人手所造而非永恒的身体，在荣耀上都远不及那不可见、非人手所造、永恒的身体。从这一比较可以想象属神身体是何等优美、辉煌、光明；并且\BibleQuote{天主为爱祂的人所预备的，是眼睛未曾看见，耳朵未曾听见，人心未曾想到的}。\BibleRef{格前2:9}然而，我们不应怀疑，我们现在这身体的本性，可以凭天主的旨意——祂造了它如此——被提升到那（所提及之身体的）精纯、纯洁、光辉的品质，按照事物情况的需要和理性本性的功绩而定。最后，当世界需要多样性和差异性时，物质以全然的顺从，在事物的不同显现和种类上，服从造物主——作为它的主和制造者——以便祂能从中产生出天上和地上各类形式。但当万物开始奔向那一切为一的圆满，如同父与子合一，可以理性地推论：当万有成为一时，将不再有任何差异性。\par

5. 此外，那被称为死亡的最后一个敌人，据说要被消灭，为的是当死亡不再存在时，没有任何哀伤的事物留下，也没有任何敌对的事物，因为没有敌人。诚然，最后一个敌人的消灭，应理解为并非天主所造的其实体要消亡，而是其思想和敌对的意志——这并非来自天主，而是来自它自身——要被消灭。因此，它的消灭并非它不存在，而是它不再成为敌人，也不再是死亡。因为对于全能者而言，没有不可能的事，也没有什么不能恢复其造物主：祂造了一切，使他们存在，而那些为存在而造的事物，不会停止存在。也正因如此，它们容许变化和多样性，以便根据其功过，被置于更好或更糟的位置；但实体的消灭不会降临到那些天主为永恒存在而造的事物上。因为那些按照普遍意见被认为会消亡的事物，无论是我们的信德本性还是真理，都不允许我们设想它们被消灭。最后，我们的肉体被无知者和不信者认为在死后会消亡，以至于它残留的任何先前实体的痕迹都不复存在。然而，我们相信复活的人明白，死亡只带来了变化，但实体确实保留；并且，按照造物主的旨意，在指定的时间，它将被恢复生命；并且，第二次变化将在它内发生，以致那起初由尘土而来的肉体，后来被死亡分解，再次归于尘土和灰烬（经上说：\BibleQuote{因为你是由土来的，你还要归于土}），将从地上再次被举起，之后，根据内住灵魂的功过，上升到属神的身体的荣耀。\par

6. 那么，我们应当设想，当我们整个肉体的实体被带入这种状态时，所有事物都将重新恢复为统一体，而天主将成为万物中的万有。而且，必须理解这一结果是并非突然，而是缓慢且逐步达成的，因为修正和改正的过程将在数不清的无可计量的时代流逝中，在个别实例中不知不觉地进行——有些超越其他，以更快的步伐趋向完美，而另一些则紧随其后，还有一些则远远落后；这样，通过无数不可计数的、正在从敌对状态与天主和好的进步存有的等级，最后才触及那被称为死亡的最后一个敌人，为使他也能被消灭，不再成为敌人。因此，当所有理性的灵魂都已恢复到这种状态时，我们这个肉体的本性将转变为属神的身体的荣耀。因为正如我们看到理性本性的情况并非如此——有些因自己的罪而活在堕落状态，而另一些因自己的功过而被召至幸福状态——而是我们看到那些曾犯罪过的同一灵魂，在悔改并与天主和好之后，被助至幸福状态；同样，我们也应思考肉体的本性：我们现在在卑贱、可朽和软弱中使用的这个身体，与我们在不朽、能力和荣耀中所将拥有的身体并非不同的身体；而是同一个身体，当它摆脱了现在缠绕它的软弱之后，将被转变为荣耀的状态，成为属神的，以致那曾是耻辱的器皿，在洁净后成为荣耀的器皿和福乐的居所。而且，在这种状态中，我们也相信，按照造物主的旨意，它将永远保持不变，没有变化，正如宗徒的宣言所证实的，他说：\BibleQuote{我们有一栋房屋，不是人手所造的，永远在天上。}因为教会的信德不接受某些希腊哲学家的观点，即除了由四种元素构成的身体之外，还有另一个完全不同的第五身体，与我们现在的这个身体截然不同；因为既不能从圣经中找出任何最轻微的怀疑证据来支持这种观点，也不能从事物的理性推论中接受它，尤其是当圣宗徒明确宣告，并非新身体赐给从死者中复活的人，而是他们领受自己生前所拥有的同一身体，从低劣状态转变为更好的状态。因为他的话是：\BibleQuote{播种的是生灵的身体，复活起来的是属神的身体；播种的是可朽坏的，复活起来的是不可朽坏的；播种的是软弱的，复活起来的是强健的；播种的是耻辱的，复活起来的是荣耀的。}因此，正如在人内有某种进步，使他起初是属生灵的，不明白属于天主圣神的事，借着教导达到成为属神的人的阶段，并判断一切，而他自己却不受任何人判断；同样，关于身体的状态，我们也应持守：这个现在因服务于灵魂而被称为生灵的身体，当灵魂与天主结合，与祂成为一神时（那时身体仍服务于神），将借着某种进步达到属神的状态和品质，尤其因为我们已多次指出，肉体本性被造物主如此塑造，以致它容易转入祂所愿的任何状态，或情势所需的状态。\par

7. 这一切推理的结论即是：天主创造了两种普遍的性体——一种是有形的，即可看见的性体；另一种是无形的，即不可见的性体。这两种性体允许两种不同的变化。那无形且具理性的性体在思想和目的上发生变化，因为它被赋予了\OriginalTerm{自由意志}{freedom of will}，并因此有时被发现致力于善的实践，有时则相反。然而，这种\OriginalTerm{有形性体}{corporeal nature}允许本质上的变化；因此，万物的安排者天主在塑造、制造或修整祂所愿的任何事物时，都掌握着这种物质的服事，以便有形性体可以根据事物应得的要求转化和变形为任何形状或种类；先知显然看到了这一点，当他说：\Quote{是那创造并转变万有的天主。}\par

8. 现今需要探讨的是：当天主成为万物之中的万有时，在万物的终末与圆满中，整个有形性体是否将由一种形状组成，而身体的唯一品质将是那在不可言喻的光荣中所照耀的，这光荣被视为未来\OriginalTerm{灵性的身体}{spiritual body}所拥有的。因为如果我们正确理解此事，这便是梅瑟在他书卷开头所说的：\BibleQuote{在起初天主创造了天地。}因为这是所有受造的开始；万物的终结与圆满必须回归于此开始，即为了使那天和那地成为虔诚者的居所与安息之地；以致所有圣洁的人和温良的人首先在那地获得产业，因为这是法律、先知和福音的教导。在这地中，我相信存在着梅瑟在\OriginalTerm{法律之影}{shadow of the law}下所传的崇拜的真实而活的形式；关于这崇拜，经上说：\BibleQuote{他们供奉的事是天上事物的模型和影子}——即那些在法律之下服从的人。梅瑟本人也受命：\BibleQuote{你要谨慎，照你在山上所指示的样式制造它们。}由此在我看来，如同在这地上，法律对那些借此被引向基督的人是一种启蒙教师，使他们受其教诲和训练后，能更容易地接受基督更完善的教理；同样，另一块接纳所有圣徒的地，可能首先通过真实和永恒的法律的体制来浸染和塑造他们，使他们更容易取得那些无可增添的完善的天上体制；在那里，将确实是那称为永久的福音，和那永新的、永不陈旧的\OriginalTerm{盟约}{Testament}。\par

9. 因此，据此我们应当认为，在万物的终结与复原时，那些逐渐进步、在改善的阶梯上攀升的人，将按照适当的尺度与次序到达那地及其包含的培育中，在那里他们可以为那些无可增添的更完善的体制而准备。因为在祂的代理者和仆人之后，主基督——万有之王——将亲自承受王国；即在圣德的教导之后，祂将亲自教导那些能够按祂是智慧来接纳祂的人，在他们内为王，直到祂将他们引向那已将万有屈服于自己的圣父，即当他们已能承受天主时，天主就成为他们之中的万有。于是，作为一种必然的结果，有形性体将获得那无可增添的最高的状态。至此，我们讨论了有形性体或灵性身体的品质，将判断权留给读者，由他决定何为最佳。至此，我们可以结束第三卷。\par

\SourceRecord{Translated by Frederick Crombie.；Ante-Nicene Fathers；Vol. 4.；Edited by Alexander Roberts, James Donaldson, and A. Cleveland Coxe.；Buffalo, NY: Christian Literature Publishing Co.,；1885.；<http://www.newadvent.org/fathers/04123.htm>.}

% structure: p0001 source=h1 target=section reason=page-title

\section{\WorkTitle{论原理}（第四卷）}

% structure: p0002 source=h2 target=subsection reason=relative-heading-rank; base=h2

\subsection{译自鲁菲努斯的拉丁文本}

1. 然而，在讨论如此重要的事理时，仅将决定托付给人的感官和理解，并就如亲眼所见那样论断不可见之事，是不够的；为了确立我们所提出的观点，我们必须引证圣经的见证。为了使这见证对我们仍将提出或已经陈述的内容产生确定无疑的信念，似乎有必要首先表明圣经本身是神圣的，即由天主圣神默感而成。因此，我们将尽可能简洁地从圣经本身中提取能对我们产生适当印象的证据——从梅瑟，希伯来民族的第一位立法者，以及从耶稣基督，基督信仰体系的创始者和元首的话语中（引证）。因为虽然在希腊人和蛮族中有过许多立法者，也有无数自称宣告真理的教师和哲学家，但我们不记得有哪一位立法者能在异邦民族心中产生那样的爱慕和热忱，以致他们自愿采纳他的法律，或以全部心力捍卫它们。那么，没有人能够将自以为真理的东西引进并传播给——我不说许多外邦——甚至仅是一个民族中的个人，以至于对它的认识和信仰能普及到所有人。然而，无疑立法者希望他们的法律被所有人遵守（如果可能的话），教师也希望他们自以为正确的东西能被所有人认识。但他们知道自己在人们心中绝不可能产生如此强大的力量，以致外邦人服从他们的法律或重视他们的言论，因此他们甚至不敢尝试这种努力，唯恐事业的失败给他们的行为打上不智的烙印。然而，在全世界——在整个希腊和所有外邦——有无数人已经放弃了他们祖国的法律和他们曾信以为神的神明，顺从梅瑟的法律，并归依基督为门徒，崇拜基督；而且他们这样做时，并不免激起偶像崇拜者的强烈仇恨，因而常常遭受后者的残酷折磨，有时甚至被处死。然而，他们拥抱基督的话和教导，并以全部爱慕珍藏它们。\par

2. 而且我们还可以看到，这一宗教本身如何在短时间内发展起来，藉着其敬拜者的惩罚和死亡、财产的掠夺以及他们所受的各种折磨而取得进展；这一结果更加令人惊讶，因为连其教师本人也都非有技巧之人，人数也不很多；然而这些话传遍了全世界，以至于希腊人和蛮族、智慧人和愚拙人都采纳了基督宗教的教义。由此可以毫不怀疑地推断，耶稣基督的话凭所有人的信德和能力在所有人的理智和灵魂中得胜，并非藉着人的力量或权力。因为这些结果既由祂预言，又由来自祂的神性回答所确立，从祂自己的话可以清楚看出：\Quote{你们要为了我的缘故被带到总督和君王面前，对他们和外邦人作见证。}又：\Quote{天国的福音要被传扬到万民中。}又：\Quote{到那一天，许多人要对我说：主啊，主啊，我们不是因祢的名吃过喝过，因祢的名驱过魔吗？那时我要对他们说：离开我，你们这些作恶的人，我从来不认识你们。}如果这些话确实由祂说出，但这些预言却没有实现，那么它们或许显得不真实，没有任何权威。但现在，当祂的宣告变为现实，且如此有能力、有权威地预言时，就最清楚地显示出，祂在降生成人时，将救恩的训诫交付给人类。\par

3. 那么，对此我们该说什么呢？先知们曾预先论及祂说：权杖不会离开犹大，首领不会出自他两腿之间，直到那应得者（即王国）来到，直到外邦人的期待来到。因为从历史本身，从今日清楚可见的事实，最明显的是：自基督时代以来，犹太人中再没有君王。不仅如此，所有犹太人引以为傲、并为之夸耀的事物，无论是圣殿的华美，祭台的装饰，还是大司祭的一切圣带和长袍，都一同被摧毁了。因为那预言应验了，它宣告说：\Quote{以色列子民将多日没有君王，没有首领；没有祭献，没有祭台，没有司铎，也没有神谕。}\BibleRef{欧 3:4}因此，我们运用这些证言，来反驳那些似乎断言雅各伯在创世纪中论及犹大的人；他们说犹大支派中仍然存有一位首领——即他们所称的\OriginalTerm{族长}{Patriarch}，是他们的民族首领——并且他的后裔中不会缺少一位统治者，他将存留到他们自己所想象的那位基督来临之时。但如果先知的话是真实的，他说：\Quote{以色列子民将多日没有君王，没有首领；没有祭献，没有祭台，没有司铎；}\BibleRef{欧 3:4}并且，既然圣殿被毁后，不再献祭，祭台消失，司祭职也不复存在，那么最确定的是，正如经上所载：\Quote{权杖已离开犹大，首领已出自他两腿之间，直到那应得者来到。}\BibleRef{创 49:10}于是，那应得者已经来到，外邦人的期待就在祂内。这显然应验在那些通过基督而信奉天主、来自不同民族的众多人群身上。\par

4. 在《申命纪》的诗歌中，也预言性地宣告：因着先前子民的罪恶，将有一个愚昧的民族被拣选——当然，这正是藉着基督所实现的；因为经文这样说：\Quote{他们以偶像激起我的愤怒，我要以非神者激起他们的嫉妒；我要以愚昧的民族激起他们的怒气。}\BibleRef{申 32:21}因此，我们可以清楚看见，那些据称以偶像（即非神者）激起天主愤怒、以雕像惹动祂怒气的希伯来人，自身也被一个愚昧的民族所嫉妒，这民族是天主藉着耶稣基督及其门徒的来临所拣选的。因为宗徒这样说：\Quote{弟兄们，你们看看你们被召叫的情况：按肉眼来看，你们中有智慧的不多，有权势的不多，有尊贵的不多；但天主却拣选了世上愚笨的，和那不存在的，为消灭那存在的。}\BibleRef{格前 1:26-28}所以，属血肉的以色列不应自夸；因为宗徒用了这样的说法：\Quote{我说，没有任何血肉能在天主面前自夸。}\BibleRef{格前 1:29}\par

5. 此外，对于《圣咏》中关于基督的预言，我们该说什么呢？尤其是那首标题为\BibleQuote{\Quote{所爱的}之歌}的圣咏，其中说\Quote{祂的舌是灵敏书记的笔，你比世人更美}，\Quote{恩宠倾注于你唇间}？那么，恩宠倾注于祂唇间的明证就是：在短暂的时间之后——因为祂只教导了一年又几个月——全世界却充满了祂的教义和对祂宗教的信仰。于是，\Quote{在祂的日子，义人兴起，和平丰盛}，存留到底，这终结被称为\Quote{月亮被除去}；并且\Quote{祂的统治从这海到那海，从大河直到地极}。还有一个记号赐给达味家族：一位童贞女怀孕，生下了\OriginalTerm{厄玛奴耳}{Emmanuel}，这名字解释为\BibleQuote{天主与我们同在：列国啊，你们要知道，并要降服}。因为我们是外邦人，我们被征服和降服，成为祂胜利的战利品，将我们的颈项置于祂的恩宠之下。甚至祂的出生地也在米该亚的预言中预告了，他说：\BibleQuote{犹大地白冷啊，你在犹大的首领中决不是最小的：因为将有一位领袖从你那里出来，牧养我的以色列民。}\BibleRef{米 5:2}此外，达尼尔先知所预言的、直到基督领导时期的七年周也已应验。而且，那一位临近了，就是在《约伯传》中预言将要毁灭巨大怪兽的那一位，祂也赐给自己的门徒权柄践踏蛇和蝎子，以及敌人一切的能力，而不受伤害。但如果有人思想基督的宗徒们到各处旅行，作为祂的使者宣讲福音，他就会发现他们所敢于承担的超乎人力，而他们所能成就的惟独来自天主。如果我们思考，人们听到这些宗徒引入新教义时如何接待他们；或者更确切地说，当他们屡次想要消灭宗徒时，却被宗徒内在的神能所阻止，我们就会发现，其中没有任何人为力量，而完全是天主的能力和眷顾——无疑的神迹奇事，为他们的言语和教义作证。\par

6. 这些要点现已简要确立，即关于基督的天主性，以及所有关于祂的预言均得以应验。我认为，这一点也已得到证实：包含这些预言的圣经本身——即那些曾预言祂的来临、或祂教义的权能、或万民归顺祂的圣经——是天主所默感的。对此必须补充：先知预言和梅瑟法律的天主性与默感性，在基督降世之后，已得到清晰的揭示和确认。因为在它们所预言的事件应验之前，尽管它们真实且由天主默感，却无法被证明，因为它们尚未应验。但基督的来临宣告了它们的话语真实且由天主默感，尽管在此之前，那些所预言之事是否会成就，确实是不确定的。\par

若有人以应有的热忱和敬畏来思考先知的话语，可以肯定，在这样细读和审慎查考中，他会感到自己的心思和感官被一股神圣气息所触动，并承认他所读的话语并非人的言论，而是天主的言语；他将从自己的情感中体会到，这些书卷并非出于人的技巧，也非任何凡俗的雄辩，而是，可以这样说，出自一种神圣的风格。因此，基督来临的光辉，以真理之光照明了梅瑟法律，揭开了覆盖在字面上的那层面纱，并为每一个信靠祂的人解开了被言语遮蔽所隐藏的一切祝福。\par

7. 然而，要在每一实例中指出先知预言如何及何时应验，以便使那些怀疑者得以确认，这确实是一件相当费力的事，因为凡愿意更深入了解这些事的人，都可以从真理本身的记录中收集充足的证据。但如果字面的意义——超越人的理解——在第一眼看来，对于那些在神圣教义上不太精通的人而言，并未立即显现，这丝毫不值得奇怪，因为神圣之事是相当缓慢地降下（为人理解），并且随着一个人越是怀疑或不配，它们就越发隐没不见。因为尽管确信这世界所存在或发生的一切，都由天主的眷顾所安排，并且某些事件确实足够清晰地显明是在祂眷顾治理之下，但另一些事件却又如此神秘和无法理解地展现，以致于天主眷顾的计划完全隐藏起来；以至于有时有人相信某些特定事件不属于（眷顾的）计划，因为那按不可言喻的智慧施行天主眷顾之工的原则避开了他们的掌握；然而，这治理的原则并非对所有人都同样隐藏。因为即使在人类中，一个人对此思虑较少，另一个人则较多；而在地上的人，无论谁是天上的居民，都更承认祂。物体的本性以一种方式为我们所明了，树木的本性以另一种方式，动物的本性以第三种方式；灵魂的本性又以不同的方式隐藏；而理性理解的各种运动如何被眷顾所安排，在更大程度上避开了人的视野，甚至，依我看来，在很大程度上也避开了天使的视野。但正如神圣眷顾的存在不会被那些确信其存在的人所驳倒，尽管他们无法以人类心智的力量理解其运作或安排；同样，贯穿圣经整体的神圣默感性也不会被认为不存在，因为我们的理解力薄弱，无法追溯每个词中隐藏的奥秘意义，神圣智慧的宝藏隐藏在粗俗和不加修饰的言语器皿中，正如宗徒也指出时所说：\BibleQuote{我们有这宝贝放在瓦器里}，这样神圣德能的力量得以更明亮地闪耀，没有人的雄辩色彩混入教义真理之中。因为如果我们的书卷由于以修辞技巧或哲学智慧写成而使人相信，那么无疑我们的信德将被认为基于言语的艺术和人的智慧，而非基于天主的德能；然而如今众所周知，这宣讲的话语几乎被全世界众多人所接受，是因为他们明白自己的信仰并非基于人的智慧之动听言辞，而是基于圣神和德能的彰显。因此，我们既被一种天上的——不，是超越天上的——能力引领至信德和接纳，从而敬拜万有的独一造物主为我们的天主，也当竭尽全力，放弃基督初级道理的语言——那只是智慧的开端——前进到成全，以便那赐予成全者的智慧也赐予我们。因为那受托传讲这智慧的人如此应许：\BibleQuote{然而在成全的人中，我们也讲智慧，但不是这世界的智慧，也不是这世界将要灭亡的执政者的智慧}；由此他表明，我们的这种智慧在言语之美上，与这世界的智慧毫无共通之处。因此，这智慧若按那从永远隐藏、如今却借着预言圣经和我们的主救主耶稣基督的来临而显明的奥秘启示给我们，就更清晰、更完美地铭刻在我们心中。愿光荣归于祂，直到永远。阿们。\par

许多人由于不按属灵的意义，而错误地理解圣经，陷入了异端。\par

8. 在简要地陈述了这些关于圣经由圣神所默感的事项之后，似乎也有必要解释这一点，即：某些人由于不正确阅读圣经，而陷入错误的意见，因为许多人不知道为了理解圣经所应遵循的方法。总而言之，犹太人由于心硬，且由于渴望在自己眼中显得有智慧，不相信我们的主和救主，他们判断那些关于祂的言论应按照字面理解，即：祂应当以一种可感知和可见的方式向俘虏宣讲释放，并首先建造一座他们真正视为天主之城的城，同时砍断\Place{厄弗辣因}的战车和\Place{耶路撒冷}的马匹；祂还应当吃奶油和蜂蜜，以便在祂知道如何行恶之前选择善。他们还认为，曾有预言说，狼——那种四足动物——在基督来临时将与羔羊一同吃草，豹子将与山羊羔同卧，牛犊和公牛将与狮子一同吃草，并将由一个小孩子带领到牧场；牛和熊要在青草地上同卧，它们的幼崽一同喂养；狮子也要与牛同进牛棚，吃稻草。并且看到，按照历史，那些关于祂的预言中没有一件应验——他们认为基督来临的标记尤其应在这当中被看见——他们便拒绝承认我们的主耶稣基督的来临；不，相反一切人法和神法的原则，即违反预言的信仰，他们将祂钉在十字架上，因为祂自称基督之名。于是异端者读到法律上记载着：\BibleQuote{在我怒火中燃起了一把火}；以及\BibleQuote{我是上主，是嫉妒的天主，追讨父亲的罪于子孙直到三四代}；以及\BibleQuote{我后悔立了\Person{撒乌耳}为王}；以及\BibleQuote{我是上主，造和平也造灾祸}；以及\BibleQuote{城中没有灾祸不是上主所降的}；以及\BibleQuote{灾祸从上主降到了\Place{耶路撒冷}的城门}；以及\BibleQuote{上主的一个恶神扰乱\Person{撒乌耳}}；他们读到许多其他类似经文，在圣经中可以看到，他们不敢断言这些不是天主的圣经，但却认为这些是犹太人敬拜的那个造物主天主的话，而他们判断这位神只应被视为公义的，而非良善的；但救主来临，是为向我们宣告一位更完美的天主，他们声称这位天主不是世界的造物主——甚至在这点上他们中间也有不同的不一致的意见，因为当他们一旦背离对万物主宰、造物主天主的信仰，他们就沉溺于各种虚构和神话，编造出某些（幻想），断言一些可见之物由一位（神）创造，另一些不可见之物由另一位创造，这是他们自己头脑空虚且任意的建议。但是，在那些似乎被约束在教会信德内的人中，也有不少较为单纯的人，他们认为没有比造物主更大的天主，在这方面持有正确而健全的意见；然而，他们对天主所持有的看法，却是连对最不义、最残忍的人也不会持有的。\par

9. 如今，以上我们所提及的那些人对所有这些点的错误领会，其原因无他，正是他们不按属灵的意义、而按字面意义理解圣经。因此，我们将尽力，以我们有限的能力，向那些相信圣经并非人的著作、而是由圣神的感动所写、并经由天主圣父的旨意、通过他的独生子耶稣基督传给我们并托付给我们的人，指出按照正确理解方式观察事物的人看来，由耶稣基督传授给宗徒、并由他们相继传给后代——圣教会的导师——的准则与纪律是什么。如今，圣经中指示了某些神秘经纶，我想这一点是所有信徒，甚至是最简单的信徒都承认的。但这些经纶是什么，或属于何种性质，一个正直而不被骄傲恶习所胜的人，会谨慎地承认自己不知道。例如，若有人举出罗特女儿们的事例，她们似乎违反天主的律法与父亲同寝，或亚巴郎的两个妻子，或雅各伯所娶的两姐妹，或那两个使他的儿子增多的婢女，除此之外还能有什么回答呢？只能说这些是某些奥秘、属灵之事的预象，但我们不知道它们的性质。甚至当我们读到会幕的建造时，我们确信所写的描述是某些隐秘事物的预象；但要将它们与恰当的标准相配合，并展开讨论每一点，我认为极其困难，甚至是不可能的。然而，正如我所说，那描述充满奥秘，这连普通的理解力也不会忽视。但所有叙述部分，无论是关于婚姻、生儿育女、各种战争或任何其他历史，除了被认为是隐秘和神圣之事的形像和预象之外，还能是什么呢？然而，人们很少努力运用他们的理智，或者以为自己尚未真正学习就拥有了知识，结果他们从未开始拥有知识；或者如果没有缺乏愿望，至少没有缺乏导师，并且如果以虔诚和神圣的精神、怀着许多要点将由天主的启示（因为对人类感觉而言它们极其困难和晦涩）而得到揭示的希望来寻求神圣知识，那么，也许这样寻求的人将会找到可以发现的合法事物。\par

10. 但是，为了不让这种困难也许被认为只存在于先知的语言中——因为先知的文体被所有人承认充满比喻和谜语——当我们来到福音书时，我们发现了什么？那里不也隐藏着内在的、即神圣的含义，这含义唯独由那获得恩宠的人所启示，他曾说：\Quote{可是我们有基督的心，为能知道天主所恩赐我们的事。我们讲论这些事，不是用人的智慧所教的言词，而是用圣神所教的言词。}现在，若有人读到给予若望的启示，他岂不惊讶其中包含着如此大量隐藏、不可言喻的奥秘，以至于连那些无法理解所隐藏之物的人也能清楚地明白，确实隐藏着\Emph{某物}？宗徒的书信，虽在有些人看来较为明白，岂不也充满了如此深奥的意义，以致通过它们，有如通过某种小容器，无限光明的清晰似乎倾注给那些能够理解神圣智慧含义的人？因此，正因为如此，并且有许多人在今生误入歧途，我认为，没有危险地宣告某人知晓或理解那些需要知识之钥才能打开的事，是不容易的；救主曾宣称，那钥匙在于精通律法的人手中。在这里，虽然稍离主题，我认为我们应该询问那些宣称在救主来临之前，研究律法的人中没有真理的人：我们的主耶稣基督怎能说知识之钥在他们那里，因为他们手中有先知和律法的书？因为他如此说：\Quote{祸哉，你们这些法学士！你们拿走了知识之钥：你们自己不进去，那些愿意进去的，你们也加以阻止。}\par

11. 但是，正如我们开始注意到的，在我们看来，理解圣经并探究其含义的正确方法似乎是这样的：因为我们由圣经本身得到了关于我们应该对它形成什么观念的教导。在\BibleBook{}{箴言}中，我们发现有如下关于圣经考察的规则：他说：\BibleQuote{你要以三重方式将这些事描述给自己，以谋略和知识，使你能回答那些向你提出问题的人。}那么，每人都应以三重方式在自己的心中描述对神圣经卷的理解——也就是说，使所有较为单纯的人可以说是藉着\OriginalTerm{圣经的体}{body of Scripture}本身得到造就；因为我们称此为普通的历史意义：而如果有些人已经开始取得长足进步，并能看出更多的东西，\Emph{他们}则藉着\OriginalTerm{圣经的魂}{soul of Scripture}本身得到造就。再者，那些成全的人，就是类似宗徒所说的人：\BibleQuote{我们在成全的人中讲智慧，但不是这世上的智慧，也不是这世上有权有位、将要败亡之人的智慧；我们讲的，是从前所隐藏、天主奥秘的智慧，就是天主在万世以前预定使我们得荣耀的。}——所有这样的人都可以由属神的法律本身（这法律有未来美事的影儿）得到造就，如同藉着圣神。因为正如人被认为由肉体、灵魂和神体组成，同样，神圣的圣经也是如此，它是由天主的恩赐为人的救恩而赐予的；此外，我们在\WorkTitle{牧者}这本小册子中也看到这一点，该书似乎被一些人轻视，其中\Person{赫马斯}奉命写两本小册子，然后向教会的长老宣布他从圣神所学到的事。因为所写的话如下：他说：\BibleQuote{你要写两本书；一本给\Person{克肋孟}，另一本给\Person{格拉普特}。让\Person{格拉普特}劝诫寡妇和孤儿，让\Person{克肋孟}将信送到所有外地的城市，而你要向教会的长老宣布。}因此，\Person{格拉普特}奉命劝诫孤儿和寡妇，代表对字句本身的纯正理解；那些尚未配得以天主为父的年轻心灵由此受到劝诫，故被称为孤儿。寡妇则是指那些从不义的、与其违法结合的男人中退出的人；但她们仍为寡妇，因为尚未达到与天上新郎结合的地步。\Person{克肋孟}奉命将所写的信送到外地的城市，即给那些已经脱离字句的人——意思似乎是给那些灵魂，他们藉此受到建造，已开始超越身体的挂虑和肉欲的渴望；而\Person{赫马斯}本人，既从圣神受到教导，就奉命不是用书信或书籍，而是用活的声音向基督教会之长老宣布，即那些拥有成熟智慧，能够接受神性教导的人。\par

12. 的确，有一点不可忽略，即圣经中有些段落并不总能找到我们所说的这个\Quote{体}，即这个推论的历史意义，我们将在后面证明这一点，而只能理解我们所说的\Quote{魂}或\Quote{灵}。我认为这在福音书中有所指示：按照犹太人净化的方式，安放了六口石缸，每口可盛两三桶水；正如我说过的，福音的语言似乎表明，对于宗徒秘密地称为\Quote{犹太人}的那些人，他们因圣经的话得以净化——有时接受两桶，即如上所说的对\Quote{魂}或\Quote{灵}的理解；有时甚至三桶，即当阅读圣经时，\Quote{肉体}的意义，即\Quote{历史}意义，得以保存，以造就百姓。六口石缸的提及是合适的，因为关于那些被安置在世界中而得以净化的人；因为我们读到，在六天内——即完全的数目——这个世界和其中所有的事物都完成了。那么，我们提到的这第一种\Quote{历史}意义有多大的用处，已被所有信徒的众多所见证，他们以足够的信德和纯朴相信，并不需要多少论证，因为它对所有人都显而易见；而对于我们上面称之为圣经之\Quote{魂}的意义，宗徒\Person{保禄}在\BibleBook{格林多}{前书}中给了我们许多例子。我们看到这句话：\BibleQuote{牛在踹谷的时候，不可笼住它的嘴。}随后，在解释由此应理解什么诫命时，他补充说：\BibleQuote{天主岂关心牛吗？还是全然为我们说的？的确，是为我们写的，因为耕地的当存希望耕地，打场的也当存希望分享。}还有许多其他类似性质的段落，都按这种方式解释法律，为听者提供了广泛的信息。\par

13. 现在，\Quote{属神}的解释具有这样的性质：当一个人能够指出这些事是天上事物的模型和阴影时，谁是\Quote{按肉体}作犹太人的，以及法律包含着将来之事的影子，以及在圣经中可以找到的其他类似表述；或者当探究那隐藏于奥秘中的智慧——即\BibleQuote{天主在万世以前，为使我们获得光荣所预定的，这世代的首领们没有一个认识过}——是什么时；或者当宗徒从\BibleBook{出谷}{纪}或\BibleBook{户籍}{纪}中借用某些比喻，说：\BibleQuote{这些事发生在他们身上，是作为预像，并且写在经上，是为了我们这些活在末世的人}时，其意义是什么。现在，我们有机会理解那些发生在他们身上的事是哪些事的预像，因为他接着补充说：\BibleQuote{他们都饮了那随着自己的神迹的盘石，那盘石就是基督。}在另一封书信中，当他提到会幕时，引用了给\Person{梅瑟}的指示：\BibleQuote{你应当依照在山上指示给你的样式，制造（一切）。}而在写给迦拉达人的书信中，他责备某些自以为阅读法律却不理解的人，因为他们不知道所写的内容隐含喻意，就用一种责备的语气对他们说：\BibleQuote{请告诉我，你们愿意受法律管辖的人，有没有听过法律？因为记载说：亚巴郎有两个儿子，一个是婢女生的，一个是自由妇人生的。但是婢女所生的是按肉体生的，自由妇人生的却是藉恩许生的。这是寓意的说法：这两个妇人代表两个盟约。}在这里要注意一点，即宗徒在说\Quote{你们愿意受法律管辖的人，没有听过法律吗？}时是多么谨慎。\Quote{没有听过}，即你们不明白、不知道吗？在致哥罗森人的书信中，他又简要地总结并浓缩了整个法律的意义，说：\BibleQuote{所以不要让任何人在饮食上，或是在节期、月朔或安息日上判断你们，这些原是将来之事的影子。}在致希伯来人的书信中，论到那些属于割礼的人时，他说：\BibleQuote{他们敬拜的只是天上之事的模型和影子。}现在，也许通过这些例证，那些视宗徒著作由天主所默感的人对\Person{梅瑟}五书就不会再有疑惑。如果他们要求，对于其余的历史，其中所记载的事件也应当被视为是为那些被记载的人所设立的榜样，我们观察到这一点在致罗马人的书信中也有所提及，宗徒在那里从\BibleBook{列王}{纪上}举了一个例子，说：\BibleQuote{我为自己留下了七千人，是未曾向巴耳屈膝的。}\Person{保禄}理解这句话是比喻性地指那些按拣选而被称为以色列人的人，以表明基督的来临不仅对外邦人有益，而且连以色列族中也有许多人蒙召得救。\par

14. 既然情况如此，我们便以举例和范例的方式，概述我们在这些要点上应如何理解圣经，首先重申并指出这一点：圣神藉着天主的眷顾和旨意，通过祂的独生圣言的德能——祂在起初就与天主同在，是天主——光照了真理的僕役，即先知与宗徒，使他们理解那些在人间或关于人所发生的事物的奥秘。在此，我所指的\Quote{人}，是身处肉身内的灵魂，他们将这些藉着基督启示而为他们所知的奥秘，仿佛是某种人类事务，或流传某些法律规诫和命令，以比喻的方式描述出来；这并非为了让任何喜欢的人可以践踏这些阐释，而是为了让那些以全然的纯洁、清醒和警醒，致力于此类研究的人，能够借此追溯天主圣神的旨意——它可能深埋其中——以及上下文，其指向可能不同于通常言语的用法。这样，他就能分享对圣神的知识，参与天主的计划，因为灵魂若不藉着神圣智慧真理的默感，就无法达到知识的完美。因此，这些充满神圣之神的人主要论述的是关于天主，即圣父、圣子及圣神的事；然后是关于天主圣子的奥秘——圣言如何成了血肉，以及祂为何降下以至于取了奴仆的形像——正如我所说，这些是充满神圣之神的人所解释的题目。接下来，他们必然要以神圣的教导来教导必死的人，关于理性的受造物，包括天上的和地上较有福的；也要（解释）灵魂之间的差异，以及这些差异的起源；然后要说明这个世界是什么，以及为何被创造；何以产生了蔓延大地的巨大而可怕的邪恶。并且，那邪恶是只存在于这个地上，还是其他地方，这一点我们有必要从神圣教导中得知。既然圣神有意光照那些献身于真理服务的圣洁灵魂，使他们了解这些及类似的主题，那么第二个目的便是：为了那些不能或不愿付出劳作和辛劳以配得受教或认识如此重要之事的人，以普通语言，在可见事物的某种历史和叙述的掩盖下，包裹和隐藏奥秘，正如我们之前所说。于是，引入了关于可见受造物的叙述，以及第一个人的创造和形成；然后是其后裔的相继诞生，以及他后代中善人所行的一些事迹，偶尔也叙述了他们作为人所犯的某些罪行；之后也叙述了恶人所行的某些不洁或邪恶的作为。此外，也以奇妙的方式描述了战争，以及胜败的交替，通过这些，某些不可言喻的奥秘向那些懂得研究此类叙述的人显明。凭借智慧令人赞叹的纪律，真理的律法，甚至先知的律法，也被植入律法的圣经中，每一部圣经都由神圣的智慧之艺编织而成，作为属灵真理的一种覆盖和面纱；这就是我们所谓的圣经的\Quote{身体}，如此，我们所谓的由智慧之艺编织的文字覆盖，也能对许多人产生建树和益处，而其他人则不会从中受益。\par

15. 但是，假如在这层覆盖（即这段历史）的所有事例中，法律的逻辑联系和次序被保留下来，那么当我们如此连续地拥有圣经的意义时，我们当然不会相信其中还包含着任何超出表面所指示的东西；因此，天主的智慧刻意使某些对历史意义的绊脚石或中断发生，通过在其中（叙事中）引入某些不可能和不协调之事；这样，叙事本身的中断就像插上门闩一样，成为读者的障碍，使他拒绝承认通向普通意义的道路；而当我们被排除和阻挡在这条道路之外时，我们可能被召回另一条道路的开端，以便通过进入一条狭窄的小径，上升到更高更崇高的道路，从而展现天主智慧的无限广度。然而，我们不可忽略这一点：既然圣神的主要目的是保持属灵意义的一致性，无论是在应当做的事还是已经完成的事中，如果祂发现那些按历史已经发生的事件可以适应属灵意义，祂就会以一种叙事风格将两者交织在一起，总是更深地隐藏那隐藏的意义；但是，当历史叙述无法适应事件的属灵连贯性时，祂有时插入一些从未发生或不可能发生的事；有时也插入可能发生但未发生的事：祂有时用寥寥数语这样做，这些字句若按\Quote{字面的}意义理解，似乎不能包含真理；有时则通过插入许多话。我们在法律部分经常发现这种情况，在\Quote{字面的}诫命中，有许多明显有用的事，但也有极多的事完全看不出任何有用之处，有时甚至是被认为不可能的事。这一切，正如我们所说，都是圣神所做的，为的是让那些表面上的事件既不能是真的，也不能是有用的，从而引导我们去探究更深隐藏的真理，并在我们相信由祂所默感的圣经中确定一种与天主相称的意义。\par

16. 圣神不仅对那些直到基督来临才写成的圣经这样处理；而且，由于祂是同一圣神，并由唯一的天主所发出，祂对福音作者和宗徒们也同样这样处理。因为即使是祂默感他们写的那些叙事，也不是没有祂的智慧之助而写成的，其性质我们已在上面解释过。因此，在他们中也混杂了不少事，叙事的历史次序被打断和分裂，读者的注意力可能因事件的不可能性而被召回，以探究内在意义。但是，为了让我们的意思通过事实本身得到确认，让我们查考圣经段落。请问，哪个有理智的人会认为这个说法合适：说第一日、第二日、第三日，其中也提到晚上和早晨，却没有太阳、月亮和星辰——第一日甚至没有天空？又有谁如此无知，竟以为天主好像农夫一样，在\Place{厄登}东方的乐园里种植树木，并在其中放置一棵生命树，即一棵可见、可摸的树，好叫任何人用肉身的牙齿吃它就能获得生命，并再吃另一棵树，就能知道善恶？我想，没有人会怀疑：天主在下午在乐园中行走，\Person{亚当}藏在一棵树下的说法，在圣经中是比喻性地叙述的，为的是借此暗示某种奥秘意义。\Person{加音}离开主的面，显然会使细心的读者询问：什么是天主的临在，人如何能离开它。但是，为了不使我们眼前的任务超出适当限度，任何愿意的人都很容易从圣经中收集那些确实被记录为已发生的事，但根据历史记述却不能相信它合理而恰当地发生过的事。在福音中，圣经叙事的同样风格大量出现，例如，当魔鬼被说将\Person{耶稣}放在一座高山上，好从那里向他展示世上万国及其荣华。这如何能按字面发生呢？无论是\Person{耶稣}被魔鬼领上一座高山，还是后者向他展示世上万国（好像它们躺在他肉身的眼目之下，并且紧邻一座山），即波斯人、西古提人和印度人的国度？或者他如何能显示这些国度的君王是怎样被人荣耀的呢？凡留心阅读福音的人，都会发现许多类似的事例，并会观察到：在那些看似按字面记录的叙事中，插入和交织了一些不能按历史接受、却可在属灵意义上接受的事。\par

17. 在包含诫命的段落中，也发现了类似的情况。因为在法律中，梅瑟奉命在第八天毁灭每一个未受割礼的男性，这是极不合宜的；因为如果照历史记载法律被执行，那么就必须命令惩罚那些未给子女行割礼的父母，以及那些幼童的乳母。现在圣经的宣告是：\BibleQuote{未受割礼的男性，即未受割礼的人，应从人民中铲除。} 如果我们要探讨法律中的不可能之事，我们发现一种被称为山羊鹿的动物，它不可能存在，但梅瑟却命令在洁净牲畜之列可以吃它；还有一种狮鹫，没有人记得或听说过它能被人力所制，但立法者却禁止用作食物。关于著名的安息日遵守，他也这样说：\BibleQuote{你们各人要坐在自己的住处；在安息日不可有人离开自己的地方。} 这条诫命字面是无法遵守的；因为没有人能整天坐着不移动。关于这些点，那些属于割礼的人，以及所有认为圣经的含义仅限于字面的人，认为不应研究山羊鹿、狮鹫和秃鹫；他们根据某些传统来源编造一些空洞琐碎的故事，声称每个人的地方被计算为两千肘以内。另一些人，其中包括撒玛黎雅人多西特奥斯，虽然批评这种解释，但他们自己却提出更可笑的东西，即每个人在安息日必须保持他当时所处的姿势、地方或位置直到晚上；即如果在坐着，就要整天坐着，如果在躺着，就要整天躺着。此外，\BibleQuote{在安息日不可担担子} 这条诫命在我看来是不可能的。因为犹太医生们，如圣宗徒所说，因这些（规定）而转向了无数的虚构故事，说若穿无钉的鞋子不算负担，但穿有钉的鞋子就算负担；若用单肩担担子，他们认为是负担，但若用双肩，他们则宣称不是负担。\par

18. 现在，如果我们对福音书进行类似的考察，那么字面理解\BibleQuote{在路上不要向任何人请安}的诫命，岂不是显得荒谬？然而有些头脑简单的人，竟认为这是我们的救主赐给祂宗徒的命令！再者，这样一条命令，尤其是在那些冬季严寒、霜冻冰封的地区，怎么可能遵守呢？即：要\BibleQuote{不要有两件内衣，也不要穿鞋}？还有，当有人打右脸时，要连左脸也转过来给他打，因为用\Emph{右}手打人的人，打的是\Emph{左}脸。福音中还有一条诫命必须被视为不可能，即：若右眼\BibleQuote{使你犯罪}，就挖出来；即使我们假定是指肉体的眼睛，又怎会显得合宜，即当两只眼睛都有视力时，\Quote{犯罪}的责任要归到一只眼睛上，而且是右眼？谁又能被认为是无罪的，如果他对自己下手，犯下了最严重的罪行？也许保禄宗徒的书信似乎超越这些。因为当他说：\BibleQuote{有人受割礼蒙召吗？不要掩盖割礼的痕迹。} 这是什么意思？首先，仔细考虑，这话似乎不是针对他当时所讨论的主题，因为这篇讲论是关于婚姻和贞洁的诫命；因此这些话对于这样的主题来说，似乎是一个不必要的添加。其次，如果一个人为了避免割礼引起的不雅观而能掩盖割礼的痕迹，那又有什么不妥呢？第三，这根本是不可能的。\par

我们所有这些陈述的目的，是要表明圣神（祂屈尊赐予我们圣经）的设计，是要我们不仅因字句本身，也不是因其中一切而得益——我们看到这常常是不可能和矛盾的；因为那样不仅会导致荒谬，还会导致不可能；而是要我们理解，某些事件被编织在\Quote{可见的}历史中，当按内在意义思考和理解时，它们会颁布一条对人有益且配得天主的法律。\par

19. 然而，不可有人抱持这样的怀疑，即我们认为圣经中的任何记载并非真实，因为我们怀疑其中所记述的某些事件并未发生；或者认为法律的规诫没有一条应按字面理解，因为我们考虑其中某些规诫——按其本性或可能性之要求——是不可能遵守的；或者我们不相信那些关于救主的预言，在感官可感知的方式上已经应验；或者认为祂的诫命不应按字面遵守。因此，我们要回答：既然我们明显持此意见，多数情况下历史事实可以而且应当予以保留。因为谁能否认\Person{亚巴郎}葬在\Place{赫贝龙}的双洞中，连同\Person{依撒格}、\Person{雅各伯}以及他们各自的妻子？谁会怀疑\Place{舍根}已赐给\Person{若瑟}作为产业？或者\Place{耶路撒冷}是\Place{犹大}的京城，\Person{撒罗满}在此为天主建造了圣殿？——还有其他无数记载。因为按历史意义成立的段落，远多于那些纯属属灵意义的段落。再者，谁不主张\BibleQuote{应孝敬你的父亲和你的母亲，好使你幸福，}这条诫命，无需任何属灵意义，其本身就足以成立，且为遵守者所必需？尤其是宗徒\Person{保禄}也曾以同样话语重申而予以确认。至于禁令：\BibleQuote{不可奸淫，}\BibleQuote{不可偷盗，}\BibleQuote{不可作假见证，}以及其他同类诫命，又有何可说的？关于福音中所吩咐的规诫，无疑许多必须按字面遵守，例如我主说：\BibleQuote{我却对你们说：你们总不可发誓；}祂又说：\BibleQuote{凡注视妇女，有意贪恋她的，他已在心里奸淫了她；}还有宗徒\Person{保禄}著作中的劝诫：\BibleQuote{劝诫不守规矩的人，抚慰怯懦的人，扶持软弱的人，对众人要忍耐，}以及其他许多。然而我毫不怀疑，细心的读者在许多情况下会犹豫：某段历史是否能视为字面真实？某条规诫是否应按字面遵守？因此，必须付出大量的辛劳和努力，直到每位读者以敬畏之心明白，他是在处理圣经中所载的神圣言词，而非人的言语。\par

20. 因此，我们认为应当恰当地、一贯地坚持的圣经理解如下：圣经宣告，在地上有一个民族被天主所拣选，这个民族有几个名称：有时整个民族称为\Place{以色列}，有时称为\Person{雅各伯}；它被乃巴特之子\Person{雅洛贝罕}分为两部分；在他统治下形成的十个支派称为\Place{以色列}，而剩下的两个支派（与它们联合的还有\Person{肋未}支派和出自\Person{达味}王族的后裔）称为\Place{犹大}。该民族所拥有的整个国土，是他们从天主所领受的，称为\Place{犹大}，其中设有京城\Place{耶路撒冷}；它被称为京城，因为它仿佛是许多城市的母亲，这些城市的名字你会在圣经其他书卷中时常看到，但它们被集中在农的儿子\Person{若苏厄}的\BibleBook{若苏厄}{书}中的一个目录里。\par

21. 既然如此，宗徒为了使我们的理解在一定程度上提升并超脱尘世，在某处说：\BibleQuote{按血肉看，这就是以色列；}他无疑是指另有另一个以色列，不是按血肉，而是按圣神。又在另一处说：\BibleQuote{凡出于以色列的，并不都是以色列人。}\par

22. 我们既由祂教导，知道有一个按肉体的以色列，也有一个按圣神的以色列。当救主说：\Quote{我被派遣，无非是为以色列家迷失的羊。}\BibleRef{玛 15:24}，我们并不像那些留恋尘世事物的人那样理解这些话，即像\OriginalTerm{贫穷派}{Ebionites}，他们从其名称本身便得出\Quote{贫穷}的称谓（因为\Quote{Ebion}希伯来语意为\Quote{贫穷}）；而是理解到存在一类称为\Quote{以色列}的灵魂种族，正如这名称本身的解释所指明的：因为以色列被解释为\OriginalTerm{心智}{mind}，或\OriginalTerm{见天主的人}{man seeing God}。\newline{} 宗徒又对耶路撒冷作了类似的启示，说：\Quote{那在上耶路撒冷是自由的，她是我们的母亲。}\BibleRef{迦 4:26}。在另一封书信中他说：\Quote{你们却来到了熙雍山，和永生天主的城，天上的耶路撒冷，以及千万天使的集会，和那些已被登录在天上的首生者的教会。}\BibleRef{希 12:22-23}。\newline{} 因此，若这世上有某些灵魂称为以色列，天上有一个城称为耶路撒冷，那么那些据说属于以色列民族的城邑，便以天上的耶路撒冷为其京城；据此，我们应将那些关于犹大全体的预言（我们认为先知们也在某些奥秘叙述中论及了这犹大全体），无论它们关乎犹太或耶路撒冷，或关乎任何种类的入侵（圣史记载这些事曾发生在犹太或耶路撒冷），都理解为与此相符。\newline{} 那么，凡关于耶路撒冷所述或所预言的一切，若我们接受保禄的话如同基督在他内说话，就必须理解为按照他对那城（他称之为天上的耶路撒冷）以及所有那些被称为圣地的城邑（耶路撒冷是其中的京城）的看法所说的。因为我们应当推定，正是从这些城邑，救主为提升我们达到更高的理解层次，应许那些妥善管理祂所托付金钱的人，将有权管理十座或五座城。\newline{} 因此，若关于犹太、耶路撒冷、犹大、以色列、雅各伯的预言，若不按肉性理解，便意味着某些神圣奥秘，那么当然，那些关于埃及或埃及人、巴比伦和巴比伦人、漆冬和漆冬人的预言，也不应理解为论及位于地上的埃及，或地上的巴比伦、提洛、漆冬。厄则克耳先知关于埃及王法郎的预言，也不能适用于任何似乎曾统治埃及的人，正如经文本身的性质所表明的那样。同样，关于提洛王子的言论，也不能理解为任何提洛人或提洛王。我们怎么能接受在许多圣经段落，尤其是依撒意亚中关于拿步高的记载，是论及一个人的呢？因为那被说成\Quote{从天上坠落}，或称为\Quote{晨星}，或\Quote{清晨升起}的，并不是一个人。\BibleRef{依 14:12}。\newline{} 至于厄则克耳中关于埃及的预言，例如它将在四十年内被毁灭，以致人的脚踪不再存于其中，并且它将遭受如此荒废，以致全地人的血高至膝盖，我不知道有理智的人会将其指为那个靠近埃塞俄比亚的地上埃及。但让我们看看是否更合适地按以下方式理解：即如同有天上的耶路撒冷和犹太，并且无疑有一个民族居住其中，称为以色列；同样，在这些地方附近，可能存在某些被称为埃及、巴比伦、提洛、漆冬的地区，而这些地方的王子以及居住其中的灵魂（如果有的话）被称为埃及人、巴比伦人、提洛人、漆冬人。从这些人那里，根据他们在那里的生活方式，似乎会产生某种被掳，因为他们从更高更好的境况，从犹太坠落到巴比伦或埃及，或者被分散到其他国家。\par

23. 因为正如那些因着众人共同的死亡离开此世的人，按其行为与功过——依其所堪受的——有的被安置在称为\Quote{地狱}的地方，有的在亚巴郎的怀抱中，或在不同处所或住所；同样，从那些地方，仿佛在那里死去（若可用此说法），他们从\Quote{上界}下来到此\Quote{地狱}。因为我认为，因着这区别，圣经称此灵魂由今世被引往的\Quote{地狱}为\Quote{下界的地狱}，正如\BibleBook{}{圣咏集}所说：\BibleQuote{你曾从最深的阴府中救出了我的灵魂}。因此，每一位降世为人者，按其功过，或按其在那里的地位，被注定生于此世，在不同的国家、不同的民族、不同的生活方式中，或带着不同种类的病弱，或出于虔诚的父母或不虔诚的父母；因此有时会发生一个以色列人降生于西徐亚人中，而一个贫穷的埃及人被带至犹太地。然而我们的救主来是为召集以色列家迷失的羊；但许多以色列人未接受祂的教导，于是外邦人蒙召。由此看来，那些向各民族发出的预言应被归于灵魂及其不同的天上住所。甚至那些关于以色列民族、耶路撒冷或犹太地遭受这国或那国攻击而说发生的事件叙述，在许多情况下不能理解为实际发生过，而更适合于那些居住在那所谓要过去的天空中的灵魂民族，或甚至现今被认为是其居民者。\par

若有人要求我们从圣经中就此给出明确清晰的陈述，我们必须回答：圣神的旨意是在那些看似记述历史事件的部分，更倾向于遮盖和隐藏其含义；例如在那些说到下到埃及或被掳到巴比伦的段落，或在这些国家中有些人被置于极度卑微并受其主人奴役；而另一些人在这些被掳之地却受尊敬和重视，以至于居高位掌大权，并被派管理省份——所有这些，如我们所说，在圣经的叙述中被隐藏和遮盖，因为\BibleQuote{天国好像藏在田里的宝贝；人找到了，就把它藏起来，高兴地去卖掉他所有的一切，买了那块田}。借此比喻，请思考是否指出：圣经的土壤和表层，即字面意义，就是那田，充满各种植物和花朵；而那更深奥的\Quote{属神}意义，正是隐藏在智慧与知识中的宝贝，圣神借依撒意亚称之为黑暗、无形和隐藏的宝藏；为寻得这些，需要天主的助佑：因为唯有天主能打破那封锁隐藏它们的铜门，粉碎那阻止接近一切写在和藏在\BibleBook{}{创世纪}中关于各种灵魂、种子和世系——它们或与以色列紧密相连或与他的后裔远隔——的铁栓和门闩；以及那七十个灵魂下到埃及，后来在那地成为天上的众多星辰。但并非他们全是这世界的光——\BibleQuote{因为凡是以色列人不都是真以色列人}——他们从七十个灵魂增长为一个重要的民族，如同\BibleQuote{海边的沙，不可胜数}。\par

% structure: p0030 source=h2 target=subsection reason=relative-heading-rank; base=h2

\subsection{译自希腊文}

\Emph{此希腊文译本有意逐字直译，以更清楚显示原文与鲁菲努斯意译之间的差异。}\par

1. 既然在我们研究如此重要的事理时，不满足于普遍的意见和可见之事的清楚证据，我们便额外采取那些被我们认为是神圣著作的见证，即所谓\WorkTitle{旧约}和所谓\WorkTitle{新约}，来证明我们的陈述，并努力以理性坚固我们的信德；又因我们尚未论及圣经是神圣的，那么，我们且来，仿佛以概要的方式，讨论有关它们的几点，并陈明那些使我们视之为神圣著作的理由。在利用这些著作本身的言语和其中所展示的事理之前，我们必须作出以下关于\Person{梅瑟}和\Person{耶稣基督}的陈述：\Person{梅瑟}是希伯来人的立法者，而\Person{耶稣基督}是根据基督教的救恩教义的引入者。因为，虽然在\Place{希腊}人和蛮族人中有过许多立法者，也有宣布自称是真理的见解的教师，但我们未曾听说有哪一位立法者能够使其他民族热心接受他的言语；虽然那些自称就真理作哲学思考的人提出了大量看似逻辑论证的工具，却没有一个人能够将他所认为的真理印刻在其他民族，甚至在一个民族中任何值得提及的若干人身上。然而，立法者们不仅愿意（若可能）将那些看来是好的法律强加于全人类，而且教师们也愿意将他们认为的真理传遍整个世界。但由于他们无法召唤其他语言和许多民族的人来遵守他们的法律并接受他们的教导，他们根本没有尝试这样做，不无明智地考虑到这种事发生在他们身上的不可能性。然而，整个\Place{希腊}，以及我们世界中的蛮族地区，却有无数的热忱者，他们离弃了祖传的法律和既定的神明，去遵守\Person{梅瑟}的法律并追随\Person{耶稣基督}的言语；尽管那些坚持\Person{梅瑟}法律的人被拜偶像者憎恨，而那些接受\Person{耶稣基督}言语的人还面临着死亡的威胁。\par

2. 此外，如果我们观察言语在短短几年内变得多么强大，尽管那些承认基督信仰的人遭遇了阴谋，其中一些人因此被杀，另一些人失去了财产，而且尽管传道的人数很少，它还是传遍了整个世界，以至于\Place{希腊}人和蛮族人、有智慧的和愚拙的，都投身于藉着\Person{耶稣}的敬拜，那么我们毫不困难地说，这结果超乎任何人力，因为\Person{耶稣}以全部权威和说服力教导，祂的言语不可战胜；因此我们可以正当地将祂的那些话视为神谕，例如：\BibleQuote{你们要为了我的缘故被带到总督和君王面前，对他们和外邦人作证。}\BibleRef{玛 10:18} 以及\BibleQuote{在那一天，许多人要对我说：主啊，主啊，我们不是藉着你的名吃过喝过，藉着你的名赶过魔鬼吗？我要对他们说：离开我吧，你们这些作恶的人，我从来不认识你们。}\BibleRef{玛 7:22-23} 现在，或许曾经有可能，祂说出这些话是徒然的，以致它们不真实；但当那以如此大权威传讲的话\Emph{已经}成就时，就表明天主真的降生成人，并将救恩的教义传给了人。\par

3. 又何必提及有关基督的预言：那时，当那所保留的（即王国）的那位来到时，权杖必离开犹大，与从祂两腿之间出来的领袖；并且外邦人的期待将居住在这地上？因为从历史和今天所见的情况清楚显明，从\Person{耶稣}的时代起，不再有任何被称为犹太王的人；所有他们引以为傲的犹太制度——我指的是有关圣殿、祭台、祭祀礼仪的供献和大司祭礼服的安排——都已被摧毁。因为那预言得到了应验，它说：\BibleQuote{以色列的子孙将要坐守许多日子，没有君王，没有首领，没有祭献，没有祭台，没有司祭，没有神谕。}\BibleRef{欧 3:4} 我们使用这些预言来回答那些对\Person{雅各伯}在\WorkTitle{创世纪}中对\Person{犹大}说的话感到困惑的人，他们声称那民族首领，出于犹大支派，是人民的统治者，并且他的后裔必不断绝，直到他们想象的那位基督来临。但倘若\BibleQuote{以色列的子孙将要坐守许多日子，没有君王，或首领，或祭台，或司祭，或神谕；}\BibleRef{欧 3:4} 并且，既然圣殿已被摧毁，不再有祭献、祭台、司祭，那就显明权杖\Emph{已经}离开了犹大，领袖已从祂两腿之间出来。既然预言宣告：\BibleQuote{权杖必不离开犹大，与从祂两腿之间出来的领袖，直到那所保留的那位来到，}\BibleRef{创 49:10} 那就显明那所保留的——外邦人的期待——的那位已经来到。这一点从通过\Person{耶稣基督}而相信天主的外邦人多得不可胜数可以清楚看出。\par

4. 在\BibleBook{}{申命纪}的颂歌中，也先知性地宣告：因先前百姓的罪，要拣选一个愚拙的民族，这除了藉耶稣基督外无人能成就。祂说：\BibleQuote{因为他们以那不是天主的，惹我嫉妒；以他们的偶像，激我愤怒。我也要以那不是子民的，惹他们嫉妒；以愚拙的民族，激他们愤怒。}如今可以极其清楚地明白：那些被称为以那不是天主的惹天主嫉妒、以偶像激祂愤怒的希伯来人，自己也被那不是子民的——即愚拙的民族所激怒；这民族就是天主藉耶稣基督及其门徒的来临所拣选的。我们实在看见\BibleQuote{我们蒙召，按肉体不是多有智慧的人，不是多有能力的，不是多有尊贵的；但天主却拣选了世上愚拙的，叫有智慧的羞愧；又拣选了世上软弱的，叫那强壮的羞愧；也拣选了世上卑贱的、被人轻视的，以及那无有的，为要废掉那有的。}并且按肉体的以色列——宗徒称之为\BibleQuote{肉体}——不可在天主面前自夸。\par

5. 关于\BibleBook{}{圣咏}中论基督的预言，我们还能说什么呢？有一首诗歌的题注是\BibleQuote{为所爱者}，其舌头被称作\BibleQuote{敏捷文士的笔，比世人更美}，因为\BibleQuote{恩宠倾注在祢唇上}？恩宠倾注在祂唇上的明证就是：虽然祂教导的时期很短——祂教导了大约一年零几个月——但世界却充满了祂的教导，以及藉祂所建立的崇拜。因为在祂的日子，\BibleQuote{正义与和平丰盛}，这和平一直持续到终局，即被称为月亮被除去的时候；祂继续\BibleQuote{管辖从这海到那海，从大河直到地极}。给达味家室一个记号：因为童贞玛利亚怀孕生子，给祂起名叫厄玛奴耳，即\BibleQuote{天主与我们同在}；正如同一先知所说，预言应验了：\BibleQuote{天主与我们同在；万民啊，你们要知道，并且屈服；你们有能力的，要被打败。}因为我们这些外邦人已经被打败、被征服了，被祂教导的恩宠所俘获。祂的出生地点也在\BibleBook{米该亚}{}的预言中被预言了：\BibleQuote{至于你，白冷，犹大地啊，你在犹大诸领袖中绝非最小的；因为将由你出生一位领袖，祂将牧养我的百姓以色列。}按照\BibleBook{达尼尔}{}，七十周直到领袖基督来临。祂来了，按照\BibleBook{约伯}{}，祂制伏了大鱼，并赐给祂的真门徒权柄践踏蛇和蝎子，以及仇敌一切的能力，而不受丝毫伤害。还要注意耶稣所派遣的宗徒们普遍地前来宣讲福音，就会看到这事业超乎人力，且命令来自天主。如果我们查考：人们如何听到新教义和奇异言论，就顺服于这些导师，尽管他们有意谋害，却因眷顾这些导师的神力而被征服，我们就会相信他们也行过奇迹，天主以神迹、奇事和种种异能来证实他们的话语。\par

6. 当我们这样简要地证明基督的天主性，并运用关于祂的预言时，同时也证明了那些预言祂的著作是受天主启示的；那些宣告祂来临和祂教导的文献是带着一切权能和权威赐下的，正因如此，它们在外邦人中获得了拣选。我们还必须说：预言的天主性以及梅瑟法律的属灵性质，在基督来临后显现出来。因为基督来临之前，不可能完全展示古圣经受天主启示的明显证据；而祂的来临使那些可能怀疑律法和先知不是神圣的人，清晰地确信它们是由天上的恩宠所写成的。凡仔细用心阅读先知话语的人，因着阅读本身就感受到其中蕴含的天主性痕迹，必因自己的情感而被引导，相信这些被视为天主话语的文字不是人的作品。此外，梅瑟法律中所含的光，曾被帕子遮蔽，如今在耶稣来临之时，帕子被除去，那包含在字句中的影儿所指向的福分，逐渐展现（给人认识）。\par

7. 现在枚举关于各未来事件的最古老预言将显得繁琐，为使怀疑者因它们的神性而受感动，放下一切犹豫和分心，将自己全灵魂奉献于天主的言语。但如果圣经的每一部分，那超越人类思想的因素似乎并不呈现给无学识者，那一点也不足为奇；因为就那统摄全世界的天主眷顾的工作而言，有些显示出它们是最清楚的天主眷顾之作，而另一些则如此隐藏，以致似乎为对那以不可言喻的技巧和权能安排万有的天主提供不信的根据。因为眷顾统治者的巧妙计划，在属于地上的事物中，并不如太阳、月亮和星辰的情况那样明显；在与人类事件相关的事上，也不如动物灵魂和身体方面那么清楚——动物冲动、幻象、本性的目的和理由，以及它们的身体结构，都被留意这些事物的人仔细查明了。但是正如眷顾的道理丝毫没有因那些不被理解的事而在诚心接受它的人眼中被削弱，同样，圣经的神性，它遍及全部圣经，也并未因我们软弱的无能无法在每一表达中发现隐藏在平凡而不吸引人的措辞中教义的隐藏光辉而丧失。因为我们有这宝贝放在瓦器里，为要使天主能力的光荣得以彰显，并使人不以为它是出于我们这些不过是凡人。因为如果圣经书卷中所有人之常情的陈旧论证方法成功地产生了确信，那么我们的信德就会被认为建立在人的智慧上，而不是在天主的能力上；但现在每一个举目仰望的人都清楚看到，言语和宣讲之所以在群众中得胜，\Quote{不是用智慧动人的言辞，而是用圣神和能力的明证。}因此，既然一种属天的甚至超天的能力迫使我们钦崇唯一的造物主，就让我们离开基督开端之道，即初步道理，努力前进，达到成全，好叫那讲给成全人的智慧也讲给我们。因为那拥有这智慧的人应许在成全人中讲智慧，却是另一种智慧，不同于这世界的智慧和这世界统治者的智慧，这智慧是归于无有的。这智慧将清晰地铭刻在我们身上，并产生对那永世所隐藏之奥秘的启示，这奥秘如今却借着先知性圣经和我们主、救主\Person{耶稣基督}的显现而被彰显了。愿光荣归于祂，直到永远。阿们。\par

8. 如此简略地论及圣经的默感之后，有必要继续考虑它们应被阅读和理解的方式，因为由于未曾被大众发现应如何查考神圣文献的方法，许多错误已经犯下了。因为无论是心硬的人，还是属于受割损的无知者，都不相信我们的救主，以为他们遵循着关于祂的先知预言的说法，却没有以他们感官可觉察的方式察觉祂已向俘虏宣告了自由，也没有察觉祂已建造了他们真正认为是天主的城的那座城，也没有剪除\Quote{以法莲的战车和耶路撒冷的战马}，也没有吃奶油和蜂蜜，并且在知道或偏爱恶事之前，已经选择了善事。他们进而认为预言了狼——四足动物——要与羔羊一同吃食，豹要与山羊羔一同躺卧，牛犊、公牛和狮子要一同吃食，由一个小孩引领；牛和熊要一同牧放，它们的幼崽一同长大，狮子要像牛一样吃草。看到这些事没有一件在我们认为是基督的那位降临时明显应验，他们就不接受我们的主耶稣；反而因祂不当地自称基督而把祂钉在十字架上。而那些属于异端教派的人读到这些话：\Quote{在我怒火中燃起了火；}和\Quote{我是忌邪的天主，追讨父亲的罪孽，直到三四代；}和\Quote{我后悔立了撒乌耳为王；}和\Quote{我是造和平、降灾祸的天主；}以及，此外，\Quote{城中没有一件邪恶是上主没有做的；}和\Quote{灾祸从上主那里降到耶路撒冷的城门上；}以及\Quote{从主那里来的邪魔搅扰撒乌耳；}以及无数类似段落——他们不敢不信这些是天主的圣经；却相信它们是\OriginalTerm{德缪哥}{Demiurge}的话，犹太人敬拜这位德缪哥。他们认为，既然德缪哥是一位不完全、不仁慈的天主，那么救主来就是宣告一位更完满的神祇，他们说这位神不是德缪哥，对祂持有不同的意见；并且一旦离开了那唯一非受造的天主德缪哥，他们就投身于虚构，为自己捏造各种假说，照此想象有些事物是可见的，另一些是不可见的，这一切都是他们自己头脑的幻想。然而，确实，在那些自称属于教会的人中，较为单纯的人认为没有比德缪哥更大的神祇，他们这样想是对的，但同时他们对祂所想象的事情，是连最野蛮、最不义的人都不会相信的。\par

9. 现在，在所有先前列举的各点中，错误意见、对天主的亵渎性说法或无知断言的原因，似乎不外乎是不按圣经的属灵意义理解圣经，而仅按字面解释。因此，对于那些相信圣书不是人的作品，而是由圣神感动、按照万有之父的旨意通过耶稣基督写成的，并已流传到我们这里的人，我们必须指出在我们看来（正确的）解释方法，因为我们紧守按照宗徒传承的耶稣基督天上教会的标准。现在，所有——甚至是最简单的那些持守圣言的人——都相信圣经中揭示了某些奥秘的经纶；但这些是什么，坦诚谦逊的人承认他们不知道。那么，如果有人对罗特与他女儿们的关系、亚巴郎的两个妻子、雅各伯娶的两个姐妹以及为他生子的两个婢女感到困惑，他们只能回答：这些是我们不理解的奥秘。还有，当读到会幕装备（的描述）时，因为相信所写的是预像，他们寻求将所能理解的适用于与会幕相关的每一细节——就其相信会幕是\Emph{某物}的预像而言，他们并没有错，但有时在将以会幕为预像的那对象的描述适用于某种具体事物时，犯了不配于圣经的错误。所有被认为是讲述婚姻、生育子女、战争或任何在民众中流传的历史，他们都宣布为预像；但在每个具体实例中，部分由于他们的习惯未经充分训练，部分由于他们的仓促，有时即使一个人确实受过良好训练且眼光清晰，也由于人类发现事物方面的过度困难——关于这些（预像）的每一特定的性质并未被清楚查明。\par

10. 至于预言，我们所有人都知道它们充满了谜语和隐晦的说法，还需要多说什么呢？如果我们来到福音书，对这些福音书的准确理解，如同基督的心意，也需要赐给那位说这话的人的恩宠：\BibleQuote{但我们拥有基督的心意，好使我们知道天主白白赐给我们的一切。这些事我们也讲论，不是用人的智慧所教导的言语，而是用圣神所教导的。} 再者，谁读到启示给若望的启示，不会对其中隐藏的、连不明白所写内容的人也显而易见的不可言说的奥秘感到惊讶呢？对于一个善于探究言辞的人，宗徒书信会显得清晰易懂吗？因为即使在它们里面，也有无数最深刻的思想，这些思想（通过一个开口）流露出来，却无法迅速理解。因此，既然事情如此，而且无数人犯错误，那么在阅读（圣经）时，就不可轻易宣称自己理解了需要知识之钥匙的内容，这钥匙救主说是在律法师那里。让那些不允许在基督来临前真理就与这些人同在的人回答：吾主耶稣基督怎么说知识之钥匙在那些据称没有包含知识秘密和完美奥秘的书籍的人那里呢？因为他的话是这样：\BibleQuote{祸哉，你们律法师！因为你们拿走了知识之钥匙；你们自己不进去，那些正要进去的，你们也阻止了。}\par

11. 我们处理圣经、从其提取意义的方法，照我们看，是由圣经本身所确定的。我们在\WorkTitle{箴言}中找见撒罗满所颁布的这样一条关于圣经神圣教义的规则：\Quote{你要以三重方式描绘它们，在谋略和知识中，以真理之言回答那些向你提出它们的人。} 于是，个人应当以三重方式把圣经的思想描绘在自己的灵魂上：为使纯朴的人能因圣经的\Quote{肉身}（我们这样称呼明显的含义）而受到启迪；而那已上升一段路程的人，可以因\Quote{灵魂}而受到启迪。至于成全的人，他像宗徒所说的人，宗徒说：\Quote{我们在成全的人中讲智慧，但不是这世界的智慧，也不是这世界将要消灭的统治者的智慧；我们讲的乃是天主奥秘的智慧，那隐藏的智慧，是天主在万世以前预定为我们的光荣而预备的，}他可从具有将来美事影子的精神法律受到启迪。因为正如人由身体、灵魂和精神组成，同样，为救赎人类而由天主安排赐予的圣经也是如此。因此，我们也从一个被某些人轻视的著作——\WorkTitle{牧者}——中推导出这一点，关于命令\Person{赫尔马}写两本书，然后向教会的长老宣布他从圣神所学到的事。其言如下：\Quote{你要写两本书，一本交给\Person{克莱孟}，一本交给\Person{格拉普特}。\Person{格拉普特}要劝诫寡妇和孤儿，\Person{克莱孟}将把书送到外邦城市，而你要向教会的长老宣布。} 现在，\Person{格拉普特}劝诫寡妇和孤儿，正是指圣经的字面意义，它劝诫那些灵魂上仍是孩子、不能称天主为父、因此被称为孤儿的人，并且也劝诫那些不再有不合法的新郎、却仍是寡妇的人，因为他们还不配成为（天上）新郎的配偶；而\Person{克莱孟}，已经超越字面意义，据说将所写的送到外邦城市，好像我们要称这些人为\Quote{灵魂}，他们超脱了肉体和卑微思想的（影响）——圣神自己的门徒奉命不再用文字，而是用活生生的话，向整个天主的教会的长老们讲明，这些长老因智慧而头发灰白。\par

12. 但是，由于圣经中有某些段落根本不包含\Quote{肉身}含义，正如我们在以下（段落）中将要指出的，也有一些地方我们只应寻求圣经的\Quote{灵魂}和\Quote{精神}。或许正因如此，\BibleBook{若望}{福音}记载说，为犹太人取洁用的水缸，每口缸可盛两或三桶：这话暗指那些被宗徒称为暗中的\Quote{犹太人}的人，他们藉圣经之言得到净化，有时接受两桶，也就是说，接受\Quote{魂性}和\Quote{精神}含义；有时接受三桶，因为有些人除此之外还接受\Quote{肉身}含义，后者能产生启迪。六口水缸合理地对应在六天内造成的世界上受净化的人——六是完美的数目。那么，第一种\Quote{含义}之所以有益，在于它能传授启迪，这由众多真诚而纯朴的信徒所证实；至于那被归结于\Quote{灵魂}的解释，在\Person{保禄}的\BibleBook{格林多}{前书}中有一个例证。其言道：\Quote{你不可笼住牛嘴，而牛在打谷；}\Person{保禄}接着又说：\Quote{天主岂关心牛吗？还是完全为我们说的？的确，这是为我们写的，因为耕种的应怀着希望耕种，打谷的应怀着分享的希望。}\BibleRef{格前 9:9-10} 此外还有许多适应群众的解释在流传，它们启迪那些不能理解更深含义的人，并且具有某种相同的性质。\par

13. 但若有人能指出犹太人\Quote{按肉体}所事奉的，是怎样成为天上之事的模型和影子，以及法律所包含的，是怎样成为未来福泽的影子，这样的解释就是\Quote{属神的}。一般而言，我们必须按照宗徒的应许，探究\Quote{奥秘中的智慧，那隐藏的智慧，是天主在万世之前为义人的荣耀所预定的，这智慧是今世的首领中没有一个认识的}。同一位置，宗徒在引述出谷纪和户籍纪中提到某些事件时，说：\Quote{这些事发生在他们身上，是作为预像，并且写下来，是为警戒我们这些生活在世末的人}。他又给了我们辨别这些事是何种模型的机会，说：\Quote{因为他们饮自随后的神石，而那石就是基督}。在另一封书信中，描绘帐幕的种种时，他用词说：\Quote{你要依照在山上指示你的模型，制造一切}。此外，在致迦拉达人书中，他仿佛责备那些自以为诵读法律却不理解的人，断定那些不明白所写内容包含寓意的人是不理解的，他说：\Quote{请告诉我，你们愿意在法律之下的人，没有听见法律吗？因为记载说：亚巴郎有两个儿子：一个由婢女所生，一个由自由妇人所生。那由婢女所生的，是按肉体生的；那由自由妇人所生的，却是藉着恩许生的。这都是寓意：这两个妇人是两个盟约}，等等。现在我们必须仔细留意他所用的每一个词。他说：\Quote{你们愿意在法律之下的人}，而不是\Quote{你们在法律之下的人}；并且\Quote{你们没有\Emph{听见}法律吗？}——\Emph{听见}被理解为\Emph{领悟}和\Emph{知道}。在致哥罗森人书中，他概要地总结了全部立法的意义，说：\Quote{所以，不要有人在饮食上，或是在节期、月朔、安息日上判断你们，这些都是未来之事的影子}。此外，在致希伯来人书中，论及属于割礼的人时，他写道：\Quote{他们事奉的是天上之事的模型和影子}。从这些例证来看，那些一经接受宗徒为神圣灵感教导的人，对梅瑟五书大概不会再有什么疑问；但你想知道其余的历史是否也是按模型发生的吗？那么我们必须注意致罗马人书中引述自《列王纪》第三卷的话：\Quote{我为自己保留了七千人，他们未曾向巴耳屈膝}，保禄认为这话在意义上相当于按拣选而成为以色列人的那些人，因为基督的来临不仅惠及外邦人，也惠及天主的族类中的某些人。\par

14. 情况既然如此，我们须勾勒出我们认为的（真正）理解圣经的标记。首先必须指出：圣神藉着天主的眷顾，通过那从起初就与天主同在的圣言，光照了真理的仆役——先知和宗徒，其目的尤其在于（传达）关于人类事物的不可言喻的奥秘（现在我所说的人，是指那使用身体作为工具的灵魂），以便那能受教的人能藉着探究，专心研究字句中所含深奥意义，而参与祂旨意的一切教训。在这些与灵魂（除非藉着天主丰盛而智慧的真理，否则无法达到完美）有关的事物中，有关天主及其独生子的（道理）必须定为首要，即祂是什么性质，祂如何是天主子，祂降生至（取得）人性血肉和完全人性的原因是什么；以及这位（圣子）的作为是什么，在谁身上和何时施行。此外，在天主的教训中还必须包含关于同类的受造物，以及其他理性的受造物，包括那些神圣者和那些从真福中堕落者，连同他们堕落的原因；还有关于灵魂的差异、这些差异的起源、世界的性质及其存在的原因。我们也必须认识那蔓延大地的巨大而可怕的邪恶的起源，以及它是仅限于这大地，还是也存在于别处。当圣神光照了真理的圣仆之灵魂时，这些及类似的目标摆在祂面前，但为了那些不能承受探究如此重大之事疲劳的人，还有第二个目的：即把关于前述主题的道理，隐藏在包含叙述的词句中，这些叙述传达关于有形受造物、人的创造、最初人类的连续后代直到他们众多、以及其他一些关于义人行为的记述，和这些同为人者偶然所犯的罪，以及有罪不虔者所行的不贞与邪恶。最值得注意的是，通过战争史、胜利者和战败者，某些奥秘向那些能检验这些叙述的人显示出来。更奇妙的是，真理的法则通过成文的立法被预言出来——所有这些都按照一种真正符合天主智慧的能力，以连续的系列描述出来。因为目的是要让属神真理的覆盖物——我指圣经的\Quote{肉体}部分——在许多情况下也不是没有益处的，而是能够按照众人的能力改善他们。\par

15. 但既然，如果立法的实用性、历史的顺序和美感普遍地不言自明，我们就不会相信圣经中除了明显的内容之外还能理解任何其他东西，天主的圣言便安排在某些地方引入一些绊脚石、障碍和不可能之事，进入法律和历史之中，以免我们被语言表面上的吸引力完全吸引，从而要么完全偏离（真正的）教义，学不到任何配得上天主的东西，要么不脱离字面，就无法认识更神圣的事物。此外，我们也必须知道，其主要目的是宣告那些已发生和应当发生之事中的\Quote{属神}关联，当圣言发现按照历史发生的事可以适应这些奥秘意义时，祂便加以利用，向大众隐藏更深的意义；但在叙述超感觉事物的发展时，并没有随之发生那些已由奥秘意义指明的事件，圣经便在历史中穿插了一些实际未发生的事件，有时是根本不可能发生的，有时是可能发生但并未发生的。有时会插入少数几个在字面意义上不真实的词句，有时则插入较多。在立法方面也应注意类似的做法，其中常能找到本身有用并适合当时立法时期的内容，有时也出现看似无用的内容，此外还有为了更精明、更探究的人而记录的不可能之事，好让他们努力探究所写的内容，从而恰当地确信：在这样的题材中，必须寻求一种配得上天主的意义。\par

16. 然而，圣神不仅对（基督来临）之前（所写的圣经）如此处理；祂既是同一圣神，又（出于）同一的天主，祂对福音作者和宗徒们也做了同样的事——因为连这些著作也并非通篇都是纯净的历史事件，它们确实按字面穿插了某些事件，但实际并未发生。就连法律和诫命也并非完全传达合理的内容。因为哪个有理解力的人会认为，第一日、第二日、第三日、晚上和早晨，在没有太阳、月亮和星辰的情况下存在呢？而且第一日似乎也没有苍穹？谁会如此愚蠢，以为天主像农夫一样，在伊甸东方栽种了一个乐园，并在其中放置了一棵可见可触摸的生命树，人用身体的牙齿品尝果子就能获得生命呢？再者，有人通过咀嚼取自树的果子而享有善恶？如果天主被称为在傍晚于乐园中行走，亚当躲在一棵树下，我想没有人会怀疑这些事比喻性地指示某些奥秘，历史只是表面发生，并非字面真实。加音从天主面前出去，对于有思想的人来说，显然会引导读者探究什么是天主的面前，以及从他面前出去意味着什么。还需要说更多吗？因为那些并非完全盲目的人可以收集无数类似的例子，它们被记录为发生过，但实际上并未字面发生。不仅如此，福音书本身也充满了这类叙述；例如，魔鬼带领耶稣上到一座高山上，以便从那里将全世界万国及其荣耀指给他看。因为那些不粗心阅读这类叙述的人，谁会不谴责那些认为用肉眼——需要高山才能看到下面和邻近的区域——能看见波斯人、西徐亚人、印度人和帕提亚人的国度，以及他们的王子如何在人前受荣耀呢？细心的读者可以在福音书中发现无数其他类似段落，从而确信在按字面记录的历史中，插入了实际并未发生的情况。\par

17. 如果我们来看梅瑟的立法，可以发现许多法律若按字面遵守，要么显出其不合理，要么显出其不可能。不合理之处在于：禁止人民吃秃鹫，然而即使在最严重的饥荒中，也没有人因匮乏而被迫食用这种鸟；又命令将出生八天未受割礼的孩童从人民中铲除，如果这法律真要按字面执行，那么他们的父亲或抚养他们的人就应该被处死。圣经说：\BibleQuote{凡未受割礼的男性，若不在第八天受割礼，应从人民中剪除。}若想看到立法中包含不可能之事，让我们观察山羊鹿是一种不存在的动物，但梅瑟却命令将其作为洁净牲畜献祭；而狮鹫，从未被人类驯服，立法者却禁止食用。再者，仔细思量关于安息日的著名禁令：\BibleQuote{你们各人要坐在自己的住处：第七天不可有人离开自己的地方}，会认为按字面遵守是不可能的：因为没有任何活物能整天坐着，保持坐姿不动。因此，那些属于割礼的人，以及所有希望只显示字面意义的人，根本不研究山羊鹿、狮鹫和秃鹫这类问题，反而在某些点上胡说八道，赘言连篇，提出无味的传统；例如，关于安息日，说每人只限两千肘。另有一些人，其中包括撒玛黎雅人多西托斯，谴责这种解释，认为人在安息日处于什么位置，就要在那里待到晚上。此外，安息日不负担子也是不可能的；因此犹太教师陷入了无数荒谬，说某种鞋是担子，另一种不是；有钉子的凉鞋是担子，无钉的则不是；同样，单肩扛的是重担，双肩扛的不是。\newline{}\par

18. 如果我们转到福音书并作类似考察，还有什么比按字面理解\BibleQuote{在路上不要向任何人请安}（这被单纯的人认为是救主对宗徒的吩咐）更不合理的呢？此外，命令打右脸也是极不可信的，因为每个打人的人，除非有身体缺陷，都是用\Emph{右}手打\Emph{左}脸。福音书中关于右眼\BibleQuote{犯罪}的说法也不可能按字面理解。因为，假设眼睛能\BibleQuote{犯罪}，那么当有两只眼观看时，为何只归咎于右眼？谁因看了妇女而动了淫念，会理性地将罪过只归给右眼，并扔掉\Emph{它}？宗徒又规定说：\BibleQuote{有人受割礼而被召吗？他不要掩盖割礼的记号。}首先，任何人都看出他说这话与眼前的话题无关。当他制定关于婚姻和洁净的规诫时，怎么会显得随意插入这些话呢？其次，谁会认为一个人若可能的话，因割礼被众人视为耻辱而努力掩盖割礼的记号是错的呢？\newline{}\par

我们说了这一切，是为了表明赐予我们圣经的那神圣能力的设计，是要我们不要只接受字面所呈现的（这类东西有时在字面意义上并非真实，而是荒谬和不可能的），而是某些事物被引入实际历史和立法中，其字面意义是有用的。\par

19. 但为避免有人以为我们断言：整体上并无真实的历史，只因某一段不是；或没有任何律法是应照字面遵守的，只因某一条照字面理解是荒谬或不可能的；或关于救主的记载并非感官可感知的方式为真；或祂的任何诫命和训诫都不必遵守——我们必须回答：关于某些事，我们十分清楚历史叙述是真实的；例如：\Person{亚巴郎}被埋葬在\Place{赫贝龙}的双洞中，\Person{依撒格}和\Person{雅各伯}以及他们各自的妻子也是如此；\Place{舍根}被赐给\Person{若瑟}作为一份产业；\Place{耶路撒冷}是\Place{犹大地}的首都，天主的圣殿在那里由\Person{撒罗满}建造；以及其他无数陈述。因为就其历史意义而言为真的经文，远多于那些穿插着纯属灵性意义的经文。再者，谁不会说这条命令——\BibleQuote{应孝敬你的父亲和你的母亲，为使你能得福}——抛开一切预像意义，是有益的且应当遵守的，宗徒\Person{保禄}也曾使用过这些相同的话呢？至于（禁止之事）：\BibleQuote{你不可奸淫}、\BibleQuote{你不可杀人}、\BibleQuote{你不可偷盗}、\BibleQuote{你不可作假见证}，又何必再提？此外，福音中所含的诫命，哪些应照字面遵守、哪些不应，是毫无疑义的；例如那句：\BibleQuote{我却对你们说：凡向自己弟兄发怒的}等等。还有：\BibleQuote{我却对你们说：什么誓都不可起。}在宗徒的著作中，字面意义也应保留：\BibleQuote{劝戒那些不守规矩的人，安慰怯懦的人，扶持软弱的人，对众人要忍耐}；即使那些渴求更深意义的人，能够保留天主智慧的深奥，同时也不废弃诫命的字面意义。然而，细心（的读者）会就某些点感到疑惑，如果没有长期研究，就无法表明这段被视为字面的历史是否真实发生，以及这条律法的字面意义是否应当遵守。因此，精确的读者必须服从救主的命令\BibleQuote{查考圣经}，仔细查明字面意义在多大程度上为真，在多大程度上不可能；并尽其所能，借助类似的陈述，追查出圣经中到处散布的、不能按字面意义理解的含义。\newline{}\par

20. 既然因此，如读者将清楚看到的，按字面理解的连贯性是不可能的，而所偏爱的含义并非不可能，甚至是真的，那么我们的目标必须是掌握整个意义，它以可理解的方式将字面上不可能的叙述与不仅不不可能而且历史上真实的内容联系起来，并且被预像地理解，因为它并未字面发生。因为，关于圣经，我们的意见是：全部圣经都有\Quote{属灵}的意义，但并非全部都有\Quote{属肉体}的意义，因为属肉体的意义在许多地方被证明是不可能的。因此，谨慎的读者必须高度重视这些神圣的书籍，视之为神圣的著作；我们认为其理解方式如下：——圣经记述天主在地上拣选了一个民族，他们用几个名字称呼它。因为这个民族整体称为以色列，也称为\Person{雅各伯}。当它在\Person{乃巴特}的儿子\Person{雅洛贝罕}的时代分裂时，隶属于他的十个支派称为以色列；其余两个支派，连同\Person{肋未}支派，由\Person{达味}的子孙统治，称为犹大。这个民族的人民所居住的全部领土，由天主赐给他们，称为犹大地，其首都是\Place{耶路撒冷}——即许多城市的首都，这些城市的名字散见于许多其他（圣经）章节，但在\BibleBook{若苏厄}{书}中一并列出。\newline{}\par

21. 情形既如此，宗徒提升我们分辨的能力（超越字面），在某处说：\Quote{你们看按肉体的以色列}，仿佛有一个\Quote{按圣神的以色列}。在另一处他又说：\Quote{因为血肉的子女不是天主的子女}；\Quote{凡从以色列生的，不都是以色列人}；\Quote{外表上作犹太人的不是真犹太人，外表上肉身的割损也不是真割损；而是内心上作犹太人的才是真犹太人，内心的割损是在圣神内，不在于文字。}因为若采纳关于\Quote{内心上的犹太人}的判断，我们必须理解，正如有\Quote{有形}的犹太种族，也有\Quote{内心上的犹太人}种族，灵魂因某种奥秘的理由获得了这种高贵。此外，有许多预言关于以色列和犹大，预言将要临到他们的事。这些关于他们的承诺，在表达上卑微，没有显出崇高的思想，也没有任何与天主的应许相称的事，难道不需要奥秘的解释吗？如果\Quote{属神的}应许藉可见的记号宣告，那么接受应许的人就不是\Quote{有形}的。关于内心上的犹太人和按照\Quote{内在的人}的以色列人，这些论述对不缺乏理解力的人已经足够，我们不再停留，回到主题，说雅各伯是十二位族长的父亲，他们是民长之父，这些民长又是其他以色列人的祖先。那么，\Quote{有形的}以色列人岂不将世系追溯到民长，民长追溯到族长，族长追溯到雅各伯，再往上溯；而\Quote{属神的}以色列人（\Quote{有形的}以色列人是其预像）岂不是出自家族，家族出自支派，支派出自某一位个体，他的世系不是\Quote{有形的}而是更优越的——他亦生于依撒格，依撒格生于亚巴郎——一切追溯到亚当，而宗徒宣告亚当是基督？因为一切与天主（万有之父）有关的家族，其起源都从基督开始，基督仅次于万有的天主和父，因此他是每一个灵魂的父，正如亚当是所有人的父。如果厄娃也被宗徒喻指教会，那么由厄娃生的加音以及其后所有追溯到厄娃的人，都是教会的预像，这一点不足为奇，因为他们在卓越的意义上都出自教会。\par

22. 现在，若有关以色列及其支派和家族的陈述令我们印象深刻，当救主说：\Quote{我被派遣，只是为了以色列家迷失的羊}时，我们并不像厄彼翁派那样理解这句话，他们理解力贫乏（其名来自理性的贫穷——\OriginalTerm{贫穷}{Ebion}在希伯来语中意为\Quote{贫穷}），竟以为救主是专门来\Quote{属血肉的}以色列人那里；因为\Quote{血肉的子女不是天主的子女}。再者，宗徒论及耶路撒冷说：\Quote{那在上面的耶路撒冷是自由的，她是我们的母亲。}在另一封书信中：\Quote{你们却来到了熙雍山和永生天主的城，天上的耶路撒冷，以及千万天使的盛会，和那些登记在天上的首生者的教会。}既然以色列属于灵魂族类，天上有一座耶路撒冷城，那么以色列各城的首府就是天上的耶路撒冷，她也就是全犹太地的首府。因此，凡预言耶路撒冷并关于它所说的一切，如果我们听从保禄的话如同天主的话、如同说出智慧者的话，我们必须理解圣经是在论及天上的城，以及圣地各城范围内的整个领土。因为也许救主正是向我们指这些城，当祂对那些因善管\Quote{米纳}而获得信赖的人，指派他们治理五座或十座城时。因此，如果关于犹太地、耶路撒冷、以色列、犹大、雅各伯的预言，不被我们按\Quote{肉性}理解，而表明某些这样的奥秘（如上所述），那么关于埃及和埃及人、巴比伦和巴比伦人、提洛和提洛人、漆冬和漆冬人或其他民族的预言，也不仅是关于这些\Quote{有形}的埃及人、巴比伦人、提洛人、漆冬人说的，也是关于其\Quote{属神}的对应者。因为若有\Quote{属神的}以色列人，也就有\Quote{属神的}埃及人和巴比伦人。因为厄则克耳论及埃及王法郎的事，完全不适用于某一位统治或据说统治埃及的人，仔细考虑的人就会明白。同样，关于提洛王所说的事，不能理解为某一位统治提洛的人。许多地方（尤其是依撒意亚）论及拿步高的事，也不能用那个人来解释。因为人拿步高既没有从天坠下，也不是晨星，也没有早晨在地上出现。有理解力的人不会将厄则克耳论埃及的话——即四十年内它将成为荒芜，人的足迹不再寻见，战争如此肆虐，血漫全地，高至膝盖——解释为位于身体被太阳晒黑的厄提约丕雅人旁边的那个埃及。\par

23. 或许正如在此处的人，按众人共有的死亡而死，因在此所行之事，被按各自罪过的比例安排，获得不同的处所，若他们被认为配得上称为\OriginalTerm{阴府}{Hades}的地方；同样，在那里死去的人，可以说是下降到这个阴府，被判定应得不同的住所——或善或恶——遍及整个大地空间，并从不同种类的父母降生，以致一个以色列人有时会落在\OriginalTerm{斯基泰人}{Scythians}中间，而一个埃及人降生在\Place{犹太}。然而救主来聚集以色列家迷失的羊；但许多以色列人不听从祂的教导，外邦人就被召叫……我们认为这些要点已隐藏在历史中。因为\BibleQuote{天国好像是藏在地里的宝贝；人找到了，就把它藏起来，高兴地去卖掉他所有的一切，买了那块地。}那么，让我们注意，圣经表面、肤浅和明显的意思，是否不像一块长满各种植物的田地，而隐藏在其中、不为人所见、仿佛埋在可见植物之下的东西，是智慧和知识的宝藏；圣神藉依撒意亚称之为隐晦、不可见和隐藏的，唯有天主能打破隐藏它们的铜门，冲开门上的铁闩，为能发现\WorkTitle{创世纪}中一切关于各种真实种类、种子（可以说是灵魂的种类）的陈述，它们与以色列亲近或疏远；以及七十个灵魂下到\Place{埃及}，为在那里成为\BibleQuote{如同天上的星辰。}但由于他们中不是所有人都是世界的光——\BibleQuote{凡出自以色列的，不都是以色列人}——他们就从七十个灵魂成为\BibleQuote{如同海边的沙粒。}\par

% structure: p0056 source=h2 target=subsection reason=relative-heading-rank; base=h2

\subsection{从拉丁文}

24. 圣祖们这一下到埃及，看来是天主眷顾赐予世界的恩惠，为光照他人并教导人类，好借此帮助其他灵魂从事启蒙工作。因为与他们交谈的殊荣首先赐给了他们，因为唯有他们的民族被称为看见天主；按解释，这正是\Quote{以色列}一词的含义。现在，按照这一观点，应当接受并解释以下说法：埃及因十灾而受惩罚，以让天主的子民离开；或关于在旷野中对子民所做之事，或关于由全体子民奉献建造会幕，或关于司祭圣服、礼仪器皿，因为正如所写的，它们确实含有\Quote{天上事物的阴影与样式}。因为保禄公开论及它们说：\Quote{它们侍奉天上事物的榜样和阴影。}此外，在同一法律中，包含了在圣地生活的规诫和制度。也向那些违犯法律的人发出了即将来临的警告；此外，为那些需要净化的人规定了不同的净化方式，因为他们易受频繁的玷污，借此他们最终能到达那一次净化之后不再容许玷污的境界。子民本身被数点，但并非全部；因为孩童的灵魂尚未达到按天主命令被数点的年龄；那些不能成为他人之首、反而隶属于他人如同首领的灵魂，被称为\Quote{女人}的，当然不包括在天主所命的数点中；只有被称为\Quote{男人}的才被数点，这表明女人不能单独被数，而是包含在被称为\Quote{男人}的人中。然而，那些特别属于神圣数目的人，是准备好出去为以色列人作战、并能对抗那些公共和私人仇敌的人，圣父将这些仇敌置于坐在祂右边的圣子之下，使祂能废除一切率领者和掌权者，并借祂的这些军队——他们为天主作战，不纠缠于世俗事务——推翻仇敌的国度；他们携带信德的盾牌，挥舞智慧的武器；其中也闪耀着望德与救恩的头盔，充满天主的胸膛佩戴着光明的胸甲。在我看来，这样的士兵就是在圣经中按天主命令被数点的人所指示和预备参与这类战争的。但在这些人中，更完善和卓越的是那些连头发都被数过的人。至于那些因罪受罚、身体倒在旷野的人，似乎与那些确实有了不小进步、但因各种原因未能达到完美终点的人相似；因为他们被记载为要么抱怨、要么崇拜偶像、要么行淫、要么做了心智甚至不应思想的那类恶事。我认为以下也并非毫无奥秘意义：某些以色列人拥有许多羊群和牲畜，预先占据了适合放牧和养畜的地方，那正是希伯来人的右手在战争中所取得的第一块土地。因为他们请求摩西赐予这地区，就被约但河的水分隔，不能拥有圣地中的产业。而这条约但河，按照天上事物的样式，似乎浇灌和滋润干渴的灵魂及其相邻的官能。与此相关，甚至以下说法也不显多余：摩西确实从天主那里听到\BibleBook{}{肋未纪}所记载的，而在\BibleBook{}{申命纪}中，人民是摩西的听众，从摩西学习他们不能从天主直接听到的事。因为\BibleBook{}{申命纪}被称为第二法律，对某些人而言，这似乎传达这样的意义：当通过摩西颁布的第一法律终结后，似乎又制定了第二部立法，由摩西特别传给其继承人若苏厄，这无疑被视为我们救主的预像，通过祂的第二法律——即福音的诫命——一切得以成全。\par

25. 然而，我们必须查看是否可能暗示了更深的意义：正如在\BibleBook{}{申命纪}中，立法比最初写成的那些书卷更清晰明确地显明，同样，通过救主在屈辱状态中完成的第一次来临——祂取了奴仆的形体——是否可能指向了祂在圣父光荣中更显赫、更著名第二次来临，而在其中\BibleBook{}{申命纪}的预像得以实现？那时在天国中，所有圣人将按永远的福音生活；正如在祂现今的来临时，祂成全了那具有未来美事阴影的法律，同样，通过那未来的光荣来临，现今来临的阴影也将得以完成和完美。因为先知论及此事说：\Quote{我们口中的气息，主基督，我们对祂说，愿在祢的荫庇下我们在万民中生活}；那时，祂将更相称地把所有圣人从暂时的福音转移到永远的福音，按若望在\BibleBook{若望}{默示录}中所用的称呼——\Quote{永远的福音}。\par

26. 但是，在这一切事上，我们只需将我们的理解适应于宗教的准则，并这样思考圣神的言语：不把这语言当作虚弱的人类雄辩的华丽辞藻，而是依照圣经的论述，相信\Quote{君王的一切光荣都在内在}，并且神圣意义的宝藏被封闭在普通文字的脆弱器皿内。如果还有好奇的读者想要追问个别要点的解释，就让他同我们一起听听，宗徒保禄如何藉着探查天主\Quote{深奥事理}的圣神的帮助，力图深入神圣智慧和知识之深渊，却仍未能达到终点，可以说，未能获得透彻的认识，于是在绝望和惊叹中高呼：\Quote{啊，天主的智慧和知识的富饶是多么深奥！}而他之所以发出这呼喊，是由于对达到完全理解感到绝望，请听他自己的话：\Quote{天主的判断是多么不可探测！祂的道路是多么不可追寻！}因为他并没有说天主的判断难以发现，而是说它们完全不可探测；也没有说追寻祂的道路（仅仅）困难，而是说它们完全不可追查。因为无论一个人在探究中前进多远，无论他藉着不懈的研究取得多么大的进步，甚至得到天主的恩宠的帮助，心志蒙受光照，他也不能达到他所探寻之事的终点。任何受造的心智都不能认为自己有可能以任何方式获得完全的理解（对事物）；反而在发现了其研究的某些对象之后，它又看到还有别的对象尚待寻求。即使它成功地掌握了这些，它又会看到许多其他的对象接踵而来，必须成为研究的主题。因此，最智慧的人撒罗满，以他的智慧观看万物的本性，说：\Quote{我说：我将成为有智慧的；智慧本身却离我甚远，比原来更远；深不可测，谁能找到？}依撒意亚也知道，万物的起源不能被必死的本性所发现，甚至不能被那些虽然比人性更属神、但本身仍是被造或形成的本性所发现；他知道，起源和终结都不能为这些本性中的任何一个所发现，就说：\Quote{你们要说明先前的事，我们就知道你们是神；或者宣告后来的事，我们便看见你们是神。}因为我的希伯来老师也曾这样教导：既然万物的起源和终结除了我们的主耶稣基督和圣神之外无人能理解，所以依撒意亚在异象中只提到两个色辣芬，他们用两个翅膀遮盖天主的容貌，用两个遮盖祂的双脚，用两个飞翔，彼此交替呼喊说：\Quote{圣！圣！圣！万军的上主！整个大地充满了祢的光荣。}我们敢于宣称，只有色辣芬才将他们的翅膀放在天主的容貌上，放在祂的双脚上，其意义是：无论是圣天使的万军，或是\Quote{宝座}，或是\Quote{宰制者}，或是\Quote{率领者}，或是\Quote{掌权者}，都不能完全理解万物的起源和宇宙的界限。但我们应当理解，那些被圣神登记入册的\Quote{圣者}和\Quote{德能}，非常接近那些起源本身，并达到了别人无法达到的高度；然而，无论这些\Quote{德能}通过天主子和圣神的启示学到了什么——他们当然能学到很多，地位高的比地位低的学得更多——他们仍不可能理解一切，正如经上所载：\Quote{天主的工程，多半是隐藏的。}\BibleRef{德 16:21}因此，也应当渴望每个人按自己的力量，不断向前伸展，\Quote{忘记背后}，趋向更善的作为和更清晰的领悟与理解，藉着我们的救主耶稣基督——愿光荣归于祂，直到永远！\par

27. 凡关心真理的人，对于言词和语言不必过于在意，因为每个民族都有不同的言语习惯；而应更注意言词所传达的意义，而非传达意义的言词本身，尤其是在如此重要和困难的事情上：例如，当研究是否存在一种\Quote{本质}，其中既没有颜色、形状、触觉，也没有大小，只能被理智所理解，任何人可随意称呼它；希腊人称此为\OriginalTerm{非物质的}{\GreekText{ἀσώματον}}，即\Quote{无形体的}，而圣经宣称它是\Quote{不可见的}，因为\Person{保禄}称\Person{基督}为\BibleQuote{不可见的天主的肖像}，又说，万物都是藉着\Person{基督}创造的，\BibleQuote{看得见的和看不见的}。由此表明，在被造物中，有一些\Quote{本质}按其特性是不可见的。它们本身虽然不是\Quote{有形的}，却使用形体，而且它们自身优于任何有形体的本质。但作为万有之本源和原因的圣三之\Quote{本质}，\BibleQuote{万有出于祂，万有藉着祂，万有在祂内}，不能相信祂是形体或在形体中，而是完全非物质的。现在，关于这些点的简短论述（虽然是因主题性质而离题）就足够了，以表明有些事物的意义完全不能用任何人类语言表达，而是更多地通过单纯的领悟而非任何言语特性来认识。圣经的理解也必须遵循这一原则，以便对圣经的陈述的判断不是根据文字的粗劣，而是根据\Person{圣神}的德能，祂默感了这些文字的写作。\par

% structure: p0061 source=h2 target=subsection reason=relative-heading-rank; base=h2

\subsection{关于圣父、圣子、圣神之教义以及前文讨论的其他主题的总结}

28. 现在，在尽我们所能对所讨论的主题进行了快速思考之后，是时候以总结我们在不同地方所说的话的方式，逐一概述各个要点，并首先重申我们关于圣父、圣子、圣神的结论。\par

既然天主圣父是不可见的，并与圣子不可分离，圣子并非像某些人以为的那样，是由圣父通过\OriginalTerm{发出}{prolation}而生的。因为如果圣子是圣父的\OriginalTerm{发出}{prolation}（\OriginalTerm{发出}{prolation}一词用以表示如动物或人类通常的生育），那么，\OriginalTerm{发出者}{prolator}和\OriginalTerm{被发出者}{prolatus}都必然是有形的。因为我们并不像异端者所认为的那样，说天主的本质的一部分转变成了圣子，或者说圣子是圣父从不存在之物中，即超出祂自己的本质之外受造的，以至于曾有一段时间祂不存在；相反，摒弃一切有形的概念，我们说圣言和智慧是从不可见的、非物质的本质中受生的，没有任何形体的感受，如同一种出自理智的意志行为。此外，既然祂被称为（祂的）爱子，那么以此方式称祂为（祂的）意志之子，也不显得荒谬。不仅如此，\Person{若望}也指明\BibleQuote{天主是光}，\Person{保禄}也宣告圣子是永恒之光的辉煌。因此，正如光绝不能没有光辉而存在，同样，圣子也不能被理解为没有圣父而存在；因为祂被称为\BibleQuote{祂位格的印像}，以及圣言和智慧。那么，怎么能断言曾有一段时间祂不是圣子呢？因为那无异于说曾有一段时间祂不是真理，也不是智慧，也不是生命，尽管在这一切中祂被认为是天主圣父的圆满本质；因为这些事物不能与祂分离，甚至不能与祂的本质分开。虽然这些属性在理智上被认为是众多，但在其本性和本质上是一个，并且在它们之中有神性的圆满。我们现在使用的这个表达——\Quote{从来不曾有一个时间祂不存在}——应当有所变通地理解。因为\Quote{何时}或\Quote{从未}这些词本身具有与时间相关的含义，而关于圣父、圣子、圣神的陈述应当理解为超越一切时间、一切世代和一切永恒。因为唯独天主圣三超越不仅时间性理智甚至永恒性理智的理解；而不在圣三之内的其他事物则应以时间和世代来衡量。因此，天主的这位圣子，作为在起初就与天主同在的圣言，没有人会合乎逻辑地认为祂被包含在任何地方；也不因祂是\Quote{智慧}、或\Quote{真理}、或\Quote{生命}、或\Quote{义德}、或\Quote{圣化}、或\Quote{救赎}而被认为是受地方限制的：因为所有这些属性都不需要空间才能行动或运作，而是每个属性都应当被理解为指那些分享祂的德能和工化的人。\par

29. 现在，若有人要说，那些分有天主\Quote{圣言}，或他的\Quote{智慧}、他的\Quote{真理}、他的\Quote{生命}的人，似乎将圣言和智慧本身包含在一个地方，我们应当回答他：毫无疑问，基督就其为\Quote{圣言}或\Quote{智慧}或其他一切而言，曾在保禄内，因此他说：\BibleQuote{你们寻求基督在我内说话的凭证吗？}又说：\BibleQuote{我生活，已不是我生活，而是基督在我内生活。}那么，既然基督在保禄内，谁还会怀疑他同样在伯多禄内、在若望内、以及在每一位圣人内呢？不仅在地上的人内，也在天上的人内。因为说基督在伯多禄和保禄内，却不在总领天使弥额尔内，也不在加俾额尔内，这是荒谬的。由此清楚表明：天主子具有神性，并未被关闭在某个地方，否则他就只在那个地方，而不在其他地方。然而，由于其无形本性的尊威，他不受任何地方的限制，因此也不能认为他缺乏于任何地方。但唯一的区别是：尽管他如我们所说，在不同的人内——如伯多禄、保禄、弥额尔或加俾额尔——，但他并非同等方式存在于所有受造物中。因为他在总领天使内比在其它圣人内更充分、更清晰、并可说是更明显地临在。这从以下事实可以显明：当所有圣人都达至完美的顶峰时，他们被说成是相似或等同于天使，正如福音中所宣布的。由此可知，基督在每个人内按其功劳的大小而存在。\par

30. 简短地重申了有关天主圣三本性的这些要点之后，接下来我们也要简要地留意这一陈述，即：借着圣子创造了\BibleQuote{天上地上的一切，有形的、无形的，无论是宝座，或是宰制者，或是率领者，或是掌权者：万物都是借着他，并为了他而受造的；他在万有之先，万有都靠他而存在，他是元首。}与此一致，若望在他的福音中也说：\BibleQuote{万物是借着他造的；凡受造的，没有一样不是借着他造的。}达味在暗示整个天主圣三的奥迹参与万物创造时，说：\BibleQuote{因上主的圣言，诸天形成；因他口中的圣神，万象生成。}\par

在此之后，我们恰当地提醒（读者）论及天主独生子的身体的来临和降生成人。关于此，我们不应认为他全部的神性尊威都被限制在他瘦弱身体的范围内，以致整个天主的\Quote{圣言}、他的\Quote{智慧}、\Quote{本质的真理}及\Quote{生命}或与圣父分离，或被束缚和局限在他身体身份的狭隘中，并且不被认为在其他地方运作。相反，虔诚之人的谨慎承认应介于两者之间：既不应相信基督内缺乏任何神性，也不应相信他与无所不在的圣父的本质有任何分离。因为若翰洗者在耶稣身体不在时对群众所说的这话似乎暗示了某种类似的意思：\BibleQuote{有一位站在你们中间，你们却不认识他；他是在我以后来的，我就是解他的鞋带也不配。}因为对于他身体而言不在场的那一位，却站在那些天主子身体不在其中的人中间，这当然是不能说的。\par

31. 然而，谁也不应以为我们借此肯定天主子的神性有一部分在基督内，而其余部分则在别处或各处，这可能是那些对无形体且不可见的本质一无所知者的看法。祂在万有之中，贯通万有，超越万有，其方式正如我们上文所述，即祂被理解为或是\Quote{智慧}，或是\Quote{圣言}，或是\Quote{生命}，或是\Quote{真理}；这种理解方式无疑排除了所有空间性的限制。因此，天主子为拯救人类，决意显现于世人，并居住在他们中间，所取的不止是人的身体，如某些人所猜想，且亦取了一个灵魂——这灵魂在本质上确实与我们的灵魂相似，但在意志和能力上却与祂自己相似，并能万无一失地完成\Quote{圣言}和\Quote{智慧}的一切意愿和安排。救主在福音中清楚表明祂有一个灵魂，当祂说：\BibleQuote{没有人夺去我的生命，是我自己舍弃的；我有权舍弃生命，也有权再取回它。}又说：\BibleQuote{我的心灵忧闷得要死。}又说：\BibleQuote{现在我心灵烦乱。}因为天主的\Quote{圣言}不能被视为一个\Quote{忧闷烦乱的灵魂}，因祂以天主的权威说：\BibleQuote{我有权舍弃生命。}我们也不断言天主子在那灵魂内如同在\Person{保禄}或\Person{伯多禄}及其他圣人们的灵魂内一样——人们相信基督在他们内说话，如同在保禄内那样。但关于所有这些，我们应坚持圣经所言：\Quote{没有一人是洁净的，即使他的生命只持续了一天。}然而，在耶稣内的这个灵魂，在尚未认识恶之前，就选择了善；并且因为祂喜爱正义，憎恨邪恶，所以天主\BibleQuote{以喜乐油傅抹了祂，胜过祂的同伴。}因此，当祂以无玷的结合与天主的\Quote{圣言}联合时，祂就被以喜乐油傅抹；并且唯有这个灵魂在所有灵魂中不可能犯罪，因为它能够良好而完全地接受天主子；因此它也与祂合一，以祂的称号被称呼，被称为耶稣基督，万物藉着祂而受造。关于这灵魂，因它已接受了天主的全部智慧、真理和生命，我认为宗徒也曾说：\BibleQuote{我们的生命与基督一同藏在天主内；当基督——我们的生命显现时，我们也要与祂一同出现在光荣中。}这里所理解的、被称为藏在天主内、后来要显现的基督，还能是谁呢？岂不是那曾被记载以喜乐油傅抹的祂吗？即本质上充满天主的祂，现在被说成隐藏在天主内。正因如此，基督被立为所有信徒的榜样，因为祂始终——甚至在尚未认识任何恶之前——就选择了善，喜爱正义，憎恨邪恶，因此天主以喜乐油傅抹了祂；同样，每一个人在跌倒或犯罪之后，也应洁净自己的污秽，以祂为榜样，以祂为旅途的向导，踏上德行的陡峭之路，如此或许藉此方法，尽可能透过效法祂，得以分享天主性，正如圣经所言：\BibleQuote{那说自己相信基督的，就该照样行走，如同祂行走了一样。}\par

那么，这\Quote{圣言}和这\Quote{智慧}，藉着效法它们，我们被称为智者或理性的存有，成了\BibleQuote{对一切人成为一切，为的是要赢得一切}；并且因为它成为软弱的，故论及它说：\BibleQuote{祂虽然因软弱而被钉在十字架上，却因天主的德能而活着。}最后，对软弱的格林多人，\Person{保禄}声明他\BibleQuote{只知道耶稣基督，并被钉在十字架上的耶稣基督。}\par

32. 有些人确实希望将宗徒的以下话语应用于那从\Person{玛利亚}取得肉身的灵魂本身，即：\BibleQuote{祂虽具有天主的形象，却不以自己与天主同等为应当把持不舍的，反而空虚了自己，取了仆人的形象。}因为祂无疑藉着更好的榜样和训练，将它恢复到天主的形象，并召回它到祂所空虚自己的丰满。\par

正如藉着参与天主子，人得以成为义子；藉着参与那在天主内的智慧，人成为有智慧的；同样，藉着参与圣神，人得以成为圣洁和属神的。因为参与圣神与参与父和子的圣神是同一回事，因为天主圣三的本性是同一且无形体的。我们论及灵魂的参与所说的一切，也应同样理解为适用于天使和天上的权力，因为每一个理性的受造物都需要参与天主圣三。\par

关于这个可见世界的计划——既然通常提出的最重要问题之一是它存在的方式——我们在前面各页已尽其所能地讨论过，这是为了那些惯于在我们宗教中寻求信仰根据的人，也为了那些兴起异端问题反对我们、并且惯于议论\Quote{物质}这个他们还不理解之词的人。关于这个主题，我现在认为有必要简要提醒（读者）。\par

33. 首先应指出，我们迄今为止在正典圣经中从未发现\Quote{质料}一词用于指称据说为物体基体的那种实体。因为在依撒意亚那句话中，当论及那些被指定要受惩罚的人时，他说：\Quote{他将吞下\OriginalTerm{质料}{\GreekText{ὕλη}}，即质料，如同干草，}\Quote{质料}一词被用来代替\Quote{罪恶}。倘若\Quote{质料}一词碰巧出现在其他任何段落中，依我之见，它绝不会具有我们现在所要探求的含义，除非在名为\BibleBook{智慧}{篇}的书中，这部作品并非被所有人视为权威。然而，在那本书中，我们读到如下内容：\Quote{因为你全能的手，从无形体的质料中创造了世界，也能轻易地派遣一群熊和猛狮来惩罚他们。}确实有许多人认为，梅瑟在\BibleBook{创世}{纪}开头所用的语言本身即指明万物由以构成的质料：\Quote{在起初，天主创造了天地；地是无形可见且未经整理的：}因为梅瑟藉着\Quote{无形可见且未经整理}这些话，似乎不过是指无形体的质料。但如果这确实是质料，那么物体的原始元素显然并非不可改变。因为那些主张\OriginalTerm{原子}{atoms}（无论是不可再分割的微粒，还是分割成相等部分的微粒）或任何一种元素作为物体本质的人，不可能在万有的第一原理中恰当地将\Quote{质料}一词列入其中。因为如果他们坚持质料是每一物体的基体——一种可转化、可改变或在其所有部分中均可分的实体——那么他们不会（恰当地）主张它没有性质而存在。我们赞同他们，因为我们完全否认质料应被称为\Quote{\OriginalTerm{非受生}{unbegotten}}或\Quote{\OriginalTerm{非受造}{uncreated}}，这与我们之前的陈述一致。那时我们指出：从水、地、空气或热中，不同种类的树木产生了不同种类的果实；或者我们展示了火、空气、水、地相互转化，一个元素通过一种相互的亲缘关系分解为另一个元素；并且我们还证明了：从人或动物的食物中产生了肉体的实体，或者天然种子的水分转化为坚固的肉和骨——所有这些都证明身体的实体是可变的，并且可以从一种性质过渡到所有其他性质。\par

34. 然而，我们不可忘记，实体绝不能没有性质而存在，并且正是仅凭理智的活动，这作为物体基体且能接受性质的（实体）才被揭示为质料。诚然，有些人出于更深入探究这些主题的愿望，竟敢断言物体性无非就是性质。因为如果坚硬与柔软、热与冷、潮湿与干燥是性质，而如果当这些或这类（性质）被除去后，剩下的东西就认为是无，那么一切事物都将显得是\Quote{性质}。因此，那些作出这些主张的人也试图论证：既然所有说质料是非受造的人都承认性质是由天主创造的，那么就可以由此表明，即便按照他们，质料也并非非受造；因为性质构成一切，而所有人都毫无异议地宣称这些是由天主造的。另外，那些认为性质是从外部附加于某种基体质料之上的人，使用此类例证：例如，\Person{保禄}无疑不是沉默就是在说话，或是在守夜，或是在睡觉，或是保持某种身体姿势；因为他不是坐着就是站着，或是躺着。因为这些是属于人的\Quote{\OriginalTerm{偶性}{accidents}}，没有这些偶性，人几乎从未被找到。然而，我们关于人的概念并没有将其中任何一点作为其定义；但我们通过它们如此理解和看待他，以至于我们完全不考虑他（特定）状态的原因，无论是在守夜、睡觉、说话、保持沉默，还是任何其他必然发生在人身上的行动。因此，如果有人能设想\Person{保禄}没有所有这些可能发生的事情，那么他也同样能理解这个没有性质的基体（实体）。于是，当我们的心灵从它的概念中除去所有性质，并且仿佛只凝视于基体元素本身，全神贯注于它，丝毫不涉及实体的软硬、冷热、干湿，那么通过这种有点模拟的思考过程，它似乎就会看到完全脱离一切性质的质料。\par

35. 但或许有人会问，我们是否能从圣经中为这样的理解找到依据。我认为在\WorkTitle{圣咏}中，当先知说：\Quote{我的眼睛看见了你的不完全；} 时，就表明这种看法；先知以此更锐利的目光审视万物的首要原理，在思想和想象中只将\OriginalTerm{物质}{matter}及其\OriginalTerm{性质}{qualities}分开，从而看到了天主的\OriginalTerm{不完全}{imperfection}——这当然是藉性质的增添而得以完善的。\Person{哈诺客}在他的\WorkTitle{哈诺客书}中也这样说道：\Quote{我甚至走到了不完全之中。}我认为这句话可以作类似的理解，即先知的心灵在审视和探究所有可见之物时，一直到达了那最初的开端，在那里他看见了没有\Quote{性质}的不完全物质。因为在\WorkTitle{哈诺客书}的同一本书中记载着：\Quote{我看见了整个物质。}其意思是如此理解的，犹如他所说的：\Quote{我清楚地看见了物质的一切划分，这些划分从一体分裂为各种个体种类，无论是人、动物、天空、太阳，或世界上所有其他事物。}在论述了这些之后，我们在前文中已竭尽全力证明，所有存在之物皆由天主所造，没有一样不是被造的，除了圣父、圣子和圣神的本性；并且，本性为善的天主，渴望拥有一些祂可以施恩的对象，这些对象也因领受祂的恩惠而喜乐，便创造了配得此恩的受造物，即那些能以相称的方式接受祂的受造物，祂说他们也由祂而生作为祂的儿子。此外，祂还按数目和度量创造了万物。因为在天主面前，没有什么是无限制或无尺度的。祂的能力包含一切，而祂自身却不被任何受造之物的能力所包含，因为那本性只有自己知道。因为只有圣父认识圣子，只有圣子认识圣父，只有圣神甚至洞察天主的深奥。因此，所有受造之物，即理性受造物的数目或物质身体的度量，都被祂区分在一定的数目或度量之内；因为，既然理智性本性有必要使用身体，而这一本性因其受造的条件而显示出可改变和可转化（因为不存在的事物开始存在，正因这一情况就被说成具有可变的性质），它就不能将善或恶作为本质属性，而只能作为其存在的\OriginalTerm{偶性}{accidental attribute}。因此，正如我们所说，既然理性本性是可变的、可更改的，以致它根据自身的功过使用这样或那样性质的身体外衣，那么，由于天主预知灵魂或灵性能力会有多样性，祂就必须也创造一种身体性本性，其性质可按照造物主的意愿改变为所需的一切。并且，这种身体性本性必须与那些需要它作为外衣的事物共存：因为将永远有理性的本性需要身体外衣；因此将永远有一个身体性本性，其外衣必须由理性受造物使用，除非有人能用论证证明理性本性可以不依赖身体而生活。但这对我们的理解来说是何等困难——不，几乎不可能——我们在前文讨论各个主题时已经表明。\par

36. 我认为，若以尽可能简洁的方式重述我们对理性本性不朽的看法，并不违背我们工作的性质。凡分享某事物者，无疑与分享同一事物之他者具备相同本质与本性。例如，既然所有眼睛都分享光，那么所有分受光的眼睛便同属一个本性；然而，虽然每只眼睛都分享光，但因一只眼睛看得更清晰，另一只更模糊，并非每只眼睛都同等地分受光。同样，所有听觉都接受声音或声响，因此所有听觉同属一个本性；但每个人的听力有快慢之别，取决于其听觉的清晰与健全。现在让我们从这些感官的例证转到对理智事物的思考。每一分享理智之光的理智，无疑与每一以相似方式分享理智之光的理智同属一个本性。因此，如果天上的德能分享理智之光，即分享天主性，因为它们参与智慧与圣德；而人的灵魂也分受了同一光明与智慧，因而彼此同属一个本性与一个本质——那么，既然天上的德能是不朽坏且不死的，人的灵魂的本质也将是不死且不朽坏的。不仅如此，而且因为圣父、圣子、圣神的本性——唯独其理智之光为一切受造物所分享——是不朽坏且永恒的，那么，凡分享那永恒本性的每一实体，都理应并必然永久存在，且是不朽坏、永恒的；如此，天主的永恒的善也可在这方面被理解，即那些获得其恩惠者也是永恒的。然而，正如在上述例子中，我们观察到分享光时的差异，即观看者的视线被描述为较迟钝或较敏锐；同样，在分享圣父、圣子、圣神时，也当根据心灵的热忱或能力的程度而注意其差异。若非如此，我们便要考虑：说那能（接受）天主的理智的本质可被毁灭，是否似乎是一种不虔的行为？仿佛它感觉并理解天主的能力本身，不足以使其永远存在；尤其因为，即使理智因疏忽而偏离对天主纯洁而圆满的接受，它内仍含有某些种子，能恢复并更新至更好的理解，正如\Quote{内在的}（又被称为\Quote{理性的}人）按照\BibleQuote{那创造他的天主的肖像和模样}而更新。因此先知说：\BibleQuote{大地四极将想起上主，并归向祂；万邦的家族都要在祢面前朝拜。}\par

37. 如果任何人胆敢将本质的败坏归给那按天主的肖像和模样所造者，那么，依我之见，这种不虔的指控甚至延及天主子自身，因为圣经称祂为天主的肖像。或者，持此意见者必然攻击圣经的权威，圣经说人按天主的肖像受造；在他身上显然可发现天主肖像的痕迹，这不是通过可朽坏的身体外形，而是通过心灵的智慧、正义、节制、德行、智慧、纪律；总之，通过一切德性的全体，这些德性内在于天主的本质，并可通过人的勤勉和效法天主而进入人内；正如主在福音中所暗示的，祂说：\BibleQuote{所以你们应当慈悲，如同你们的父一样慈悲。}又说：\BibleQuote{你们应当是成全的，如同你们的父一样成全。}由此清楚显示，这一切德性永远存在于天主内，它们绝不能接近或离开祂；而人只是逐渐地、一个一个地获得它们。因此，通过这些方式，人似乎与天主有一种类似的关系；并且，既然天主知晓万事，没有任何理智性的事物本身能逃避祂的注意（因为天主圣父、祂的独生子以及圣神不仅知道他们所创造之物，也知晓自身），那么，一个理性的理智，从微小进步到伟大，从可见进步到不可见，也能达到更完善的知识。因为它居于身体内，并从感性事物本身（即有形之物）前进到理智之物。但为避免我们关于理智之物不能通过感官认知的陈述显得不妥，我们将引用\Person{撒罗满}的例证，他说：\BibleQuote{你也要寻得一种神圣的感觉。}由此他表明，理智之物不应通过身体感官去寻求，而应通过另一种他称之为\Quote{神圣的}感官。以此感官，我们必须注视我们上面枚举的每一个理性存在；以此感官，必须理解我们所说的话语，并权衡我们所写下的论述。因为天主性甚至知道我们静默中内心所转的意念。关于我们已讨论的这些事，或从它们引出的其他事，应根据上述规则形成我们的意见。\par

\SourceRecord{Translated by Frederick Crombie.；Ante-Nicene Fathers；Vol. 4.；Edited by Alexander Roberts, James Donaldson, and A. Cleveland Coxe.；Buffalo, NY: Christian Literature Publishing Co.,；1885.；<http://www.newadvent.org/fathers/04124.htm>.}
